\mysection{\PrintfSeveralArgumentsSectionName}

Now let's extend the \IT{\HelloWorldSectionName}~(\myref{sec:helloworld}) example, replacing \printf in
the \main function body with this:

\lstinputlisting[label=hw_c,style=customc]{patterns/03_printf/1.c}

% sections
\documentclass[a4paper,oneside]{book}

% http://www.tex.ac.uk/FAQ-noroom.html
\usepackage{etex}

\usepackage[table,usenames,dvipsnames]{xcolor}

\usepackage{fontspec}
% fonts
%\setmonofont{DroidSansMono}
%\setmainfont[Ligatures=TeX]{PT Sans}
%\setmainfont{DroidSans}
\setmainfont{DejaVu Sans}
\setmonofont{DejaVu Sans Mono}
\usepackage{polyglossia}
\defaultfontfeatures{Scale=MatchLowercase} % ensure all fonts have the same 1ex
\usepackage{ucharclasses}
\usepackage{csquotes}

\ifdefined\ENGLISH
%\wlog{main ENGLISH defined OK}
\setmainlanguage{english}
\setotherlanguage{russian}
\fi

\ifdefined\RUSSIAN
\setmainlanguage{russian}
%\newfontfamily\cyrillicfont{LiberationSans}
%\newfontfamily\cyrillicfonttt{LiberationMono}
%\newfontfamily\cyrillicfontsf{lmsans10-regular.otf}
\setotherlanguage{english}
\fi

\ifdefined\GERMAN
%\wlog{main GERMAN defined OK}
\setmainlanguage{german}
\setotherlanguage{english}
\fi

\ifdefined\SPANISH
\setmainlanguage{spanish}
\setotherlanguage{english}
\fi

\ifdefined\ITALIAN
\setmainlanguage{italian}
\setotherlanguage{english}
\fi

\ifdefined\BRAZILIAN
\setmainlanguage{portuges}
\setotherlanguage{english}
\fi

\ifdefined\POLISH
\setmainlanguage{polish}
\setotherlanguage{english}
\fi

\ifdefined\DUTCH
\setmainlanguage{dutch}
\setotherlanguage{english}
\fi

\ifdefined\TURKISH
\setmainlanguage{turkish}
\setotherlanguage{english}
\fi

\ifdefined\THAI
\setmainlanguage{thai}
%\usepackage[thai]{babel}
%\usepackage{fonts-tlwg}
\setmainfont[Script=Thai]{TH SarabunPSK}
\newfontfamily{\thaifont}[Script=Thai]{TH SarabunPSK}
\let\thaifonttt\ttfamily
\setotherlanguage{english}
\fi

\ifdefined\FRENCH
\setmainlanguage{french}
\setotherlanguage{english}
\fi

\ifdefined\JAPANESE
\usepackage{xeCJK}
\xeCJKallowbreakbetweenpuncts
\defaultfontfeatures{Ligatures=TeX,Scale=MatchLowercase}
\setCJKmainfont{IPAGothic}
\setCJKsansfont{IPAGothic}
\setCJKmonofont{IPAGothic}
\DeclareQuoteStyle{japanese}
  {「}
  {」}
  {『}
  {』}
\setquotestyle{japanese}
\setmainlanguage{japanese}
\setotherlanguage{english}
\fi

\usepackage{microtype}
\usepackage{fancyhdr}
\usepackage{listings}
%\usepackage{ulem} % used for \sout{}...
\usepackage{url}
\usepackage{graphicx}
\usepackage{makeidx}
\usepackage[cm]{fullpage}
%\usepackage{color}
\usepackage{fancyvrb}
\usepackage{xspace}
\usepackage{tabularx}
\usepackage{framed}
\usepackage{parskip}
\usepackage{epigraph}
\usepackage{ccicons}
\usepackage[nottoc]{tocbibind}
\usepackage{longtable}
\usepackage[footnote,printonlyused,withpage]{acronym}
\usepackage[]{bookmark,hyperref} % must be last
\usepackage[official]{eurosym}
\usepackage[usestackEOL]{stackengine}

% ************** myref
% http://tex.stackexchange.com/questions/228286/how-to-mix-ref-and-pageref#228292
\ifdefined\RUSSIAN
\newcommand{\myref}[1]{%
  \ref{#1}
  (стр.~\pageref{#1})%
  }
% FIXME: I wasn't able to force varioref to output russian text...
\else
\usepackage{varioref}
\newcommand{\myref}[1]{\vref{#1}}
\fi
% ************** myref

\usepackage{glossaries}
\usepackage{tikz}
%\usepackage{fixltx2e}
\usepackage{bytefield}

\usepackage{amsmath}
\usepackage{MnSymbol}
\undef\mathdollar

\usepackage{float}

\usepackage{shorttoc}
\usetikzlibrary{calc,positioning,chains,arrows}
\usepackage[margin=0.5in,headheight=15.5pt]{geometry}
% FIXME would be good without "margin="...
%\usepackage[headheight=15.5pt]{geometry}

%--------------------
% to prevent clashing of numbers and titles in TOC:
% https://tex.stackexchange.com/a/64124
\usepackage{tocloft}% http://ctan.org/pkg/tocloft
\makeatletter
\renewcommand{\numberline}[1]{%
  \@cftbsnum #1\@cftasnum~\@cftasnumb%
}
\makeatother
%--------------------

\newcommand{\footnoteref}[1]{\textsuperscript{\ref{#1}}}

%\definecolor{lstbgcolor}{rgb}{0.94,0.94,0.94}

% I don't know why this voodoo works, but without all-caps, it can't find LIGHT-GRAY color. WTF?
% see also: https://tex.stackexchange.com/questions/64298/error-with-xcolor-package
\definecolor{light-gray}{gray}{0.87}
\definecolor{LIGHT-GRAY}{gray}{0.87}
\definecolor{RED}{rgb}{1,0,0}
\makeindex

% TODO: split this file by languages

\newcommand{\myhrule}{\begin{center}\rule{0.5\textwidth}{.4pt}\end{center}}

% TODO: find better name:
\newcommand{\myincludegraphics}[1]{\includegraphics[width=\textwidth]{#1}}
\newcommand{\myincludegraphicsSmall}[1]{\includegraphics[width=0.3\textwidth]{#1}}

\newcommand{\myincludegraphicsSmallOrNormalForEbook}{\myincludegraphicsSmall}

% advertisement:
%\ifdefined\ENGLISH{}
%\newcommand{\mysection}[2][]{
%    \clearpage\input{ad_EN}\clearpage
%    \ifx&#1& \section{#2}
%    \else    \section[#1]{#2}
%    \fi
%}
%\else
\newcommand{\mysection}[2][]{
    \ifx&#1& \section{#2}
    \else    \section[#1]{#2}
    \fi
}
%\fi

%\newcommand*{\TT}[1]{\texttt{#1}}
% synonyms, so far
\newcommand*{\InSqBrackets}[1]{\lbrack{}#1\rbrack{}}

% too contrast text?
%\newcommand*{\TT}[1]{\colorbox{light-gray}{\texttt{#1}}}
%\newcommand*{\GTT}[1]{\colorbox{light-gray}{\texttt{#1}}}
\newcommand*{\TT}[1]{\texttt{#1}}
\newcommand*{\GTT}[1]{\texttt{#1}}

\newcommand*{\IT}[1]{\textit{#1}}
\newcommand*{\EN}[1]{\iflanguage{english}{#1}{}}

\newcommand{\HeaderColor}{\cellcolor{blue!25}}

\ifdefined\RUSSIAN{}
\newcommand*{\RU}[1]{\iflanguage{russian}{#1}{}}
\else
\newcommand*{\RU}[1]{}
\fi

\ifdefined\CHINESE{}
\newcommand*{\CN}[1]{\iflanguage{chinese}{#1}{}}
\else
\newcommand*{\CN}[1]{}
\fi

\ifdefined\SPANISH{}
\newcommand*{\ES}[1]{\iflanguage{spanish}{#1}{}}
\else
\newcommand*{\ES}[1]{}
\fi

\ifdefined\ITALIAN{}
\newcommand*{\ITA}[1]{\iflanguage{italian}{#1}{}}
\else
\newcommand*{\ITA}[1]{}
\fi

\ifdefined\BRAZILIAN{}
\newcommand*{\PTBR}[1]{\iflanguage{portuges}{#1}{}}
\else
\newcommand*{\PTBR}[1]{}
\fi

\ifdefined\POLISH{}
\newcommand*{\PL}[1]{\iflanguage{polish}{#1}{}}
\else
\newcommand*{\PL}[1]{}
\fi

\ifdefined\GERMAN{}
%\newcommand*{\DE}[1]{\iflanguage{german}{#1}{}}
\newcommand*{\DE}[1]{#1}
\else
\newcommand*{\DE}[1]{}
\fi

\ifdefined\THAI{}
\newcommand*{\THA}[1]{\iflanguage{thai}{#1}{}}
\else
\newcommand*{\THA}[1]{}
\fi

\ifdefined\DUTCH{}
\newcommand*{\NL}[1]{\iflanguage{dutch}{#1}{}}
\else
\newcommand*{\NL}[1]{}
\fi

\ifdefined\FRENCH{}
\newcommand*{\FR}[1]{\iflanguage{french}{#1}{}}
\else
\newcommand*{\FR}[1]{}
\fi

\ifdefined\JAPANESE{}
\newcommand*{\JPN}[1]{\iflanguage{japanese}{#1}{}}
\else
\newcommand*{\JPN}[1]{}
\fi

\ifdefined\TURKISH{}
\newcommand*{\TR}[1]{\iflanguage{turkish}{#1}{}}
\else
\newcommand*{\TR}[1]{}
\fi

\newcommand{\ESph}{\ES{{\color{red}Spanish text placeholder}}}
\newcommand{\PTBRph}{\PTBR{{\color{red}Brazilian Portuguese text placeholder}}}
\newcommand{\PLph}{\PL{{\color{red}Polish text placeholder}}}
\newcommand{\ITAph}{\ITA{{\color{red}Italian text placeholder}}}
\newcommand{\DEph}{\DE{{\color{red}German text placeholder}}}
\newcommand{\THAph}{\THA{{\color{red}Thai text placeholder}}}
\newcommand{\NLph}{\NL{{\color{red}Dutch text placeholder}}}
\newcommand{\FRph}{\FR{{\color{red}French text placeholder}}}
\newcommand{\JPNph}{\JPN{{\color{red}Japanese text placeholder}}}
\newcommand{\TRph}{\TR{{\color{red}Turkish text placeholder}}}

\newcommand{\TitleRE}{%
	\EN{Reverse Engineering for Beginners}%
	\RU{Reverse Engineering для начинающих}%
	\ES{Ingenier\'ia Inversa para Principiantes}%
	\PTBRph{}%
	\DE{Reverse Engineering für Einsteiger}%
	\PLph{}%
	\ITA{Reverse Engineering per Principianti}%
	\THAph{}%
        \NLph{}%
	\FR{Reverse Engineering pour Débutants}%
	\JPN{リバースエンジニアリング入門}%
	\TR{Acemiler için Tersine Mühendislik}%
}

\newcommand{\TitleUAL}{%
	\EN{Understanding Assembly Language}%
	\RU{Понимание языка ассемблера}%
	\ES{Understanding Assembly Language}%     %TBT
	\PTBR{Understanding Assembly Language}%   %TBT
	\DE{Understanding Assembly Language}%     %TBT
	\PL{Understanding Assembly Language}%     %TBT
	\ITA{Capire il linguaggio Assembly}%
	\THA{Understanding Assembly Language}%    %TBT
	\NL{Understanding Assembly Language}%     %TBT
	\FR{Comprendre le langage d'assemblage}%
	\JPN{Understanding Assembly Language}%    %TBT
}

\ifdefined\UAL
\newcommand{\TitleMain}{\TitleUAL}
\newcommand{\TitleAux}{\TitleRE}
\else
\newcommand{\TitleMain}{\TitleRE}
\newcommand{\TitleAux}{\TitleUAL}
\fi

\ifdefined\RUSSIAN
\newcommand{\AUTHOR}{Денис Юричев}
\else
\newcommand{\AUTHOR}{Dennis Yurichev}
\fi

%\newcommand{\EMAIL}{dennis(@)yurichev.com}
\newcommand{\EMAIL}{dennis@yurichev.com}

\newcommand*{\dittoclosing}{---''---}
\newcommand*{\EMDASH}{\RU{~--- }\EN{---}\ES{---}\PTBR{---}\PL{---}\ITA{---}\DE{---}\THA{---}\NL{---}\JPN{---}}
\newcommand*{\AsteriskOne}{${}^{*}$}
\newcommand*{\AsteriskTwo}{${}^{**}$}
\newcommand*{\AsteriskThree}{${}^{***}$}

% http://www.ctan.org/pkg/csquotes
\newcommand{\q}[1]{\enquote{#1}}

\newcommand{\var}[1]{\textit{#1}}

\newcommand{\ttf}{\GTT{f()}\xspace}
\newcommand{\ttfone}{\GTT{f1()}\xspace}

% FIXME get rid of
% without dot!
\newcommand{\etc}{%
	\RU{и~т.д}%
	\EN{etc}%
	\ES{etc}%
	\PTBRph{}%
	\PLph{}%
	\ITA{etc}%
	\DEph{}%
	\THAph{}%
	\NLph{}%
	\FR{etc}%
	\JPN{etc}%
}

% http://tex.stackexchange.com/questions/32160/new-line-after-paragraph
\newcommand{\myparagraph}[1]{\paragraph{#1}\mbox{}\\}
\newcommand{\mysubparagraph}[1]{\subparagraph{#1}\mbox{}\\}

\newcommand{\figname}{%
	\RU{илл.}%
	\EN{fig.}%
	\ES{fig.}%
	\PTBRph{}%
	\PLph{}%
	\ITA{fig.}%
	\DE{Abb.}%
	\THAph{}%
	\NLph{}%
	\FR{fig.}%
	\JPN{fig.}\xspace}
\newcommand{\figref}[1]{\figname{}\ref{#1}\xspace}
\newcommand{\listingname}{%
	\RU{листинг.}%
	\EN{listing.}%
	\ES{listado.}%
	\PTBRph{}%
	\PLph{}%
	\ITA{listato.}%
	\DE{Listing.}%
	\THAph{}%
	\NLph{}%
	\FR{liste.}%
	\JPN{リスト}%
\xspace}
\newcommand{\lstref}[1]{\listingname{}\ref{#1}\xspace}
\newcommand{\bitENRU}{%
	\RU{бит}%
	\EN{bit}%
	\ES{bit}%
	\PTBRph{}%
	\PLph{}%
	\ITA{bit}%
	\DE{bit}%
	\THAph{}%
	\NLph{}%
	\FR{bit}%
	\JPN{bit}%
\xspace}
% FIXME get rid of:
\newcommand{\bitsENRU}{%
	\RU{бита}%
	\EN{bits}%
	\ES{bits}%
	\PTBRph{}%
	\PLph{}%
	\ITA{bit}%
	\DE{bits}%
	\THAph{}%
	\NLph{}%
	\FR{bits}%
	\JPN{bits}%
\xspace}
\newcommand{\Sourcecode}{%
	\RU{Исходный код}%
	\EN{Source code}%
	\ES{C\'odigo fuente}%
	\PTBRph{}%
	\PLph{}%
	\ITA{Codice sorgente}%
	\DE{Quellcode}%
	\THAph{}%
	\NLph{}%
	\FR{Code source}%
	\JPN{ソースコード}%
\xspace}
\newcommand{\Seealso}{%
	\RU{См. также}%
	\EN{See also}%
	\ES{V\'ease tambi\'en}%
	\PTBRph{}%
	\PLph{}%
	\ITA{Vedi anche}%
	\DE{siehe auch}%
	\THAph{}%
	\NLph{}%
	\FR{Voir également}%
	\JPN{参照}%
\xspace}
\newcommand{\MacOSX}{Mac OS X\xspace}

% FIXME TODO non-overlapping color!
% \newcommand{\headercolor}{\cellcolor{blue!25}}
\newcommand{\headercolor}{}

\newcommand{\tableheader}{\headercolor{}%
	\RU{смещение}%
	\EN{offset}%
	\ES{offset}%
	\PLph{}%
	\ITA{offset}%
	\DE{Offset}%
	\THAph{}%
	\NLph{}%
	\FR{offset}%
	\JPN{オフセット}%
& \headercolor{}%
	\RU{описание}%
	\EN{description}%
	\ES{descripci\'on}%
	\PTBRph{}%
	\PLph{}%
	\ITA{descrizione}%
	\DE{Beschreibung}%
	\THAph{}%
	\NLph{}%
	\FR{description}%
	\JPN{記述}%
}

% FIXME: get rid of:
\newcommand{\IDA}{\ac{IDA}\xspace}

% FIXME: get rid of:
\newcommand{\tracer}{\protect\gls{tracer}\xspace}

\newcommand{\Tchar}{\IT{char}\xspace}
\newcommand{\Tint}{\IT{int}\xspace}
\newcommand{\Tbool}{\IT{bool}\xspace}
\newcommand{\Tfloat}{\IT{float}\xspace}
\newcommand{\Tdouble}{\IT{double}\xspace}
\newcommand{\Tvoid}{\IT{void}\xspace}
\newcommand{\ITthis}{\IT{this}\xspace}

\newcommand{\Ox}{\GTT{/Ox}\xspace}
\newcommand{\Obzero}{\GTT{/Ob0}\xspace}
\newcommand{\Othree}{\GTT{-O3}\xspace}

\newcommand{\oracle}{Oracle RDBMS\xspace}

\newcommand{\idevices}{iPod/iPhone/iPad\xspace}
\newcommand{\olly}{OllyDbg\xspace}

% common C functions
\newcommand{\printf}{\GTT{printf()}\xspace}
\newcommand{\puts}{\GTT{puts()}\xspace}
\newcommand{\main}{\GTT{main()}\xspace}
\newcommand{\qsort}{\GTT{qsort()}\xspace}
\newcommand{\strlen}{\GTT{strlen()}\xspace}
\newcommand{\scanf}{\GTT{scanf()}\xspace}
\newcommand{\rand}{\GTT{rand()}\xspace}


% for easier fiddling with formatting of all instructions together
\newcommand{\INS}[1]{\GTT{#1}\xspace}

% x86 instructions
\newcommand{\ADD}{\INS{ADD}}
\newcommand{\ADRP}{\INS{ADRP}}
\newcommand{\AND}{\INS{AND}}
\newcommand{\CALL}{\INS{CALL}}
\newcommand{\CPUID}{\INS{CPUID}}
\newcommand{\CMP}{\INS{CMP}}
\newcommand{\DEC}{\INS{DEC}}
\newcommand{\FADDP}{\INS{FADDP}}
\newcommand{\FCOM}{\INS{FCOM}}
\newcommand{\FCOMP}{\INS{FCOMP}}
\newcommand{\FCOMI}{\INS{FCOMI}}
\newcommand{\FCOMIP}{\INS{FCOMIP}}
\newcommand{\FUCOM}{\INS{FUCOM}}
\newcommand{\FUCOMI}{\INS{FUCOMI}}
\newcommand{\FUCOMIP}{\INS{FUCOMIP}}
\newcommand{\FUCOMPP}{\INS{FUCOMPP}}
\newcommand{\FDIVR}{\INS{FDIVR}}
\newcommand{\FDIV}{\INS{FDIV}}
\newcommand{\FLD}{\INS{FLD}}
\newcommand{\FMUL}{\INS{FMUL}}
\newcommand{\MUL}{\INS{MUL}}
\newcommand{\FSTP}{\INS{FSTP}}
\newcommand{\FDIVP}{\INS{FDIVP}}
\newcommand{\IDIV}{\INS{IDIV}}
\newcommand{\IMUL}{\INS{IMUL}}
\newcommand{\INC}{\INS{INC}}
\newcommand{\JAE}{\INS{JAE}}
\newcommand{\JA}{\INS{JA}}
\newcommand{\JBE}{\INS{JBE}}
\newcommand{\JB}{\INS{JB}}
\newcommand{\JE}{\INS{JE}}
\newcommand{\JGE}{\INS{JGE}}
\newcommand{\JG}{\INS{JG}}
\newcommand{\JLE}{\INS{JLE}}
\newcommand{\JL}{\INS{JL}}
\newcommand{\JMP}{\INS{JMP}}
\newcommand{\JNE}{\INS{JNE}}
\newcommand{\JNZ}{\INS{JNZ}}
\newcommand{\JNA}{\INS{JNA}}
\newcommand{\JNAE}{\INS{JNAE}}
\newcommand{\JNB}{\INS{JNB}}
\newcommand{\JNBE}{\INS{JNBE}}
\newcommand{\JZ}{\INS{JZ}}
\newcommand{\JP}{\INS{JP}}
\newcommand{\Jcc}{\INS{Jcc}}
\newcommand{\SETcc}{\INS{SETcc}}
\newcommand{\LEA}{\INS{LEA}}
\newcommand{\LOOP}{\INS{LOOP}}
\newcommand{\MOVSX}{\INS{MOVSX}}
\newcommand{\MOVZX}{\INS{MOVZX}}
\newcommand{\MOV}{\INS{MOV}}
\newcommand{\NOP}{\INS{NOP}}
\newcommand{\POP}{\INS{POP}}
\newcommand{\PUSH}{\INS{PUSH}}
\newcommand{\NOT}{\INS{NOT}}
\newcommand{\NOR}{\INS{NOR}}
\newcommand{\RET}{\INS{RET}}
\newcommand{\RETN}{\INS{RETN}}
\newcommand{\SETNZ}{\INS{SETNZ}}
\newcommand{\SETBE}{\INS{SETBE}}
\newcommand{\SETNBE}{\INS{SETNBE}}
\newcommand{\SUB}{\INS{SUB}}
\newcommand{\TEST}{\INS{TEST}}
\newcommand{\TST}{\INS{TST}}
\newcommand{\FNSTSW}{\INS{FNSTSW}}
\newcommand{\SAHF}{\INS{SAHF}}
\newcommand{\XOR}{\INS{XOR}}
\newcommand{\OR}{\INS{OR}}
\newcommand{\SHL}{\INS{SHL}}
\newcommand{\SHR}{\INS{SHR}}
\newcommand{\SAR}{\INS{SAR}}
\newcommand{\LEAVE}{\INS{LEAVE}}
\newcommand{\MOVDQA}{\INS{MOVDQA}}
\newcommand{\MOVDQU}{\INS{MOVDQU}}
\newcommand{\PADDD}{\INS{PADDD}}
\newcommand{\PCMPEQB}{\INS{PCMPEQB}}
\newcommand{\LDR}{\INS{LDR}}
\newcommand{\LSL}{\INS{LSL}}
\newcommand{\LSR}{\INS{LSR}}
\newcommand{\ASR}{\INS{ASR}}
\newcommand{\RSB}{\INS{RSB}}
\newcommand{\BTR}{\INS{BTR}}
\newcommand{\BTS}{\INS{BTS}}
\newcommand{\BTC}{\INS{BTC}}
\newcommand{\LUI}{\INS{LUI}}
\newcommand{\ORI}{\INS{ORI}}
\newcommand{\BIC}{\INS{BIC}}
\newcommand{\EOR}{\INS{EOR}}
\newcommand{\MOVS}{\INS{MOVS}}
\newcommand{\LSLS}{\INS{LSLS}}
\newcommand{\LSRS}{\INS{LSRS}}
\newcommand{\FMRS}{\INS{FMRS}}
\newcommand{\CMOVNE}{\INS{CMOVNE}}
\newcommand{\CMOVNZ}{\INS{CMOVNZ}}
\newcommand{\ROL}{\INS{ROL}}
\newcommand{\CSEL}{\INS{CSEL}}
\newcommand{\SLL}{\INS{SLL}}
\newcommand{\SLLV}{\INS{SLLV}}
\newcommand{\SW}{\INS{SW}}
\newcommand{\LW}{\INS{LW}}

% x86 flags

\newcommand{\ZF}{\GTT{ZF}\xspace}
\newcommand{\CF}{\GTT{CF}\xspace}
\newcommand{\PF}{\GTT{PF}\xspace}

% x86 registers

\newcommand{\AL}{\GTT{AL}\xspace}
\newcommand{\AH}{\GTT{AH}\xspace}
\newcommand{\AX}{\GTT{AX}\xspace}
\newcommand{\EAX}{\GTT{EAX}\xspace}
\newcommand{\EBX}{\GTT{EBX}\xspace}
\newcommand{\ECX}{\GTT{ECX}\xspace}
\newcommand{\EDX}{\GTT{EDX}\xspace}
\newcommand{\DL}{\GTT{DL}\xspace}
\newcommand{\ESI}{\GTT{ESI}\xspace}
\newcommand{\EDI}{\GTT{EDI}\xspace}
\newcommand{\EBP}{\GTT{EBP}\xspace}
\newcommand{\ESP}{\GTT{ESP}\xspace}
\newcommand{\RSP}{\GTT{RSP}\xspace}
\newcommand{\EIP}{\GTT{EIP}\xspace}
\newcommand{\RIP}{\GTT{RIP}\xspace}
\newcommand{\RAX}{\GTT{RAX}\xspace}
\newcommand{\RBX}{\GTT{RBX}\xspace}
\newcommand{\RCX}{\GTT{RCX}\xspace}
\newcommand{\RDX}{\GTT{RDX}\xspace}
\newcommand{\RBP}{\GTT{RBP}\xspace}
\newcommand{\RSI}{\GTT{RSI}\xspace}
\newcommand{\RDI}{\GTT{RDI}\xspace}
\newcommand*{\ST}[1]{\GTT{ST(#1)}\xspace}
\newcommand*{\XMM}[1]{\GTT{XMM#1}\xspace}

% ARM
\newcommand*{\Reg}[1]{\GTT{R#1}\xspace}
\newcommand*{\RegX}[1]{\GTT{X#1}\xspace}
\newcommand*{\RegW}[1]{\GTT{W#1}\xspace}
\newcommand*{\RegD}[1]{\GTT{D#1}\xspace}
\newcommand{\ADREQ}{\GTT{ADREQ}\xspace}
\newcommand{\ADRNE}{\GTT{ADRNE}\xspace}
\newcommand{\BEQ}{\GTT{BEQ}\xspace}

% instructions descriptions
\newcommand{\ASRdesc}{%
	\RU{арифметический сдвиг вправо}%
	\EN{arithmetic shift right}%
	\ES{desplazamiento aritm\'etico a la derecha}%
	\PTBRph{}%
	\PLph{}%
	\ITA{scorrimento aritmetico a destra}%
	\DEph{}%
	\THAph{}\NLph{}%
	\FR{décalage arithmétique vers la droite}%
	\JPNph{}%
}

% x86 registers tables
\newcommand{\RegHeaderTop}{
	\multicolumn{8}{ | c | }{
	\RU{Номер байта:}%
	\EN{Byte number:}%
	\ES{}%
	\PTBRph{}%
	\PLph{}%
	\ITA{Numero byte}%
	\DE{Byte-Nummer:}%
	\THAph{}%
	\NLph{}%
	\FR{Octet d'indice}%
	\JPN{バイトの並び順}%
	}
}

% TODO: non-overlapping color!
\newcommand{\RegHeader}{
\RU{ 7-й & 6-й & 5-й & 4-й & 3-й & 2-й & 1-й & 0-й }%
\EN{ 7th & 6th & 5th & 4th & 3rd & 2nd & 1st & 0th }%
\ES{ 7mo & 6to & 5to & 4to & 3ro & 2do & 1ro & 0 }%
\PTBRph{}%
\PLph{}%
\ITA{ 7°  & 6°  & 5°  & 4°  & 3°  & 2°  & 1°  & 0 }%
\DEph{}%
\THAph{}%
\NLph{}%
\FR{7 & 6 & 5 & 4 & 3 & 2 & 1 & 0}%
\JPN{ 第7 & 第6 & 第5 & 第4 & 第3 & 第2 & 第1 & 第0 }%
}

% FIXME навести порядок тут...
\newcommand{\RegTableThree}[5]{
\begin{center}
\begin{tabular}{ | l | l | l | l | l | l | l | l | l |}
\hline
\RegHeaderTop \\
\hline
\RegHeader \\
\hline
\multicolumn{8}{ | c | }{#1} \\
\hline
\multicolumn{4}{ | c | }{} & \multicolumn{4}{ c | }{#2} \\
\hline
\multicolumn{6}{ | c | }{} & \multicolumn{2}{ c | }{#3} \\
\hline
\multicolumn{6}{ | c | }{} & #4 & #5 \\
\hline
\end{tabular}
\end{center}
}

\newcommand{\RegTableOne}[5]{\RegTableThree{#1\textsuperscript{x64}}{#2}{#3}{#4}{#5}}

\newcommand{\RegTableTwo}[4]{
\begin{center}
\begin{tabular}{ | l | l | l | l | l | l | l | l | l |}
\hline
\RegHeaderTop \\
\hline
\RegHeader \\
\hline
\multicolumn{8}{ | c | }{#1\textsuperscript{x64}} \\
\hline
\multicolumn{4}{ | c | }{} & \multicolumn{4}{ c | }{#2} \\
\hline
\multicolumn{6}{ | c | }{} & \multicolumn{2}{ c | }{#3} \\
\hline
\multicolumn{7}{ | c | }{} & #4\textsuperscript{x64} \\
\hline
\end{tabular}
\end{center}
}

\newcommand{\RegTableFour}[4]{
\begin{center}
\begin{tabular}{ | l | l | l | l | l | l | l | l | l |}
\hline
\RegHeaderTop \\
\hline
\RegHeader \\
\hline
\multicolumn{8}{ | c | }{#1} \\
\hline
\multicolumn{4}{ | c | }{} & \multicolumn{4}{ c | }{#2} \\
\hline
\multicolumn{6}{ | c | }{} & \multicolumn{2}{ c | }{#3} \\
\hline
\multicolumn{7}{ | c | }{} & #4 \\
\hline
\end{tabular}
\end{center}
}

\newcommand{\ReturnAddress}{%
	\RU{Адрес возврата}%
	\EN{Return Address}%
	\ES{Direcci\'on de Retorno}%
	\NL{Return Adres}%
	\PTBR{Endereço de retorno}%
	\DE{Rücksprungadresse}%
	\ITA{Indirizzo di Ritorno}%
	\PLph{}%
	\FR{Adresse de retour}%
	\JPN{リターンアドレス}%
}

\newcommand{\localVariable}{%
	\RU{локальная переменная}%
	\EN{local variable}%
	\PTBR{Variável local}%
	\FR{variable locale}%
	\DE{lokale Variable}%
	\ITA{variabile locale}%
	\JPN{ローカル変数}%
}
\newcommand{\savedValueOf}{%
	\RU{сохраненное значение}%
	\EN{saved value of}%
	\PTBR{Valor salvo de}%
	\FR{valeur enregistrée de}%
	\ITA{valore memorizzato di}%
	\JPN{saved value of}%
}

% was in common_phrases.tex
% for index
\newcommand{\GrepUsage}{\RU{Использование grep}\EN{grep usage}\PTBR{Uso do grep}\ES{Uso de grep}\FR{Utilisation de grep}\JPN{grep usage}\ITA{uso di grep}}
\newcommand{\SyntacticSugar}{\RU{Синтаксический сахар}\EN{Syntactic Sugar}\PTBR{Açúcar sintático}\ES{Azúcar sintáctica}\FR{Sucre syntaxique}\JPN{糖衣構文}\ITA{Syntactic Sugar}}
\newcommand{\CompilerAnomaly}{\RU{Аномалии компиляторов}\EN{Compiler's anomalies}\PTBR{Anomalias do compilador}\ES{Anomalías del compilador}\FR{Anomalies du compilateur}\ITA{Anomalie del compilatore}\JPN{コンパイラアノマリ}}
\newcommand{\CLanguageElements}{\RU{Элементы языка Си}\EN{C language elements}\DE{C Sprachelemente}\PTBR{Elementos da linguagem C}\ES{Elementos del lenguaje C}\FR{Eléments du langage C}\ITA{Elementi del linguaggio C}\JPN{C言語の要素}}
\newcommand{\CStandardLibrary}{\RU{Стандартная библиотека Си}\EN{C standard library}\PTBR{Biblioteca padrão C}\ES{Librería estándar C}\FR{Bibliothèque standard C}\ITA{Libreria C standard}\JPN{C標準ライブラリ}}
\newcommand{\Instructions}{\RU{Инструкции}\EN{Instructions}\PTBR{Instruções}\ES{Instrucciones}\FR{Instructions}\ITA{Istruzioni}\JPN{命令}}
\newcommand{\Pseudoinstructions}{\RU{Псевдоинструкции}\EN{Pseudoinstructions}\PTBR{Pseudo-instruções}\ES{Pseudo-instrucciones}\FR{Pseudo-instructions}\ITA{Pseudo-istruzioni}\JPN{疑似命令}}
\newcommand{\Prefixes}{\RU{Префиксы}\EN{Prefixes}\PTBR{Prefixos}\ES{Prefijos}\FR{Préfixes}\ITA{Prefissi}\JPN{プリフィックス}}

\newcommand{\Flags}{\RU{Флаги}\EN{Flags}\PTBR{Flags}\ES{Flags}\FR{Flags}\ITA{Flag}\JPN{フラグ}}
\newcommand{\Registers}{\RU{Регистры}\EN{Registers}\PTBR{Registradores}\ES{Registros}\FR{Registres}\ITA{Registri}\JPN{レジスタ}}
\newcommand{\registers}{\RU{регистры}\EN{registers}\PTBR{registradores}\ES{registros}\FR{registres}\ITA{registri}\JPN{レジスタ}}
\newcommand{\Stack}{\RU{Стек}\EN{Stack}\DE{Stack}\PTBR{Pilha}\ES{Pila}\FR{Pile}\ITA{Stack}\JPN{スタック}}
\newcommand{\Recursion}{\RU{Рекурсия}\EN{Recursion}\PTBR{Recursividade}\ES{Recursión}\FR{Récursivité}\ITA{Recursione}\JPN{再帰}}
\newcommand{\RAM}{\RU{ОЗУ}\EN{RAM}\PTBR{RAM}\ES{RAM}\FR{RAM}\ITA{RAM}\JPN{RAM}}
\newcommand{\ROM}{\RU{ПЗУ}\EN{ROM}\PTBR{ROM}\ES{ROM}\FR{ROM}\ITA{ROM}\JPN{ROM}}
\newcommand{\Pointers}{\RU{Указатели}\EN{Pointers}\PTBR{Ponteiros}\ES{Apuntadores}\FR{Pointeurs}\ITA{Puntatori}\JPN{ポインタ}}
\newcommand{\BufferOverflow}{\RU{Переполнение буфера}\EN{Buffer Overflow}\PTBR{Buffer Overflow}\ES{Desbordamiento de buffer}\FR{Débordement de tampon}\ITA{Buffer Overflow}\JPN{バッファオーバーフロー}}

% DE: also "Zusammenfassung"
\newcommand{\Conclusion}{\RU{Вывод}\EN{Conclusion}\DE{Fazit}\PTBR{Conclusão}\ES{Conclusión}\FR{Conclusion}\ITA{Conclusione}\JPN{結論}}

\newcommand{\Exercise}{\RU{Упражнение}\EN{Exercise}\PTBR{Exercício}\ES{Ejercicio}\FR{Exercice}\ITA{Esercizio}\JPN{練習問題}\xspace}
\newcommand{\Exercises}{\RU{Упражнения}\EN{Exercises}\PTBR{Exercícios}\ES{Ejercicios}\FR{Exercices}\ITA{Esercizi}\DE{Übungen}\JPN{練習問題}\xspace}
\newcommand{\Arrays}{\RU{Массивы}\EN{Arrays}\PTBR{Matriz}\ES{Matriz}\FR{Tableaux}\ITA{Array}\JPN{配列}\DE{Arrays}}
\newcommand{\Cpp}{\RU{Си++}\EN{C++}\PTBR{C++}\ES{C++}\FR{C++}\ITA{C++}\DE{C++}\JPN{C++}\xspace}
\newcommand{\CCpp}{\RU{Си/Си++}\EN{C/C++}\PTBR{C/C++}\ES{C/C++}\FR{C/C++}\ITA{C/C++}\DE{C/C++}\JPN{C/C++}\xspace}
\newcommand{\NonOptimizing}{\RU{Неоптимизирующий}\EN{Non-optimizing}\PTBR{Sem otimização}\ES{Sin optimización}\FR{sans optimisation}\ITA{Senza ottimizzazione}\JPN{非最適化}\xspace}
\newcommand{\Optimizing}{\RU{Оптимизирующий}\EN{Optimizing}\PTBR{Com otimização}\ES{Con optimización}\FR{avec optimisation}\ITA{Con ottimizzazione}\JPN{最適化}\xspace}
\newcommand{\NonOptimizingKeilVI}{\NonOptimizing Keil 6/2013\xspace}
\newcommand{\OptimizingKeilVI}{\Optimizing Keil 6/2013\xspace}
\newcommand{\NonOptimizingXcodeIV}{\NonOptimizing Xcode 4.6.3 (LLVM)\xspace}
\newcommand{\OptimizingXcodeIV}{\Optimizing Xcode 4.6.3 (LLVM)\xspace}
\newcommand{\ARMMode}{\RU{Режим ARM}\EN{ARM mode}\PTBR{Modo ARM}\ES{Modo ARM}\FR{Mode ARM}\ITA{Modalità ARM}\JPN{ARMモード}\xspace}
\newcommand{\ThumbMode}{\RU{Режим Thumb}\EN{Thumb mode}\PTBR{Modo Thumb}\ES{Modo Thumb}\FR{Mode Thumb}\ITA{Modalità Thumb}\JPN{Thumbモード}\xspace}
\newcommand{\ThumbTwoMode}{\RU{Режим Thumb-2}\EN{Thumb-2 mode}\PTBR{Modo Thumb-2}\ES{Modo Thumb-2}\FR{Mode Thumb-2}\ITA{Modalità Thumb-2}\JPN{Thumb-2モード}\xspace}
\newcommand{\AndENRU}{\RU{и}\EN{and}\PTBR{e}\ES{y}\FR{et}\JPN{and}\xspace}
\newcommand{\OrENRU}{\RU{или}\EN{or}\PTBR{ou}\ES{o}\FR{ou}\JPN{or}\xspace}
\newcommand{\InENRU}{\RU{в}\EN{in}\PTBR{em}\ES{en}\FR{dans}\JPN{in}\xspace}
\newcommand{\ForENRU}{\RU{для}\EN{for}\PTBR{para}\ES{para}\FR{pour}\JPN{for}\xspace}
\newcommand{\LineENRU}{\RU{строка}\EN{line}\PTBR{linha}\ES{línea}\FR{ligne}\JPN{line}\xspace}

\newcommand{\DataProcessingInstructionsFootNote}{%
	\RU{Эти инструкции также называются \q{data processing instructions}}%
	\EN{These instructions are also called \q{data processing instructions}}%
	\PTBR{Estas intruções também são chamadas \q{data processing instructions}}%
	\ES{Estas instrucciones también son llamadas \q{instrucciones de procesamiento de datos}}%
	\FR{Ces instructions sont également appelées \q{instructions de traitement de données}}%
  \EN{Queste istruzioni sono anche chiamate \q{istruzioni di processamento dei dati}}%
	\JPN{この命令は \q{データプロセス命令} とも呼ばれます}%
}

% for .bib files
\newcommand{\AlsoAvailableAs}{\RU{Также доступно здесь:}\EN{Also available as}\PTBR{Também disponível como}\ES{También disponible como}\FR{Aussi disponible en}\EN{Disponibile anche come}\JPN{以下で利用可能}\xspace}

% section names
\newcommand{\ShiftsSectionName}{\RU{Сдвиги}\EN{Shifts}\PTBR{Shifts}\ES{Desplazamientos}\FR{Décalages}\JPN{シフト}}
\newcommand{\SignedNumbersSectionName}{\RU{Представление знака в числах}\EN{Signed number representations}\ES{Representaci\'on de n\'umeros con signo}\FR{Représentations des nombres signés}\ITA{Rappresentazione di numeri con segno}\DE{Darstellung vorzeichenbehafteter Zahlen}\JPN{符号付き整数表現}}
\newcommand{\HelloWorldSectionName}{\RU{Hello, world!}\EN{Hello, world!}\ES{!`Hola, mundo!}\FR{Hello, world!}\ITA{Hello, world!}\DE{Hallo, Welt!}\JPN{ハローワールド!}}
\newcommand{\SwitchCaseDefaultSectionName}{switch()/case/default}
\newcommand{\PrintfSeveralArgumentsSectionName}{printf() \RU{с несколькими аргументами}\EN{with several arguments}\PTBR{com vários argumentos}\ES{con varios argumentos}\ITA{con più argomenti}\FR{avec plusieurs arguments}\JPN{引数を取って}}
\newcommand{\BitfieldsChapter}{\RU{Работа с отдельными битами}\EN{Manipulating specific bit(s)}\PTBR{Manipulando bit(s) específicos}\ES{Manipulando bit(s) específicos}\FR{Manipulation de bits spécifiques}\ITA{Manipolando dei bit speifici}\JPN{特定のビットを操作する}\DE{Manipulieren einzelner Bits}}
\newcommand{\ArithOptimizations}{%
	\RU{Замена одних арифметических инструкций на другие}%
	\EN{Replacing arithmetic instructions to other ones}%
	\PTBR{Substituição de instruções aritiméticas por outras}%
	\ES{Substituición de instrucciones aritméticas por otras}%
	\FR{Remplacement de certaines instructions arithmétiques par d'autres}%
  \ITA{Sostituzione di istruzioni aritmetiche con altre}%
	\JPN{算術命令を他の命令に置換する}%
	\DE{Ersetzen von arithmetischen Operationen}%
	}
\newcommand{\FPUChapterName}{\RU{Работа с FPU}\EN{Floating-point unit}\PTBR{Unidade de Ponto flutuante}\ES{Unidad de punto flotante}\FR{Unité à virgule flottante}\ITA{Floating-point unit}\JPN{フローティングポイントユニット}\DE{Gleitkommaeinheit}}
\newcommand{\MoreAboutStrings}{\RU{Еще кое-что о строках}\EN{More about strings}\PTBRph{}\ESph{}\FR{Plus d'information sur les chaînes}\ITA{Maggiori informazioni sulle stringhe}\JPN{文字列に関する加筆}\DE{Mehr über Zeichenketten}}
\newcommand{\DivisionByMultSectionName}{\RU{Деление используя умножение}\EN{Division using multiplication}\PTBR{Divisão por 9}\ES{División entre 9}\FR{Division par la multiplication}\ITA{Divisione utilizzando la moltiplicazione}\JPN{乗法を使用した除算}}
\newcommand{\Answer}{\RU{Ответ}\EN{Answer}\PTBR{Responda}\ES{Respuesta}\FR{Réponse}\ITA{Risposta}\JPN{解答}}
\newcommand{\WhatThisCodeDoes}{\RU{Что делает этот код}\EN{What does this code do}\PTBR{O que este código faz}\ES{?`Qu\'e hace este código}?\FR{Que fait ce code ?}\ITA{Cosa fa qusto codice}\JPN{このコードは何をしている}}
\newcommand{\WorkingWithFloatAsWithStructSubSubSectionName}{%
\RU{Работа с типом float как со структурой}\EN{Handling float data type as a structure}\PTBR{Trabalhando com o tipo float como uma estrutura}\ES{Trabajando con el tipo float como una estructura}\FR{Travailler avec le type float comme une structure}\ITA{Utilizzare i tipi di dato float come una struttura}\JPN{フロート型のデータを構造体として扱う}}

\newcommand{\MinesweeperWinXPExampleChapterName}{\RU{Сапёр}\EN{Minesweeper}\PTBR{Campo minado}\ES{Buscaminas}\FR{Démineur}\JPN{マインスイーパー} (Windows XP)}

\newcommand{\StructurePackingSectionName}{\RU{Упаковка полей в структуре}\EN{Fields packing in structure}\PTBR{Organização de campos na estrutura}\ES{Organización de campos en la estructura}\FR{Organisation des champs dans la structure}\ITA{Organizzazione dei campi in una struttura}\JPN{フィールドを構造体にパッキングする}}
\newcommand{\StructuresChapterName}{\RU{Структуры}\EN{Structures}\PTBR{Estruturas}\ES{Estructuras}\FR{Structures}\ITA{Strutture}\JPN{構造体}\DE{Strukturen}}
\newcommand{\PICcode}{\RU{адресно-независимый код}\EN{position-independent code}\PTBR{código independente de posição}\ES{código independiente de la posición}\FR{code indépendant de la position}\ITA{codice indipendente dalla posizione}\DE{positionsabhängiger Code}\JPN{位置独立コード}}
\newcommand{\CapitalPICcode}{\RU{Адресно-независимый код}\EN{Position-independent code}\PTBR{Código independente de posição}\ES{Código independiente de lá posición}\FR{Code indépendant de la position}\ITA{Codice indipendente dalla posizione}\DE{Positionsabhängiger Code}\JPN{位置独立コード}}
\newcommand{\Loops}{\RU{Циклы}\EN{Loops}\DE{Schleifen}\PTBR{Laços}\ES{Bucles}\FR{Boucles}\ITA{Cicli}\JPN{ループ}}

% C
\newcommand{\PostIncrement}{\RU{Пост-инкремент}\EN{Post-increment}\PTBR{Pós-incremento}\ES{Post-incremento}\FR{Post-incrémentation}\ITA{Post-incremento}\JPN{後置インクリメント}}
\newcommand{\PostDecrement}{\RU{Пост-декремент}\EN{Post-decrement}\PTBR{Pós-decremento}\ES{Post-decremento}\FR{Post-décrémentation}\ITA{Post-decremento}\JPN{後置デクリメント}}
\newcommand{\PreIncrement}{\RU{Пре-инкремент}\EN{Pre-increment}\PTBR{Pré-incremento}\ES{Pre-incremento}\FR{Pré-incrémentation}\ITA{Pre-incremento}\JPN{前置インクリメント}}
\newcommand{\PreDecrement}{\RU{Пре-декремент}\EN{Pre-decrement}\PTBR{Pré-decremento}\ES{Pre-decremento}\FR{Pré-décrémentation}\ITA{Pre-decremento}\JPN{前置デクリメント}}

% MIPS
\newcommand{\GlobalPointer}{\RU{Глобальный указатель}\EN{Global Pointer}\PTBR{Ponteiro Global}\ES{Apuntador Global}\FR{Pointeur Global}\ITA{Puntatore Globale}\JPN{グローバルポインタ}}

\newcommand{\garbage}{\RU{мусор}\EN{garbage}\PTBR{Lixo}\ES{Basura}\FR{déchets}\ITA{garbage}\JPN{ガーベッジ}}
\newcommand{\IntelSyntax}{\RU{Синтаксис Intel}\EN{Intel syntax}\PTBR{Sintaxe Intel}\ES{Sintaxis Intel}\FR{Syntaxe Intel}\ITA{Sintassi Intel}\JPN{インテル構文}}
\newcommand{\ATTSyntax}{\RU{Синтаксис AT\&T}\EN{AT\&T syntax}\PTBR{Sintaxe AT\&T}\ES{Sintaxis AT\&T}\FR{Syntaxe AT\&T}\ITA{Sintassi AT\&T}\JPN{AT\&T構文}}
\newcommand{\randomNoise}{\RU{случайный шум}\EN{random noise}\PTBR{Ruído aleatório}\ES{Ruido aleatorio}\FR{bruit aléatoire}\ITA{rumore casuale}\JPN{ランダムノイズ}}
\newcommand{\Example}{\RU{Пример}\EN{Example}\PTBR{Exemplo}\ES{Ejemplo}\FR{Exemple}\ITA{Esempio}\JPN{例}}
\newcommand{\argument}{\RU{аргумент}\EN{argument}\PTBR{argumento}\ES{argumento}\FR{argument}\ITA{argomento}\JPN{引数}}
\newcommand{\MarkedInIDAAs}{\RU{маркируется в \IDA как}\EN{marked in \IDA as}\PTBR{Marcado no \IDA como}\ES{Marcado en \IDA como}\FR{marqué dans \IDA comme}\ITA{marcato in IDA come}\JPN{\IDA にマークする}}
\newcommand{\stepover}{\RU{сделать шаг, не входя в функцию}\EN{step over}\PTBR{passar por cima}\ES{pasar por encima}\FR{enjamber}\ITA{step over}\JPN{ステップオーバー}}
\newcommand{\ShortHotKeyCheatsheet}{\RU{Краткий справочник горячих клавиш}\EN{Hot-keys cheatsheet}\PTBR{Cheatsheet de teclas de atalho}\ES{Cheatsheet de teclas de acceso rápido}\FR{Anti-sèche des touches de raccourci}\ITA{Elenco delle scorciatoie da tastiera}\JPN{ホットキーチートシート}}

\newcommand{\assemblyOutput}{\RU{вывод на ассемблере}\EN{assembly output}\PTBR{saída do assembly}\ES{salida del ensamblador}\FR{résultat en sortie de l'assembleur}\ITA{risultato dell'assembly}\JPN{アセンブリ出力}}

% was in common_URLS.tex:
\newcommand{\URLWPDA}{%
	\RU{\href{http://go.yurichev.com/17012}{Wikipedia: Выравнивание данных}}%
	\EN{\href{http://go.yurichev.com/17013}{Wikipedia: Data structure alignment}}%
	\ES{\href{URL}{\ESph{}}}%
	\PL{\href{URL}{\PLph{}}}%
	\PTBR{\href{URL}{\PTBRph{}}}%
	\ITA{\href{https://en.wikipedia.org/wiki/Data_structure_alignment}{Wikipedia: Data structure alignment}}%
	\DE{\href{URL}{\DEph{}}}%
	\FR{\href{https://fr.wikipedia.org/wiki/Alignement_en_m\%C3\%A9moire}{Wikipedia: Alignement en mémoire}}%
	\JPN{\href{http://go.yurichev.com/17013}{Wikipedia: Data structure alignment}}%
}

\newcommand{\OracleTablesName}{oracle tables\xspace}
\newcommand{\oracletables}{\OracleTablesName\footnote{\href{http://go.yurichev.com/17014}{yurichev.com}}\xspace}

\newcommand{\BGREPURL}{\href{http://go.yurichev.com/17017}{GitHub}}
\newcommand{\FNMSDNROTxURL}{\footnote{\href{http://go.yurichev.com/17018}{MSDN}}}

\newcommand{\YurichevIDAIDCScripts}{http://go.yurichev.com/17019}

% ML prefix is for multi-lingual words and sentences:
\newcommand{\MLHeap}{\EN{Heap}\RU{Куча}\PTBR{Heap}\FR{Heap}\ITA{Heap}\JPN{ヒープ}}
\newcommand{\MLStack}{\EN{Stack}\RU{Стэк}\PTBR{Pilha}\FR{Pile}\ITA{Stack}\JPN{スタック}}
\newcommand{\MLStackOverflow}{\EN{Stack overflow}\RU{Переполнение стека}\PTBR{\PTBRph{}}\FR{Débordement de pile}\DEph{}\ITA{Stack overflow}\JPN{スタックオーバーフロー}\PTBR{\PTBRph{}}}
\newcommand{\MLStartOfHeap}{\RU{Начало кучи}\EN{Start of heap}\PTBR{começo da heap}\FR{Début du heap}\ITA{Inizio dell'heap}\JPN{ヒープの開始}}
\newcommand{\MLStartOfStack}{\RU{Вершина стека}\EN{Start of stack}\PTBR{começo da pilha}\FR{Début de la pile}\ITA{Inizio dello stack}\JPN{スタックの開始}}
\newcommand{\MLinputA}{\RU{вход А}\EN{input A}\FR{entrée A}\ITA{input A}\JPN{入力A}}
\newcommand{\MLinputB}{\RU{вход Б}\EN{input B}\FR{entrée B}\ITA{input A}\JPN{入力B}}
\newcommand{\MLoutput}{\RU{выход}\EN{output}\FR{sortie}\ITA{output}\JPN{出力}}

% sources: books, etc
\newcommand{\TAOCPvolI}{Donald E. Knuth, \IT{The Art of Computer Programming}, Volume 1, 3rd ed., (1997)}
\newcommand{\TAOCPvolII}{Donald E. Knuth, \IT{The Art of Computer Programming}, Volume 2, 3rd ed., (1997)}
\newcommand{\Russinovich}{Mark Russinovich, \IT{Microsoft Windows Internals}}
\newcommand{\Schneier}{Bruce Schneier, \IT{Applied Cryptography}, (John Wiley \& Sons, 1994)}
\newcommand{\AgnerFog}{Agner Fog, \IT{The microarchitecture of Intel, AMD and VIA CPUs}, (2016)}
\newcommand{\AgnerFogCPP}{Agner Fog, \IT{Optimizing software in C++} (2015)}
\newcommand{\AgnerFogCC}{Agner Fog, \IT{Calling conventions} (2015)}
\newcommand{\JavaBook}{[Tim Lindholm, Frank Yellin, Gilad Bracha, Alex Buckley, \IT{The Java(R) Virtual Machine Specification / Java SE 7 Edition}]
\footnote{\AlsoAvailableAs \url{https://docs.oracle.com/javase/specs/jvms/se7/jvms7.pdf}; \url{http://docs.oracle.com/javase/specs/jvms/se7/html/}}}

\newcommand{\ARMPCS}{\InSqBrackets{\IT{Procedure Call Standard for the ARM 64-bit Architecture (AArch64)}, (2013)}\footnote{\AlsoAvailableAs \url{http://go.yurichev.com/17287}}}

\newcommand{\IgorSkochinsky}{[Igor Skochinsky, \IT{Compiler Internals: Exceptions and RTTI}, (2012)] \footnote{\AlsoAvailableAs \url{http://go.yurichev.com/17294}}}

\newcommand{\PietrekSEH}{[Matt Pietrek, \IT{A Crash Course on the Depths of Win32\texttrademark{} Structured Exception Handling}, (1997)]\footnote{\AlsoAvailableAs \url{http://go.yurichev.com/17293}}}

\newcommand{\PietrekPE}{Matt Pietrek, \IT{An In-Depth Look into the Win32 Portable Executable File Format}, (2002)]}

\newcommand{\PietrekPEURL}{\PietrekPE\footnote{\AlsoAvailableAs \url{http://go.yurichev.com/17318}}}

\newcommand{\RitchieDevC}{[Dennis M. Ritchie, \IT{The development of the C language}, (1993)]\footnote{\AlsoAvailableAs \url{http://go.yurichev.com/17264}}}

\newcommand{\RitchieThompsonUNIX}{[D. M. Ritchie and K. Thompson, \IT{The UNIX Time Sharing System}, (1974)]\footnote{\AlsoAvailableAs \url{http://go.yurichev.com/17270}}}

\newcommand{\DrepperTLS}{[Ulrich Drepper, \IT{ELF Handling For Thread-Local Storage}, (2013)]\footnote{\AlsoAvailableAs \url{http://go.yurichev.com/17272}}}

\newcommand{\DrepperMemory}{[Ulrich Drepper, \IT{What Every Programmer Should Know About Memory}, (2007)]\footnote{\AlsoAvailableAs \url{http://go.yurichev.com/17341}}}

\newcommand{\AlephOne}{[Aleph One, \IT{Smashing The Stack For Fun And Profit}, (1996)]\footnote{\AlsoAvailableAs \url{http://go.yurichev.com/17266}}}

\ifdefined\RUSSIAN
\newcommand{\KRBook}{Брайан Керниган, Деннис Ритчи, \IT{Язык программирования Си}, второе издание, (1988, 2009)}
\newcommand{\CppOneOneStd}{Стандарт Си++11}
\newcommand{\CNotes}{Денис Юричев, \IT{Заметки о языке программирования Си/Си++}}
\newcommand{\TAOCP}{Дональд Э. Кнут, \IT{Искусство программирования}}
\else
\newcommand{\KRBook}{Brian W. Kernighan, Dennis M. Ritchie, \IT{The C Programming Language}, 2ed, (1988)}
\newcommand{\CppOneOneStd}{C++11 standard}
\newcommand{\CNotes}{Dennis Yurichev, \IT{C/C++ programming language notes}}
\newcommand{\TAOCP}{Donald E. Knuth, \IT{The Art of Computer Programming}}
\fi

\newcommand{\RobPikePractice}{Brian W. Kernighan, Rob Pike, \IT{Practice of Programming}, (1999)}

\newcommand{\ARMSixFourRef}{\IT{ARM Architecture Reference Manual, ARMv8, for ARMv8-A architecture profile}, (2013)}
\newcommand{\ARMSixFourRefURL}{\InSqBrackets{\ARMSixFourRef}\footnote{\AlsoAvailableAs \url{http://yurichev.com/mirrors/ARMv8-A_Architecture_Reference_Manual_(Issue_A.a).pdf}}}

\newcommand{\SysVABI}{[Michael Matz, Jan Hubicka, Andreas Jaeger, Mark Mitchell, \IT{System V Application Binary Interface. AMD64 Architecture Processor Supplement}, (2013)]
\footnote{\AlsoAvailableAs \url{https://software.intel.com/sites/default/files/article/402129/mpx-linux64-abi.pdf}}}

\newcommand{\IOSABI}{\InSqBrackets{\IT{iOS ABI Function Call Guide}, (2010)}\footnote{\AlsoAvailableAs \url{http://go.yurichev.com/17276}}}

\newcommand{\CNineNineStd}{\IT{ISO/IEC 9899:TC3 (C C99 standard)}, (2007)}

\newcommand{\TCPPPL}{Bjarne Stroustrup, \IT{The C++ Programming Language, 4th Edition}, (2013)}

\newcommand{\AMDOptimization}{\InSqBrackets{\IT{Software Optimization Guide for AMD Family 16h Processors}, (2013)}}
% note = "\AlsoAvailableAs \url{http://go.yurichev.com/17285}",

\newcommand{\IntelOptimization}{\InSqBrackets{\IT{Intel® 64 and IA-32 Architectures Optimization Reference Manual}, (2014)}}
% note = "\AlsoAvailableAs \url{http://go.yurichev.com/17342}",

\newcommand{\TAOUP}{Eric S. Raymond, \IT{The Art of UNIX Programming}, (2003)}
% note = "\AlsoAvailableAs \url{http://go.yurichev.com/17277}",

\newcommand{\HenryWarren}{Henry S. Warren, \IT{Hacker's Delight}, (2002)}

\newcommand{\ParashiftCPPFAQ}{Marshall Cline, \IT{C++ FAQ}}
% note = "\AlsoAvailableAs \url{http://go.yurichev.com/17291}",

\newcommand{\ARMCookBook}{Advanced RISC Machines Ltd, \IT{The ARM Cookbook}, (1994)}
% note = "\AlsoAvailableAs \url{http://go.yurichev.com/17273}",

\newcommand{\ARMSevenRef}{\IT{ARM(R) Architecture Reference Manual, ARMv7-A and ARMv7-R edition}, (2012)}

\newcommand{\MAbrash}{Michael Abrash, \IT{Graphics Programming Black Book}, 1997}

\newcommand{\radare}{rada.re}

\newcommand{\RedditHiringThread}{\href{https://www.reddit.com/r/ReverseEngineering/comments/4sbd11/rreverseengineerings_2016_triannual_hiring_thread/}{2016}}
\newcommand{\NetsecHiringThread}{\href{https://www.reddit.com/r/netsec/comments/552rz1/rnetsecs_q4_2016_information_security_hiring/}{2016}}

\newcommand{\FNURLREDDIT}{\footnote{\href{http://go.yurichev.com/17027}{reddit.com/r/ReverseEngineering/}}}

\EN{\newcommand{\AcronymsUsed}{Acronyms Used}

\newcommand{\TitleRE}{Reverse Engineering for Beginners}

\newcommand{\TitleUAL}{Understanding Assembly Language}

\newcommand{\AUTHOR}{Dennis Yurichev}

\newcommand{\etc}{etc}

\newcommand{\figname}{fig.\xspace}
\newcommand{\listingname}{listing.\xspace}
\newcommand{\bitENRU}{bit\xspace}
% FIXME get rid of:
\newcommand{\bitsENRU}{bits\xspace}
\newcommand{\Sourcecode}{Source code\xspace}
\newcommand{\Seealso}{See also\xspace}
\newcommand{\tableheader}{\headercolor{} offset & \headercolor{} description}
% instructions descriptions
\newcommand{\ASRdesc}{arithmetic shift right}

% x86 registers tables
\newcommand{\RegHeaderTop}{ \multicolumn{8}{ | c | }{ Byte number: } }
% TODO: non-overlapping color!
\newcommand{\RegHeader}{ 7th & 6th & 5th & 4th & 3rd & 2nd & 1st & 0th }
\newcommand{\ReturnAddress}{Return Address}

\newcommand{\localVariable}{local variable}

\newcommand{\savedValueOf}{saved value of}

% for index
\newcommand{\GrepUsage}{grep usage}
\newcommand{\SyntacticSugar}{Syntactic Sugar}
\newcommand{\CompilerAnomaly}{Compiler's anomalies}
\newcommand{\CLanguageElements}{C language elements}
\newcommand{\CStandardLibrary}{C standard library}
\newcommand{\Instructions}{Instructions}
\newcommand{\Pseudoinstructions}{Pseudoinstructions}
\newcommand{\Prefixes}{Prefixes}

\newcommand{\Flags}{Flags}
\newcommand{\Registers}{Registers}
\newcommand{\registers}{registers}
\newcommand{\Stack}{Stack}
\newcommand{\Recursion}{Recursion}
\newcommand{\RAM}{RAM}
\newcommand{\ROM}{ROM}
\newcommand{\Pointers}{Pointers}
\newcommand{\BufferOverflow}{Buffer Overflow}

% DE: also "Zusammenfassung"
\newcommand{\Conclusion}{Conclusion}

\newcommand{\Exercise}{Exercise}
\newcommand{\Exercises}{Exercises}
\newcommand{\Arrays}{Arrays}
\newcommand{\Cpp}{C++\xspace}
\newcommand{\CCpp}{C/C++\xspace}
\newcommand{\NonOptimizing}{Non-optimizing\xspace}
\newcommand{\Optimizing}{Optimizing\xspace}
\newcommand{\ARMMode}{ARM mode\xspace}
\newcommand{\ThumbMode}{Thumb mode\xspace}
\newcommand{\ThumbTwoMode}{Thumb-2 mode\xspace}
\newcommand{\AndENRU}{and\xspace}
\newcommand{\OrENRU}{or\xspace}
\newcommand{\InENRU}{in\xspace}
\newcommand{\ForENRU}{for\xspace}
\newcommand{\LineENRU}{line\xspace}

\newcommand{\DataProcessingInstructionsFootNote}{These instructions are also called \q{data processing instructions}}
% for .bib files
\newcommand{\AlsoAvailableAs}{Also available as\xspace}

% section names
\newcommand{\ShiftsSectionName}{Shifts}
\newcommand{\SignedNumbersSectionName}{Signed number representations}
\newcommand{\HelloWorldSectionName}{Hello, world!}
\newcommand{\SwitchCaseDefaultSectionName}{switch()/case/default}
\newcommand{\PrintfSeveralArgumentsSectionName}{printf() with several arguments}
\newcommand{\BitfieldsChapter}{Manipulating specific bit(s)}
\newcommand{\ArithOptimizations}{Replacing arithmetic instructions to other ones}
\newcommand{\FPUChapterName}{Floating-point unit}
\newcommand{\MoreAboutStrings}{More about strings}
\newcommand{\DivisionByMultSectionName}{Division using multiplication}
\newcommand{\Answer}{Answer}
\newcommand{\WhatThisCodeDoes}{What does this code do}
\newcommand{\WorkingWithFloatAsWithStructSubSubSectionName}{Handling float data type as a structure}

\newcommand{\MinesweeperWinXPExampleChapterName}{Minesweeper (Windows XP)}
\newcommand{\StructurePackingSectionName}{Fields packing in structure}
\newcommand{\StructuresChapterName}{Structures}
\newcommand{\PICcode}{position-independent code}
\newcommand{\CapitalPICcode}{Position-independent code}
\newcommand{\Loops}{Loops}

% C
\newcommand{\PostIncrement}{Post-increment}
\newcommand{\PostDecrement}{Post-decrement}
\newcommand{\PreIncrement}{Pre-increment}
\newcommand{\PreDecrement}{Pre-decrement}

% MIPS
\newcommand{\GlobalPointer}{Global Pointer}

\newcommand{\garbage}{garbage}
\newcommand{\IntelSyntax}{Intel syntax}
\newcommand{\ATTSyntax}{AT\&T syntax}
\newcommand{\randomNoise}{random noise}
\newcommand{\Example}{Example}
\newcommand{\argument}{argument}
\newcommand{\MarkedInIDAAs}{marked in \IDA as}
\newcommand{\stepover}{step over}
\newcommand{\ShortHotKeyCheatsheet}{Hot-keys cheatsheet}

\newcommand{\assemblyOutput}{assembly output}
% was in common_URLS.tex:
\newcommand{\URLWPDA}{\href{http://go.yurichev.com/17013}{Wikipedia: Data structure alignment}}

% ML prefix is for multi-lingual words and sentences:
\newcommand{\MLHeap}{Heap}
\newcommand{\MLStack}{Stack}
\newcommand{\MLStackOverflow}{Stack overflow}
\newcommand{\MLStartOfHeap}{Start of heap}
\newcommand{\MLStartOfStack}{Start of stack}
\newcommand{\MLinputA}{input A}
\newcommand{\MLinputB}{input B}
\newcommand{\MLoutput}{output}

}
\RU{\newcommand{\AcronymsUsed}{Список принятых сокращений}

\newcommand{\TitleRE}{Reverse Engineering для начинающих}

\newcommand{\TitleUAL}{Понимание языка ассемблера}

\newcommand{\AUTHOR}{Денис Юричев}

\newcommand{\figname}{илл.\xspace}
\newcommand{\listingname}{листинг.\xspace}
\newcommand{\bitENRU}{бит\xspace}
% FIXME get rid of:
\newcommand{\bitsENRU}{бита\xspace}
\newcommand{\Sourcecode}{Исходный код}
\newcommand{\Seealso}{См. также\xspace}
\newcommand{\tableheader}{\headercolor{} смещение & \headercolor{} описание }
% instructions descriptions
\newcommand{\ASRdesc}{арифметический сдвиг вправо}

% x86 registers tables
\newcommand{\RegHeaderTop}{ \multicolumn{8}{ | c | }{ Номер байта: } }
% TODO: non-overlapping color!
\newcommand{\RegHeader}{ 7-й & 6-й & 5-й & 4-й & 3-й & 2-й & 1-й & 0-й }
\newcommand{\ReturnAddress}{Адрес возврата}

\newcommand{\localVariable}{локальная переменная}

\newcommand{\savedValueOf}{сохраненное значение}

% for index
\newcommand{\GrepUsage}{Использование grep}
\newcommand{\SyntacticSugar}{Синтаксический сахар}
\newcommand{\CompilerAnomaly}{Аномалии компиляторов}
\newcommand{\CLanguageElements}{Элементы языка Си}
\newcommand{\CStandardLibrary}{Стандартная библиотека Си}
\newcommand{\Instructions}{Инструкции}
\newcommand{\Pseudoinstructions}{Псевдоинструкции}
\newcommand{\Prefixes}{Префиксы}

\newcommand{\Flags}{Флаги}
\newcommand{\Registers}{Регистры}
\newcommand{\registers}{регистры}
\newcommand{\Stack}{Стек}
\newcommand{\Recursion}{Рекурсия}
\newcommand{\RAM}{ОЗУ}
\newcommand{\ROM}{ПЗУ}
\newcommand{\Pointers}{Указатели}
\newcommand{\BufferOverflow}{Переполнение буфера}

\newcommand{\Conclusion}{Вывод}

\newcommand{\Exercise}{Упражнение}
\newcommand{\Exercises}{Упражнения}
\newcommand{\Arrays}{Массивы}
\newcommand{\Cpp}{Си++}
\newcommand{\CCpp}{Си/Си++}
\newcommand{\NonOptimizing}{Неоптимизирующий}
\newcommand{\Optimizing}{Оптимизирующий}
\newcommand{\ARMMode}{Режим ARM}
\newcommand{\ThumbMode}{Режим Thumb\xspace}
\newcommand{\ThumbTwoMode}{Режим Thumb-2\xspace}
\newcommand{\AndENRU}{и\xspace}
\newcommand{\OrENRU}{или\xspace}
\newcommand{\InENRU}{в\xspace}
\newcommand{\ForENRU}{для\xspace}
\newcommand{\LineENRU}{строка\xspace}

\newcommand{\DataProcessingInstructionsFootNote}{Эти инструкции также называются \q{data processing instructions}}

% for .bib files
\newcommand{\AlsoAvailableAs}{Также доступно здесь:\xspace}

% section names
\newcommand{\ShiftsSectionName}{Сдвиги}
\newcommand{\SignedNumbersSectionName}{Представление знака в числах}
\newcommand{\HelloWorldSectionName}{Hello, world!}
\newcommand{\SwitchCaseDefaultSectionName}{switch()/case/default}
\newcommand{\PrintfSeveralArgumentsSectionName}{printf() с несколькими аргументами}
\newcommand{\BitfieldsChapter}{Работа с отдельными битами}
\newcommand{\ArithOptimizations}{Замена одних арифметических инструкций на другие}

\newcommand{\FPUChapterName}{Работа с FPU}
\newcommand{\MoreAboutStrings}{Еще кое-что о строках}
\newcommand{\DivisionByMultSectionName}{Деление используя умножение}
\newcommand{\Answer}{Ответ}
\newcommand{\WhatThisCodeDoes}{Что делает этот код}
\newcommand{\WorkingWithFloatAsWithStructSubSubSectionName}{Работа с типом float как со структурой}

\newcommand{\MinesweeperWinXPExampleChapterName}{Сапёр (Windows XP)}
\newcommand{\StructurePackingSectionName}{Упаковка полей в структуре}
\newcommand{\StructuresChapterName}{Структуры}
\newcommand{\PICcode}{адресно-независимый код}
\newcommand{\CapitalPICcode}{Адресно-независимый код}
\newcommand{\Loops}{Циклы}

% C
\newcommand{\PostIncrement}{Пост-инкремент}
\newcommand{\PostDecrement}{Пост-декремент}
\newcommand{\PreIncrement}{Пре-инкремент}
\newcommand{\PreDecrement}{Пре-декремент}

% MIPS
\newcommand{\GlobalPointer}{Глобальный указатель}

\newcommand{\garbage}{мусор}
\newcommand{\IntelSyntax}{Синтаксис Intel}
\newcommand{\ATTSyntax}{Синтаксис AT\&T}
\newcommand{\randomNoise}{случайный шум}
\newcommand{\Example}{Пример}
\newcommand{\argument}{аргумент}
\newcommand{\MarkedInIDAAs}{маркируется в \IDA как}
\newcommand{\stepover}{сделать шаг, не входя в функцию}
\newcommand{\ShortHotKeyCheatsheet}{Краткий справочник горячих клавиш}

\newcommand{\assemblyOutput}{вывод на ассемблере}
% was in common_URLS.tex:
\newcommand{\URLWPDA}{\href{http://go.yurichev.com/17012}{Wikipedia: Выравнивание данных}}

% ML prefix is for multi-lingual words and sentences:
\newcommand{\MLHeap}{Куча}
\newcommand{\MLStack}{Стэк}
\newcommand{\MLStackOverflow}{Переполнение стека}
\newcommand{\MLStartOfHeap}{Начало кучи}
\newcommand{\MLStartOfStack}{Вершина стека}
\newcommand{\MLinputA}{вход А}
\newcommand{\MLinputB}{вход Б}
\newcommand{\MLoutput}{выход}

}
\FR{\newcommand{\AcronymsUsed}{Acronymes utilisés}

\newcommand{\TitleRE}{Reverse Engineering pour Débutants}

\newcommand{\TitleUAL}{Comprendre le langage d'assemblage}

\newcommand{\AUTHOR}{Dennis Yurichev}

\newcommand*{\EMDASH}{---}

\newcommand{\etc}{etc}

\newcommand{\figname}{fig.\xspace}
\newcommand{\listingname}{listado.\xspace}
\newcommand{\bitENRU}{bit\xspace}
% FIXME get rid of:
\newcommand{\bitsENRU}{bits\xspace}
\newcommand{\Sourcecode}{Code source\xspace}
\newcommand{\Seealso}{Voir également\xspace}
\newcommand{\tableheader}{\headercolor{} offset & \headercolor{} description }
% instructions descriptions
\newcommand{\ASRdesc}{décalage arithmétique vers la droite}

% x86 registers tables
\newcommand{\RegHeaderTop}{ \multicolumn{8}{ | c | }{ Octet d'indice } }
% TODO: non-overlapping color!
\newcommand{\RegHeader}{7 & 6 & 5 & 4 & 3 & 2 & 1 & 0}

\newcommand{\ReturnAddress}{Adresse de retour}

\newcommand{\localVariable}{variable locale}

\newcommand{\savedValueOf}{valeur enregistrée de}

% for index
\newcommand{\GrepUsage}{Utilisation de grep}
\newcommand{\SyntacticSugar}{Sucre syntaxique}
\newcommand{\CompilerAnomaly}{Anomalies du compilateur}
\newcommand{\CLanguageElements}{Eléments du langage C}
\newcommand{\CStandardLibrary}{Bibliothèque standard C}
\newcommand{\Instructions}{Instructions}
\newcommand{\Pseudoinstructions}{Pseudo-instructions}
\newcommand{\Prefixes}{Préfixes}

\newcommand{\Flags}{Flags}
\newcommand{\Registers}{Registres}
\newcommand{\registers}{registres}
\newcommand{\Stack}{Pile}
\newcommand{\Recursion}{Récursivité}
\newcommand{\RAM}{RAM}
\newcommand{\ROM}{ROM}
\newcommand{\Pointers}{Pointeurs}
\newcommand{\BufferOverflow}{Débordement de tampon}

\newcommand{\Conclusion}{Conclusion}

\newcommand{\Exercise}{Exercice}
\newcommand{\Exercises}{Exercices\xspace}
\newcommand{\Arrays}{Tableaux}
\newcommand{\Cpp}{C++\xspace}
\newcommand{\CCpp}{C/C++\xspace}
\newcommand{\NonOptimizing}{sans optimisation\xspace}
\newcommand{\Optimizing}{avec optimisation\xspace}
\newcommand{\ARMMode}{Mode ARM\xspace}
\newcommand{\ThumbMode}{Mode Thumb\xspace}
\newcommand{\ThumbTwoMode}{Mode Thumb-2\xspace}
\newcommand{\AndENRU}{et\xspace}
\newcommand{\OrENRU}{ou\xspace}
\newcommand{\InENRU}{dans\xspace}
\newcommand{\ForENRU}{pour\xspace}
\newcommand{\LineENRU}{ligne\xspace}

\newcommand{\DataProcessingInstructionsFootNote}{Ces instructions sont également appelées \q{instructions de traitement de données}}

% for .bib files
\newcommand{\AlsoAvailableAs}{Aussi disponible en\xspace}

% section names
\newcommand{\ShiftsSectionName}{Décalages}
\newcommand{\SignedNumbersSectionName}{Représentations des nombres signés}
\newcommand{\HelloWorldSectionName}{Hello, world!}
\newcommand{\SwitchCaseDefaultSectionName}{switch()/case/default}
\newcommand{\PrintfSeveralArgumentsSectionName}{printf() avec plusieurs arguments}
\newcommand{\BitfieldsChapter}{Manipulation de bits spécifiques}
\newcommand{\ArithOptimizations}{Remplacement de certaines instructions arithmétiques par d'autres}

\newcommand{\FPUChapterName}{Unité à virgule flottante}
\newcommand{\MoreAboutStrings}{Plus d'information sur les chaînes}
\newcommand{\DivisionByMultSectionName}{Division par la multiplication}
\newcommand{\Answer}{Réponse}
\newcommand{\WhatThisCodeDoes}{Que fait ce code ?}
\newcommand{\WorkingWithFloatAsWithStructSubSubSectionName}{Travailler avec le type float comme une structure}

\newcommand{\MinesweeperWinXPExampleChapterName}{Démineur (Windows XP)}
\newcommand{\StructurePackingSectionName}{Organisation des champs dans la structure}
\newcommand{\StructuresChapterName}{Structures}
\newcommand{\PICcode}{code indépendant de la position}
\newcommand{\CapitalPICcode}{Code indépendant de la position}
\newcommand{\Loops}{Boucles}

% C
\newcommand{\PostIncrement}{Post-incrémentation}
\newcommand{\PostDecrement}{Post-décrémentation}
\newcommand{\PreIncrement}{Pré-incrémentation}
\newcommand{\PreDecrement}{Pré-décrémentation}

% MIPS
\newcommand{\GlobalPointer}{Pointeur Global}

\newcommand{\garbage}{déchets}
\newcommand{\IntelSyntax}{Syntaxe Intel}
\newcommand{\ATTSyntax}{Syntaxe AT\&T}
\newcommand{\randomNoise}{bruit aléatoire}
\newcommand{\Example}{Exemple}
\newcommand{\argument}{argument}
\newcommand{\MarkedInIDAAs}{marqué dans \IDA comme}
\newcommand{\stepover}{enjamber}
\newcommand{\ShortHotKeyCheatsheet}{Anti-sèche des touches de raccourci}

\newcommand{\assemblyOutput}{résultat en sortie de l'assembleur}
% was in common_URLS.tex:
\newcommand{\URLWPDA}{\href{https://fr.wikipedia.org/wiki/Alignement_en_m\%C3\%A9moire}{Wikipedia: Alignement en mémoire}}

% ML prefix is for multi-lingual words and sentences:
\newcommand{\MLHeap}{Heap}
\newcommand{\MLStack}{Pile}
\newcommand{\MLStackOverflow}{Débordement de pile}
\newcommand{\MLStartOfHeap}{Début du heap}
\newcommand{\MLStartOfStack}{Début de la pile}
\newcommand{\MLinputA}{entrée A}
\newcommand{\MLinputB}{entrée B}
\newcommand{\MLoutput}{sortie}

}
\DE{\newcommand{\AcronymsUsed}{Verwendete Abkürzungen}

\newcommand{\TitleRE}{Reverse Engineering für Einsteiger}

% TBT?
\newcommand{\TitleUAL}{Understanding Assembly Language}

\newcommand{\AUTHOR}{Dennis Yurichev}
\newcommand{\etc}{\DEph{}}

\newcommand{\figname}{Abb.\xspace}

\newcommand{\listingname}{\DE{Listing.}\xspace}

\newcommand{\bitENRU}{bit\xspace}

% FIXME get rid of:
\newcommand{\bitsENRU}{bits\xspace}
\newcommand{\Sourcecode}{Quellcode\xspace}
\newcommand{\Seealso}{siehe auch\xspace}
\newcommand{\tableheader}{\headercolor{}Offset & \headercolor{} Beschreibung}
% instructions descriptions
\newcommand{\ASRdesc}{\DEph{}}

% x86 registers tables
\newcommand{\RegHeaderTop}{ \multicolumn{8}{ | c | }{ Byte-Nummer: }}

% TODO: non-overlapping color!
\newcommand{\RegHeader}{\DEph{}}

\newcommand{\ReturnAddress}{Rücksprungadresse}

\newcommand{\localVariable}{lokale Variable}

\newcommand{\savedValueOf}{\DEph{}}

% for index
\newcommand{\GrepUsage}{\DEph}
\newcommand{\SyntacticSugar}{\DEph{}}
\newcommand{\CompilerAnomaly}{\DEph{}}
\newcommand{\CLanguageElements}{C Sprachelemente}
\newcommand{\CStandardLibrary}{\DEph{}}
\newcommand{\Instructions}{\DEph{}}
\newcommand{\Pseudoinstructions}{\DEph{}}
\newcommand{\Prefixes}{\DEph{}}

\newcommand{\Flags}{\DEph{}}
\newcommand{\Registers}{\DEph{}}
\newcommand{\registers}{\DEph{}}
\newcommand{\Stack}{Stack}
\newcommand{\Recursion}{\DEph{}}
\newcommand{\RAM}{RAM}
\newcommand{\ROM}{ROM}
\newcommand{\Pointers}{\DEph{}}
\newcommand{\BufferOverflow}{\DEph{}}

% DE: also "Zusammenfassung"
\newcommand{\Conclusion}{Fazit}

\newcommand{\Exercise}{\DEph{}\xspace}
\newcommand{\Exercises}{Übungen\xspace}
\newcommand{\Arrays}{Arrays}
\newcommand{\Cpp}{C++\xspace}
\newcommand{\CCpp}{C/C++\xspace}
\newcommand{\NonOptimizing}{\DEph{}\xspace}
\newcommand{\Optimizing}{\DEph{}\xspace}
\newcommand{\ARMMode}{\DEph{}\xspace}
\newcommand{\ThumbMode}{\DEph{}\xspace}
\newcommand{\ThumbTwoMode}{\DEph{}\xspace}
\newcommand{\AndENRU}{\DEph{}\xspace}
\newcommand{\OrENRU}{\DEph{}\xspace}
\newcommand{\InENRU}{\DEph{}\xspace}
\newcommand{\ForENRU}{\DEph{}\xspace}
\newcommand{\LineENRU}{\DEph{}\xspace}

\newcommand{\DataProcessingInstructionsFootNote}{\DEph{}}
% for .bib files
\newcommand{\AlsoAvailableAs}{\DEph{}\xspace}

% section names
\newcommand{\ShiftsSectionName}{\DEph{}}
\newcommand{\SignedNumbersSectionName}{\DEph{}}
\newcommand{\HelloWorldSectionName}{Hallo, Welt!}
\newcommand{\SwitchCaseDefaultSectionName}{switch()/case/default}
\newcommand{\PrintfSeveralArgumentsSectionName}{\DEph{}}
\newcommand{\BitfieldsChapter}{Manipulieren einzelner Bits}
\newcommand{\ArithOptimizations}{Ersetzen von arithmetischen Operationen}
\newcommand{\FPUChapterName}{Gleitkommaeinheit}
\newcommand{\MoreAboutStrings}{Mehr über Zeichenketten}
\newcommand{\DivisionByMultSectionName}{\DEph{}}
\newcommand{\Answer}{\DEph{}}
\newcommand{\WhatThisCodeDoes}{\DEph{}}
\newcommand{\WorkingWithFloatAsWithStructSubSubSectionName}{\DEph{}}

\newcommand{\MinesweeperWinXPExampleChapterName}{\DEph{} (Windows XP)}
\newcommand{\StructurePackingSectionName}{\DEph{}}
\newcommand{\StructuresChapterName}{\DEph{}}
\newcommand{\PICcode}{positionsabhängiger Code}
\newcommand{\CapitalPICcode}{Positionsabhängiger Code}
\newcommand{\Loops}{Schleifen}

% C
\newcommand{\PostIncrement}{\DEph{}}
\newcommand{\PostDecrement}{\DEph{}}
\newcommand{\PreIncrement}{\DEph{}}
\newcommand{\PreDecrement}{\DEph{}}

% MIPS
\newcommand{\GlobalPointer}{\DEph{}}

\newcommand{\garbage}{\DEph{}}
\newcommand{\IntelSyntax}{\DEph{}}
\newcommand{\ATTSyntax}{\DEph{}}
\newcommand{\randomNoise}{\DEph{}}
\newcommand{\Example}{\DEph{}}
\newcommand{\argument}{\DEph{}}
\newcommand{\MarkedInIDAAs}{\DEph{}}
\newcommand{\stepover}{\DEph{}}
\newcommand{\ShortHotKeyCheatsheet}{\DEph{}}

\newcommand{\assemblyOutput}{\DEph{}}
% was in common_URLS.tex:
\newcommand{\URLWPDA}{\href{http://go.yurichev.com/17013}{Wikipedia: Data structure alignment}}

% ML prefix is for multi-lingual words and sentences:
\newcommand{\MLHeap}{\DEph{}}
\newcommand{\MLStack}{\DEph{}}
\newcommand{\MLStackOverflow}{\DEph{}}
\newcommand{\MLStartOfHeap}{\DEph{}}
\newcommand{\MLStartOfStack}{\DEph{}}
\newcommand{\MLinputA}{\DEph{}}
\newcommand{\MLinputB}{\DEph{}}
\newcommand{\MLoutput}{\DEph{}}

}
\JA{\newcommand{\AcronymsUsed}{頭字語}

\newcommand{\TitleRE}{リバースエンジニアリング入門}

%TBT
\newcommand{\TitleUAL}{Understanding Assembly Language}

\newcommand{\AUTHOR}{Dennis Yurichev}

\newcommand{\figname}{fig.\xspace}
\newcommand{\listingname}{リスト}

\newcommand{\bitENRU}{bit\xspace}
% FIXME get rid of:
\newcommand{\bitsENRU}{bits\xspace}
\newcommand{\Sourcecode}{ソースコード\xspace}
\newcommand{\Seealso}{参照\xspace}
\newcommand{\tableheader}{\headercolor{} オフセット & \headercolor{} 記述 }
% instructions descriptions
\newcommand{\ASRdesc}{\JPNph{}}

% x86 registers tables
\newcommand{\RegHeaderTop}{ \multicolumn{8}{ | c | }{ バイトの並び順 } }
% TODO: non-overlapping color!
\newcommand{\RegHeader}{ 第7 & 第6 & 第5 & 第4 & 第3 & 第2 & 第1 & 第0 }

\newcommand{\ReturnAddress}{リターンアドレス}

\newcommand{\localVariable}{ローカル変数}

\newcommand{\savedValueOf}{saved value of}

% for index
\newcommand{\GrepUsage}{\JPNph{}}
\newcommand{\SyntacticSugar}{糖衣構文}
\newcommand{\CompilerAnomaly}{コンパイラアノマリ}
\newcommand{\CLanguageElements}{C言語の要素}
\newcommand{\CStandardLibrary}{C標準ライブラリ}
\newcommand{\Instructions}{命令}
\newcommand{\Pseudoinstructions}{疑似命令}
\newcommand{\Prefixes}{プリフィックス}

\newcommand{\Flags}{フラグ}
\newcommand{\Registers}{レジスタ}
\newcommand{\registers}{レジスタ}
\newcommand{\Stack}{スタック}
\newcommand{\Recursion}{再帰}
\newcommand{\RAM}{RAM}
\newcommand{\ROM}{ROM}
\newcommand{\Pointers}{ポインタ}
\newcommand{\BufferOverflow}{バッファオーバーフロー}

% DE: also "Zusammenfassung"
\newcommand{\Conclusion}{結論}

\newcommand{\Exercise}{練習問題\xspace}
\newcommand{\Exercises}{練習問題\xspace}
\newcommand{\Arrays}{配列}
\newcommand{\Cpp}{C++\xspace}
\newcommand{\CCpp}{C/C++\xspace}
\newcommand{\NonOptimizing}{非最適化\xspace}
\newcommand{\Optimizing}{最適化\xspace}
\newcommand{\ARMMode}{ARMモード\xspace}
\newcommand{\ThumbMode}{Thumbモード\xspace}
\newcommand{\ThumbTwoMode}{Thumb-2モード\xspace}
\newcommand{\AndENRU}{\JPNph{}}
\newcommand{\OrENRU}{\JPNph{}}
\newcommand{\InENRU}{\JPNph{}}
\newcommand{\ForENRU}{\JPNph{}}
\newcommand{\LineENRU}{\JPNph{}}

\newcommand{\DataProcessingInstructionsFootNote}{この命令は \q{データプロセス命令} とも呼ばれます}

% for .bib files
\newcommand{\AlsoAvailableAs}{以下で利用可能\xspace}

% section names
\newcommand{\ShiftsSectionName}{シフト}
\newcommand{\SignedNumbersSectionName}{符号付き整数表現}
\newcommand{\HelloWorldSectionName}{ハローワールド!}
\newcommand{\SwitchCaseDefaultSectionName}{switch()/case/default}
\newcommand{\PrintfSeveralArgumentsSectionName}{printf() 引数を取って}
\newcommand{\BitfieldsChapter}{特定のビットを操作する}
\newcommand{\ArithOptimizations}{算術命令を他の命令に置換する}

\newcommand{\FPUChapterName}{フローティングポイントユニット}
\newcommand{\MoreAboutStrings}{文字列に関する加筆}
\newcommand{\DivisionByMultSectionName}{乗法を使用した除算}
\newcommand{\Answer}{解答}
\newcommand{\WhatThisCodeDoes}{このコードは何をしている}
\newcommand{\WorkingWithFloatAsWithStructSubSubSectionName}{フロート型のデータを構造体として扱う}

\newcommand{\MinesweeperWinXPExampleChapterName}{マインスイーパー (Windows XP)}
\newcommand{\StructurePackingSectionName}{フィールドを構造体にパッキングする}
\newcommand{\StructuresChapterName}{構造体}
\newcommand{\PICcode}{位置独立コード}
\newcommand{\CapitalPICcode}{位置独立コード}
\newcommand{\Loops}{ループ}

% C
\newcommand{\PostIncrement}{後置インクリメント}
\newcommand{\PostDecrement}{後置デクリメント}
\newcommand{\PreIncrement}{前置インクリメント}
\newcommand{\PreDecrement}{前置デクリメント}

% MIPS
\newcommand{\GlobalPointer}{グローバルポインタ}

\newcommand{\garbage}{ガーベッジ}
\newcommand{\IntelSyntax}{インテル構文}
\newcommand{\ATTSyntax}{AT\&T構文}
\newcommand{\randomNoise}{ランダムノイズ}
\newcommand{\Example}{例}
\newcommand{\argument}{引数}
\newcommand{\MarkedInIDAAs}{\IDA にマークする}
\newcommand{\stepover}{ステップオーバー}
\newcommand{\ShortHotKeyCheatsheet}{ホットキーチートシート}

\newcommand{\assemblyOutput}{アセンブリ出力}
% was in common_URLS.tex:
\newcommand{\URLWPDA}{\href{https://en.wikipedia.org/wiki/Data_structure_alignment}{Wikipedia: Data structure alignment}}

% ML prefix is for multi-lingual words and sentences:
\newcommand{\MLHeap}{ヒープ}
\newcommand{\MLStack}{スタック}
\newcommand{\MLStackOverflow}{スタックオーバーフロー}
\newcommand{\MLStartOfHeap}{ヒープの開始}
\newcommand{\MLStartOfStack}{スタックの開始}
\newcommand{\MLinputA}{入力A}
\newcommand{\MLinputB}{入力B}
\newcommand{\MLoutput}{出力}

}
\IT{\newcommand{\AcronymsUsed}{Acronimi utilizzati}

\newcommand{\TitleRE}{Reverse Engineering per Principianti}

\newcommand{\TitleUAL}{Capire il linguaggio Assembly}

\newcommand{\AUTHOR}{Dennis Yurichev}

\newcommand{\figname}{fig.\xspace}
\newcommand{\listingname}{listato.\xspace}
% FIXME get rid of:
\newcommand{\bitsENRU}{bit\xspace}
\newcommand{\Sourcecode}{Codice sorgente\xspace}
\newcommand{\Seealso}{Vedi anche\xspace}
\newcommand{\tableheader}{\headercolor{} offset & \headercolor{} descrizione}

% instructions descriptions
\newcommand{\ASRdesc}{scorrimento aritmetico a destra}

% x86 registers tables
\newcommand{\RegHeaderTop}{ \multicolumn{8}{ | c | }{ Numero byte } }
% TODO: non-overlapping color!
\newcommand{\RegHeader}{ 7°  & 6°  & 5°  & 4°  & 3°  & 2°  & 1°  & 0 }
\newcommand{\ReturnAddress}{Indirizzo di Ritorno}

\newcommand{\localVariable}{variabile locale}

\newcommand{\savedValueOf}{valore memorizzato di}
% for index
\newcommand{\GrepUsage}{uso di grep}
\newcommand{\SyntacticSugar}{Syntactic Sugar}
\newcommand{\CompilerAnomaly}{Anomalie del compilatore}
\newcommand{\CLanguageElements}{Elementi del linguaggio C}
\newcommand{\CStandardLibrary}{Libreria C standard}
\newcommand{\Instructions}{Istruzioni}
\newcommand{\Pseudoinstructions}{Pseudo-istruzioni}
\newcommand{\Prefixes}{Prefissi}

\newcommand{\Flags}{Flag}
\newcommand{\Registers}{Registri}
\newcommand{\registers}{registri}
\newcommand{\Stack}{Stack}
\newcommand{\Recursion}{Recursione}
\newcommand{\RAM}{RAM}
\newcommand{\ROM}{ROM}
\newcommand{\Pointers}{Puntatori}
\newcommand{\BufferOverflow}{Buffer Overflow}

\newcommand{\Conclusion}{Conclusione}

\newcommand{\Exercise}{Esercizio}
\newcommand{\Exercises}{Esercizi\xspace}
\newcommand{\Arrays}{Array}
\newcommand{\Cpp}{C++\xspace}
\newcommand{\CCpp}{C/C++\xspace}
\newcommand{\NonOptimizing}{Senza ottimizzazione\xspace}
\newcommand{\Optimizing}{Con ottimizzazione\xspace}
\newcommand{\ARMMode}{Modalità ARM\xspace}
\newcommand{\ThumbMode}{Modalità Thumb\xspace}
\newcommand{\ThumbTwoMode}{Modalità Thumb-2\xspace}
\newcommand{\AndENRU}{\ITph{}}
\newcommand{\OrENRU}{\ITph{}}
\newcommand{\InENRU}{\ITph{}}
\newcommand{\ForENRU}{\ITph{}}
\newcommand{\LineENRU}{\ITph{}}

\newcommand{\DataProcessingInstructionsFootNote}{\ITph{}}

% for .bib files
\newcommand{\AlsoAvailableAs}{\ITph{}}

% section names
\newcommand{\ShiftsSectionName}{\ITph{}}
\newcommand{\SignedNumbersSectionName}{Rappresentazione di numeri con segno}
\newcommand{\HelloWorldSectionName}{Hello, world!}
\newcommand{\SwitchCaseDefaultSectionName}{switch()/case/default}
\newcommand{\PrintfSeveralArgumentsSectionName}{printf() con più argomenti}
\newcommand{\BitfieldsChapter}{Manipolando dei bit speifici}
\newcommand{\ArithOptimizations}{Sostituzione di istruzioni aritmetiche con altre}
\newcommand{\FPUChapterName}{Floating-point unit}
\newcommand{\MoreAboutStrings}{Maggiori informazioni sulle stringhe}
\newcommand{\DivisionByMultSectionName}{Divisione utilizzando la moltiplicazione}
\newcommand{\Answer}{Risposta}
\newcommand{\WhatThisCodeDoes}{Cosa fa qusto codice}
\newcommand{\WorkingWithFloatAsWithStructSubSubSectionName}{Utilizzare i tipi di dato float come una struttura}

\newcommand{\MinesweeperWinXPExampleChapterName}{\ITph{} (Windows XP)}
\newcommand{\StructurePackingSectionName}{Organizzazione dei campi in una struttura}
\newcommand{\StructuresChapterName}{Strutture}
\newcommand{\PICcode}{codice indipendente dalla posizione}
\newcommand{\CapitalPICcode}{Codice indipendente dalla posizione}
\newcommand{\Loops}{Cicli}

% C
\newcommand{\PostIncrement}{Post-incremento}
\newcommand{\PostDecrement}{Post-decremento}
\newcommand{\PreIncrement}{Pre-incremento}
\newcommand{\PreDecrement}{Pre-decremento}

% MIPS
\newcommand{\GlobalPointer}{Puntatore Globale}

\newcommand{\garbage}{garbage}
\newcommand{\IntelSyntax}{Sintassi Intel}
\newcommand{\ATTSyntax}{Sintassi AT\&T}
\newcommand{\randomNoise}{rumore casuale}
\newcommand{\Example}{Esempio}
\newcommand{\argument}{argomento}
\newcommand{\MarkedInIDAAs}{marcato in IDA come}
\newcommand{\stepover}{step over}
\newcommand{\ShortHotKeyCheatsheet}{Elenco delle scorciatoie da tastiera}

\newcommand{\assemblyOutput}{risultato dell'assembly}
% was in common_URLS.tex:
\newcommand{\URLWPDA}{\href{https://en.wikipedia.org/wiki/Data_structure_alignment}{Wikipedia: Data structure alignment}}

% ML prefix is for multi-lingual words and sentences:
\newcommand{\MLHeap}{Heap}
\newcommand{\MLStack}{Stack}
\newcommand{\MLStackOverflow}{Stack overflow}
\newcommand{\MLStartOfHeap}{Inizio dell'heap}
\newcommand{\MLStartOfStack}{Inizio dello stack}
\newcommand{\MLinputA}{input A}
\newcommand{\MLinputB}{input A}
\newcommand{\MLoutput}{output}

}
\PL{\newcommand{\TitleRE}{\PLph{}}
\newcommand{\TitleUAL}{Understanding Assembly Language}%     %TBT

\newcommand{\AUTHOR}{Dennis Yurichev}

\newcommand*{\EMDASH}{---}

\newcommand{\tableheader}{\headercolor{} \PLph{} & \headercolor{} \PLph{} }

% x86 registers tables
\newcommand{\RegHeaderTop}{ \multicolumn{8}{ | c | }{ \PLph{} } }
}
\ES{\newcommand{\AcronymsUsed}{Acr\'onimos utilizados}

\newcommand{\TitleRE}{Ingenier\'ia Inversa para Principiantes}

% TBT?
\newcommand{\TitleUAL}{Understanding Assembly Language}

\newcommand{\AUTHOR}{Dennis Yurichev}

\newcommand*{\EMDASH}{---}

\newcommand{\etc}{etc}

\newcommand{\figname}{fig.\xspace}
\newcommand{\listingname}{listado.\xspace}
\newcommand{\bitENRU}{bit}\xspace}
% FIXME get rid of:
\newcommand{\bitsENRU}{bits}\xspace}
\newcommand{\Sourcecode}{C\'odigo fuente\xspace}
\newcommand{\Seealso}{V\'ease tambi\'en\xspace}
\newcommand{\tableheader}{\headercolor{} offset & \headercolor{} descripci\'on }
% instructions descriptions
\newcommand{\ASRdesc}{desplazamiento aritm\'etico a la derecha}

% x86 registers tables
\newcommand{\RegHeaderTop}{ \multicolumn{8}{ | c | }{ \ESph{} } }
% TODO: non-overlapping color!
\newcommand{\RegHeader}{ 7mo & 6to & 5to & 4to & 3ro & 2do & 1ro & 0 }
\newcommand{\ReturnAddress}{ Direcci\'on de Retorno }

\newcommand{\localVariable}{\ESph{}}

\newcommand{\savedValueOf}{\ESph{}}

% for index
\newcommand{\GrepUsage}{Uso de grep}
\newcommand{\SyntacticSugar}{Azúcar sintáctica}
\newcommand{\CompilerAnomaly}{Anomalías del compilador}
\newcommand{\CLanguageElements}{Elementos del lenguaje C}
\newcommand{\CStandardLibrary}{Librería estándar C}
\newcommand{\Instructions}{Instrucciones}
\newcommand{\Pseudoinstructions}{Pseudo-instrucciones}
\newcommand{\Prefixes}{Prefijos}

\newcommand{\Flags}{Flags}
\newcommand{\Registers}{Registros}
\newcommand{\registers}{registros}
\newcommand{\Stack}{Pila}
\newcommand{\Recursion}{Recursión}
\newcommand{\RAM}{RAM}
\newcommand{\ROM}{ROM}
\newcommand{\Pointers}{Apuntadores}
\newcommand{\BufferOverflow}{Desbordamiento de buffer}

% DE: also "Zusammenfassung"
\newcommand{\Conclusion}{Conclusión}

\newcommand{\Exercise}{Ejercicio\xspace}
\newcommand{\Exercises}{Ejercicios\xspace}
\newcommand{\Arrays}{Matriz}
\newcommand{\Cpp}{C++}
\newcommand{\CCpp}{C/C++}
\newcommand{\NonOptimizing}{Sin optimización}
\newcommand{\Optimizing}{Con optimización}
\newcommand{\ARMMode}{Modo ARM\xspace}
\newcommand{\ThumbMode}{Modo Thumb\xspace}
\newcommand{\ThumbTwoMode}{Modo Thumb-2\xspace}
\newcommand{\AndENRU}{y\xspace}
\newcommand{\OrENRU}{o\xspace}
\newcommand{\InENRU}{en\xspace}
\newcommand{\ForENRU}{para\xspace}
\newcommand{\LineENRU}{línea\xspace}

\newcommand{\DataProcessingInstructionsFootNote}{Estas instrucciones también son llamadas \q{instrucciones de procesamiento de datos}}

% for .bib files
\newcommand{\AlsoAvailableAs}{También disponible como\xspace}

% section names
\newcommand{\ShiftsSectionName}{Desplazamientos}
\newcommand{\SignedNumbersSectionName}{Representaci\'on de n\'umeros con signo}
\newcommand{\HelloWorldSectionName}{!`Hola, mundo!}
\newcommand{\SwitchCaseDefaultSectionName}{switch()/case/default}
\newcommand{\PrintfSeveralArgumentsSectionName}{printf() con varios argumentos}
\newcommand{\BitfieldsChapter}{Manipulando bit(s) específicos}
\newcommand{\ArithOptimizations}{Substituición de instrucciones aritméticas por otras}
	
\newcommand{\FPUChapterName}{Unidad de punto flotante}
\newcommand{\MoreAboutStrings}{\ESph{}}
\newcommand{\DivisionByMultSectionName}{División entre 9}
\newcommand{\Answer}{Respuesta}
\newcommand{\WhatThisCodeDoes}{?`Qu\'e hace este código?}
\newcommand{\WorkingWithFloatAsWithStructSubSubSectionName}{Trabajando con el tipo float como una estructura}

\newcommand{\MinesweeperWinXPExampleChapterName}{Buscaminas (Windows XP)}
\newcommand{\StructurePackingSectionName}{Organización de campos en la estructura}
\newcommand{\StructuresChapterName}{Estructuras}
\newcommand{\PICcode}{código independiente de la posición}
\newcommand{\CapitalPICcode}{Código independiente de lá posición}
\newcommand{\Loops}{Bucles}

% C
\newcommand{\PostIncrement}{Post-incremento}
\newcommand{\PostDecrement}{Post-decremento}
\newcommand{\PreIncrement}{Pre-incremento}
\newcommand{\PreDecrement}{Pre-decremento}

% MIPS
\newcommand{\GlobalPointer}{Apuntador Global}

\newcommand{\garbage}{Basura}
\newcommand{\IntelSyntax}{Sintaxis Intel}
\newcommand{\ATTSyntax}{Sintaxis AT\&T}
\newcommand{\randomNoise}{Ruido aleatorio}
\newcommand{\Example}{Ejemplo}
\newcommand{\argument}{argumento}
\newcommand{\MarkedInIDAAs}{Marcado en \IDA como}
\newcommand{\stepover}{pasar por encima}
\newcommand{\ShortHotKeyCheatsheet}{Cheatsheet de teclas de acceso rápido}

\newcommand{\assemblyOutput}{salida del ensamblador}
% was in common_URLS.tex:
\newcommand{\URLWPDA}{\href{http://go.yurichev.com/17013}{Wikipedia: Data structure alignment}}

% ML prefix is for multi-lingual words and sentences:
\newcommand{\MLHeap}{\ESph{}}
\newcommand{\MLStack}{\ESph{}}
\newcommand{\MLStackOverflow}{\ESph{}}
\newcommand{\MLStartOfHeap}{\ESph{}}
\newcommand{\MLStartOfStack}{\ESph{}}
\newcommand{\MLinputA}{\ESph{}}
\newcommand{\MLinputB}{\ESph{}}
\newcommand{\MLoutput}{\ESph{}}

}
\NL{\newcommand{\AcronymsUsed}{Gebruikte afkortingen}

\newcommand{\TitleRE}{\NLph{}}

\newcommand{\TitleUAL}{\NLph{}}

\newcommand{\AUTHOR}{Dennis Yurichev}

% x86 registers tables
\newcommand{\RegHeaderTop}{ \multicolumn{8}{ | c | }{ \NLph{} } }
% TODO: non-overlapping color!
% TBT
\newcommand{\RegHeader}{ 7th & 6th & 5th & 4th & 3rd & 2nd & 1st & 0th }

\newcommand{\ReturnAddress}{Return Adres}
\newcommand{\SoftwareCracking}{\NLph{}}

}
\TR{\newcommand{\TitleRE}{Acemiler için Tersine Mühendislik}

}
\PTBR{\newcommand{\ReturnAddress}{Endereço de retorno}

\newcommand{\localVariable}{Variável local}

\newcommand{\savedValueOf}{Valor salvo de}

% for index
\newcommand{\GrepUsage}{Uso do grep}
\newcommand{\SyntacticSugar}{Açúcar sintático}
\newcommand{\CompilerAnomaly}{Anomalias do compilador}
\newcommand{\CLanguageElements}{Elementos da linguagem C}
\newcommand{\CStandardLibrary}{Biblioteca padrão C}
\newcommand{\Instructions}{Instruções}
\newcommand{\Pseudoinstructions}{Pseudo-instruções}
\newcommand{\Prefixes}{Prefixos}

\newcommand{\Flags}{Flags}
\newcommand{\Registers}{Registradores}
\newcommand{\registers}{registradores}
\newcommand{\Stack}{Pilha}
\newcommand{\Recursion}{Recursividade}
\newcommand{\RAM}{RAM}
\newcommand{\ROM}{ROM}
\newcommand{\Pointers}{Ponteiros}
\newcommand{\BufferOverflow}{Buffer Overflow}

% DE: also "Zusammenfassung"
\newcommand{\Conclusion}{Conclusão}

\newcommand{\Exercise}{Exercício}
\newcommand{\Exercises}{Exercícios}
\newcommand{\Arrays}{Matriz}
\newcommand{\Cpp}{C++\xspace}
\newcommand{\CCpp}{C/C++\xspace}
\newcommand{\NonOptimizing}{Sem otimização\xspace}
\newcommand{\Optimizing}{Com otimização\xspace}
\newcommand{\ARMMode}{Modo ARM\xspace}
\newcommand{\ThumbMode}{Modo Thumb\xspace}
\newcommand{\ThumbTwoMode}{Modo Thumb-2\xspace}
\newcommand{\AndENRU}{e\xspace}
\newcommand{\OrENRU}{ou\xspace}
\newcommand{\InENRU}{em\xspace}
\newcommand{\ForENRU}{para\xspace}
\newcommand{\LineENRU}{linha\xspace}

\newcommand{\DataProcessingInstructionsFootNote}{Estas intruções também são chamadas \q{data processing instructions}}

% for .bib files
\newcommand{\AlsoAvailableAs}{Também disponível como\xspace}

% section names
\newcommand{\ShiftsSectionName}{Shifts}
\newcommand{\SwitchCaseDefaultSectionName}{switch()/case/default}
\newcommand{\BitfieldsChapter}{Manipulando bit(s) específicos}
\newcommand{\ArithOptimizations}{Substituição de instruções aritiméticas por outras}
\newcommand{\FPUChapterName}{Unidade de Ponto flutuante}
\newcommand{\DivisionByMultSectionName}{Divisão por 9}
\newcommand{\Answer}{Responda}
\newcommand{\WhatThisCodeDoes}{O que este código faz}
\newcommand{\WorkingWithFloatAsWithStructSubSubSectionName}{Trabalhando com o tipo float como uma estrutura}

\newcommand{\MinesweeperWinXPExampleChapterName}{Campo minado (Windows XP)}
\newcommand{\StructurePackingSectionName}{Organização de campos na estrutura}
\newcommand{\StructuresChapterName}{Estruturas}
\newcommand{\PICcode}{código independente de posição}
\newcommand{\CapitalPICcode}{Código independente de posição}
\newcommand{\Loops}{Laços}

% C
\newcommand{\PostIncrement}{Pós-incremento}
\newcommand{\PostDecrement}{Pós-decremento}
\newcommand{\PreIncrement}{Pré-incremento}
\newcommand{\PreDecrement}{Pré-decremento}

% MIPS
\newcommand{\GlobalPointer}{Ponteiro Global}

\newcommand{\garbage}{Lixo}
\newcommand{\IntelSyntax}{Sintaxe Intel}
\newcommand{\ATTSyntax}{Sintaxe AT\&T}
\newcommand{\randomNoise}{Ruído aleatório}
\newcommand{\Example}{Exemplo}
\newcommand{\argument}{argumento}
\newcommand{\MarkedInIDAAs}{Marcado no \IDA como}
\newcommand{\stepover}{passar por cima}
\newcommand{\ShortHotKeyCheatsheet}{Cheatsheet de teclas de atalho}

\newcommand{\assemblyOutput}{saída do assembly}
% was in common_URLS.tex:
\newcommand{\URLWPDA}{\href{http://go.yurichev.com/17013}{Wikipedia: Data structure alignment}}

% ML prefix is for multi-lingual words and sentences:
\newcommand{\MLHeap}{Heap}
\newcommand{\MLStack}{Pilha}
\newcommand{\MLStackOverflow}{\PTBRph{}}
\newcommand{\MLStartOfHeap}{começo da heap}
\newcommand{\MLStartOfStack}{começo da pilha}
\newcommand{\SoftwareCracking}{\PTBRph{}}

}

\ifdefined\UAL
\newcommand{\TitleMain}{\TitleUAL}
\newcommand{\TitleAux}{\TitleRE}
\else
\newcommand{\TitleMain}{\TitleRE}
\newcommand{\TitleAux}{\TitleUAL}
\fi

\EN{\newglossaryentry{tail call}
{
  name=tail call,
  description={It is when the compiler (or interpreter) transforms the recursion
  (\IT{tail recursion}) into an iteration for efficiency}
}

\newglossaryentry{endianness}
{
  name=endianness,
  description={Byte order}
}

\newglossaryentry{caller}
{
  name=caller,
  description={A function calling another}
}

\newglossaryentry{callee}
{
  name=callee,
  description={A function being called by another}
}

\newglossaryentry{debuggee}
{
  name=debuggee,
  description={A program being debugged}
}

\newglossaryentry{leaf function}
{
  name=leaf function,
  description={A function which does not call any other function}
}

\newglossaryentry{link register}
{
  name=link register,
  description=(RISC) {A register where the return address is usually stored.
  This makes it possible to call leaf functions without using the stack, i.e., faster}
}

\newglossaryentry{anti-pattern}
{
  name=anti-pattern,
  description={Generally considered as bad practice}
}

\newglossaryentry{stack pointer}
{
  name=stack pointer,
  description={A register pointing to a place in the stack}
}

\newglossaryentry{decrement}
{
  name=decrement,
  description={Decrease by 1}
}

\newglossaryentry{increment}
{
  name=increment,
  description={Increase by 1}
}

\newglossaryentry{loop unwinding}
{
  name=loop unwinding,
  description={It is when a compiler, instead of generating loop code for $n$ iterations, generates just $n$ copies of the
  loop body, in order to get rid of the instructions for loop maintenance}
}

\newglossaryentry{register allocator}
{
  name=register allocator,
  description=
  {The part of the compiler that assigns CPU registers to local variables}
}

\newglossaryentry{quotient}
{
  name=quotient,
  description={Division result}
}

\newglossaryentry{product}
{
  name=product,
  description={Multiplication result}
}

\newglossaryentry{NOP}
{
  name=NOP,
  description={\q{no operation}, idle instruction}
}

\newglossaryentry{POKE}
{
  name=POKE,
  description={BASIC language instruction for writing a byte at a specific address}
}

\newglossaryentry{keygenme}
{
  name=keygenme,
  description={A program which imitates software protection,
  for which one needs to make a key/license generator}
} % TODO clarify: A software which generate key/license value to bypass sotfware protection?

\newglossaryentry{dongle}
{
  name=dongle,
  description={Dongle is a small piece of hardware connected to LPT printer port (in past) or to USB}
}

\newglossaryentry{thunk function}
{
  name=thunk function,
  description={Tiny function with a single role: call another function}
}

\newglossaryentry{user mode}
{
  name=user mode,
  description={A restricted CPU mode in which it all application software code is executed. cf. \gls{kernel mode}}
}

\newglossaryentry{kernel mode}
{
  name=kernel mode,
  description={A restrictions-free CPU mode in which the OS kernel and drivers execute. cf. \gls{user mode}}
}

\newglossaryentry{Windows NT}
{
  name=Windows NT,
  description={Windows NT, 2000, XP, Vista, 7, 8, 10}
}

\newglossaryentry{atomic operation}
{
  name=atomic operation,
  description={
  \q{$\alpha{}\tau{}o\mu{}o\varsigma{}$}
  %\q{atomic}
  stands for \q{indivisible} in Greek, so an atomic operation is guaranteed not
  to be interrupted by other threads}
}

% to be proofreaded (begin)
\newglossaryentry{NaN}
{
  name=NaN,
  description={not a number: 
  	a special cases for floating point numbers, usually signaling about errors}
}

\newglossaryentry{basic block}
{
  name=basic block,
  description={
	a group of 
	instructions that do not have jump/branch instructions, and also don't have
	jumps inside the block from the outside.
	In \IDA it looks just like as a list of instructions without empty lines}
}

\newglossaryentry{NEON}
{
  name=NEON,
  description={\ac{AKA} \q{Advanced SIMD}---\ac{SIMD} from ARM}
}

\newglossaryentry{reverse engineering}
{
  name=reverse engineering,
  description={act of understanding how the thing works, sometimes in order to clone it}
}

\newglossaryentry{compiler intrinsic}
{
  name=compiler intrinsic,
  description={A function specific to a compiler which is not an usual library function.
	The compiler generates a specific machine code instead of a call to it.
	Often, it's a pseudofunction for a specific \ac{CPU} instruction. Read more:}
 (\myref{sec:compiler_intrinsic})
}

\newglossaryentry{heap}
{
  name=heap,
  description={usually, a big chunk of memory provided by the \ac{OS} so that applications can divide it by themselves as they wish.
  malloc()/free() work with the heap}
}

\newglossaryentry{name mangling}
{
  name=name mangling,
  description={used at least in \Cpp, where the compiler needs to encode the name of class, method and argument types in one string,
  which will become the internal name of the function. You can read more about it here: \myref{namemangling}}
}

\newglossaryentry{xoring}
{
  name=xoring,
  description={often used in the English language, which implying applying the \ac{XOR} operation}
}

\newglossaryentry{security cookie}
{
  name=security cookie,
  description={A random value, different at each execution. You can read more about it here:
  \myref{subsec:BO_protection}}
}

\newglossaryentry{tracer}
{
  name=tracer,
  description={My own simple debugging tool. You can read more about it here: \myref{tracer}}
}

\newglossaryentry{GiB}
{
  name=GiB,
  description={Gibibyte: $2^{30}$ or 1024 mebibytes or 1073741824 bytes}
}

\newglossaryentry{CP/M}
{
  name=CP/M,
  description={Control Program for Microcomputers: a very basic disk \ac{OS} used before MS-DOS}
}

\newglossaryentry{stack frame}
{
  name=stack frame,
  description={A part of the stack that contains information specific to the current function:
  local variables, function arguments, \ac{RA}, etc.}
}

\newglossaryentry{jump offset}
{
  name=jump offset,
  description={a part of the JMP or Jcc instruction's opcode, 
  to be added to the address
  of the next instruction, and this is how the new \ac{PC} is calculated. May be negative as well}
}

\newglossaryentry{integral type}
{
  name=integral data type,
  description={usual numbers, but not a real ones. may be used for passing variables of boolean data type and enumerations}
}

\newglossaryentry{real number}
{
  name={real number},
  description={numbers which may contain a dot. this is \Tfloat and \Tdouble in \CCpp}
}

\newglossaryentry{PDB}
{
  name=PDB,
  description={(Win32) Debugging information file, usually just function names, but sometimes also function
  arguments and local variables names}
}

\newglossaryentry{NTAPI}
{
  name=NTAPI,
  description={\ac{API} available only in the Windows NT line.  Largely not documented by Microsoft}
}

\newglossaryentry{stdout}
{
  name=stdout,
  description={standard output}
}

\newglossaryentry{word}
{
  name=word,
  description={data type fitting in \ac{GPR}.
  In the computers older than PCs, 
  the memory size was often measured in words rather than bytes.}
}

\newglossaryentry{arithmetic mean}
{
  name=arithmetic mean,
  description={a sum of all values divided by their count}
}
\newglossaryentry{padding}
{
  name=padding,
  description={
  \IT{Padding} in English language means to stuff a pillow with something
  to give it a desired (bigger) form.
  In computer science, padding means to add more bytes to a block so it will have desired size, like $2^n$ bytes.}
}

}
\RU{\newglossaryentry{tail call}
{
  name=хвостовая рекурсия,
  description={Это когда компилятор или интерпретатор превращает рекурсию 
  (с которой возможно это проделать, т.е. \IT{хвостовую}) в итерацию для эффективности}
}

\newglossaryentry{endianness}
{
  name=endianness,
  description={Порядок байт}
}

\newglossaryentry{caller}
{
  name=caller,
  description={Функция вызывающая другую функцию}
}

\newglossaryentry{callee}
{
  name=callee,
  description={Вызываемая функция}
}

\newglossaryentry{debuggee}
{
  name=debuggee,
  description={Отлаживаемая программа}
}

\newglossaryentry{leaf function}
{
  name=leaf function,
  description={Функция не вызывающая больше никаких функций}
}

\newglossaryentry{link register}
{
  name=link register,
  description=(RISC) {Регистр в котором обычно записан адрес возврата.
  Это позволяет вызывать leaf-функции без использования стека, т.е. быстрее}
}

\newglossaryentry{anti-pattern}
{
  name=anti-pattern,
  description={Нечто широко известное как плохое решение}
}

\newglossaryentry{stack pointer}
{
  name=указатель стека,
  description={Регистр указывающий на место в стеке}
}

\newglossaryentry{decrement}
{
  name=декремент,
  description={Уменьшение на 1}
}

\newglossaryentry{increment}
{
  name=инкремент,
  description={Увеличение на 1}
}

\newglossaryentry{loop unwinding}
{
  name=loop unwinding,
  description={Это когда вместо организации цикла на $n$ итераций, компилятор генерирует $n$ копий тела
  цикла, для экономии на инструкциях, обеспечивающих сам цикл}
}

\newglossaryentry{register allocator}
{
  name=register allocator,
  description={Функция компилятора распределяющая локальные переменные по регистрам процессора}
}

\newglossaryentry{quotient}
{
  name=частное,
  description={Результат деления}
}

\newglossaryentry{product}
{
  name={произведение},
  description={Результат умножения}
}

\newglossaryentry{NOP}
{
  name=NOP,
  description={\q{no operation}, холостая инструкция}
}

\newglossaryentry{POKE}
{
  name=POKE,
  description={Инструкция языка BASIC записывающая байт по определенному адресу}
}

\newglossaryentry{keygenme}
{
  name=keygenme,
  description={Программа, имитирующая защиту вымышленной программы, для которой нужно сделать 
  генератор ключей/лицензий}
} % TODO clarify: A software which generate key/license value to bypass sotfware protection?

\newglossaryentry{dongle}
{
  name=dongle,
  description={Небольшое устройство подключаемое к LPT-порту для принтера (в прошлом) или к USB}
}

\newglossaryentry{thunk function}
{
  name=thunk function,
  description={Крохотная функция делающая только одно: вызывающая другую функцию}
}

\newglossaryentry{user mode}
{
  name=user mode,
  description={Режим CPU с ограниченными возможностями в котором он исполняет прикладное ПО. ср.}
}

\newglossaryentry{kernel mode}
{
  name=kernel mode,
  description={Режим CPU с неограниченными возможностями в котором он исполняет ядро OS и драйвера. ср.}
}

\newglossaryentry{Windows NT}
{
  name=Windows NT,
  description={Windows NT, 2000, XP, Vista, 7, 8, 10}
}

\newglossaryentry{atomic operation}
{
  name=atomic operation,
  description={
  \q{$\alpha{}\tau{}o\mu{}o\varsigma{}$}
  %\q{atomic}
  означает \q{неделимый} в греческом языке, так что атомарная операция ---
  это операция которая гарантированно не будет прервана другими тредами}
}

% to be proofreaded (begin)
\newglossaryentry{NaN}
{
  name=NaN,
  description={
  	не число: специальные случаи чисел с плавающей запятой, 
  	обычно сигнализирующие об ошибках
  }
}

\newglossaryentry{basic block}
{
  name=basic block,
  description={
  	группа инструкций, не имеющая инструкций переходов,
	а также не имеющая переходов в середину блока извне.
	В \IDA он выглядит как просто список инструкций без строк-разрывов
  }
}

\newglossaryentry{NEON}
{
  name=NEON,
  description={\ac{AKA} \q{Advanced SIMD} --- от ARM}
}

\newglossaryentry{reverse engineering}
{
  name=reverse engineering,
  description={процесс понимания как устроена некая вещь, иногда, с целью клонирования оной}
}

\newglossaryentry{compiler intrinsic}
{
  name=compiler intrinsic,
  description={Специфичная для компилятора функция не являющаяся обычной библиотечной функцией.
	Компилятор вместо её вызова генерирует определенный машинный код.
	Нередко, это псевдофункции для определенной инструкции \ac{CPU}. Читайте больше:}
}

\newglossaryentry{heap}
{
  name=heap,
  description={(куча) обычно, большой кусок памяти предоставляемый \ac{OS}, так что прикладное ПО может делить его
  как захочет. malloc()/free() работают с кучей}
}

\newglossaryentry{name mangling}
{
  name=name mangling,
  description={применяется как минимум в \Cpp, где компилятору нужно закодировать имя класса,
  метода и типы аргументов в одной
  строке, которая будет внутренним именем функции. читайте также здесь: \myref{namemangling}}
}

\newglossaryentry{xoring}
{
  name=xoring,
  description={нередко применяемое в английском языке, означает применение операции 
  \ac{XOR}}
}

\newglossaryentry{security cookie}
{
  name=security cookie,
  description={Случайное значение, разное при каждом исполнении. Читайте больше об этом тут}
}

\newglossaryentry{tracer}
{
  name=tracer,
  description={Моя простейшая утилита для отладки. Читайте больше об этом тут: \myref{tracer}}
}

\newglossaryentry{GiB}
{
  name=GiB,
  description={Гибибайт: $2^{30}$ или 1024 мебибайт или 1073741824 байт}
}

\newglossaryentry{CP/M}
{
  name=CP/M,
  description={Control Program for Microcomputers: очень простая дисковая \ac{OS} использовавшаяся перед MS-DOS}
}

\newglossaryentry{stack frame}
{
  name=stack frame,
  description={Часть стека, в которой хранится информация, связанная с текущей функцией: локальные переменные,
  аргументы функции, \ac{RA}, итд.}
}

\newglossaryentry{jump offset}
{
  name=jump offset,
  description={Часть опкода JMP или Jcc инструкции, просто прибавляется к адресу следующей инструкции,
  и так вычисляется новый \ac{PC}. Может быть отрицательным}
}

\newglossaryentry{integral type}
{
  name=интегральный тип данных,
  description={обычные числа, но не вещественные. могут использоваться для передачи булевых типов и перечислений (enumerations)}
}

\newglossaryentry{real number}
{
  name=вещественное число,
  description={числа, которые могут иметь точку. в \CCpp это \Tfloat и \Tdouble}
}

\newglossaryentry{PDB}
{
  name=PDB,
  description={(Win32) Файл с отладочной информацией, обычно просто имена функций, 
  но иногда имена аргументов функций и локальных переменных}
}

\newglossaryentry{NTAPI}
{
  name=NTAPI,
  description={\ac{API} доступное только в линии Windows NT. 
  Большей частью не документировано Microsoft-ом}
}

\newglossaryentry{stdout}
{
  name=stdout,
  description={standard output}
}

\newglossaryentry{word}
{
  name=word,
  description={(слово) тип данных помещающийся в \ac{GPR}. 
  В компьютерах старше персональных, память часто измерялась не в байтах, 
  а в словах}
}

\newglossaryentry{arithmetic mean}
{
  name=среднее арифметическое,
  description={сумма всех значений, разделенная на их количество}
}
\newglossaryentry{padding}
{
  name=padding,
  description={\IT{Padding} в английском языке означает набивание подушки чем-либо для придания ей желаемой (большей)
  формы. В компьютерных науках, \IT{padding} означает добавление к блоку дополнительных байт, чтобы он имел нужный
  размер, например, $2^n$ байт.}
}

}
\FR{\newglossaryentry{tail call}
{
  name={tail call},
  description={C'est lorsque le compilateur (ou l'interpréteur) transforme la récursion (ce qui est possible: \emph{tail recursion})
  en une itération pour l'efficacité: \href{http://go.yurichev.com/17105}{wikipedia}}
}

\newglossaryentry{endianness}
{
  name=endianness,
  description={Ordre des octets: \myref{sec:endianness}}
}

\newglossaryentry{caller}
{
  name=caller,
  description={Une fonction en appelant une autre}
}

\newglossaryentry{callee}
{
  name=callee,
  description={Une fonction appelée par une autre}
}

\newglossaryentry{debuggee}
{
  name=debuggee,
  description={Un programme en train d'être débogué}
}

\newglossaryentry{leaf function}
{
  name=leaf function,
  description={Une fonction qui n'appelle pas d'autre fonction}
}

\newglossaryentry{link register}
{
  name=link register,
  description={(RISC) Un registre où l'adresse de retour est en général stockée. Ceci permet
  d'appeler une fonction leaf sans utiliser la pile, i.e, plus rapidemment}
}

\newglossaryentry{anti-pattern}
{
  name=anti-pattern,
  description={En général considéré comme une mauvaise pratique}
}

\newglossaryentry{stack pointer}
{
  name=pointeur de pile,
  description={Un registre qui pointe dans la pile}
}

\newglossaryentry{decrement}
{
  name=décrémenter,
  description={Décrémenter de 1}
}

\newglossaryentry{increment}
{
  name=incrémenter,
  description={Incrémenter de 1}
}

\newglossaryentry{loop unwinding}
{
  name=loop unwinding,
  description={C'est lorsqu'un compilateur, au lieu de générer du code pour une boucle de
  $n$ itérations, génère juste $n$ copies du corps de la boucle, afin de supprimer
  les instructions pour la gestion de la boucle}
}

\newglossaryentry{register allocator}
{
  name=register allocator,
  description={La partie du compilateur qui assigne des registes du CPU aux variables locales}
}

\newglossaryentry{quotient}
{
  name=quotient,
  description={Résultat de la division}
}

\newglossaryentry{product}
{
  name=produit,
  description={Résultat d'une multiplication}
}

\newglossaryentry{NOP}
{
  name=NOP,
  description={\q{no operation}, instruction ne faisant rien}
}

\newglossaryentry{POKE}
{
  name=POKE,
  description={instruction du langage BASIC pour écrire un octet à une adresse spécifique}
}

\newglossaryentry{keygenme}
{
  name=keygenme,
  description={Un programme qui imite la protection des logiciels pour lesquels on a besoin d'un générateur de clef/licence}
} % TODO clarify: A software which generate key/license value to bypass sotfware protection?

\newglossaryentry{dongle}
{
  name=dongle,
  description={Un dongle est un petit périphérique se connectant sur un port d'imprimante LPT (par le passé) ou USB.
  Sa fonction est similaire au tokens de sécurité, il a de la mémoire et, parfois, un algorithme secret de (crypto-)hachage.}
}

\newglossaryentry{thunk function}
{
  name=thunk function,
  description={Minuscule fonction qui a un seul rôle: appeler une autre fonction}
}

\newglossaryentry{user mode}
{
  name=user mode,
  description={Un mode CPU restreint dans lequel le code de toutes les applications est exécuté. cf. \gls{kernel mode}}
}

\newglossaryentry{kernel mode}
{
  name=kernel mode,
  description={Un mode CPU sans restriction dans lequel le noyau de l'OS et les drivers sont exécutés. cf. \gls{user mode}}
}

\newglossaryentry{Windows NT}
{
  name=Windows NT,
  description={Windows NT, 2000, XP, Vista, 7, 8, 10}
}

\newglossaryentry{atomic operation}
{
  name=atomic operation,
  description={
  \q{$\alpha{}\tau{}o\mu{}o\varsigma{}$}
  %\q{atomic}
  signifie \q{indivisible} en grec, donc il est garantie qu'une opération atomique ne sera pas interrompue par d'autres threads}
}

% to be proofreaded (begin)
\newglossaryentry{NaN}
{
  name=NaN,
  description={pas un nombre: un cas particulier pour les nombres à virgule flottante, indiquant généralement une erreur}
}

\newglossaryentry{NEON}
{
  name=NEON,
  description={\ac{AKA} \q{Advanced SIMD} --- \ac{SIMD} de ARM}
}

\newglossaryentry{reverse engineering}
{
  name=reverse engineering,
  description={action d'examiner et de comprendre comment quelque chose fonctionne, parfois dans le but de le reproduire}
}

\newglossaryentry{compiler intrinsic}
{
  name=compiler intrinsic,
  description={Une foncion spécifique d'un compilateur, qui n'est pas une fonction usuelle de bibliothèque.
    Le compilateur génère du code machine spécifique au lieu d'un appel à celui-ci.
    Souvent il s'agit d'une pseudo-fonction pour une instruction \ac{CPU} spécifique. Lire plus:
    (\myref{sec:compiler_intrinsic})}
}

\newglossaryentry{heap}
{
  name={tas},
  description={Généralement c'est un gros bout de mémoire fournit par l'\ac{OS} et utilisé par
  les applications pour le diviser comme elles le souhaitent. malloc()/free() fonctionnent en utilisant le tas}
}

\newglossaryentry{name mangling}
{
  name=name mangling,
  description={utilisé au moins en \Cpp, où le compilateur doit encoder le nom de la classe, la méthode et le type des arguments dans une chaîne,
  qui devient le nom interne de la fonction. Vous pouvez en lire plus à ce propos ici: \myref{namemangling}}
}

\newglossaryentry{xoring}
{
  name=xoring,
  description={souvent utilisé en anglais, qui signifie appliquer l'opération \ac{XOR}}
}

\newglossaryentry{security cookie}
{
  name=security cookie,
  description={Une valeur aléatoire, différente à chaque exécution. Vous pouvez en lire plus à ce propos ici: \myref{subsec:BO_protection}}
}

\newglossaryentry{tracer}
{
  name=tracer,
  description={Mon propre outil de debugging. Vous pouvez en lire plus à son propos ici: \myref{tracer}}
}

\newglossaryentry{GiB}
{
  name=GiB,
  description={Gibioctet: $2^{30}$ or 1024 mebioctets ou 1073741824 octets}
}

\newglossaryentry{CP/M}
{
  name=CP/M,
  description={Control Program for Microcomputers:
  un \ac{OS} de disque très basique utilisé avant MS-DOS}
}

\newglossaryentry{stack frame}
{
  name=stack frame,
  description={Une partie de la pile qui contient des informations spécifiques à la fonction courante:
  variables locales, arguments de la fonction, \ac{RA}, etc.}
}

\newglossaryentry{jump offset}
{
  name=jump offset,
  description={une partie de l'opcode de l'instruction JMP ou Jcc, qui doit être ajoutée à l'adresse de l'instruction suivante,
  et c'est ainsi que le nouveau \ac{PC} est calculé. Peut être négatif}
}

\newglossaryentry{integral type}
{
  name={type de donnée intégral},
  description={nombre usuel, mais pas un réel. peut être utilisé pour passer des variables de type booléen et des énumérations}
}

\newglossaryentry{real number}
{
  name={nombre réel},
  description={nombre qui peut contenir un point. comme \Tfloat et \Tdouble en \CCpp}
}

\newglossaryentry{PDB}
{
  name=PDB,
  description={(Win32) Fichier contenant des informations de débogage, en général seulement les noms des fonctions,
  mais aussi parfois les arguments des fonctions et le nom des variables locales}
}

\newglossaryentry{NTAPI}
{
  name=NTAPI,
  description={\ac{API} disponible seulement dans la série de Windows NT. Très peu documentée par Microsoft}
}

\newglossaryentry{stdout}
{
  name=stdout,
  description={standard output, sortie standard}
}

\newglossaryentry{word}
{
  name=word,
  description={Dans les ordinateurs plus vieux que les PCs, la taille de la mémoire était
  souvent mesurée en mots plutôt qu'en octet}
}

\newglossaryentry{arithmetic mean}
{
  name={moyenne arithmétique},
  description={la somme de toutes les valeurs, divisé par leur nombre}
}
\newglossaryentry{padding}
{
  name=padding,
  description={\emph{Padding} en anglais signifie rembourrer un oreiller, un matelas, etc. avec quelque chose afin de lui donner la forme désirée.
  En informatique, padding signifie ajouter des octets à un bloc, afin qu'il ait une certaine taille, comme $2^n$ octets.}
}

\newglossaryentry{basic block}
{
  name=basic block,
  description={
	un groupe d'instructions qui n'a pas d'instruction de saut/branchement,
	et donc n'as pas de saut de l'intérieur du bloc vers l'extérieur.
	Dans \IDA il ressemble à une liste d'instructions sans ligne vide}
}
}
\DE{\newglossaryentry{caller}
{
  name=caller,
  description={aufrufende Funktion}
}

\newglossaryentry{callee}
{
  name=callee,
  description={aufgerufene Funktion}
}

\newglossaryentry{stack pointer}
{
  name={Stapel-Zeiger},
  description={Ein Register das auf eine Stelle im Stack zeigt}
}

\newglossaryentry{product}
{
  name={Produkt},
  description={Ergebnis einer Multiplikation}
}

\newglossaryentry{GiB}
{
  name={GiB},
  description={\DEph{}}
}

\newglossaryentry{increment}
{
  name={\DEph},
  description={\DEph{}}
}

\newglossaryentry{decrement}
{
  name={\DEph},
  description={\DEph{}}
}

\newglossaryentry{stdout}
{
  name={\DEph},
  description={\DEph{}}
}

\newglossaryentry{endianness}
{
  name={\DEph},
  description={\DEph{}}
}

\newglossaryentry{thunk function}
{
  name={\DEph},
  description={\DEph{}}
}

\newglossaryentry{leaf function}
{
  name={\DEph},
  description={\DEph{}}
}

\newglossaryentry{heap}
{
  name={\DEph},
  description={\DEph{}}
}

\newglossaryentry{link register}
{
  name={\DEph},
  description={\DEph{}}
}

\newglossaryentry{anti-pattern}
{
  name={\DEph},
  description={\DEph{}}
}

\newglossaryentry{stack frame}
{
  name={\DEph},
  description={\DEph{}}
}

\newglossaryentry{jump offset}
{
  name={\DEph},
  description={\DEph{}}
}

\newglossaryentry{loop unwinding}
{
  name={\DEph},
  description={\DEph{}}
}

\newglossaryentry{tracer}
{
  name={\DEph},
  description={\DEph{}}
}

\newglossaryentry{register allocator}
{
  name={\DEph},
  description={\DEph{}}
}

\newglossaryentry{quotient}
{
  name={\DEph},
  description={\DEph{}}
}

\newglossaryentry{real number}
{
  name={\DEph},
  description={\DEph{}}
}

\newglossaryentry{NaN}
{
  name={\DEph},
  description={\DEph{}}
}

\newglossaryentry{Windows NT}
{
  name={\DEph},
  description={\DEph{}}
}

\newglossaryentry{word}
{
  name={\DEph},
  description={\DEph{}}
}

\newglossaryentry{PDB}
{
  name={\DEph},
  description={\DEph{}}
}

\newglossaryentry{name mangling}
{
  name={\DEph},
  description={\DEph{}}
}

\newglossaryentry{padding}
{
  name={\DEph},
  description={\DEph{}}
}

\newglossaryentry{NOP}
{
  name={\DEph},
  description={\DEph{}}
}

\newglossaryentry{POKE}
{
  name={\DEph},
  description={\DEph{}}
}

\newglossaryentry{xoring}
{
  name={\DEph},
  description={\DEph{}}
}

\newglossaryentry{atomic operation}
{
  name={\DEph},
  description={\DEph{}}
}

\newglossaryentry{basic block}
{
  name={\DEph},
  description={\DEph{}}
}

\newglossaryentry{reverse engineering}
{
  name={\DEph},
  description={\DEph{}}
}

\newglossaryentry{compiler intrinsic}
{
  name={\DEph},
  description={\DEph{}}
}


}
\IT{\newglossaryentry{caller}
{
  name=caller,
  description={\ITph{}}
}

\newglossaryentry{callee}
{
  name=callee,
  description={\ITph{}}
}

\newglossaryentry{stack pointer}
{
  name={\ITph{}},
  description={\ITph{}}
}

\newglossaryentry{product}
{
  name={\ITph{}},
  description={\ITph{}}
}

\newglossaryentry{GiB}
{
  name={GiB},
  description={\ITph{}}
}

\newglossaryentry{increment}
{
  name={\ITph},
  description={\ITph{}}
}

\newglossaryentry{decrement}
{
  name={\ITph},
  description={\ITph{}}
}

\newglossaryentry{stdout}
{
  name={\ITph},
  description={\ITph{}}
}

\newglossaryentry{endianness}
{
  name={\ITph},
  description={\ITph{}}
}

\newglossaryentry{thunk function}
{
  name={\ITph},
  description={\ITph{}}
}

\newglossaryentry{leaf function}
{
  name={\ITph},
  description={\ITph{}}
}

\newglossaryentry{heap}
{
  name={\ITph},
  description={\ITph{}}
}

\newglossaryentry{link register}
{
  name={\ITph},
  description={\ITph{}}
}

\newglossaryentry{anti-pattern}
{
  name={\ITph},
  description={\ITph{}}
}

\newglossaryentry{stack frame}
{
  name={\ITph},
  description={\ITph{}}
}

\newglossaryentry{jump offset}
{
  name={\ITph},
  description={\ITph{}}
}

\newglossaryentry{loop unwinding}
{
  name={\ITph},
  description={\ITph{}}
}

\newglossaryentry{tracer}
{
  name={\ITph},
  description={\ITph{}}
}

\newglossaryentry{register allocator}
{
  name={\ITph},
  description={\ITph{}}
}

\newglossaryentry{quotient}
{
  name={\ITph},
  description={\ITph{}}
}

\newglossaryentry{real number}
{
  name={\ITph},
  description={\ITph{}}
}

\newglossaryentry{NaN}
{
  name={\ITph},
  description={\ITph{}}
}

\newglossaryentry{Windows NT}
{
  name={\ITph},
  description={\ITph{}}
}

\newglossaryentry{word}
{
  name={\ITph},
  description={\ITph{}}
}

\newglossaryentry{PDB}
{
  name={\ITph},
  description={\ITph{}}
}

\newglossaryentry{name mangling}
{
  name={\ITph},
  description={\ITph{}}
}

\newglossaryentry{padding}
{
  name={\ITph},
  description={\ITph{}}
}

\newglossaryentry{NOP}
{
  name={\ITph},
  description={\ITph{}}
}

\newglossaryentry{POKE}
{
  name={\ITph},
  description={\ITph{}}
}

\newglossaryentry{xoring}
{
  name={\ITph},
  description={\ITph{}}
}

\newglossaryentry{atomic operation}
{
  name={\ITph},
  description={\ITph{}}
}

\newglossaryentry{basic block}
{
  name={\ITph},
  description={\ITph{}}
}

\newglossaryentry{reverse engineering}
{
  name={\ITph},
  description={\ITph{}}
}

\newglossaryentry{compiler intrinsic}
{
  name={\ITph},
  description={\ITph{}}
}


}
\JA{\newglossaryentry{tail call}
{
  name=tail call,
  description={コンパイラ（またはインタプリタ）が再帰（\emph{末尾再帰}が可能なとき）を効率のため、反復に変換するときです}
}

\newglossaryentry{endianness}
{
  name=endianness,
  description={バイトオーダー: \myref{sec:endianness}}
}

\newglossaryentry{caller}
{
  name=caller,
  description={呼び出し元の関数}
}

\newglossaryentry{callee}
{
  name=callee,
  description={呼び出された関数}
}

\newglossaryentry{debuggee}
{
  name=debuggee,
  description={デバッグされるプログラム}
}

\newglossaryentry{leaf function}
{
  name=leaf function,
  description={他の関数から呼び出されない関数}
}

\newglossaryentry{link register}
{
  name=link register,
  description=(RISC) リターンアドレスが保存されるレジスタ。これはleaf functionをスタックを使わずに呼び出すのを可能にする
}

\newglossaryentry{anti-pattern}
{
  name=anti-pattern,
  description=
  {一般に、よくないと考えられるやり方}
}

\newglossaryentry{stack pointer}
{
  name={スタックポインタ},
  description=
  {スタックの場所を示すレジスタ}
}

\newglossaryentry{decrement}
{
  name={デクリメント},
  description={1の減算}
}

\newglossaryentry{increment}
{
  name={インクリメント},
  description={1の加算}
}

\newglossaryentry{loop unwinding}
{
  name=loop unwinding,
  description={
  $n$回のイテレーションのループコードをコンパイラが生成する代わりに、ループボディを$n$回コピーするコードを生成する。
  ループに使用する命令を削除するため}
}

\newglossaryentry{register allocator}
{
  name=register allocator,
  description=
  {CPUレジスタをローカル変数に割り当てるコンパイラの機構}
}

\newglossaryentry{quotient}
{
  name=商,
  description={除算結果}
}

\newglossaryentry{product}
{
  name=積,
  description={乗算結果}
}

\newglossaryentry{NOP}
{
  name=NOP,
  description={何もしない命令}
}

\newglossaryentry{POKE}
{
  name=POKE,
  description={特定のアドレスにバイトを書き込むための基本的な言語命令}
}

\newglossaryentry{keygenme}
{
  name=keygenme,
  description={キー/ライセンスジェネレータのようなソフトウェア保護を模倣するプログラム}
} % TODO clarify: A software which generate key/license value to bypass sotfware protection?

\newglossaryentry{dongle}
{
  name=dongle,
  description={LPTプリンタポートやUSBにさす小さなハードウェア。セキュリティトークンと似て、メモリを持ち、時には秘密の(暗号学的)ハッシュアルゴリズムを持つ}
}

\newglossaryentry{thunk function}
{
  name=thunk function,
  description={単一の役割だけ持つ小さな関数：他の関数を呼び出す等}
}

\newglossaryentry{user mode}
{
  name=user mode,
  description={アプリケーションソフトウェアが実行される制限されたCPUモード \gls{kernel mode}}
}

\newglossaryentry{kernel mode}
{
  name=kernel mode,
  description={OSカーネルやドライバが実行される制限のないCPUモード \gls{user mode}}
}

\newglossaryentry{Windows NT}
{
  name=Windows NT,
  description={Windows NT, 2000, XP, Vista, 7, 8, 10}
}

\newglossaryentry{atomic operation}
{
  name=atomic operation,
  description={
  \q{$\alpha{}\tau{}o\mu{}o\varsigma{}$}
  %\q{atomic}
  ギリシャ語で分割不可能を意味し、処理は他のスレッドに割り込まれないことが保証される
  }
}

% to be proofreaded (begin)
\newglossaryentry{NaN}
{
  name=NaN,
  description=
    {非数：float型の数の特殊なケースで、エラーが通知される}
}

\newglossaryentry{basic block}
{
  name=basic block,
  description=
    {jump/branch命令を持たない命令群で、ブロックの外側から内側へのjumpも持たない}
}

\newglossaryentry{NEON}
{
  name=NEON,
  description={\ac{AKA} \q{Advanced SIMD} --- \ac{SIMD} 拡張、ARM における}
}

\newglossaryentry{reverse engineering}
{
  name=reverse engineering,
  description={時にはクローンするため、どうやって動いているのかを理解しようとする行為}
}

\newglossaryentry{compiler intrinsic}
{
  name=compiler intrinsic,
  description={
    通常のライブラリ関数ではない、コンパイラ特有の関数。コンパイラは特定の機械語を生成する。しばしば、特定のCPU命令のための疑似関数 (\myref{sec:compiler_intrinsic})
  }
}

\newglossaryentry{heap}
{
  name=ヒープ,
  description={\ac{OS}が提供する大きなメモリの塊のことで、アプリケーションが好きなように分割することができる。malloc()/free()を呼び出して使用する}
}

\newglossaryentry{name mangling}
{
  name=name mangling,
  description={コンパイラがクラス、メソッド、および引数型の名前を1つの文字列にエンコードする必要がある少なくとも \CCpp で使用され、
  関数の内部名になります。 以下で詳細を読むことができます: \myref{namemangling}}
}

\newglossaryentry{xoring}
{
  name=xoring,
  description={英語圏でしばしばみられ、\ac{XOR}操作を適用する意味になる}
}

\newglossaryentry{security cookie}
{
  name=security cookie,
  description={
  ランダムな値で、実行の度に異なった値になります。以下で詳細を読むことができます: \myref{subsec:BO_protection}}
}

\newglossaryentry{tracer}
{
  name=tracer,
  description={シンプルなデバッグツールです。以下で詳細を読むことができます: \myref{tracer}}
}

\newglossaryentry{GiB}
{
  name=GiB,
  description={ギガバイト：$2^{30}$ または1024メガバイトまたは1073741824バイト}
}

\newglossaryentry{CP/M}
{
  name=CP/M,
  description={Control Program for Microcomputers: 以前使用されていたとても基本的なディスク\ac{OS}です MS-DOS}
}

\newglossaryentry{stack frame}
{
  name=stack frame,
  description=
  {現在の関数に固有の情報（ローカル変数、関数の引数、\ac{RA}など）を含むスタックの一部}
}

\newglossaryentry{jump offset}
{
  name=jump offset,
  description=
  {JMP命令またはJcc命令のオペコードの一部を次の命令のアドレスに追加する必要があります。これが新しい\ac{PC}の計算方法です。 負となる場合もあります}
}

\newglossaryentry{integral type}
{
  name=整数型,
  description=
  {通常の数字ですが、実際の数字はありません。 ブールと列挙型の変数を渡すために使用できます}
}

\newglossaryentry{real number}
{
  name=実数,
  description={
  小数点以下を含む可能性のある数字。 これは \CCpp で \Tfloat と \Tdouble です
  }
}

\newglossaryentry{PDB}
{
  name=PDB,
  description={(Win32) 
  デバッグ情報ファイルで、通常は関数名だけではなく、関数の引数とローカル変数名を含む}
}

\newglossaryentry{NTAPI}
{
  name=NTAPI,
  description=
  {APIはWindows NT系列でのみ使用できます。 Microsoftは大部分で文書化していません}
}

\newglossaryentry{stdout}
{
  name=stdout,
  description={standard output}
}

\newglossaryentry{word}
{
  name=word,
  description=
  {PCよりも昔のコンピュータでは、メモリサイズはバイトではなくワードで測定されることがしばしばでした}
}

\newglossaryentry{arithmetic mean}
{
  name={算術平均},
  description=
  {すべての値の合計を数で割った値}
}
\newglossaryentry{padding}
{
  name=padding,
  description=
  {英語で\emph{パディング}とは、（より大きな）望ましい形状にするために、あるものを枕に詰めることを意味します。
  コンピュータサイエンスでは、パディングとはブロックにバイトを追加することで、$2^n$バイトのようなサイズにすることを意味します。}
}

}

\makeglossaries

\hypersetup{
    colorlinks=true,
    allcolors=blue,
    pdfauthor={\AUTHOR},
    pdftitle={\TitleMain}
    }

%\ifdefined\RUSSIAN
\newcommand{\LstStyle}{\ttfamily\small}
%\else
%\newcommand{\LstStyle}{\ttfamily}
%\fi

% inspired by http://prismjs.com/
\definecolor{digits}{RGB}{0,0,0}
\definecolor{bg}{RGB}{255,255,255}
%\definecolor{bg}{RGB}{255,252,250}
\definecolor{col1}{RGB}{154,20,150}
\definecolor{col2}{RGB}{112,128,144}
\definecolor{col3}{RGB}{10,120,180}
\definecolor{col4}{RGB}{106,164,108}


\lstset{
    %backgroundcolor=\color{lstbgcolor},
    %backgroundcolor=\color{light-gray},
    backgroundcolor=\color{bg},
    basicstyle=\LstStyle,
    breaklines=true,
    %prebreak=\raisebox{0ex}[0ex][0ex]{->},
    %postbreak=\raisebox{0ex}[0ex][0ex]{->},
    prebreak=\raisebox{0ex}[0ex][0ex]{\ensuremath{\rhookswarrow}},
    postbreak=\raisebox{0ex}[0ex][0ex]{\ensuremath{\rcurvearrowse\space}},
    frame=single,
    columns=fullflexible,keepspaces,
    escapeinside=§§,
    inputencoding=utf8
}

\definecolor{digits}{RGB}{0,0,0}
\definecolor{colkw}{RGB}{0,0,0}

\definecolor{colstr}{RGB}{0,0,0}
\definecolor{colid}{RGB}{0,0,0}

\definecolor{colcomment}{RGB}{130,130,130}

\definecolor{colbg}{RGB}{255,255,255}

\lstdefinestyle{custompy}{
  texcl=true,
  language=Python,
  showstringspaces=false,
  backgroundcolor=\color{colbg},
  keywordstyle=\color{colkw},
  commentstyle=\color{colcomment},
  identifierstyle=\color{colid},
  stringstyle=\color{colstr}
}

\lstdefinestyle{customc}{
  texcl=true,
  language=C,
  showstringspaces=false,
  backgroundcolor=\color{colbg},
  keywordstyle=\color{colkw},
  commentstyle=\color{colcomment},
  identifierstyle=\color{colid},
  stringstyle=\color{colstr}
}

\lstdefinestyle{custommath}{
  texcl=true,
  language=Mathematica,
  showstringspaces=false,
  backgroundcolor=\color{colbg},
  keywordstyle=\color{colkw},
  commentstyle=\color{colcomment},
  identifierstyle=\color{colid},
  stringstyle=\color{colstr}
}

\lstdefinestyle{customjava}{
  texcl=true,
  language=Java,
  showstringspaces=false,
  backgroundcolor=\color{colbg},
  keywordstyle=\color{colkw},
  commentstyle=\color{colcomment},
  identifierstyle=\color{colid},
  stringstyle=\color{colstr}
}

\lstdefinestyle{customasmx86}{
  texcl=true,
  morecomment=[l]{;},
  morestring=[b]",
  morestring=[d]',
  showstringspaces=false,
  backgroundcolor=\color{colbg},
  keywordstyle=\color{colkw},
  commentstyle=\color{colcomment},
  identifierstyle=\color{colid},
  stringstyle=\color{colstr}
}

\lstdefinestyle{customasmMIPS}{
  texcl=true,
  morecomment=[l]{;},
  morecomment=[l]{\#},
  morestring=[b]",
  morestring=[d]',
  showstringspaces=false,
  backgroundcolor=\color{colbg},
  keywordstyle=\color{colkw},
  commentstyle=\color{colcomment},
  identifierstyle=\color{colid},
  stringstyle=\color{colstr}
}

\lstdefinestyle{customasmARM}{
  texcl=true,
  morecomment=[l]{;},
  morestring=[b]",
  morestring=[d]',
  showstringspaces=false,
  backgroundcolor=\color{colbg},
  keywordstyle=\color{colkw},
  commentstyle=\color{colcomment},
  identifierstyle=\color{colid},
  stringstyle=\color{colstr}
}

\lstdefinestyle{customasmPPC}{
  texcl=true,
  morecomment=[l]{;},
  morecomment=[l]{\#},
  morestring=[b]",
  morestring=[d]',
  showstringspaces=false,
  backgroundcolor=\color{colbg},
  keywordstyle=\color{colkw},
  commentstyle=\color{colcomment},
  identifierstyle=\color{colid},
  stringstyle=\color{colstr}
}



\ifdefined\RUSSIAN
\renewcommand\lstlistingname{Листинг}
\renewcommand\lstlistlistingname{Листинг}
\fi

\DeclareMathSizes{12}{30}{16}{12}%

% see also:
% http://tex.stackexchange.com/questions/129225/how-can-i-get-get-makeindex-to-ignore-capital-letters
% http://tex.stackexchange.com/questions/18336/correct-sorting-of-index-entries-containing-macros
\def\myindex#1{\expandafter\index\expandafter{#1}}

\begin{document}

% fancyhdr
\pagestyle{fancy}
\setlength{\headheight}{13pt}
\fancyhead[R]{} % suppress chapter name

\VerbatimFootnotes

\frontmatter

%\RU{% To translators: don't bother to translate this... english-only version.

\begin{center}
\LARGE{} Это моя собственная доска объявления \normalsize{}
\end{center}

Мои дорогие читатели! Время от времени, у меня появляются вопросы, и я не знаю, кого (или где) спросить.
Или я просто ленив...
Поможете мне?

\myhrule{}

Что использовать на нерутованном Андроиде как SSH-сервер, который умеет писать на внешнюю micro-SD?
SSHDroid нормально работал, но теперь он устарел...
Например, когда Total Commander пишет на micro-SD, появляется окно с вопросом, разрешить или нет.
В случае с виденными мною SSH-серверами этого что-то не происходит.

\myhrule{}

Win32-процесс А запущен.
Процесс Б аттачится к нему как отладчик, либо открывает его используя OpenProcess().
ReadProcessMemory() работает, но не работает, если пытается читать незакоммиченные страницы процесса А.

Проблема: как заставить менеджер памяти Windows закоммитить страницу в процессе А из userland-а процесса Б?
Я могу всунуть в процесс А инструкцию чтения, запустить её, страница закоммитится, но это не решение.

\myhrule{}

Если знаете что-то, пожалуйста помогите мне: \EMAIL{}

}
%\EN{% To translators: don't bother to translate this... english-only version.

\begin{center}
\LARGE{} This is my own bulletin board \normalsize{}
\end{center}

My dear readers! From time to time, I have questions, I don't know who (or where) to ask.
Or I just lazy...
Can you please help me?

\myhrule{}

What do you use on unrooted Android device as a SSH server, that can write to external micro-SD card?
I was happy with SSHDroid, but now it's outdated...
For instance, when Total Commander writes to micro-SD, a dialog box appears, whether to allow this or not.
This is not a case with SSH-servers that I saw.

\myhrule{}

A win32 process A is running.
Process B is attaching to it as a debugger, or opens it using OpenProcess().
ReadProcessMemory() works OK, but fails if it tries to read uncommited memory pages of process A.

The problem: how to force the Windows Memory Manager to commit a page in process A from userland of process B?
I can inject a read instruction into process A, run it, and the page would be committed, but this is not the solution.

\myhrule{}

If you know something, please help me: \EMAIL{}

}
%\DE{\include{1st_page_DE}}
%\ES{\include{1st_page_ES}}
%\CN{\include{1st_page_CN}}

\RU{\begin{titlepage}
\vspace*{\fill}

\begin{flushright}
\emph{Эта книга посвящается \\
компьютеру Intertec Superbrain II}
\end{flushright}

\vspace*{\fill}
\end{titlepage}

}
\EN{\begin{titlepage}
\vspace*{\fill}

\begin{flushright}
\emph{This book is dedicated to \\
Intertec Superbrain II computer}

\end{flushright}

\vspace*{\fill}
\end{titlepage}

}
\FR{\begin{titlepage}
\vspace*{\fill}

\begin{flushright}
\emph{Ce livre est dédié à \\
l'ordinateur Intertec Superbrain II}
\end{flushright}

\vspace*{\fill}
\end{titlepage}

}
\JA{\begin{titlepage}
\vspace*{\fill}

\begin{flushright}
\emph{この本は \\
Intertec Superbrain II computer専用です}

\end{flushright}

\vspace*{\fill}
\end{titlepage}

}
\DE{\begin{titlepage}
\vspace*{\fill}

\begin{flushright}
\emph{Diese Buch ist gewidmet: \\
Intertec Superbrain II computer}

\end{flushright}

\vspace*{\fill}
\end{titlepage}

}
\IT{\begin{titlepage}
\vspace*{\fill}

\begin{flushright}
\emph{Questo libro è dedicato al \\
Intertec Superbrain II computer}

\end{flushright}

\vspace*{\fill}
\end{titlepage}
}

\EN{\begin{titlepage}

\input{cover}
\end{titlepage}

\newpage

\begin{center}
\vspace*{\fill}
{\LARGE \TitleMain}

\bigskip

{\large (\TitleAux)}

\bigskip
\bigskip
Why two titles? Read here: \myref{TwoTitles}.

\vspace*{\fill}

{\large \AUTHOR}

{\large \TT{<\EMAIL>}}
\vspace*{\fill}
\vfill

\ccbysa

\textcopyright 2013-2019, \AUTHOR. 

This work is licensed under the Creative Commons Attribution-ShareAlike 4.0 International (CC BY-SA 4.0) license.
To view a copy of this license, visit \url{https://creativecommons.org/licenses/by-sa/4.0/}.

Text version ({\large \today}).

The latest version (and Russian edition) of this text is accessible at \href{http://go.yurichev.com/17009}{beginners.re}.

The cover was made by Andy Nechaevsky: \href{http://go.yurichev.com/17023}{facebook}.

\end{center}
}
\RU{\begin{titlepage}

\input{cover}
\end{titlepage}

\newpage

\begin{center}
\vspace*{\fill}
{\LARGE \TitleMain}

\bigskip

{\large (\TitleAux)}

\bigskip
\bigskip
Почему два названия? Читайте здесь: \myref{TwoTitles}.

\vspace*{\fill}

{\large \AUTHOR}

{\large \TT{<\EMAIL>}}
\vspace*{\fill}
\vfill

\ccbysa

\textcopyright 2013-2019, \AUTHOR. 

Это произведение доступно по лицензии Creative Commons «Attribution-ShareAlike 4.0 International» (CC BY-SA 4.0).
Чтобы увидеть копию этой лицензии, посетите \url{https://creativecommons.org/licenses/by-sa/4.0/}.

Версия этого текста ({\large \today}).

Самая новая версия текста (а также англоязычная версия) доступна на сайте \href{http://go.yurichev.com/17009}{beginners.re}.

Обложка нарисована Андреем Нечаевским: \href{http://go.yurichev.com/17023}{facebook}.

\end{center}
}
\ES{\begin{titlepage}

\input{cover}
\end{titlepage}

\newpage

\begin{center}
\vspace*{\fill}
{\LARGE \TitleMain}

\bigskip

{\large (\TitleAux)}

\bigskip
\bigskip
Why two titles? Read here: \myref{TwoTitles}. % TBT

\vspace*{\fill}

{\large \AUTHOR}

{\large \TT{<\EMAIL>}}
\vspace*{\fill}
\vfill

\ccbysa

\textcopyright 2013-2019, \AUTHOR. 

Esta obra est\'a bajo una Licencia Creative Commons ``Attribution-ShareAlike 4.0 International'' (CC BY-SA 4.0)
Para ver una copia de esta licencia, visita \url{https://creativecommons.org/licenses/by-sa/4.0/}.

Versi\'on del texto ({\large \today}).

La \'ultima versi\'on (as\'i como las versiones en ingl\'es y ruso) de este texto est\'a disponible en
\href{http://go.yurichev.com/17009}{beginners.re}.

La portada fue hecha por Andy Nechaevsky: \href{http://go.yurichev.com/17023}{facebook}.

\end{center}
}
\IT{\begin{titlepage}

\input{cover}
\end{titlepage}

\newpage

\begin{center}
\vspace*{\fill}
{\LARGE \TitleMain}

\bigskip

{\large (\TitleAux)}

\bigskip
\bigskip
Perchè due titoli? Leggi qua: \myref{TwoTitles}.

\vspace*{\fill}

{\large \AUTHOR}

{\large \TT{<\EMAIL>}}
\vspace*{\fill}
\vfill

\ccbysa

\textcopyright 2013-2019, \AUTHOR.

Questo lavoro è rilasciato sotto licenza Creative Commons Attribution-ShareAlike 4.0 International (CC BY-SA 4.0).
Per vedere la copia di questa licenza visita il sito \url{https://creativecommons.org/licenses/by-sa/4.0/}.

Teston ({\large \today}).

L'ultima versione (e l'edizione Russa) del testo sono accessibili al seguente sito: \href{http://go.yurichev.com/17009}{beginners.re}.

La copertina è stata creata da Andy Nechaevsky: \href{http://go.yurichev.com/17023}{facebook}.

\end{center}
}
\FR{\begin{titlepage}

\input{cover}
\end{titlepage}

\newpage

\begin{center}
\vspace*{\fill}
{\LARGE \TitleMain}

\bigskip

{\large (\TitleAux)}

\bigskip
\bigskip
Pourquoi deus titres? Lire ici: \myref{TwoTitles}.

\vspace*{\fill}

{\large \AUTHOR}

{\large \TT{<\EMAIL>}}
\vspace*{\fill}
\vfill

\ccbysa

\textcopyright 2013-2019, \AUTHOR. 

Ce travail est sous licence Creative Commons Attribution-ShareAlike 4.0 International (CC BY-SA 4.0).
Pour voir une copie de cette licence, rendez vous sur \url{https://creativecommons.org/licenses/by-sa/4.0/}.

Version du texte ({\large \today}).

La dernière version (et édition en russe) de ce texte est accessible sur \href{http://go.yurichev.com/17009}{beginners.re}.

La couverture a été réalisée par Andy Nechaevsky: \href{http://go.yurichev.com/17023}{facebook}.

\end{center}
}
\DE{\begin{titlepage}

\input{cover}
\end{titlepage}

\newpage

\begin{center}
\vspace*{\fill}
{\LARGE \TitleMain}

\bigskip

{\large (\TitleAux)}

\bigskip
\bigskip
Why two titles? Read here: \myref{TwoTitles}. % TBT

\vspace*{\fill}

{\large \AUTHOR}

{\large \TT{<\EMAIL>}}
\vspace*{\fill}
\vfill

\ccbysa

\textcopyright 2013-2019, \AUTHOR. 

Diese Arbeit ist lizenziert unter der Creative Commons Attribution-ShareAlike 4.0 International (CC BY-SA 4.0) Lizenz.
Um eine Kopie dieser Lizenz zu lesen, besuchen Sie \url{https://creativecommons.org/licenses/by-sa/4.0/}.

Text-Version ({\large \today}).

Die aktuellste Version (und eine russiche Ausgabe) dieses Textes ist auf \href{http://go.yurichev.com/17009}{beginners.re} verfügbar.
Eine Version für den E-Book-Leser ist dort ebenfalls erhältlich.

Das Cover wurde erstellt von Andy Nechaevsky: \href{http://go.yurichev.com/17023}{Facebook}.

\end{center}
}
\JA{\begin{titlepage}

\input{cover}
\end{titlepage}

\newpage

\begin{center}
\vspace*{\fill}
{\LARGE \TitleMain}

\bigskip

{\large (\TitleAux)}

\bigskip
\bigskip
どうしてタイトルが2つあるの？ 参照: \myref{TwoTitles}

\vspace*{\fill}

{\large \AUTHOR}

{\large \TT{<\EMAIL>}}
\vspace*{\fill}
\vfill

\ccbysa

\textcopyright 2013-2019, \AUTHOR. 

この作品は、クリエイティブ・コモンズの表示 - 継承 4.0 国際(CC BY-SA 4.0)の下でライセンスされています。
このライセンスのコピーを表示するには、\url{https://creativecommons.org/licenses/by-sa/4.0/}を訪れてください。

Text version ({\large \today}).

この本の最新版(そしてロシア版)は以下で見ることができます: \href{http://go.yurichev.com/17009}{beginners.re}

カバーはAndy Nechaevskyが制作しました: \href{http://go.yurichev.com/17023}{facebook}

\end{center}
}

\EN{\vspace*{\fill}

\Huge Call for translators!

\normalsize

\bigskip
\bigskip
\bigskip

You may want to help me with translating this work into languages other than English and Russian.
Just send me any piece of translated text (no matter how short) and I'll put it into my LaTeX source code.

\href{https://github.com/DennisYurichev/RE-for-beginners/blob/master/Translation.md}{Read here}.

We already have something in \href{https://beginners.re/RE4B-DE.pdf}{German},
\href{https://beginners.re/RE4B-FR.pdf}{French}, a bit in
\href{https://yurichev.com/tmp/RE4B-IT-partial.pdf}{Italian},
\href{https://yurichev.com/tmp/RE4B-PTBR-lite2.pdf}{Portuguese} and
\href{https://yurichev.com/tmp/RE4B-PL.pdf}{Polish}.

% Add new languages?

Speed isn't important, because this is an open-source project, after all.
Your name will be mentioned as a project contributor.
Korean, Chinese, and Persian languages are reserved by publishers.
English and Russian versions I do by myself, but my English is still that horrible, so I'm very grateful for any notes about grammar, etc.
Even my Russian is flawed, so I'm grateful for notes about Russian text as well!%

So do not hesitate to contact me: \GTT{\EMAIL}.

\vspace*{\fill}
\vfill
}
\RU{\input{call_for_translators_RU}}
\DE{\input{call_for_translators_DE}}
\FR{\input{call_for_translators_FR}}
\IT{\vspace*{\fill}

\Huge Cerchiamo traduttori!

\normalsize

\bigskip
\bigskip
\bigskip

Puoi aiutare a tradurre questo progetto in linguaggi diversi dall'Inglese ed il Russo.
Basta inviarmi un qualsiasi pezzo di testo tradotto (non importa quanto è lungo) e lo aggiungerò al mio codice sorgente scritto in LaTeX.

\href{https://github.com/DennisYurichev/RE-for-beginners/blob/master/Translation.md}{Leggi qui}.

Abbiamo già traduzioni in \href{https://beginners.re/RE4B-DE.pdf}{Tedesco},
\href{https://beginners.re/RE4B-FR.pdf}{Francese}, e qualcosa in
\href{https://yurichev.com/tmp/RE4B-IT-partial.pdf}{Italiano},
\href{https://yurichev.com/tmp/RE4B-PTBR-lite2.pdf}{Portoghese} e
\href{https://yurichev.com/tmp/RE4B-PL.pdf}{Polacco}.

La velocità non è importante, perchè è un progetto Open Source, dopo tutto.
Il tuo nome sarà menzionato come Contributore del progetto.
Coreano, Cinese e Persiano sono linguaggi reservati ai publisher.
Le versioni in Inglese e Russo le traduco da solo, ma il mio Inglese è ancora terribile, quindi vi sono grato ad ogni correzione sulla mia grammatica, ecc...
Anche il mio Russo non è ancora perfetto, quindi sono felice di ricevere correzioni anche per il Russo!%

Non esitate a contattarmi: \GTT{\EMAIL}.

\vspace*{\fill}
\vfill
}
\PL{\input{call_for_translators_PL}}
\JA{\vspace*{\fill}

\Huge 翻訳者求む！

\normalsize

\bigskip
\bigskip
\bigskip

この作品を英語とロシア語以外の言語に翻訳するのを手伝ってください。
どのように翻訳されたテキストを私に送っても（どれほど短くても）、私はLaTeXのソースコードに入れます。

\href{https://github.com/DennisYurichev/RE-for-beginners/blob/master/Translation.md}{ここを読んでください}.

いくつかは既にあります。\href{https://beginners.re/RE4B-DE.pdf}{ドイツ語},
\href{https://beginners.re/RE4B-FR.pdf}{フランス語}, 少しですが
\href{https://yurichev.com/tmp/RE4B-IT-partial.pdf}{イタリア語},
\href{https://yurichev.com/tmp/RE4B-PTBR-lite2.pdf}{ポルトガル語} と
\href{https://yurichev.com/tmp/RE4B-PL.pdf}{ポーランド語}です。

スピードは重要ではありません。なぜなら、これはオープンソースプロジェクトなのですから。
あなたの名前はプロジェクト寄稿者として言及されます。
韓国語、中国語、ペルシャ語は出版社によって予約されています。
英語とロシア語のバージョンは自分でやっていますが、私の英語はまだひどいので、文法などに関するメモにはとても感謝しています。
私のロシア語にも欠陥があるので、ロシア語のテキストについての注釈にも感謝しています！

だから私に連絡するのをためらうことはありません: \GTT{\EMAIL}

\vspace*{\fill}
\vfill
}



\shorttoc{%
    \RU{Краткое оглавление}%
    \EN{Abridged contents}%
    \ES{Contenidos abreviados}%
    \PTBRph{}%
    \DE{Inhaltsverzeichnis (gekürzt)}%
    \PLph{}%
    \IT{Sommario}%
    \THAph{}\NLph{}%
    \FR{Contenus abrégés}%
    \JA{簡略版}
    \TR{İçindekiler}
}{0}

\tableofcontents
\cleardoublepage

\cleardoublepage
\EN{\section*{Preface}

\subsection*{What is with two titles?}
\label{TwoTitles}

The book was named ``Reverse Engineering for Beginners'' in 2014-2018, but I always suspected this makes readership too narrow.

Infosec people know about ``reverse engineering'', but I've rarely hear the ``assembler'' word from them.

Likewise, the ``reverse engineering'' term is somewhat cryptic to a general audience of programmers, but they know about ``assembler''.

In July 2018, as an experiment, I've changed the title to ``Assembly Language for Beginners''
and posted the link to Hacker News website\footnote{\url{https://news.ycombinator.com/item?id=17549050}}, and the book was received generally well.

So let it be, the book now has two titles.

However, I've changed the second title to ``Understanding Assembly Language'', because someone already written ``Assembly Language for Beginners'' book.
Also, people say ``for Beginners'' sounds a bit sarcastic for the book of \textasciitilde{}1000 pages.

Two books differ only by title, filename (UAL-XX.pdf versus RE4B-XX.pdf), URL and couple of first pages.

\subsection*{About reverse engineering}

There are several popular meanings of the term \q{\gls{reverse engineering}}:

1) The reverse engineering of software; researching compiled programs

2) The scanning of 3D structures and the subsequent digital manipulation required in order to duplicate them

3) Recreating \ac{DBMS} structure

This book is about the first meaning.

\subsection*{Prerequisites}

Basic knowledge of the C \ac{PL}.
Recommended reading: \myref{CCppBooks}.

\subsection*{Exercises and tasks}

\dots
can be found at: \url{http://challenges.re}.

\subsection*{About the author}
\begin{tabularx}{\textwidth}{ l X }

\raisebox{-\totalheight}{
\includegraphics[scale=0.60]{Dennis_Yurichev.jpg}
}

&
Dennis Yurichev is an experienced reverse engineer and programmer.
He can be contacted by email: \textbf{\EMAIL{}}. % or Skype: \textbf{dennis.yurichev}.

% FIXME: no link. \tablefootnote doesn't work
\end{tabularx}

% subsections:
\input{praise}
\input{thanks}
\input{FAQ_EN}

\subsection*{How to learn programming}

Many people keep asking about it.

There is no ``royal road'', but there are quite efficient ways.

From my own experience, this is just: solving exercises from:

\begin{itemize}
\item \KRBook
\item Harold Abelson, Gerald Jay Sussman, Julie Sussman -- Structure and Interpretation of Computer Programs
\item \TAOCP
\item Niklaus Wirth's books
\item \RobPikePractice
\end{itemize}

... in pure C and LISP.
You may never use these programming languages in future at all.
Almost all commercial programmers don’t. But C and LISP coding experience will help enormously in long run.

Also, you can skip reading books themselves,
just skim them whenever you feel you need to understand something you missing for the exercise you currently solving.

This may take years at best, or a lifetime, but still this is way faster than to rush between fads.

The success of these books probably related to the fact that their authors are teachers
and all this material has been honed on students first.

As of LISP, I personally would recommend Racket (Scheme dialect). But this is matter of taste, anyway.

Some people say assembly language understanding is also very helpful, even if you will never use it.
This is true.
But this is a way for the most dedicated geeks, and it can be postponed at start.

Also, self-taught people (including author of these lines) often has the problem of trying too hard on hard problems,
skipping easy ones.
This is a great mistake.
Compare to sport or music -- no one starts at 100kg weights, or Paganini's Caprices.
I would say -- you can try to tackle a problem if you can outline its solution in your mind.

\input{knuth_baby_steps}
( Donald E. Knuth -- \url{https://theorydish.blog/2018/02/01/donald-knuth-on-doing-research/} )

Good luck!

\subsection*{About the Korean translation}

In January 2015, the Acorn publishing company (\href{http://www.acornpub.co.kr}{www.acornpub.co.kr}) in South Korea did a huge amount of work in translating and publishing
this book (as it was in August 2014) into Korean.

It's available now at \href{http://go.yurichev.com/17343}{their website}.

\iffalse
\begin{figure}[H]
\centering
\includegraphics[scale=0.3]{acorn_cover.jpg}
\end{figure}
\fi

The translator is Byungho Min (\href{http://go.yurichev.com/17344}{twitter/tais9}).
The cover art was done by the artistic Andy Nechaevsky, a friend of the author:
\href{http://go.yurichev.com/17023}{facebook/andydinka}.
Acorn also holds the copyright to the Korean translation.

So, if you want to have a \emph{real} book on your shelf in Korean and
want to support this work, it is now available for purchase.

\subsection*{About the Persian/Farsi translation}

In 2016 the book was translated by Mohsen Mostafa Jokar (who is also known to Iranian community for his translation of Radare manual\footnote{\url{http://rada.re/get/radare2book-persian.pdf}}).
It is available on the publisher’s website\footnote{\url{http://goo.gl/2Tzx0H}} (Pendare Pars).

Here is a link to a 40-page excerpt: \url{https://beginners.re/farsi.pdf}.

National Library of Iran registration information: \url{http://opac.nlai.ir/opac-prod/bibliographic/4473995}.

\subsection*{About the Chinese translation}

In April 2017, translation to Chinese was completed by Chinese PTPress. They are also the Chinese translation copyright holders.

 The Chinese version is available for order here: \url{http://www.epubit.com.cn/book/details/4174}. A partial review and history behind the translation can be found here: \url{http://www.cptoday.cn/news/detail/3155}.

The principal translator is Archer, to whom the author owes very much. He was extremely meticulous (in a good sense) and reported most of the known mistakes and bugs, which is very important in literature such as this book.
The author would recommend his services to any other author!

The guys from \href{http://www.antiy.net/}{Antiy Labs} has also helped with translation. \href{http://www.epubit.com.cn/book/onlinechapter/51413}{Here is preface} written by them.

}
\RU{\section*{Предисловие}

\subsection*{Почему два названия?}
\label{TwoTitles}

В 2014-2018 книга называлась ``Reverse Engineering для начинающих'', но я всегда подозревал что это слишком сужает аудиторию.

Люди от инфобезопасности знают о ``reverse engineering'', но я от них редко слышу слово ``ассемблер''.

Точно также, термин ``reverse engineering'' слишком незнакомый для общей аудитории программистов, но они знают про ``ассемблер''.

В июле 2018, для эксперимента, я заменил название на ``Assembly Language for Beginners''
и запостил ссылку на сайт Hacker News\footnote{\url{https://news.ycombinator.com/item?id=17549050}}, и книгу приняли, в общем, хорошо.

Так что, пусть так и будет, у книги будет два названия.

Хотя, я поменял второе название на ``Understanding Assembly Language'' (``Понимание языка ассемблера''), потому что кто-то уже написал книгу ``Assembly Language for Beginners''.
Также, люди говорят что ``для начинающих'' уже звучит немного саркастично для книги объемом в \textasciitilde{}1000 страниц.

Книги отличаются только названием, именем файла (UAL-XX.pdf и RE4B-XX.pdf), URL-ом и парой первых страниц.

\subsection*{О reverse engineering}

У термина \q{\gls{reverse engineering}} несколько популярных значений:
1) исследование скомпилированных
программ;
2) сканирование трехмерной модели для последующего копирования;
3) восстановление структуры СУБД.

Настоящая книга связана с первым значением.

\subsection*{Желательные знания перед началом чтения}

Очень желательно базовое знание \ac{PL} Си.
Рекомендуемые материалы: \myref{CCppBooks}.

\subsection*{Упражнения и задачи}

\dots 
все перемещены на отдельный сайт: \url{http://challenges.re}.

\subsection*{Об авторе}
\begin{tabularx}{\textwidth}{ l X }

\raisebox{-\totalheight}{
\includegraphics[scale=0.60]{Dennis_Yurichev.jpg}
}

&
Денис Юричев~--- опытный reverse engineer и программист.
С ним можно контактировать по емейлу: \textbf{\EMAIL{}}. % или по Skype: \textbf{dennis.yurichev}.

% FIXME: no link. \tablefootnote doesn't work
\end{tabularx}

% subsections:
\input{praise}
\input{thanks}
\input{FAQ_RU}

\subsection*{Как научиться программированию}

Многие люди спрашивают об этом.

Легких путей нет, но есть пути быстрые и эффективные.

Из моего опыта, это просто: решать задачи из:

\begin{itemize}
\item \KRBook
\item Харольд Абельсон, Джеральд Сассман -- Структура и интерпретация компьютерных программ
\item \TAOCP
\item Книги Никласа Вирта
\item \RobPikePractice
\end{itemize}

... на чистом Си и Лиспе.
В будущем, возможно, вы не будете использовать эти языки вообще.
Большинство коммерческих программистов не используют. Но опыт программирования на Си и Лиспе помогает очень сильно в 
долгосрочной перспективе.

Также, вы можете не читать сами книги, просто листайте, тогда, когда вы чувствуете, что вы не можете чего-то
понять для решения задачи.

В лучшем случае, это занимает годы, или всю жизнь, но это все же намного быстрее, чем метаться между очередными модными
штуками.

Успех этих книг вероятно связан с тем, что их авторы -- преподаватели, и весь материал был в начале отточен на студентах.

Насчет Лиспа, я бы рекомендовал Racket (диалект Scheme). Но это дело вкуса, так или иначе.

Некоторые люди говорят, что понимание ассемблера очень помогает, даже если вы не будете использовать его.
Это верно.
Но это уже путь для наиболее посвященных гиков, и для начала это можно отложить на будущее.

Также, очень частая проблема самоучек (включая автора этих строк), это то, что они слишком часто хватаются
за слишком трудные проблемы, пропуская более легкие.
Это большая ошибка.
Сравните со спортом или музыкой -- никто не начинает с попыток поднимать сразу 100 кг,
или же сразу пытаться играть Капризы Паганини.
Я бы сказал так -- браться за задачу можно тогда, когда вы зараннее, в уме, можете примерно прикинуть решение.

\input{knuth_baby_steps}
( Дональд Э. Кнут -- \url{https://theorydish.blog/2018/02/01/donald-knuth-on-doing-research/} )

Удачи!

\subsection*{О переводе на корейский язык}

В январе 2015, издательство Acorn в Южной Корее сделало много работы в переводе 
и издании моей книги (по состоянию на август 2014) на корейский язык.
Она теперь доступна на \href{http://go.yurichev.com/17343}{их сайте}.

\iffalse
\begin{figure}[H]
\centering
\includegraphics[scale=0.3]{acorn_cover.jpg}
\end{figure}
\fi

Переводил Byungho Min (\href{http://go.yurichev.com/17344}{twitter/tais9}).
Обложку нарисовал мой хороший знакомый художник Андрей Нечаевский
\href{http://go.yurichev.com/17023}{facebook/andydinka}.
Они также имеют права на издание книги на корейском языке.
Так что если вы хотите иметь \emph{настоящую} книгу на полке на корейском языке и
хотите поддержать мою работу, вы можете купить её.

\subsection*{О переводе на персидский язык (фарси)}

В 2016 году книга была переведена Mohsen Mostafa Jokar (который также известен иранскому сообществу по переводу руководства Radare\footnote{\url{http://rada.re/get/radare2book-persian.pdf}}).
Книга доступна на сайте издательства\footnote{\url{http://goo.gl/2Tzx0H}} (Pendare Pars).

Первые 40 страниц: \url{https://beginners.re/farsi.pdf}.

Регистрация книги в Национальной Библиотеке Ирана: \url{http://opac.nlai.ir/opac-prod/bibliographic/4473995}.

\subsection*{О переводе на китайский язык}

В апреле 2017, перевод на китайский был закончен китайским издательством PTPress. Они также имеют права на издание книги на китайском языке.

Она доступна для заказа здесь: \url{http://www.epubit.com.cn/book/details/4174}. Что-то вроде рецензии и история о переводе: \url{http://www.cptoday.cn/news/detail/3155}.

Основным переводчиком был Archer, перед которым я теперь в долгу.
Он был крайне дотошным (в хорошем смысле) и сообщил о большинстве известных ошибок и баг, что крайне важно для литературы вроде этой книги.
Я буду рекомендовать его услуги всем остальным авторам!

Ребята из \href{http://www.antiy.net/}{Antiy Labs} также помогли с переводом. \href{http://www.epubit.com.cn/book/onlinechapter/51413}{Здесь предисловие} написанное ими.

}
\ES{% TODO to be synced with EN version
\section*{Pr\'ologo}

Existen muchos significados populares para el t\'ermino \q{\gls{reverse engineering}}:
1) La ingenier\'ia inversa de software: la investigaci\'on de programas compilados;
2) El escaneo de estructuras 3D y la manipulaci\'on digital subsecuente requerida para duplicarlas;
3) La recreaci\'on de la estructura de un \ac{DBMS}.
Este libro es acerca del primer significado.

\subsection*{Ejercicios y tareas}

\dots 
fueron movidos al sitio web: \url{http://challenges.re}.

\subsection*{Sobre el autor}
\begin{tabularx}{\textwidth}{ l X }

\raisebox{-\totalheight}{
\includegraphics[scale=0.60]{Dennis_Yurichev.jpg}
}

&
Dennis Yurichev es un reverser y programador experimentado.
Puede ser contactado por email: \textbf{\EMAIL{}}.

% FIXME: no link. \tablefootnote doesn't work
\end{tabularx}

% subsections:
\input{praise}
\input{thanks}
\input{FAQ_ES}

\subsection*{Acerca de la traducci\'on al Coreano}

En enero del 2015, la editorial Acorn (\href{http://www.acornpub.co.kr}{www.acornpub.co.kr}) en Corea del Sur realiz\'o una enorme cantidad de trabajo
traduciendo y publicando mi libro (como era en agosto del 2014) en Coreano.
Ahora se encuentra disponible en
\href{http://go.yurichev.com/17343}{su sitio web}.

\iffalse
\begin{figure}[H]
\centering
\includegraphics[scale=0.3]{acorn_cover.jpg}
\end{figure}
\fi

El traductor es Byungho Min (\href{http://go.yurichev.com/17344}{twitter/tais9}).
El arte de la portada fue hecho por mi art\'istico amigo, Andy Nechaevsky
\href{http://go.yurichev.com/17023}{facebook/andydinka}.
Ellos tambi\'en poseen los derechos de autor de la traducci\'on al coreano.
As\'i que, si quieren tener un libro \emph{real} en coreano en su estante
y quieren apoyar mi trabajo, ya se encuentra disponible a la venta.

%\subsection*{About the Persian/Farsi translation}
%TBT

}
\NL{% TODO to be synced with EN version
\section*{Voorwoord}

Er zijn verschillende populaire betekenissen voor de term \q{\gls{reverse engineering}}:
1) Reverse engineeren van software: gecompileerde programma\'s onderzoeken;
2) Scannen van 3D structuren en de onderliggende digitale bewerkingen om deze te kunnen dupliceren;
3) Het nabootsen van een \ac{DBMS} structuur.
Dit boek gaat over de eerste betekenis.

\subsection*{Oefeningen en opdrachten}

\dots 
zijn allen verplaatst naar de website: \url{http://challenges.re}.

\subsection*{Over de auteur}
\begin{tabularx}{\textwidth}{ l X }

\raisebox{-\totalheight}{
\includegraphics[scale=0.60]{Dennis_Yurichev.jpg}
}

&
Dennis Yurichev is een ervaren reverse engineer en programmeur.
Je kan hem contacteren via email: \textbf{\EMAIL{}}.

% FIXME: no link. \tablefootnote doesn't work
\end{tabularx}

% subsections:
\input{praise}
\input{thanks}
%\input{FAQ_NL} % to be translated

% {\RU{Целевая аудитория}\EN{Target audience}}

\subsection*{Over de Koreaanse vertaling}

In Januari 2015 heeft de Acorn uitgeverij (\href{http://www.acornpub.co.kr}{www.acornpub.co.kr}) in Zuid Korea een enorme hoeveelheid werk verricht in het vertalen en uitgeven
van mijn boek (zoals het was in augustus 2014) in het Koreaans.

Het is nu beschikbaar op
\href{http://go.yurichev.com/17343}{hun website}

\iffalse
\begin{figure}[H]
\centering
\includegraphics[scale=0.3]{acorn_cover.jpg}
\end{figure}
\fi

De vertaler is Byungho Min (\href{http://go.yurichev.com/17344}{twitter/tais9}).
De cover art is verzorgd door mijn artistieke vriend, Andy Nechaevsky
\href{http://go.yurichev.com/17023}{facebook/andydinka}.
Zij bezitten ook de auteursrechten voor de Koreaanse vertaling.
Dus, als je een \emph{echt} boek op je kast wil in het Koreaans en je
wil mijn werk steunen, is het nu beschikbaar voor verkoop.

%\subsection*{About the Persian/Farsi translation}
%TBT

}
\IT{\section*{Prefazione}

\subsection*{Da cosa derivano i due titoli?}
\label{TwoTitles}

Il libro era chiamato ``Reverse Engineering for Beginners'' nel periodo 2014-2018, ma ho sempre sospettato che questo restringesse troppo i potenziali lettori.

Nel campo Infosec le persone conoscono il ``reverse engineering'', ma raramente ho sentito la parola ``assembler'' da parte loro.

Similmente, il termine ``reverse engineering'' è in qualche modo criptico per il resto dei programmatori, ma sanno cos'è l'``assembler''.

A luglio 2018, per esperimento, ho cambiato il titolo in ``Assembly Language for Beginners''
e postato il link sul sito Hacker News \footnote{\url{https://news.ycombinator.com/item?id=17549050}}, ed il libro ha avuto un buon successo.

Quindi è così, il libro adesso ha due titoli.

Tuttavia, ho modificato il secondo titolo in ``Understanding Assembly Language'', perchè qualcuno aveva già scritto un libro ``Assembly Language for Beginners''.
Inoltre lla gente dice che ``for Beginners'' sembra un po' sarcastico per un libro di \textasciitilde{}1000 pagine.

I due libri sono differenti solo per il titolo, il nome del file (UAL-XX.pdf versus RE4B-XX.pdf), l'URL ed un paio di pagine iniziali.

\subsection*{Sul reverse engineering}

Esistono diversi significati per il termine \q{\gls{reverse engineering}}:

1) Il reverse engineering del software; riguardo la ricerca su programmi compilati

2) La scansione di strutture 3D e la successiva manipolazione digitale necessaria alla loro riproduzione

3) Ricreare strutture in \ac{DBMS}

Questo libro riguarda il primo significato.

\subsection*{Prerequisiti}

Conoscenza di base del C \ac{PL}.
Letture raccomandate: \myref{CCppBooks}.

\subsection*{Esercizi e compiti}

\dots
possono essere trovati su: \url{http://challenges.re}.

\subsection*{L'Autore}
\begin{tabularx}{\textwidth}{ l X }

\raisebox{-\totalheight}{
\includegraphics[scale=0.60]{Dennis_Yurichev.jpg}
}

&
Dennis Yurichev è un reverse engineer e programmatore.
Può essere contattato via mail: \textbf{\EMAIL{}}. % o Skype: \textbf{dennis.yurichev}.

% FIXME: no link. \tablefootnote doesn't work
\end{tabularx}

% subsections:
\input{praise}
\input{thanks}
\input{FAQ_IT}

\subsection*{Come imparare a programmare}

Molte persone se lo chiedono.

Non esiste una via unica preferenziale, ma esistono alcuni modi efficaci.

Per la mia esperienza, basta risolvere gli esercizi su:

\begin{itemize}
\item \KRBook
\item Harold Abelson, Gerald Jay Sussman, Julie Sussman -- Structure and Interpretation of Computer Programs
\item \TAOCP
\item Niklaus Wirth's books
\item \RobPikePractice
\end{itemize}

... in C puro e LISP.
Potresti non utilizzare mai questi linguaggi in futuro.
La maggiorparte dei programmatori commerciali non lo fanno. Ma l'esperienza nella programmazione in C e LISP ti aiuterà enormemente a lungo termine.

Inoltre, puoi saltare la lettura dei libri in se,
salta semplicemente alle sezioni in cui pensi di aver bisogno per capire qualcosa che ti manca per risolvere l'esercizio che stai facendo.

Può richiedere anni nel migliore dei casi, o una vita, ma è comunque più veloce che correre dietro alle mode del momento.

Il successo di questi libri dipende probabilmente dal fatto che i loro autori sono insegnanti
e tutto il materiale è stato verificato sugli studenti.

Riguardo al LISP, personalmente consiglierei Racket (un dialetto). Ma questa è una questione di preferenze.

Secondo alcune persone anche la comprensione del linguaggio assembly è molto utile, anche se non lo utilizzi mai.
Questo è vero.
Ma questa è una strada per i geek più dedicati, e inizialmente può essere rimandata.

Inoltre, gli autodidatti (incluso l'autore di queste righe) spesso hanno la caratteristica di perdere troppo tempo sui problemi più complessi,
saltando quelli più semplici.
Questo è un grande errore.
Pensa allo sport o alla musica -- nessuno inizia con sollvare pesi di 100kg, o a suonare i Capricci di Paganini.
Io direi che -- puoi tentare di affrontare un problema se sei in grado di abbozzare mentalmente una soluzione.

\input{knuth_baby_steps}
( Donald E. Knuth -- \url{https://theorydish.blog/2018/02/01/donald-knuth-on-doing-research/} )

Buona fortuna!

\subsection*{La traduzione in Coreano}

A gennaio 2015, la Acorn publishing company (\href{http://www.acornpub.co.kr}{www.acornpub.co.kr}) in Corea del Sud ha compiuto un enorme lavoro traducendo e pubblicando
questo libro (aggiornato ad agosto 2014) in Coreano.

Adesso è disponibile sul \href{http://go.yurichev.com/17343}{loro sito}.

\iffalse
\begin{figure}[H]
\centering
\includegraphics[scale=0.3]{acorn_cover.jpg}
\end{figure}
\fi

Il traduttore è Byungho Min (\href{http://go.yurichev.com/17344}{twitter/tais9}).
La copertina è stata creata dall'artistico Andy Nechaevsky, un amico dell'autore:
\href{http://go.yurichev.com/17023}{facebook/andydinka}.
Acorn detiene inoltre i diritti della traduzione in Coreano.

Quindi se vuoi un \emph{vero} libro in Coreano nella tua libreria e
vuoi supportare questo lavoro, è disponibile per l'acquisto.

\subsection*{La traduzione in Persiano/Farsi}

Nel 2016 il libro è stato tradotto da Mohsen Mostafa Jokar (che è anche conosciuto nell comunità iraniana per la traduzione del manuale di Radare \footnote{\url{http://rada.re/get/radare2book-persian.pdf}}).
Potete trovarlo sul sito dell'editore\footnote{\url{http://goo.gl/2Tzx0H}} (Pendare Pars).

Qua c'è il link ad un estratto di 40 pagine: \url{https://beginners.re/farsi.pdf}.

Informazioni nella Libreria Nazionale dell'Iran: \url{http://opac.nlai.ir/opac-prod/bibliographic/4473995}.

\subsection*{La traduzione Cinese}

Ad aprile 2017, la traduzione in Cinese è stata completata da Chinese PTPress. Possiedono inoltre i diritti della traduzione in Cinese.

La versione cinese può essere ordinata a questo indirizzo: \url{http://www.epubit.com.cn/book/details/4174}. Una recensione parziale, con informazioni sulla traduzione è disponibile qua: \url{http://www.cptoday.cn/news/detail/3155}.

Il traduttore principale è Archer, al quale l'autore deve molto. E' stato estremamente meticolosos (in senso buono) ed ha segnalato buona parte degli errori e bug, il che è molto importante in un libro come questo.
L'autore raccomanderebbe i suoi servizi a qualsiasi altro autore!

I ragazzi di \href{http://www.antiy.net/}{Antiy Labs} hanno inoltre aiutato nella traduzione. \href{http://www.epubit.com.cn/book/onlinechapter/51413}{Qua c'è la prefazione} scritta da loro.
}
\DE{\section*{Vorwort}

\subsection*{Warum zwei Titel?}
\label{TwoTitles}

Dieses Buch hieß von 2014-2018 ``Reverse Engineering for Beginners'', jedoch hatte ich immer die Befürchtung, dass es den Leserkreis zu sehr einengen würde.

Infosec Leute kennen sich mit ``Reverse Engineering'' aus, jedoch hörte ich selten das Wort ``Assembler'' von Ihnen.

Desweiteren ist der Begriff ``Reverse Engineering'' etwas zu kryptisch für den Großteil von Programmierern, diesen ist jedoch ``Assembler'' geläufig.

Im Juli 2018 änderte ich als Experiment den Titel zu ``Assembly Language for Beginners'' und veröffentlichte den Link auf der Hacker News-Website\footnote{\url{https://news.ycombinator.com/item?id=17549050}}. Das Buch kam allgemein gut an.

Aus diesem Grund hat das Buch nun zwei Titel.

Ich habe den zweiten Titel zu ``Understanding Assembly Language'' geändert, da es bereits eine Erscheinung mit dem Titel ``Assembly Language for Beginners'' gab. Einige Leute sind der Meinung, dass ``for Beginners'' etwas sarkastisch klingt, für ein Buch mit \textasciitilde{}1000 Seiten.

Die beiden Bücher unterscheiden sich lediglich im Titel, dem Dateinamen (UAL-XX.pdf beziehungsweise RE4B-XX.pdf), URL und ein paar der einleitenden Seiten.

\subsection*{Über Reverse Engineering}

Es gibt verschiedene verbreitete Interpretationen des Begriffs Reverse Engineering:\\
1) Reverse Engineering von Software: Rückgewinnung des Quellcodes bereits kompilierter Programme;\\
2) Das Erfassen von 3D Strukturen und die digitalen Manipulationen die zur Duplizierung notwendig sind;\\
3) Nachbilden von \ac{DBMS}-Strukturen.\\
Dieses Buch behandelt die erste Interpretation.

\subsection*{Voraussetzungen}

Grundlegende Kenntnisse der Programmiersprache C.
Empfohlene Literatur: \myref{CCppBooks}.

\subsection*{Übungen und Aufgaben}
\dots 
befinden sich nun alle auf der Website: \url{http://challenges.re}.

\subsection*{Über den Autor}
\begin{tabularx}{\textwidth}{ l X }

\raisebox{-\totalheight}{
\includegraphics[scale=0.60]{Dennis_Yurichev.jpg}
}

&
Dennis Yurichev ist ein erfahrener Reverse Engineer und Programmierer.
Er kann per E-Mail kontaktiert werden: \textbf{\EMAIL{}}.

% FIXME: no link. \tablefootnote doesn't work
\end{tabularx}

% subsections:
\input{praise}
\input{thanks}
\input{FAQ_DE}

\subsection*{Programmieren lernen}

Viele Leute stellen hierzu Fragen.

Es gibt keinen ``goldenen Weg'', aber einige effiziente Möglichkeiten.

Aus meiner eigenen Erfahrung ist dies: lösen von Übungen aus:

\begin{itemize}
	\item \KRBook
	\item Harold Abelson, Gerald Jay Sussman, Julie Sussman -- Structure and Interpretation of Computer Programs
	\item \TAOCP
	\item Niklaus Wirth's books
	\item \RobPikePractice
\end{itemize}

... in C und LISP.
Möglicherweise benutzt du diese Programmiersprachen in der Zukunft nicht mehr, zumindest die meisten professionellen Programmierer tun es nicht. Aber Programmiererfahrung in C und LISP wird langfristig eine wichtige Bedeutung haben.

Das Lesen der Bücher kann übersprungen und die Bücher nur zu Rate gezogen werden wenn das Gefühl aufkommt, dass einige Grundkenntnisse zum Lösen der Aufgaben fehlen.

Dies mag Jahre dauern, vielleicht auch ein Leben lang. Trotzdem ist es noch schneller als Modeerscheinungen hinterherzulaufen.

Der Erfolg dieser Bücher ist vielleicht in der Tatsache begründet, dass die Autoren Lehrer sind und das komplette Material zunächst von Studenten genutzt wurde.

Bezüglich LISP, empfehle ich persönlich Racket (einen Scheme-Dialekt). Dies ist aber stark abhängig vom Geschmack jedes Einzelnen.

Manche sagen, dass das Verständnis von Assembler auch hilfreich ist, wenn man es niemals nutzen wird.
Das ist wahr.
Aber dies ist der Weg für die meisten ``geeks'' und kann erst mal aufgeschoben werden.

Selbstlerner (inklusive dem Autoren dieses Buchs) stehen oft vor dem Problem, sich erst mit schwierigen Aufgaben zu beschäftigen, die einfachen jedoch zu überspringen.
Das ist ein großer Fehler.
Als Vergleich mit Sport oder Musik -- niemand startet mit 100Kg-Gewichten oder Paganini's Caprices.
Ich würde sagen -- man kann die Lösung eines Problems angehen, wenn man diese schon grob im Kopf skizzieren kann.

\input{knuth_baby_steps}
( Donald E. Knuth -- \url{https://theorydish.blog/2018/02/01/donald-knuth-on-doing-research/} )

Viel Glück!

\subsection*{Über die koreanische Übersetzung}

Im Januar 2015 hat die Acorn Publishing Company (\href{http://www.acornpub.co.kr}{www.acornpub.co.kr}) in Süd-Korea
viel Aufwand in die Übersetzung und Veröffentlichung meines Buchs ins Koreanische (mit Stand August 2014) investiert.

Es ist jetzt unter dieser \href{http://go.yurichev.com/17343}{Webseite} verfügbar.

\iffalse
\begin{figure}[H]
\centering
\includegraphics[scale=0.3]{acorn_cover.jpg}
\end{figure}
\fi

Der Übersetzer ist Byungho Min (\href{http://go.yurichev.com/17344}{twitter/tais9}).
Die Cover-Gestaltung wurde von meinem künstlerisch begabten Freund Andy Nechaevsky erstellt:
\href{http://go.yurichev.com/17023}{facebook/andydinka}.
Die Acorn Publishing Company besetzt die Urheberrechte an der koreanischen Übersetzung.

Wenn Sie also ein \emph{echtes} Buch in Ihrem Buchregal auf koreanisch haben und 
mich bei meiner Arbeit unterstützen wollen, können Sie das Buch nun kaufen.

\subsection*{Über die persische Übersetzung (Farsi)}

In 2016 wurde das Buch von Mohsen Mostafa Jokar übersetzt, der in der iranischen Community
auch für die Übersetzung des Radare
Handbuchs\footnote{\url{http://rada.re/get/radare2book-persian.pdf}} bekannt ist).
Es ist auf der Homepage des Verlegers\footnote{\url{http://goo.gl/2Tzx0H}} (Pendare Pars)
verfügbar.

Hier ist ein Link zu einem 40-seitigen Auszug: \url{https://beginners.re/farsi.pdf}.

National Library of Iran registration information: \url{http://opac.nlai.ir/opac-prod/bibliographic/4473995}.

\subsection*{Über die chinesische Übersetzung}

Im April 2017 wurde die chinesische Übersetzung von Chinese PTPress fertiggestellt, bei denen
auch das Copyright liegt.

Die chinesische Version kann hier bestellt werden: \url{http://www.epubit.com.cn/book/details/4174}.
Ein Auszug und die Geschichte der Übersetzung kann hier gefunden werden: \url{http://www.cptoday.cn/news/detail/3155}.

Der Hauptübersetzer ist Archer, dem der Autor sehr iel verdankt.
Archer war extrem akribisch (im positiven Sinne) und meldete
den Großteil der bekannten Fehler, was für Literatur wie dieses Buch extrem wichtig ist.
Der Autor empfiehlt diesen Service jedem anderen Autor!

Die Mitarbeiter von \href{http://www.antiy.net/}{Antiy Labs} halfen ebenfalls bei der Übersetzung.
\href{http://www.epubit.com.cn/book/onlinechapter/51413}{Hier ist das Vorwort} von ihnen.
}
\FR{\section*{Préface}

\subsection*{C'est quoi ces deux titres?}
\label{TwoTitles}

Le livre a été appelé ``Reverse Engineering for Beginners'' en 2014-2018, mais j'ai
toujours suspecté que ça rendait son audience trop réduite.

Les gens de l'infosec connaissent le ``reverse engineering'', mais j'ai rarement
entendu le mot ``assembleur'' de leur part.

De même, le terme ``reverse engineering'' est quelque peu cryptique pour une audience
générale de programmeurs, mais qui ont des connaissances à propos de l'``assembleur''.

En juillet 2018, à titre d'expérience, j'ai changé le titre en ``Assembly Language for Beginners''
et publié le lien sur le site Hacker News\footnote{\url{https://news.ycombinator.com/item?id=17549050}},
et le livre a été plutôt bien accueilli.

Donc, c'est ainsi que le livre a maintenant deux titres.

Toutefois, j'ai changé le second titre à ``Understanding Assembly Language'', car
quelqu'un a déjà écrit le livre ``Assembly Language for Beginners''.
De même, des gens disent que ``for Beginners'' sonne sarcastique pour un livre de
\textasciitilde{}1000 pages.

Les deux livres diffèrent seulement par le titre, le nom du fichier (UAL-XX.pdf versus
RE4B-XX.pdf), l'URL et quelques-une des première pages.

\subsection*{À propos de la rétro-ingénierie}

Il existe plusieurs définitions pour l'expression \q{ingénierie inverse ou rétro-ingénierie \gls{reverse engineering}} :

1) L'ingénierie inverse de logiciels : examiner des programmes compilés;

2) Le balayage des structures en 3D et la manipulation numérique nécessaire afin de les reproduire;

3) Recréer une structure de base de données.

Ce livre concerne la première définition.

\subsection*{Prérequis}

Connaissance basique du C \ac{PL}.
Il est recommandé de lire: \myref{CCppBooks}.

\subsection*{Exercices et tâches}

\dots
ont été déplacés sur un site différent : \url{http://challenges.re}.

\subsection*{A propos de l'auteur}
\begin{tabularx}{\textwidth}{ l X }

\raisebox{-\totalheight}{
\includegraphics[scale=0.60]{Dennis_Yurichev.jpg}
}

&
Dennis Yurichev est un ingénieur expérimenté en rétro-ingénierie et un programmeur.
Il peut être contacté par email : \textbf{\EMAIL{}}.

% FIXME: no link. \tablefootnote doesn't work
\end{tabularx}

% subsections:
\input{praise}
\input{thanks}
\input{FAQ_FR}

\subsection*{Comment apprendre à programmer}

Beaucoup de gens me l'ont demandé.

Il n'y a pas de ``voie royale'', mais il y a quelques chemins efficaces.

De ma propre expérience, ceci est juste: résoudre les exercices de:

\begin{itemize}
\item \KRBook
\item Harold Abelson, Gerald Jay Sussman, Julie Sussman -- Structure and Interpretation of Computer Programs
\item \TAOCP
\item Niklaus Wirth's books
\item \RobPikePractice
\end{itemize}

... en pur C et LISP.
Vous pourriez ne jamais utiliser du tout ces langages de programmation dans le futur.
Presqu'aucun des programmeurs de métier le les utilisent. Mais l'expérience du développement
en C et en LISP aide énormément sur la long terme.

Vous pouvez également passer leur lecture elle-même, parcourez-les seulement lorsque
vous sentez que vous devez comprendre quelque chose que vous ne comprenez pas dans
l'exercice que vous faites.

Ceci peu prendre des années, ou une vie entière, mais c'est toujours plus rapide
que de se précipiter entre les modes.

Le succès de ces livres est sans doute lié au fait que les auteurs sont des enseignants
et tous ce matériel a d'abord été amélioré avec les étudiants.

Pour LISP, je recommande personnellement Racket (dialecte Scheme). Mais c'est une
affaire de goût.

Certaines personnes disent que comprendre le langage d'assemblage est très utile,
même si vous ne l'utiliserez jamais.
Ceci est vrai.
Mais ceci est une voie pour les geeks les plus persévérants, et ça peut être mis de
côté au début.

De même, les autodidactes (incluant l'auteur de ces lignes) ont souvent le problème
de travailler dur sur des problèmes difficiles, passant sur les plus faciles.
Ceci est une grosse erreur.
Comparez avec le sport ou le musique -- personne ne commence à un poids de 100kg,
ou Paganini's Caprices.
Je dirais -- essayez de résoudre un problème si vous pouvez en esquisser la solution
dans votre tête.

\input{knuth_baby_steps}
( Donald E. Knuth -- \url{https://theorydish.blog/2018/02/01/donald-knuth-on-doing-research/} )

Bonne chance!

\subsection*{À propos de la traduction en Coréen}

En Janvier 2015, la maison d'édition Acorn (\href{http://www.acornpub.co.kr}{www.acornpub.co.kr}) en Corée du Sud a réalisé un énorme travail en traduisant et en publiant mon livre (dans son état en Août 2014) en Coréen.

Il est désormais disponible sur \href{http://go.yurichev.com/17343}{leur site web}.

\iffalse
\begin{figure}[H]
\centering
\includegraphics[scale=0.3]{acorn_cover.jpg}
\end{figure}
\fi

Le traducteur est Byungho Min (\href{http://go.yurichev.com/17344}{twitter/tais9}).
L'illustration de couverture a été réalisée l'artiste, Andy Nechaevsky, un ami de l'auteur:
\href{http://go.yurichev.com/17023}{facebook/andydinka}.
Ils détiennent également les droits d'auteurs sur la traduction coréenne.

Donc si vous souhaitez avoir un livre \emph{réel} en coréen sur votre étagère et que vous souhaitez soutenir ce travail, il est désormais disponible à l'achat.

\subsection*{Á propos de la traduction en Farsi/Perse}

En 2016, ce livre a été traduit par Mohsen Mostafa Jokar (qui est aussi connu dans
la communauté iranienne pour sa traduction du manuel de Radare\footnote{\url{http://rada.re/get/radare2book-persian.pdf}}).
Il est disponible sur le site web de l'éditeur\footnote{\url{http://goo.gl/2Tzx0H}}
(Pendare Pars).

Extrait de 40 pages: \url{https://beginners.re/farsi.pdf}.

Enregistrement du livre à la Bibliothèque Nationale d'Iran: \url{http://opac.nlai.ir/opac-prod/bibliographic/4473995}.

\subsection*{Á propos de la traduction en Chinois}

En avril 2017, la traduction en Chinois a été terminée par Chinese PTPress. Ils sont
également les détenteurs des droits de la traduction en Chinois.

La version chinoise est disponible à l'achat ici: \url{http://www.epubit.com.cn/book/details/4174}.
Une revue partielle et l'historique de la traduction peut être trouvé ici: \url{http://www.cptoday.cn/news/detail/3155}.


Le traducteur principal est Archer, à qui je dois beaucoup. Il a été très méticuleux
(dans le bon sens du terme) et a signalé la plupart des erreurs et bugs connus, ce
qui est très important dans le genre de littérature de ce livre.
Je recommanderais ses services à tout autre auteur!

Les gens de \href{http://www.antiy.net/}{Antiy Labs} ont aussi aidé pour la traduction.
\href{http://www.epubit.com.cn/book/onlinechapter/51413}{Voici la préface} écrite par eux.

}
\PL{% TODO to be synced with EN version
\section*{Słowo wstępne}

Termin \q{\gls{reverse engineering}} ma kilka popularnych definicji:
1) analiza skompilowanych programów; 2) skanowanie modelu w 3D, żeby następnie go skopiować;
3) odzyskanie struktury SZBD. Nasza książka będzie powiązana z tą pierwszą definicją.

\subsection*{Pożądana wiedza}

Bardzo pożądane, aczkolwiek nie niezbędne są podstawowe umiętności programowania w \ac{PL} C.
Polecane materiały: \myref{CCppBooks}.

\subsection*{Ćwiczenia i zadania}

\dots 
wszystkie są na osobnej stronie: \url{http://challenges.re}.

\subsection*{O autorze}
\begin{tabularx}{\textwidth}{ l X }

\raisebox{-\totalheight}{
\includegraphics[scale=0.60]{Dennis_Yurichev.jpg}
}

&
Dennis Yurichev~--- reverse engineer i programista z dużym doświadczeniem.
Można się z nim skontaktować mailowo: \textbf{\EMAIL{}}.

% FIXME: no link. \tablefootnote doesn't work
\end{tabularx}

% subsections:
\input{praise}
\input{thanks}
\input{FAQ_PL}

\subsection*{O tłumaczeniu na język koreański}

W styczniu 2015, wydawnictwo Acorn w Korei Południowej zrobiło kawał dobrej roboty, żeby przetłumaczyć i wydać moją książkę  (stanem na siepień 2014) w języku koreańskim.
Jest ona teraz dostępna na \href{http://go.yurichev.com/17343}{ich stronie}.

\iffalse
\begin{figure}[H]
\centering
\includegraphics[scale=0.3]{acorn_cover.jpg}
\end{figure}
\fi

Tłumaczył Byungho Min (\href{http://go.yurichev.com/17344}{twitter/tais9}).
Okładka została narysowana przez mojego dobrego znajomego Andy Nechaevsky
\href{http://go.yurichev.com/17023}{facebook/andydinka}.
Oni również mają prawa do wydania książki w języku koreańskim.
Jakbyście chcieli mieć \emph{prawdziwą} książkę w języku koreańskim i chcielibyście wesprzeć moją pracę, możecie ją kupić.

\subsection*{O tłumaczeniu na język perski (farsi)}

W roku 2016 książka była przetłumaczona przez Mohsen Mostafa Jokar.
Książka jest dostępna na stronie wydawnictwa \footnote{\url{http://goo.gl/2Tzx0H}} (Pendare Pars).

Pierwsze 40 stron: \url{https://beginners.re/farsi.pdf}.

Pozycja książki w Nacjonalnej bibliotece Iranu: \url{http://opac.nlai.ir/opac-prod/bibliographic/4473995}.

\subsection*{O tłumaczeniu na język chiński}

W kwietniu 2017, wydawnictwo PTPress skończyło tłumaczenie mojej książki na j. chiński. Oni również mają prawa do wydania książki w języku chińskim.

Można ją zamówić tutaj: \url{http://www.epubit.com.cn/book/details/4174}. Coś w stylu recenzji: \url{http://www.cptoday.cn/news/detail/3155}.

Głównym tłumaczem był Archer, przed którym mam teraz dług widzięczności.
Był bardzo dociekliwy i znalazł w książce sporo bugów i błędów, co jest krytycznie ważne dla książki open-source.
Będę polecał go również innym autorom!

Chłopaki z \href{http://www.antiy.net/}{Antiy Labs} również pomogli z tłumaczeniem. \href{http://www.epubit.com.cn/book/onlinechapter/51413}{Tutaj słowo wstępne} napisane przez nich.


}
%\CN{% !TEX program = XeLaTeX
% !TEX encoding = UTF-8
\documentclass[UTF8,nofonts]{ctexart}
\setCJKmainfont[BoldFont=STHeiti,ItalicFont=STKaiti]{STSong}
\setCJKsansfont[BoldFont=STHeiti]{STXihei}
\setCJKmonofont{STFangsong}

\begin{document}

%daveti: translated on Dec 28, 2016
%NOTE: above is needed for MacTex.

\section*{前言 Preface}

术语\q{\gls{逆向工程 reverse engineering}}有好多种流行的说法:
1) 软件逆向工程: 研究编译器产生的程序;
2) 3D结构扫描和随后的用来复制该结构的数字模型重建;
3) Recreating \ac{DBMS} structure.
This book is about the first meaning.

\subsection*{前提条件 Prerequisites}

基本的C \ac{语言 PL}知识.
推荐阅读: \myref{CCppBooks}.

\subsection*{练习和项目 Exercises and tasks}

\dots 
都已经转移到了一个单独的网站: \url{http://challenges.re}.

\subsection*{关于作者 About the author}
\begin{tabularx}{\textwidth}{ l X }

\raisebox{-\totalheight}{
\includegraphics[scale=0.60]{Dennis_Yurichev.jpg}
}

&
Dennis Yurichev是一个经验丰富的逆向工程师和程序员.
可以通过邮件联系他: \textbf{\EMAIL{}}.

% FIXME: no link. \tablefootnote doesn't work
\end{tabularx}

% subsections:
\input{praise}
\input{thanks}
\input{FAQ_CN}

\subsection*{关于韩语翻译 About the Korean translation}

2015一月，南韩的Acorn出版社 (\href{http://www.acornpub.co.kr}{www.acornpub.co.kr}) 做了大量的翻译工作并且最终把这本书（2014年8月的版本）韩语版出版。

现在可以从这里买到\href{http://go.yurichev.com/17343}{出版社网站}.

\iffalse
\begin{figure}[H]
\centering
\includegraphics[scale=0.3]{acorn_cover.jpg}
\end{figure}
\fi

译者是Byungho Min (\href{http://go.yurichev.com/17344}{twitter/tais9}).
封面依然由我的艺术家哥们Andy Nechaevsky设计:
\href{http://go.yurichev.com/17023}{facebook/andydinka}.
他们共同持有韩语翻译版的版权。

所以，如果你想有一本\emph{真正的 real}韩语版本书而且顺便支持我的工作，请购买该书。

\end{document}

%\subsection*{About the Persian/Farsi translation}
%TBT

}
\JPN{\input{preface_JPN}}


\mainmatter

% only chapters here!
\ifdefined\SPANISH
\chapter{Patrones de código}
\fi % SPANISH

\ifdefined\GERMAN
\chapter{Code-Muster}
\fi % GERMAN

\ifdefined\ENGLISH
\chapter{Code Patterns}
\fi % ENGLISH

\ifdefined\ITALIAN
\chapter{Pattern di codice}
\fi % ITALIAN

\ifdefined\RUSSIAN
\chapter{Образцы кода}
\fi % RUSSIAN

\ifdefined\BRAZILIAN
\chapter{Padrões de códigos}
\fi % BRAZILIAN

\ifdefined\THAI
\chapter{รูปแบบของโค้ด}
\fi % THAI

\ifdefined\FRENCH
\chapter{Pattern de code}%meilleur?
\fi % FRENCH

\ifdefined\POLISH
\chapter{\PLph{}}
\fi % POLISH

\ifdefined\JAPANESE
\chapter{コードパターン}
\fi % JAPANESE

\ifdefined\TURKISH
\chapter{Kod Şablonları}
\fi % Turkish

% sections
\EN{\input{patterns/patterns_opt_dbg_EN}}
\ES{\input{patterns/patterns_opt_dbg_ES}}
\IT{\input{patterns/patterns_opt_dbg_IT}}
\PTBR{\input{patterns/patterns_opt_dbg_PTBR}}
\RU{\input{patterns/patterns_opt_dbg_RU}}
\THA{\input{patterns/patterns_opt_dbg_THA}}
\DE{\input{patterns/patterns_opt_dbg_DE}}
\FR{\input{patterns/patterns_opt_dbg_FR}}
\PL{\input{patterns/patterns_opt_dbg_PL}}
\JA{\input{patterns/patterns_opt_dbg_JA}}

\RU{\mysection{Некоторые базовые понятия}}
\EN{\mysection{Some basics}}
\DE{\mysection{Einige Grundlagen}}
\FR{\mysection{Quelques bases}}
\ES{\mysection{Algunos conceptos basicos}}
\IT{\mysection{Alcune basi teoriche}}
\PTBR{\mysection{\PTBRph{}}}
\THA{\mysection{\THAph{}}}
\PL{\mysection{\PLph{}}}
\TR{\mysection{Temeller}}
\JA{\mysection{基本的な事柄}}

% sections:
\EN{\input{patterns/intro_CPU_ISA_EN}}
\ES{\input{patterns/intro_CPU_ISA_ES}}
\IT{\input{patterns/intro_CPU_ISA_IT}}
\PTBR{\input{patterns/intro_CPU_ISA_PTBR}}
\RU{\input{patterns/intro_CPU_ISA_RU}}
\DE{\input{patterns/intro_CPU_ISA_DE}}
\FR{\input{patterns/intro_CPU_ISA_FR}}
\PL{\input{patterns/intro_CPU_ISA_PL}}
\JA{\input{patterns/intro_CPU_ISA_JA}}
\TR{\input{patterns/intro_CPU_ISA_TR}}

\EN{\input{patterns/numeral_EN}}
\RU{\input{patterns/numeral_RU}}
\IT{\input{patterns/numeral_IT}}
\DE{\input{patterns/numeral_DE}}
\FR{\input{patterns/numeral_FR}}
\PL{\input{patterns/numeral_PL}}
\ES{\input{patterns/numeral_ES}}
\JA{\input{patterns/numeral_JA}}

% chapters
\input{patterns/00_empty/main}
\input{patterns/011_ret/main}
\input{patterns/01_helloworld/main}
\input{patterns/015_prolog_epilogue/main}
\input{patterns/02_stack/main}
\input{patterns/03_printf/main}
\input{patterns/04_scanf/main}
\input{patterns/05_passing_arguments/main}
\input{patterns/06_return_results/main}
\input{patterns/061_pointers/main}
\input{patterns/065_GOTO/main}
\input{patterns/07_jcc/main}
\input{patterns/075_cracking/main}
\input{patterns/08_switch/main}
\input{patterns/09_loops/main}
\input{patterns/10_strings/main}
\input{patterns/11_arith_optimizations/main}
\input{patterns/12_FPU/main}
\input{patterns/13_arrays/main}
\input{patterns/135_angband/main}
\input{patterns/14_bitfields/main}
\EN{\input{patterns/145_LCG/main_EN}}
\RU{\input{patterns/145_LCG/main_RU}}
\DE{\input{patterns/145_LCG/main_DE}}
\FR{\input{patterns/145_LCG/main_FR}}
\JA{\input{patterns/145_LCG/main_JA}}
\input{patterns/15_structs/main}
\input{patterns/17_unions/main}
\input{patterns/18_pointers_to_functions/main}
\input{patterns/185_64bit_in_32_env/main}

\EN{\input{patterns/19_SIMD/main_EN}}
\RU{\input{patterns/19_SIMD/main_RU}}
\DE{\input{patterns/19_SIMD/main_DE}}
\FR{\input{patterns/19_SIMD/main_FR}}
\JA{\input{patterns/19_SIMD/main_JA}}

\EN{\input{patterns/20_x64/main_EN}}
\RU{\input{patterns/20_x64/main_RU}}
\DE{\input{patterns/20_x64/main_DE}}
\FR{\input{patterns/20_x64/main_FR}}
\JA{\input{patterns/20_x64/main_JA}}

\EN{\input{patterns/205_floating_SIMD/main_EN}}
\RU{\input{patterns/205_floating_SIMD/main_RU}}
\DE{\input{patterns/205_floating_SIMD/main_DE}}
\FR{\input{patterns/205_floating_SIMD/main_FR}}
\JA{\input{patterns/205_floating_SIMD/main_JA}}

\EN{\input{patterns/ARM/main_EN}}
\RU{\input{patterns/ARM/main_RU}}
\DE{\input{patterns/ARM/main_DE}}
\FR{\input{patterns/ARM/main_FR}}
\JA{\input{patterns/ARM/main_JA}}

\input{patterns/MIPS/main}

% FIXME: why chapter?! it should be section
\chapter{%
	\RU{Важные фундаментальные вещи}%
	\EN{Important fundamentals}%
	\ES{Fundamentos importantes}%
	\PTBRph{}%
	\DE{Wichtige Grundlagen}%
	\FR{Fondamentaux importants}%
	\PLph{}%
	\ITph{}%
	\JPNph{}%
}

%\clearpage
\begin{center}
\vspace*{\fill}

\begin{figure}[H]
\centering
\includegraphics[width=0.7\textwidth]{cover2.jpg}
\end{figure}

\vspace*{\fill}
\end{center}

\clearpage

% sections
\EN{\input{fundamentals/data_types_and_numbers_EN}}
\RU{\input{fundamentals/data_types_and_numbers_RU}}
\FR{\input{fundamentals/data_types_and_numbers_FR}}

\EN{\input{fundamentals/signed_numbers_EN}}
\RU{\input{fundamentals/signed_numbers_RU}}
\ES{\input{fundamentals/signed_numbers_ES}}
\DE{\input{fundamentals/signed_numbers_DE}}
\FR{\input{fundamentals/signed_numbers_FR}}

\EN{\input{fundamentals/int_overflow_EN}}
\RU{\input{fundamentals/int_overflow_RU}}
\FR{\input{fundamentals/int_overflow_FR}}

\EN{\input{fundamentals/AND_EN}}\RU{\input{fundamentals/AND_RU}}%
\FR{\input{fundamentals/AND_FR}}

\EN{\input{fundamentals/AND_OR_as_SUB_ADD_EN}}\RU{\input{fundamentals/AND_OR_as_SUB_ADD_RU}}%
\FR{\input{fundamentals/AND_OR_as_SUB_ADD_FR}}

\EN{\input{fundamentals/XOR_EN}}\RU{\input{fundamentals/XOR_RU}}%
\FR{\input{fundamentals/XOR_FR}}

\EN{\input{fundamentals/AND_OR_XOR_as_MOV_EN}}
\ES{\input{fundamentals/AND_OR_XOR_as_MOV_ES}}
\RU{\input{fundamentals/AND_OR_XOR_as_MOV_RU}}
\FR{\input{fundamentals/AND_OR_XOR_as_MOV_FR}}

\EN{\input{fundamentals/POPCNT_EN}}\RU{\input{fundamentals/POPCNT_RU}}%
\FR{\input{fundamentals/POPCNT_FR}}

\EN{\input{fundamentals/endianness_EN}}
\ES{\input{fundamentals/endianness_ES}}
\RU{\input{fundamentals/endianness_RU}}
\FR{\input{fundamentals/endianness_FR}}

\EN{\input{fundamentals/memory_EN}}
\ES{\input{fundamentals/memory_ES}}
\RU{\input{fundamentals/memory_RU}}
\FR{\input{fundamentals/memory_FR}}

\EN{\input{fundamentals/CPU_EN}}
\ES{\input{fundamentals/CPU_ES}}
\RU{\input{fundamentals/CPU_RU}}
\FR{\input{fundamentals/CPU_FR}}

\EN{\input{fundamentals/hash_EN}}
\ES{\input{fundamentals/hash_ES}}
\RU{\input{fundamentals/hash_RU}}
\FR{\input{fundamentals/hash_FR}}


\chapter{\EN{Slightly more advanced examples}\RU{Более сложные примеры}\DE{Fortgeschrittenere Beispiele}\FR{Exemples un peu plus avancés}}

% sections here:

\newcommand{\CURPATH}{advanced/030_dbl_neg}
\EN{\input{\CURPATH/main_EN}}
\RU{\input{\CURPATH/main_RU}}
\FR{\input{\CURPATH/main_FR}}

\renewcommand{\CURPATH}{advanced/050_strstr}
\EN{\input{\CURPATH/main_EN}}\DE{\input{\CURPATH/main_DE}}\RU{\input{\CURPATH/main_RU}}\FR{\input{\CURPATH/main_FR}}

\renewcommand{\CURPATH}{advanced/100_fahrenheit}
\EN{\input{\CURPATH/main_EN}}\RU{\input{\CURPATH/main_RU}}\FR{\input{\CURPATH/main_FR}}

\renewcommand{\CURPATH}{advanced/102_fib}
\EN{\input{\CURPATH/main}}\RU{\input{\CURPATH/main}}\FR{\input{\CURPATH/main}}

\renewcommand{\CURPATH}{advanced/110_CRC32}
\EN{\input{\CURPATH/main_EN}}\RU{\input{\CURPATH/main_RU}}\FR{\input{\CURPATH/main_FR}}

\renewcommand{\CURPATH}{advanced/111_netmask}
\EN{\input{\CURPATH/main_EN}}\RU{\input{\CURPATH/main_RU}}\FR{\input{\CURPATH/main_FR}}

\renewcommand{\CURPATH}{advanced/115_loop_iterators}
\EN{\input{\CURPATH/main_EN}}\RU{\input{\CURPATH/main_RU}}\FR{\input{\CURPATH/main_FR}}

\renewcommand{\CURPATH}{advanced/117_duff_device}
\EN{\input{\CURPATH/main_EN}}\RU{\input{\CURPATH/main_RU}}\FR{\input{\CURPATH/main_FR}}

\renewcommand{\CURPATH}{advanced/120_division_by_mult}
\EN{\input{\CURPATH/main}}\RU{\input{\CURPATH/main}}\FR{\input{\CURPATH/main}}

\renewcommand{\CURPATH}{advanced/125_atoi}
\EN{\input{\CURPATH/main_EN}}\RU{\input{\CURPATH/main_RU}}\FR{\input{\CURPATH/main_FR}}

\renewcommand{\CURPATH}{advanced/127_inline_function}
\EN{\input{\CURPATH/main_EN}}\RU{\input{\CURPATH/main_RU}}\FR{\input{\CURPATH/main_FR}}

\renewcommand{\CURPATH}{advanced/130_C99_restrict}
\EN{\input{\CURPATH/main_EN}}\RU{\input{\CURPATH/main_RU}}\FR{\input{\CURPATH/main_FR}}

\renewcommand{\CURPATH}{advanced/135_abs_branchless}
\EN{\input{\CURPATH/main_EN}}\RU{\input{\CURPATH/main_RU}}\FR{\input{\CURPATH/main_FR}}

\renewcommand{\CURPATH}{advanced/170_variadic_functions}
\EN{\input{\CURPATH/main_EN}}\RU{\input{\CURPATH/main_RU}}\FR{\input{\CURPATH/main_FR}}

\renewcommand{\CURPATH}{advanced/200_string_trim}
\EN{\input{\CURPATH/main_EN}}\RU{\input{\CURPATH/main_RU}}\FR{\input{\CURPATH/main_FR}}

\renewcommand{\CURPATH}{advanced/250_toupper}
\EN{\input{\CURPATH/main_EN}}\RU{\input{\CURPATH/main_RU}}\FR{\input{\CURPATH/main_FR}}

\renewcommand{\CURPATH}{advanced/310_obfuscation}
\EN{\input{\CURPATH/main_EN}}\RU{\input{\CURPATH/main_RU}}\FR{\input{\CURPATH/main_FR}}

\renewcommand{\CURPATH}{advanced/350_cpp}
\EN{\input{\CURPATH/main}}\RU{\input{\CURPATH/main}}\FR{\input{\CURPATH/main}}

\renewcommand{\CURPATH}{advanced/370_neg_arrays}
\EN{\input{\CURPATH/main_EN}}\RU{\input{\CURPATH/main_RU}}\FR{\input{\CURPATH/main_FR}}

\renewcommand{\CURPATH}{advanced/380_FAT12}
\EN{\input{\CURPATH/main_EN}}\FR{\input{\CURPATH/main_FR}}

\EN{\input{advanced/450_more_ptrs/main_EN}}\RU{\input{advanced/450_more_ptrs/main_RU}}\FR{\input{advanced/450_more_ptrs/main_FR}}

\EN{\input{advanced/500_loop_optimizations/main_EN}}\RU{\input{advanced/500_loop_optimizations/main_RU}}\FR{\input{advanced/500_loop_optimizations/main_FR}}

\EN{\input{advanced/550_more_structs/main_EN}}\RU{\input{advanced/550_more_structs/main_RU}}\FR{\input{advanced/550_more_structs/main_FR}}

\EN{\input{advanced/600_memmove/main_EN}}\RU{\input{advanced/600_memmove/main_RU}}\FR{\input{advanced/600_memmove/main_FR}}

\EN{\input{advanced/625_setjmp/main_EN}}\RU{\input{advanced/625_setjmp/main_RU}}\FR{\input{advanced/625_setjmp/main_FR}}

\EN{\input{advanced/650_stack/main}}\RU{\input{advanced/650_stack/main}}\FR{\input{advanced/650_stack/main}}

\EN{\input{advanced/700_openmp/main_EN}}\RU{\input{advanced/700_openmp/main_RU}}\FR{\input{advanced/700_openmp/main_FR}}

\EN{\input{advanced/750_heisenbug/main_EN}}\FR{\input{advanced/750_heisenbug/main_FR}}

\renewcommand{\CURPATH}{advanced/760_no_return}
\EN{\input{advanced/760_no_return/main_EN}}

\renewcommand{\CURPATH}{advanced/800_win16}
\EN{\input{\CURPATH/main}}\RU{\input{\CURPATH/main}}\FR{\input{\CURPATH/main_FR}}



\EN{\documentclass[a4paper,oneside]{book}

% http://www.tex.ac.uk/FAQ-noroom.html
\usepackage{etex}

\usepackage[table,usenames,dvipsnames]{xcolor}

\usepackage{fontspec}
% fonts
%\setmonofont{DroidSansMono}
%\setmainfont[Ligatures=TeX]{PT Sans}
%\setmainfont{DroidSans}
\setmainfont{DejaVu Sans}
\setmonofont{DejaVu Sans Mono}
\usepackage{polyglossia}
\defaultfontfeatures{Scale=MatchLowercase} % ensure all fonts have the same 1ex
\usepackage{ucharclasses}
\usepackage{csquotes}

\ifdefined\ENGLISH
%\wlog{main ENGLISH defined OK}
\setmainlanguage{english}
\setotherlanguage{russian}
\fi

\ifdefined\RUSSIAN
\setmainlanguage{russian}
%\newfontfamily\cyrillicfont{LiberationSans}
%\newfontfamily\cyrillicfonttt{LiberationMono}
%\newfontfamily\cyrillicfontsf{lmsans10-regular.otf}
\setotherlanguage{english}
\fi

\ifdefined\GERMAN
%\wlog{main GERMAN defined OK}
\setmainlanguage{german}
\setotherlanguage{english}
\fi

\ifdefined\SPANISH
\setmainlanguage{spanish}
\setotherlanguage{english}
\fi

\ifdefined\ITALIAN
\setmainlanguage{italian}
\setotherlanguage{english}
\fi

\ifdefined\BRAZILIAN
\setmainlanguage{portuges}
\setotherlanguage{english}
\fi

\ifdefined\POLISH
\setmainlanguage{polish}
\setotherlanguage{english}
\fi

\ifdefined\DUTCH
\setmainlanguage{dutch}
\setotherlanguage{english}
\fi

\ifdefined\TURKISH
\setmainlanguage{turkish}
\setotherlanguage{english}
\fi

\ifdefined\THAI
\setmainlanguage{thai}
%\usepackage[thai]{babel}
%\usepackage{fonts-tlwg}
\setmainfont[Script=Thai]{TH SarabunPSK}
\newfontfamily{\thaifont}[Script=Thai]{TH SarabunPSK}
\let\thaifonttt\ttfamily
\setotherlanguage{english}
\fi

\ifdefined\FRENCH
\setmainlanguage{french}
\setotherlanguage{english}
\fi

\ifdefined\JAPANESE
\usepackage{xeCJK}
\xeCJKallowbreakbetweenpuncts
\defaultfontfeatures{Ligatures=TeX,Scale=MatchLowercase}
\setCJKmainfont{IPAGothic}
\setCJKsansfont{IPAGothic}
\setCJKmonofont{IPAGothic}
\DeclareQuoteStyle{japanese}
  {「}
  {」}
  {『}
  {』}
\setquotestyle{japanese}
\setmainlanguage{japanese}
\setotherlanguage{english}
\fi

\usepackage{microtype}
\usepackage{fancyhdr}
\usepackage{listings}
%\usepackage{ulem} % used for \sout{}...
\usepackage{url}
\usepackage{graphicx}
\usepackage{makeidx}
\usepackage[cm]{fullpage}
%\usepackage{color}
\usepackage{fancyvrb}
\usepackage{xspace}
\usepackage{tabularx}
\usepackage{framed}
\usepackage{parskip}
\usepackage{epigraph}
\usepackage{ccicons}
\usepackage[nottoc]{tocbibind}
\usepackage{longtable}
\usepackage[footnote,printonlyused,withpage]{acronym}
\usepackage[]{bookmark,hyperref} % must be last
\usepackage[official]{eurosym}
\usepackage[usestackEOL]{stackengine}

% ************** myref
% http://tex.stackexchange.com/questions/228286/how-to-mix-ref-and-pageref#228292
\ifdefined\RUSSIAN
\newcommand{\myref}[1]{%
  \ref{#1}
  (стр.~\pageref{#1})%
  }
% FIXME: I wasn't able to force varioref to output russian text...
\else
\usepackage{varioref}
\newcommand{\myref}[1]{\vref{#1}}
\fi
% ************** myref

\usepackage{glossaries}
\usepackage{tikz}
%\usepackage{fixltx2e}
\usepackage{bytefield}

\usepackage{amsmath}
\usepackage{MnSymbol}
\undef\mathdollar

\usepackage{float}

\usepackage{shorttoc}
\usetikzlibrary{calc,positioning,chains,arrows}
\usepackage[margin=0.5in,headheight=15.5pt]{geometry}
% FIXME would be good without "margin="...
%\usepackage[headheight=15.5pt]{geometry}

%--------------------
% to prevent clashing of numbers and titles in TOC:
% https://tex.stackexchange.com/a/64124
\usepackage{tocloft}% http://ctan.org/pkg/tocloft
\makeatletter
\renewcommand{\numberline}[1]{%
  \@cftbsnum #1\@cftasnum~\@cftasnumb%
}
\makeatother
%--------------------

\newcommand{\footnoteref}[1]{\textsuperscript{\ref{#1}}}

%\definecolor{lstbgcolor}{rgb}{0.94,0.94,0.94}

% I don't know why this voodoo works, but without all-caps, it can't find LIGHT-GRAY color. WTF?
% see also: https://tex.stackexchange.com/questions/64298/error-with-xcolor-package
\definecolor{light-gray}{gray}{0.87}
\definecolor{LIGHT-GRAY}{gray}{0.87}
\definecolor{RED}{rgb}{1,0,0}
\makeindex

\input{macros}
\EN{\input{macro_lang_EN}}
\RU{\input{macro_lang_RU}}
\FR{\input{macro_lang_FR}}
\DE{\input{macro_lang_DE}}
\JA{\input{macro_lang_JP}}
\IT{\input{macro_lang_IT}}
\PL{\input{macro_lang_PL}}
\ES{\input{macro_lang_ES}}
\NL{\input{macro_lang_NL}}
\TR{\input{macro_lang_TR}}
\PTBR{\input{macro_lang_PTBR}}

\ifdefined\UAL
\newcommand{\TitleMain}{\TitleUAL}
\newcommand{\TitleAux}{\TitleRE}
\else
\newcommand{\TitleMain}{\TitleRE}
\newcommand{\TitleAux}{\TitleUAL}
\fi

\EN{\input{glossary_EN}}
\RU{\input{glossary_RU}}
\FR{\input{glossary_FR}}
\DE{\input{glossary_DE}}
\IT{\input{glossary_IT}}
\JA{\input{glossary_JA}}

\makeglossaries

\hypersetup{
    colorlinks=true,
    allcolors=blue,
    pdfauthor={\AUTHOR},
    pdftitle={\TitleMain}
    }

%\ifdefined\RUSSIAN
\newcommand{\LstStyle}{\ttfamily\small}
%\else
%\newcommand{\LstStyle}{\ttfamily}
%\fi

% inspired by http://prismjs.com/
\definecolor{digits}{RGB}{0,0,0}
\definecolor{bg}{RGB}{255,255,255}
%\definecolor{bg}{RGB}{255,252,250}
\definecolor{col1}{RGB}{154,20,150}
\definecolor{col2}{RGB}{112,128,144}
\definecolor{col3}{RGB}{10,120,180}
\definecolor{col4}{RGB}{106,164,108}


\lstset{
    %backgroundcolor=\color{lstbgcolor},
    %backgroundcolor=\color{light-gray},
    backgroundcolor=\color{bg},
    basicstyle=\LstStyle,
    breaklines=true,
    %prebreak=\raisebox{0ex}[0ex][0ex]{->},
    %postbreak=\raisebox{0ex}[0ex][0ex]{->},
    prebreak=\raisebox{0ex}[0ex][0ex]{\ensuremath{\rhookswarrow}},
    postbreak=\raisebox{0ex}[0ex][0ex]{\ensuremath{\rcurvearrowse\space}},
    frame=single,
    columns=fullflexible,keepspaces,
    escapeinside=§§,
    inputencoding=utf8
}

\input{syntax_color}

\ifdefined\RUSSIAN
\renewcommand\lstlistingname{Листинг}
\renewcommand\lstlistlistingname{Листинг}
\fi

\DeclareMathSizes{12}{30}{16}{12}%

% see also:
% http://tex.stackexchange.com/questions/129225/how-can-i-get-get-makeindex-to-ignore-capital-letters
% http://tex.stackexchange.com/questions/18336/correct-sorting-of-index-entries-containing-macros
\def\myindex#1{\expandafter\index\expandafter{#1}}

\begin{document}

% fancyhdr
\pagestyle{fancy}
\setlength{\headheight}{13pt}
\fancyhead[R]{} % suppress chapter name

\VerbatimFootnotes

\frontmatter

%\RU{\include{1st_page_RU}}
%\EN{\include{1st_page_EN}}
%\DE{\include{1st_page_DE}}
%\ES{\include{1st_page_ES}}
%\CN{\include{1st_page_CN}}

\RU{\include{dedication_RU}}
\EN{\include{dedication_EN}}
\FR{\include{dedication_FR}}
\JA{\include{dedication_JA}}
\DE{\include{dedication_DE}}
\IT{\include{dedication_IT}}

\include{page_after_cover}
\include{call_for_translators}

\shorttoc{%
    \RU{Краткое оглавление}%
    \EN{Abridged contents}%
    \ES{Contenidos abreviados}%
    \PTBRph{}%
    \DE{Inhaltsverzeichnis (gekürzt)}%
    \PLph{}%
    \IT{Sommario}%
    \THAph{}\NLph{}%
    \FR{Contenus abrégés}%
    \JA{簡略版}
    \TR{İçindekiler}
}{0}

\tableofcontents
\cleardoublepage

\cleardoublepage
\include{preface}

\mainmatter

\include{parts}

\EN{\include{appendix/appendix}}
\RU{\include{appendix/appendix}}
\DE{\include{appendix/appendix}}
\FR{\include{appendix/appendix}}
\IT{\include{appendix/appendix}}
\include{acronyms}

\bookmarksetup{startatroot}

\clearpage
\phantomsection
\addcontentsline{toc}{chapter}{%
    \RU{Глоссарий}%
    \EN{Glossary}%
    \ES{Glosario}%
    \PTBRph{}%
    \DE{Glossar}%
    \PLph{}%
    \IT{Glossario}%
    \THAph{}\NLph{}%
    \FR{Glossaire}%
    \JA{用語}
    \TR{Bolum}
}
\printglossaries

\clearpage
\phantomsection
\printindex

\end{document}
}
\RU{\documentclass[a4paper,oneside]{book}

% http://www.tex.ac.uk/FAQ-noroom.html
\usepackage{etex}

\usepackage[table,usenames,dvipsnames]{xcolor}

\usepackage{fontspec}
% fonts
%\setmonofont{DroidSansMono}
%\setmainfont[Ligatures=TeX]{PT Sans}
%\setmainfont{DroidSans}
\setmainfont{DejaVu Sans}
\setmonofont{DejaVu Sans Mono}
\usepackage{polyglossia}
\defaultfontfeatures{Scale=MatchLowercase} % ensure all fonts have the same 1ex
\usepackage{ucharclasses}
\usepackage{csquotes}

\ifdefined\ENGLISH
%\wlog{main ENGLISH defined OK}
\setmainlanguage{english}
\setotherlanguage{russian}
\fi

\ifdefined\RUSSIAN
\setmainlanguage{russian}
%\newfontfamily\cyrillicfont{LiberationSans}
%\newfontfamily\cyrillicfonttt{LiberationMono}
%\newfontfamily\cyrillicfontsf{lmsans10-regular.otf}
\setotherlanguage{english}
\fi

\ifdefined\GERMAN
%\wlog{main GERMAN defined OK}
\setmainlanguage{german}
\setotherlanguage{english}
\fi

\ifdefined\SPANISH
\setmainlanguage{spanish}
\setotherlanguage{english}
\fi

\ifdefined\ITALIAN
\setmainlanguage{italian}
\setotherlanguage{english}
\fi

\ifdefined\BRAZILIAN
\setmainlanguage{portuges}
\setotherlanguage{english}
\fi

\ifdefined\POLISH
\setmainlanguage{polish}
\setotherlanguage{english}
\fi

\ifdefined\DUTCH
\setmainlanguage{dutch}
\setotherlanguage{english}
\fi

\ifdefined\TURKISH
\setmainlanguage{turkish}
\setotherlanguage{english}
\fi

\ifdefined\THAI
\setmainlanguage{thai}
%\usepackage[thai]{babel}
%\usepackage{fonts-tlwg}
\setmainfont[Script=Thai]{TH SarabunPSK}
\newfontfamily{\thaifont}[Script=Thai]{TH SarabunPSK}
\let\thaifonttt\ttfamily
\setotherlanguage{english}
\fi

\ifdefined\FRENCH
\setmainlanguage{french}
\setotherlanguage{english}
\fi

\ifdefined\JAPANESE
\usepackage{xeCJK}
\xeCJKallowbreakbetweenpuncts
\defaultfontfeatures{Ligatures=TeX,Scale=MatchLowercase}
\setCJKmainfont{IPAGothic}
\setCJKsansfont{IPAGothic}
\setCJKmonofont{IPAGothic}
\DeclareQuoteStyle{japanese}
  {「}
  {」}
  {『}
  {』}
\setquotestyle{japanese}
\setmainlanguage{japanese}
\setotherlanguage{english}
\fi

\usepackage{microtype}
\usepackage{fancyhdr}
\usepackage{listings}
%\usepackage{ulem} % used for \sout{}...
\usepackage{url}
\usepackage{graphicx}
\usepackage{makeidx}
\usepackage[cm]{fullpage}
%\usepackage{color}
\usepackage{fancyvrb}
\usepackage{xspace}
\usepackage{tabularx}
\usepackage{framed}
\usepackage{parskip}
\usepackage{epigraph}
\usepackage{ccicons}
\usepackage[nottoc]{tocbibind}
\usepackage{longtable}
\usepackage[footnote,printonlyused,withpage]{acronym}
\usepackage[]{bookmark,hyperref} % must be last
\usepackage[official]{eurosym}
\usepackage[usestackEOL]{stackengine}

% ************** myref
% http://tex.stackexchange.com/questions/228286/how-to-mix-ref-and-pageref#228292
\ifdefined\RUSSIAN
\newcommand{\myref}[1]{%
  \ref{#1}
  (стр.~\pageref{#1})%
  }
% FIXME: I wasn't able to force varioref to output russian text...
\else
\usepackage{varioref}
\newcommand{\myref}[1]{\vref{#1}}
\fi
% ************** myref

\usepackage{glossaries}
\usepackage{tikz}
%\usepackage{fixltx2e}
\usepackage{bytefield}

\usepackage{amsmath}
\usepackage{MnSymbol}
\undef\mathdollar

\usepackage{float}

\usepackage{shorttoc}
\usetikzlibrary{calc,positioning,chains,arrows}
\usepackage[margin=0.5in,headheight=15.5pt]{geometry}
% FIXME would be good without "margin="...
%\usepackage[headheight=15.5pt]{geometry}

%--------------------
% to prevent clashing of numbers and titles in TOC:
% https://tex.stackexchange.com/a/64124
\usepackage{tocloft}% http://ctan.org/pkg/tocloft
\makeatletter
\renewcommand{\numberline}[1]{%
  \@cftbsnum #1\@cftasnum~\@cftasnumb%
}
\makeatother
%--------------------

\newcommand{\footnoteref}[1]{\textsuperscript{\ref{#1}}}

%\definecolor{lstbgcolor}{rgb}{0.94,0.94,0.94}

% I don't know why this voodoo works, but without all-caps, it can't find LIGHT-GRAY color. WTF?
% see also: https://tex.stackexchange.com/questions/64298/error-with-xcolor-package
\definecolor{light-gray}{gray}{0.87}
\definecolor{LIGHT-GRAY}{gray}{0.87}
\definecolor{RED}{rgb}{1,0,0}
\makeindex

\input{macros}
\EN{\input{macro_lang_EN}}
\RU{\input{macro_lang_RU}}
\FR{\input{macro_lang_FR}}
\DE{\input{macro_lang_DE}}
\JA{\input{macro_lang_JP}}
\IT{\input{macro_lang_IT}}
\PL{\input{macro_lang_PL}}
\ES{\input{macro_lang_ES}}
\NL{\input{macro_lang_NL}}
\TR{\input{macro_lang_TR}}
\PTBR{\input{macro_lang_PTBR}}

\ifdefined\UAL
\newcommand{\TitleMain}{\TitleUAL}
\newcommand{\TitleAux}{\TitleRE}
\else
\newcommand{\TitleMain}{\TitleRE}
\newcommand{\TitleAux}{\TitleUAL}
\fi

\EN{\input{glossary_EN}}
\RU{\input{glossary_RU}}
\FR{\input{glossary_FR}}
\DE{\input{glossary_DE}}
\IT{\input{glossary_IT}}
\JA{\input{glossary_JA}}

\makeglossaries

\hypersetup{
    colorlinks=true,
    allcolors=blue,
    pdfauthor={\AUTHOR},
    pdftitle={\TitleMain}
    }

%\ifdefined\RUSSIAN
\newcommand{\LstStyle}{\ttfamily\small}
%\else
%\newcommand{\LstStyle}{\ttfamily}
%\fi

% inspired by http://prismjs.com/
\definecolor{digits}{RGB}{0,0,0}
\definecolor{bg}{RGB}{255,255,255}
%\definecolor{bg}{RGB}{255,252,250}
\definecolor{col1}{RGB}{154,20,150}
\definecolor{col2}{RGB}{112,128,144}
\definecolor{col3}{RGB}{10,120,180}
\definecolor{col4}{RGB}{106,164,108}


\lstset{
    %backgroundcolor=\color{lstbgcolor},
    %backgroundcolor=\color{light-gray},
    backgroundcolor=\color{bg},
    basicstyle=\LstStyle,
    breaklines=true,
    %prebreak=\raisebox{0ex}[0ex][0ex]{->},
    %postbreak=\raisebox{0ex}[0ex][0ex]{->},
    prebreak=\raisebox{0ex}[0ex][0ex]{\ensuremath{\rhookswarrow}},
    postbreak=\raisebox{0ex}[0ex][0ex]{\ensuremath{\rcurvearrowse\space}},
    frame=single,
    columns=fullflexible,keepspaces,
    escapeinside=§§,
    inputencoding=utf8
}

\input{syntax_color}

\ifdefined\RUSSIAN
\renewcommand\lstlistingname{Листинг}
\renewcommand\lstlistlistingname{Листинг}
\fi

\DeclareMathSizes{12}{30}{16}{12}%

% see also:
% http://tex.stackexchange.com/questions/129225/how-can-i-get-get-makeindex-to-ignore-capital-letters
% http://tex.stackexchange.com/questions/18336/correct-sorting-of-index-entries-containing-macros
\def\myindex#1{\expandafter\index\expandafter{#1}}

\begin{document}

% fancyhdr
\pagestyle{fancy}
\setlength{\headheight}{13pt}
\fancyhead[R]{} % suppress chapter name

\VerbatimFootnotes

\frontmatter

%\RU{\include{1st_page_RU}}
%\EN{\include{1st_page_EN}}
%\DE{\include{1st_page_DE}}
%\ES{\include{1st_page_ES}}
%\CN{\include{1st_page_CN}}

\RU{\include{dedication_RU}}
\EN{\include{dedication_EN}}
\FR{\include{dedication_FR}}
\JA{\include{dedication_JA}}
\DE{\include{dedication_DE}}
\IT{\include{dedication_IT}}

\include{page_after_cover}
\include{call_for_translators}

\shorttoc{%
    \RU{Краткое оглавление}%
    \EN{Abridged contents}%
    \ES{Contenidos abreviados}%
    \PTBRph{}%
    \DE{Inhaltsverzeichnis (gekürzt)}%
    \PLph{}%
    \IT{Sommario}%
    \THAph{}\NLph{}%
    \FR{Contenus abrégés}%
    \JA{簡略版}
    \TR{İçindekiler}
}{0}

\tableofcontents
\cleardoublepage

\cleardoublepage
\include{preface}

\mainmatter

\include{parts}

\EN{\include{appendix/appendix}}
\RU{\include{appendix/appendix}}
\DE{\include{appendix/appendix}}
\FR{\include{appendix/appendix}}
\IT{\include{appendix/appendix}}
\include{acronyms}

\bookmarksetup{startatroot}

\clearpage
\phantomsection
\addcontentsline{toc}{chapter}{%
    \RU{Глоссарий}%
    \EN{Glossary}%
    \ES{Glosario}%
    \PTBRph{}%
    \DE{Glossar}%
    \PLph{}%
    \IT{Glossario}%
    \THAph{}\NLph{}%
    \FR{Glossaire}%
    \JA{用語}
    \TR{Bolum}
}
\printglossaries

\clearpage
\phantomsection
\printindex

\end{document}
}
\DE{\documentclass[a4paper,oneside]{book}

% http://www.tex.ac.uk/FAQ-noroom.html
\usepackage{etex}

\usepackage[table,usenames,dvipsnames]{xcolor}

\usepackage{fontspec}
% fonts
%\setmonofont{DroidSansMono}
%\setmainfont[Ligatures=TeX]{PT Sans}
%\setmainfont{DroidSans}
\setmainfont{DejaVu Sans}
\setmonofont{DejaVu Sans Mono}
\usepackage{polyglossia}
\defaultfontfeatures{Scale=MatchLowercase} % ensure all fonts have the same 1ex
\usepackage{ucharclasses}
\usepackage{csquotes}

\ifdefined\ENGLISH
%\wlog{main ENGLISH defined OK}
\setmainlanguage{english}
\setotherlanguage{russian}
\fi

\ifdefined\RUSSIAN
\setmainlanguage{russian}
%\newfontfamily\cyrillicfont{LiberationSans}
%\newfontfamily\cyrillicfonttt{LiberationMono}
%\newfontfamily\cyrillicfontsf{lmsans10-regular.otf}
\setotherlanguage{english}
\fi

\ifdefined\GERMAN
%\wlog{main GERMAN defined OK}
\setmainlanguage{german}
\setotherlanguage{english}
\fi

\ifdefined\SPANISH
\setmainlanguage{spanish}
\setotherlanguage{english}
\fi

\ifdefined\ITALIAN
\setmainlanguage{italian}
\setotherlanguage{english}
\fi

\ifdefined\BRAZILIAN
\setmainlanguage{portuges}
\setotherlanguage{english}
\fi

\ifdefined\POLISH
\setmainlanguage{polish}
\setotherlanguage{english}
\fi

\ifdefined\DUTCH
\setmainlanguage{dutch}
\setotherlanguage{english}
\fi

\ifdefined\TURKISH
\setmainlanguage{turkish}
\setotherlanguage{english}
\fi

\ifdefined\THAI
\setmainlanguage{thai}
%\usepackage[thai]{babel}
%\usepackage{fonts-tlwg}
\setmainfont[Script=Thai]{TH SarabunPSK}
\newfontfamily{\thaifont}[Script=Thai]{TH SarabunPSK}
\let\thaifonttt\ttfamily
\setotherlanguage{english}
\fi

\ifdefined\FRENCH
\setmainlanguage{french}
\setotherlanguage{english}
\fi

\ifdefined\JAPANESE
\usepackage{xeCJK}
\xeCJKallowbreakbetweenpuncts
\defaultfontfeatures{Ligatures=TeX,Scale=MatchLowercase}
\setCJKmainfont{IPAGothic}
\setCJKsansfont{IPAGothic}
\setCJKmonofont{IPAGothic}
\DeclareQuoteStyle{japanese}
  {「}
  {」}
  {『}
  {』}
\setquotestyle{japanese}
\setmainlanguage{japanese}
\setotherlanguage{english}
\fi

\usepackage{microtype}
\usepackage{fancyhdr}
\usepackage{listings}
%\usepackage{ulem} % used for \sout{}...
\usepackage{url}
\usepackage{graphicx}
\usepackage{makeidx}
\usepackage[cm]{fullpage}
%\usepackage{color}
\usepackage{fancyvrb}
\usepackage{xspace}
\usepackage{tabularx}
\usepackage{framed}
\usepackage{parskip}
\usepackage{epigraph}
\usepackage{ccicons}
\usepackage[nottoc]{tocbibind}
\usepackage{longtable}
\usepackage[footnote,printonlyused,withpage]{acronym}
\usepackage[]{bookmark,hyperref} % must be last
\usepackage[official]{eurosym}
\usepackage[usestackEOL]{stackengine}

% ************** myref
% http://tex.stackexchange.com/questions/228286/how-to-mix-ref-and-pageref#228292
\ifdefined\RUSSIAN
\newcommand{\myref}[1]{%
  \ref{#1}
  (стр.~\pageref{#1})%
  }
% FIXME: I wasn't able to force varioref to output russian text...
\else
\usepackage{varioref}
\newcommand{\myref}[1]{\vref{#1}}
\fi
% ************** myref

\usepackage{glossaries}
\usepackage{tikz}
%\usepackage{fixltx2e}
\usepackage{bytefield}

\usepackage{amsmath}
\usepackage{MnSymbol}
\undef\mathdollar

\usepackage{float}

\usepackage{shorttoc}
\usetikzlibrary{calc,positioning,chains,arrows}
\usepackage[margin=0.5in,headheight=15.5pt]{geometry}
% FIXME would be good without "margin="...
%\usepackage[headheight=15.5pt]{geometry}

%--------------------
% to prevent clashing of numbers and titles in TOC:
% https://tex.stackexchange.com/a/64124
\usepackage{tocloft}% http://ctan.org/pkg/tocloft
\makeatletter
\renewcommand{\numberline}[1]{%
  \@cftbsnum #1\@cftasnum~\@cftasnumb%
}
\makeatother
%--------------------

\newcommand{\footnoteref}[1]{\textsuperscript{\ref{#1}}}

%\definecolor{lstbgcolor}{rgb}{0.94,0.94,0.94}

% I don't know why this voodoo works, but without all-caps, it can't find LIGHT-GRAY color. WTF?
% see also: https://tex.stackexchange.com/questions/64298/error-with-xcolor-package
\definecolor{light-gray}{gray}{0.87}
\definecolor{LIGHT-GRAY}{gray}{0.87}
\definecolor{RED}{rgb}{1,0,0}
\makeindex

\input{macros}
\EN{\input{macro_lang_EN}}
\RU{\input{macro_lang_RU}}
\FR{\input{macro_lang_FR}}
\DE{\input{macro_lang_DE}}
\JA{\input{macro_lang_JP}}
\IT{\input{macro_lang_IT}}
\PL{\input{macro_lang_PL}}
\ES{\input{macro_lang_ES}}
\NL{\input{macro_lang_NL}}
\TR{\input{macro_lang_TR}}
\PTBR{\input{macro_lang_PTBR}}

\ifdefined\UAL
\newcommand{\TitleMain}{\TitleUAL}
\newcommand{\TitleAux}{\TitleRE}
\else
\newcommand{\TitleMain}{\TitleRE}
\newcommand{\TitleAux}{\TitleUAL}
\fi

\EN{\input{glossary_EN}}
\RU{\input{glossary_RU}}
\FR{\input{glossary_FR}}
\DE{\input{glossary_DE}}
\IT{\input{glossary_IT}}
\JA{\input{glossary_JA}}

\makeglossaries

\hypersetup{
    colorlinks=true,
    allcolors=blue,
    pdfauthor={\AUTHOR},
    pdftitle={\TitleMain}
    }

%\ifdefined\RUSSIAN
\newcommand{\LstStyle}{\ttfamily\small}
%\else
%\newcommand{\LstStyle}{\ttfamily}
%\fi

% inspired by http://prismjs.com/
\definecolor{digits}{RGB}{0,0,0}
\definecolor{bg}{RGB}{255,255,255}
%\definecolor{bg}{RGB}{255,252,250}
\definecolor{col1}{RGB}{154,20,150}
\definecolor{col2}{RGB}{112,128,144}
\definecolor{col3}{RGB}{10,120,180}
\definecolor{col4}{RGB}{106,164,108}


\lstset{
    %backgroundcolor=\color{lstbgcolor},
    %backgroundcolor=\color{light-gray},
    backgroundcolor=\color{bg},
    basicstyle=\LstStyle,
    breaklines=true,
    %prebreak=\raisebox{0ex}[0ex][0ex]{->},
    %postbreak=\raisebox{0ex}[0ex][0ex]{->},
    prebreak=\raisebox{0ex}[0ex][0ex]{\ensuremath{\rhookswarrow}},
    postbreak=\raisebox{0ex}[0ex][0ex]{\ensuremath{\rcurvearrowse\space}},
    frame=single,
    columns=fullflexible,keepspaces,
    escapeinside=§§,
    inputencoding=utf8
}

\input{syntax_color}

\ifdefined\RUSSIAN
\renewcommand\lstlistingname{Листинг}
\renewcommand\lstlistlistingname{Листинг}
\fi

\DeclareMathSizes{12}{30}{16}{12}%

% see also:
% http://tex.stackexchange.com/questions/129225/how-can-i-get-get-makeindex-to-ignore-capital-letters
% http://tex.stackexchange.com/questions/18336/correct-sorting-of-index-entries-containing-macros
\def\myindex#1{\expandafter\index\expandafter{#1}}

\begin{document}

% fancyhdr
\pagestyle{fancy}
\setlength{\headheight}{13pt}
\fancyhead[R]{} % suppress chapter name

\VerbatimFootnotes

\frontmatter

%\RU{\include{1st_page_RU}}
%\EN{\include{1st_page_EN}}
%\DE{\include{1st_page_DE}}
%\ES{\include{1st_page_ES}}
%\CN{\include{1st_page_CN}}

\RU{\include{dedication_RU}}
\EN{\include{dedication_EN}}
\FR{\include{dedication_FR}}
\JA{\include{dedication_JA}}
\DE{\include{dedication_DE}}
\IT{\include{dedication_IT}}

\include{page_after_cover}
\include{call_for_translators}

\shorttoc{%
    \RU{Краткое оглавление}%
    \EN{Abridged contents}%
    \ES{Contenidos abreviados}%
    \PTBRph{}%
    \DE{Inhaltsverzeichnis (gekürzt)}%
    \PLph{}%
    \IT{Sommario}%
    \THAph{}\NLph{}%
    \FR{Contenus abrégés}%
    \JA{簡略版}
    \TR{İçindekiler}
}{0}

\tableofcontents
\cleardoublepage

\cleardoublepage
\include{preface}

\mainmatter

\include{parts}

\EN{\include{appendix/appendix}}
\RU{\include{appendix/appendix}}
\DE{\include{appendix/appendix}}
\FR{\include{appendix/appendix}}
\IT{\include{appendix/appendix}}
\include{acronyms}

\bookmarksetup{startatroot}

\clearpage
\phantomsection
\addcontentsline{toc}{chapter}{%
    \RU{Глоссарий}%
    \EN{Glossary}%
    \ES{Glosario}%
    \PTBRph{}%
    \DE{Glossar}%
    \PLph{}%
    \IT{Glossario}%
    \THAph{}\NLph{}%
    \FR{Glossaire}%
    \JA{用語}
    \TR{Bolum}
}
\printglossaries

\clearpage
\phantomsection
\printindex

\end{document}
}
\FR{\documentclass[a4paper,oneside]{book}

% http://www.tex.ac.uk/FAQ-noroom.html
\usepackage{etex}

\usepackage[table,usenames,dvipsnames]{xcolor}

\usepackage{fontspec}
% fonts
%\setmonofont{DroidSansMono}
%\setmainfont[Ligatures=TeX]{PT Sans}
%\setmainfont{DroidSans}
\setmainfont{DejaVu Sans}
\setmonofont{DejaVu Sans Mono}
\usepackage{polyglossia}
\defaultfontfeatures{Scale=MatchLowercase} % ensure all fonts have the same 1ex
\usepackage{ucharclasses}
\usepackage{csquotes}

\ifdefined\ENGLISH
%\wlog{main ENGLISH defined OK}
\setmainlanguage{english}
\setotherlanguage{russian}
\fi

\ifdefined\RUSSIAN
\setmainlanguage{russian}
%\newfontfamily\cyrillicfont{LiberationSans}
%\newfontfamily\cyrillicfonttt{LiberationMono}
%\newfontfamily\cyrillicfontsf{lmsans10-regular.otf}
\setotherlanguage{english}
\fi

\ifdefined\GERMAN
%\wlog{main GERMAN defined OK}
\setmainlanguage{german}
\setotherlanguage{english}
\fi

\ifdefined\SPANISH
\setmainlanguage{spanish}
\setotherlanguage{english}
\fi

\ifdefined\ITALIAN
\setmainlanguage{italian}
\setotherlanguage{english}
\fi

\ifdefined\BRAZILIAN
\setmainlanguage{portuges}
\setotherlanguage{english}
\fi

\ifdefined\POLISH
\setmainlanguage{polish}
\setotherlanguage{english}
\fi

\ifdefined\DUTCH
\setmainlanguage{dutch}
\setotherlanguage{english}
\fi

\ifdefined\TURKISH
\setmainlanguage{turkish}
\setotherlanguage{english}
\fi

\ifdefined\THAI
\setmainlanguage{thai}
%\usepackage[thai]{babel}
%\usepackage{fonts-tlwg}
\setmainfont[Script=Thai]{TH SarabunPSK}
\newfontfamily{\thaifont}[Script=Thai]{TH SarabunPSK}
\let\thaifonttt\ttfamily
\setotherlanguage{english}
\fi

\ifdefined\FRENCH
\setmainlanguage{french}
\setotherlanguage{english}
\fi

\ifdefined\JAPANESE
\usepackage{xeCJK}
\xeCJKallowbreakbetweenpuncts
\defaultfontfeatures{Ligatures=TeX,Scale=MatchLowercase}
\setCJKmainfont{IPAGothic}
\setCJKsansfont{IPAGothic}
\setCJKmonofont{IPAGothic}
\DeclareQuoteStyle{japanese}
  {「}
  {」}
  {『}
  {』}
\setquotestyle{japanese}
\setmainlanguage{japanese}
\setotherlanguage{english}
\fi

\usepackage{microtype}
\usepackage{fancyhdr}
\usepackage{listings}
%\usepackage{ulem} % used for \sout{}...
\usepackage{url}
\usepackage{graphicx}
\usepackage{makeidx}
\usepackage[cm]{fullpage}
%\usepackage{color}
\usepackage{fancyvrb}
\usepackage{xspace}
\usepackage{tabularx}
\usepackage{framed}
\usepackage{parskip}
\usepackage{epigraph}
\usepackage{ccicons}
\usepackage[nottoc]{tocbibind}
\usepackage{longtable}
\usepackage[footnote,printonlyused,withpage]{acronym}
\usepackage[]{bookmark,hyperref} % must be last
\usepackage[official]{eurosym}
\usepackage[usestackEOL]{stackengine}

% ************** myref
% http://tex.stackexchange.com/questions/228286/how-to-mix-ref-and-pageref#228292
\ifdefined\RUSSIAN
\newcommand{\myref}[1]{%
  \ref{#1}
  (стр.~\pageref{#1})%
  }
% FIXME: I wasn't able to force varioref to output russian text...
\else
\usepackage{varioref}
\newcommand{\myref}[1]{\vref{#1}}
\fi
% ************** myref

\usepackage{glossaries}
\usepackage{tikz}
%\usepackage{fixltx2e}
\usepackage{bytefield}

\usepackage{amsmath}
\usepackage{MnSymbol}
\undef\mathdollar

\usepackage{float}

\usepackage{shorttoc}
\usetikzlibrary{calc,positioning,chains,arrows}
\usepackage[margin=0.5in,headheight=15.5pt]{geometry}
% FIXME would be good without "margin="...
%\usepackage[headheight=15.5pt]{geometry}

%--------------------
% to prevent clashing of numbers and titles in TOC:
% https://tex.stackexchange.com/a/64124
\usepackage{tocloft}% http://ctan.org/pkg/tocloft
\makeatletter
\renewcommand{\numberline}[1]{%
  \@cftbsnum #1\@cftasnum~\@cftasnumb%
}
\makeatother
%--------------------

\newcommand{\footnoteref}[1]{\textsuperscript{\ref{#1}}}

%\definecolor{lstbgcolor}{rgb}{0.94,0.94,0.94}

% I don't know why this voodoo works, but without all-caps, it can't find LIGHT-GRAY color. WTF?
% see also: https://tex.stackexchange.com/questions/64298/error-with-xcolor-package
\definecolor{light-gray}{gray}{0.87}
\definecolor{LIGHT-GRAY}{gray}{0.87}
\definecolor{RED}{rgb}{1,0,0}
\makeindex

\input{macros}
\EN{\input{macro_lang_EN}}
\RU{\input{macro_lang_RU}}
\FR{\input{macro_lang_FR}}
\DE{\input{macro_lang_DE}}
\JA{\input{macro_lang_JP}}
\IT{\input{macro_lang_IT}}
\PL{\input{macro_lang_PL}}
\ES{\input{macro_lang_ES}}
\NL{\input{macro_lang_NL}}
\TR{\input{macro_lang_TR}}
\PTBR{\input{macro_lang_PTBR}}

\ifdefined\UAL
\newcommand{\TitleMain}{\TitleUAL}
\newcommand{\TitleAux}{\TitleRE}
\else
\newcommand{\TitleMain}{\TitleRE}
\newcommand{\TitleAux}{\TitleUAL}
\fi

\EN{\input{glossary_EN}}
\RU{\input{glossary_RU}}
\FR{\input{glossary_FR}}
\DE{\input{glossary_DE}}
\IT{\input{glossary_IT}}
\JA{\input{glossary_JA}}

\makeglossaries

\hypersetup{
    colorlinks=true,
    allcolors=blue,
    pdfauthor={\AUTHOR},
    pdftitle={\TitleMain}
    }

%\ifdefined\RUSSIAN
\newcommand{\LstStyle}{\ttfamily\small}
%\else
%\newcommand{\LstStyle}{\ttfamily}
%\fi

% inspired by http://prismjs.com/
\definecolor{digits}{RGB}{0,0,0}
\definecolor{bg}{RGB}{255,255,255}
%\definecolor{bg}{RGB}{255,252,250}
\definecolor{col1}{RGB}{154,20,150}
\definecolor{col2}{RGB}{112,128,144}
\definecolor{col3}{RGB}{10,120,180}
\definecolor{col4}{RGB}{106,164,108}


\lstset{
    %backgroundcolor=\color{lstbgcolor},
    %backgroundcolor=\color{light-gray},
    backgroundcolor=\color{bg},
    basicstyle=\LstStyle,
    breaklines=true,
    %prebreak=\raisebox{0ex}[0ex][0ex]{->},
    %postbreak=\raisebox{0ex}[0ex][0ex]{->},
    prebreak=\raisebox{0ex}[0ex][0ex]{\ensuremath{\rhookswarrow}},
    postbreak=\raisebox{0ex}[0ex][0ex]{\ensuremath{\rcurvearrowse\space}},
    frame=single,
    columns=fullflexible,keepspaces,
    escapeinside=§§,
    inputencoding=utf8
}

\input{syntax_color}

\ifdefined\RUSSIAN
\renewcommand\lstlistingname{Листинг}
\renewcommand\lstlistlistingname{Листинг}
\fi

\DeclareMathSizes{12}{30}{16}{12}%

% see also:
% http://tex.stackexchange.com/questions/129225/how-can-i-get-get-makeindex-to-ignore-capital-letters
% http://tex.stackexchange.com/questions/18336/correct-sorting-of-index-entries-containing-macros
\def\myindex#1{\expandafter\index\expandafter{#1}}

\begin{document}

% fancyhdr
\pagestyle{fancy}
\setlength{\headheight}{13pt}
\fancyhead[R]{} % suppress chapter name

\VerbatimFootnotes

\frontmatter

%\RU{\include{1st_page_RU}}
%\EN{\include{1st_page_EN}}
%\DE{\include{1st_page_DE}}
%\ES{\include{1st_page_ES}}
%\CN{\include{1st_page_CN}}

\RU{\include{dedication_RU}}
\EN{\include{dedication_EN}}
\FR{\include{dedication_FR}}
\JA{\include{dedication_JA}}
\DE{\include{dedication_DE}}
\IT{\include{dedication_IT}}

\include{page_after_cover}
\include{call_for_translators}

\shorttoc{%
    \RU{Краткое оглавление}%
    \EN{Abridged contents}%
    \ES{Contenidos abreviados}%
    \PTBRph{}%
    \DE{Inhaltsverzeichnis (gekürzt)}%
    \PLph{}%
    \IT{Sommario}%
    \THAph{}\NLph{}%
    \FR{Contenus abrégés}%
    \JA{簡略版}
    \TR{İçindekiler}
}{0}

\tableofcontents
\cleardoublepage

\cleardoublepage
\include{preface}

\mainmatter

\include{parts}

\EN{\include{appendix/appendix}}
\RU{\include{appendix/appendix}}
\DE{\include{appendix/appendix}}
\FR{\include{appendix/appendix}}
\IT{\include{appendix/appendix}}
\include{acronyms}

\bookmarksetup{startatroot}

\clearpage
\phantomsection
\addcontentsline{toc}{chapter}{%
    \RU{Глоссарий}%
    \EN{Glossary}%
    \ES{Glosario}%
    \PTBRph{}%
    \DE{Glossar}%
    \PLph{}%
    \IT{Glossario}%
    \THAph{}\NLph{}%
    \FR{Glossaire}%
    \JA{用語}
    \TR{Bolum}
}
\printglossaries

\clearpage
\phantomsection
\printindex

\end{document}
}

\EN{\chapter{Finding important/interesting stuff in the code}

Minimalism it is not a prominent feature of modern software.

\myindex{\Cpp!STL}

But not because the programmers are writing a lot, but because a lot of libraries are commonly linked statically
to executable files.
If all external libraries were shifted into an external DLL files, the world would be different.
(Another reason for C++ are the \ac{STL} and other template libraries.)

\newcommand{\FOOTNOTEBOOST}{\footnote{\url{http://go.yurichev.com/17036}}}
\newcommand{\FOOTNOTELIBPNG}{\footnote{\url{http://go.yurichev.com/17037}}}

Thus, it is very important to determine the origin of a function, if it is from standard library or 
well-known library (like Boost\FOOTNOTEBOOST, libpng\FOOTNOTELIBPNG),
or if it is related to what we are trying to find in the code.

It is just absurd to rewrite all code in \CCpp to find what we're looking for.

One of the primary tasks of a reverse engineer is to find quickly the code he/she needs.

\myindex{\GrepUsage}

The \IDA disassembler allow us to search among text strings, byte sequences and constants.
It is even possible to export the code to .lst or .asm text files and then use \TT{grep}, \TT{awk}, etc.

When you try to understand what some code is doing, this easily could be some open-source library like libpng.
So when you see some constants or text strings which look familiar, it is always worth to \emph{google} them.
And if you find the opensource project where they are used, 
then it's enough just to compare the functions.
It may solve some part of the problem.

For example, if a program uses XML files, the first step may be determining which
XML library is used for processing, since the standard (or well-known) libraries are usually used
instead of self-made one.

\myindex{SAP}
\myindex{Windows!PDB}

For example, the author of these lines once tried to understand how the compression/decompression of network packets works in SAP 6.0. 
It is a huge software, but a detailed .\gls{PDB} with debugging information is present, 
and that is convenient.
He finally came to the idea that one of the functions, that was called \emph{CsDecomprLZC}, was doing the decompression of network packets.
Immediately he tried to google its name and he quickly found the function was used in MaxDB
(it is an open-source SAP project) \footnote{More about it in relevant section~(\myref{sec:SAPGUI})}.

\url{http://www.google.com/search?q=CsDecomprLZC}

Astoundingly, MaxDB and SAP 6.0 software shared likewise code for the compression/decompression of network packets.

\input{digging_into_code/identification/exec_EN}

% binary files might be also here

\mysection{Communication with outer world (function level)}
It's often advisable to track function arguments and return values in debugger or \ac{DBI}.
For example, the author once tried to understand meaning of some obscure function, which happens to be incorrectly
implemented bubble sort\footnote{\url{https://yurichev.com/blog/weird_sort/}}.
(It worked correctly, but slower.)
Meanwhile, watching inputs and outputs of this function helps instantly to understand what it does.

Often, when you see division by multiplication (\myref{sec:divisionbymult}),
but forgot all details about its mechanics, you can just observe input
and output and quickly find divisor.

% sections:
\input{digging_into_code/communication_win32_EN}
\input{digging_into_code/strings_EN}
\input{digging_into_code/assert_EN}
\input{digging_into_code/constants_EN}
\input{digging_into_code/instructions_EN}
\input{digging_into_code/suspicious_code_EN}
\input{digging_into_code/magic_numbers_tracing_EN}
\input{digging_into_code/loops_EN}
\input{digging_into_code/binary/main_EN}
\input{digging_into_code/snapshots_comparing_EN}
\input{digging_into_code/ISA_detect_EN}

\mysection{Other things}

\subsection{General idea}

A reverse engineer should try to be in programmer's shoes as often as possible. 
To take his/her viewpoint and ask himself, how would one solve some task the specific case.

\subsection{Order of functions in binary code}

All functions located in a single .c or .cpp-file are compiled into corresponding object (.o) file.
Later, linker puts all object files it needs together, not changing order or functions in them.
As a consequence, if you see two or more consecutive functions, it means, that they were placed together
in a single source code file (unless you're on border of two object files, of course.)
This means these functions have something in common, that they are from the same \ac{API} level, from same library, etc.

\subsection{Tiny functions}

Tiny functions like empty functions (\myref{empty_func})
or function which returns just ``true'' (1) or ``false'' (0) (\myref{ret_val_func}) are very common,
and almost all decent compilers tend to put only one such function into resulting executable code even if there were several
similar functions in source code.
So, whenever you see a tiny function consisting just of \TT{mov eax, 1 / ret}
which is referenced (and can be called) from many places,
which are seems unconnected to each other, this may be a result of such optimization.%

\subsection{\Cpp}

\ac{RTTI}~(\myref{RTTI})-data may be also useful for \Cpp class identification.

}
\RU{\chapter{Поиск в коде того что нужно}

Современное ПО, в общем-то, минимализмом не отличается.

\myindex{\Cpp!STL}
Но не потому, что программисты слишком много пишут, 
а потому что к исполняемым файлам обыкновенно прикомпилируют все подряд библиотеки. 
Если бы все вспомогательные библиотеки всегда выносили во внешние DLL, мир был бы иным.
(Еще одна причина для Си++ --- \ac{STL} и прочие библиотеки шаблонов.)

\newcommand{\FOOTNOTEBOOST}{\footnote{\url{http://go.yurichev.com/17036}}}
\newcommand{\FOOTNOTELIBPNG}{\footnote{\url{http://go.yurichev.com/17037}}}

Таким образом, очень полезно сразу понимать, какая функция из стандартной библиотеки или 
более-менее известной (как Boost\FOOTNOTEBOOST, libpng\FOOTNOTELIBPNG), 
а какая --- имеет отношение к тому что мы пытаемся найти в коде.

Переписывать весь код на \CCpp, чтобы разобраться в нем, безусловно, не имеет никакого смысла.

Одна из важных задач reverse engineer-а это быстрый поиск в коде того что собственно его интересует.

\myindex{\GrepUsage}
Дизассемблер \IDA позволяет делать поиск как минимум строк, последовательностей байт, констант.
Можно даже сделать экспорт кода в текстовый файл .lst или .asm и затем натравить на него \TT{grep}, \TT{awk}, итд.

Когда вы пытаетесь понять, что делает тот или иной код, это запросто может быть какая-то 
опенсорсная библиотека вроде libpng. Поэтому, когда находите константы, или текстовые строки, которые 
выглядят явно знакомыми, всегда полезно их \emph{погуглить}.
А если вы найдете искомый опенсорсный проект где это используется, 
то тогда будет достаточно будет просто сравнить вашу функцию с ней. 
Это решит часть проблем.

К примеру, если программа использует какие-то XML-файлы, первым шагом может быть
установление, какая именно XML-библиотека для этого используется, ведь часто используется какая-то
стандартная (или очень известная) вместо самодельной.

\myindex{SAP}
\myindex{Windows!PDB}
К примеру, автор этих строк однажды пытался разобраться как происходит компрессия/декомпрессия сетевых пакетов в SAP 6.0. 
Это очень большая программа, но к ней идет подробный .\gls{PDB}-файл с отладочной информацией, и это очень удобно. 
Он в конце концов пришел к тому что одна из функций декомпрессирующая пакеты называется CsDecomprLZC(). 
Не сильно раздумывая, он решил погуглить и оказалось, что функция с таким же названием имеется в MaxDB
(это опен-сорсный проект SAP) \footnote{Больше об этом в соответствующей секции~(\myref{sec:SAPGUI})}.

\url{http://www.google.com/search?q=CsDecomprLZC}

Каково же было мое удивление, когда оказалось, что в MaxDB используется точно такой же алгоритм, 
скорее всего, с таким же исходником.

\input{digging_into_code/identification/exec_RU}
% binary files might be also here

\mysection{Связь с внешним миром (на уровне функции)}
Очень желательно следить за аргументами ф-ции и возвращаемыми значениями, в отладчике или \ac{DBI}.
Например, автор этих строк однажды пытался понять значение некоторой очень запутанной ф-ции, которая, как потом оказалось,
была неверно реализованной пузырьковой сортировкой\footnote{\url{https://yurichev.com/blog/weird_sort/}}.
(Она работала правильно, но медленнее.)
В то же время, наблюдение за входами и выходами этой ф-ции помогает мгновенно понять, что она делает.

Часто, когда вы видите деление через умножение (\myref{sec:divisionbymult}),
но забыли все детали о том, как оно работает, вы можете просто наблюдать за входом и выходом, и так быстро найти делитель.

% sections:
\input{digging_into_code/communication_win32_RU}
\input{digging_into_code/strings_RU}
\input{digging_into_code/assert_RU}
\input{digging_into_code/constants_RU}
\input{digging_into_code/instructions_RU}
\input{digging_into_code/suspicious_code_RU}
\input{digging_into_code/magic_numbers_tracing_RU}
\input{digging_into_code/loops_RU}
\input{digging_into_code/binary/main_RU}
\input{digging_into_code/snapshots_comparing_RU}
\input{digging_into_code/ISA_detect_RU}

\mysection{Прочее}

\subsection{Общая идея}

Нужно стараться как можно чаще ставить себя на место программиста и задавать себе вопрос, 
как бы вы сделали ту или иную вещь в этом случае и в этой программе.

\subsection{Порядок функций в бинарном коде}

Все функции расположеные в одном .c или .cpp файле компилируются в соответствующий объектный (.o) файл.
Линкер впоследствии складывает все нужные объектные файлы вместе, не меняя порядок ф-ций в них.
Как следствие, если вы видите в коде две или более идущих подряд ф-ций, то это означает, что и в исходном коде они 
были расположены в одном и том же файле (если только вы не на границе двух объектных файлов, конечно).
Это может означать, что эти ф-ции имеют что-то общее между собой, что они из одного слоя \ac{API}, из одной библиотеки, итд.%

\subsection{Крохотные функции}

Крохотные ф-ции, такие как пустые ф-ции (\myref{empty_func})
или ф-ции возвращающие только ``true'' (1) или ``false'' (0) (\myref{ret_val_func}) очень часто встречаются,
и почти все современные компиляторы, как правило, помещают только одну такую ф-цию в исполняемый код,
даже если в исходном их было много одинаковых.
Так что если вы видите ф-цию состояющую только из \TT{mov eax, 1 / ret}, которая может вызываться из разных мест,
которые, судя по всему, друг с другом никак не связаны, это может быть результат подобной оптимизации.

\subsection{\Cpp}

\ac{RTTI}~(\myref{RTTI})-информация также может быть полезна для идентификации 
классов в \Cpp.

}
\DE{\chapter{Finden von wichtigen / interessanten Stellen im Code}

Minimalismus ist kein beliebtes Feature moderner Software.

\myindex{\Cpp!STL}

Aber nicht weil die Programmierer so viel Code schreiben, sondern weil die Libaries
allgemein statisch zu ausf\"uhrbaren Dateien gelinkt werden. Wenn alle externen
Libraries in externe DLL Dateien verschoben werden w\"urden, w\"are die Welt ein
anderer Ort. (Ein weiterer Grund f\"ur C++ sind die \ac{STL} und andere Template-Libraries.)

\newcommand{\FOOTNOTEBOOST}{\footnote{\url{http://go.yurichev.com/17036}}}
\newcommand{\FOOTNOTELIBPNG}{\footnote{\url{http://go.yurichev.com/17037}}}

Deshalb ist es sehr wichtig den Ursprung einer Funktion zu bestimmen, wenn die
Funktion aus einer Standard-Library oder aus einer sehr bekannten Library stammt
(wie z.B Boost\FOOTNOTEBOOST, libpng\FOOTNOTELIBPNG), oder ob die Funktion sich
auf das bezieht was wir im Code versuchen zu finden.

Es ist ein wenig absurd s\"amtlichen Code in \CCpp neu zu schreiben, um das zu
finden was wir suchen.

Eine der Hauptaufgaben eines Reverse Enigneers ist es schnell Code zu finden den
er/sie sucht.

\myindex{\GrepUsage}

Der \IDA-Disassembler erlaubt es durch Textstrings, Byte-Sequenzen und Konstanten
zu suchen.  Es ist sogar m\"oglich den Code in .lst oder .asm Text Dateien zu
exportieren und diese mit \TT{grep}, \TT{awk}, etc. zu untersuchen.

Wenn man versucht zu verstehen wie ein bestimmter Code funktioniert, kann auch
eine einfache Open-Source-Library wie libpng als Beispiel dienen.
Wenn man also eine Konstante oder Textstrings findet die vertraut erscheinen, ist
es immer einen Versuch wert diese zu \emph{google}n .
Und wenn man ein Opensource Projekt findet in dem diese Funktion benutzt wird, 
reicht es meist aus diese Funktionen miteinander zu vergleichen.
Es k\"onnte helfen Teile des Problems zu l\"osen.

% When you try to understand what some code is doing, this easily could be some open-source library like libpng.
% So when you see some constants or text strings which look familiar, it is always worth to \emph{google} them.
% And if you find the opensource project where they are used, 
% then it's enough just to compare the functions.
% It may solve some part of the problem.

Zum Beispiel, wenn ein Programm XML Dateien benutzt, w\"are der erste Schritt zu ermitteln welche
XML-Library benutzt wird f\"ur die Verarbeitung, da die Standard (oder am weitesten verbreitete) libraries
normal benutzt werden anstatt selbst geschriebene librarys.

\myindex{SAP}
\myindex{Windows!PDB}

Zum Beispiel, der Autor dieser Zeilen wollte verstehen wie die Kompression/Dekompression von Netzwerkpaketen in SAP 6.0 funktioniert.
SAP ist ein gewaltiges St\"uck Software, aber detaillierte -\gls{PDB} Dateien mit Debug Informationen sind vorhanden, was sehr praktisch 
ist. Der Autor hat schließlich eine Ahnung gehabt, das eine Funktion genannt \emph{CsDecomprLZC} die Dekompression der Netzwerkpakete \"ubernahm.
Er hat nach dem Namen der Funktion auf google gesucht und ist schnell zum schluss gekommen das diese Funktion in 
MaxDB benutzt wurde (Das ist ein Open-Source SAP Projekt) \footnote{Mehr dar\"uber in der relevanten Sektion~(\myref{sec:SAPGbUI})}. 

\url{http://www.google.com/search?q=CsDecomprLZC}

Erstaunlich, das MaxDB und die SAP 6.0 Software den selben Code geteilt haben f\"ur die Kompression/Dekompression der Netzwerkpakete.

\input{digging_into_code/identification/exec_DE} 

\mysection{Kommunikation mit der außen Welt (Funktion Level)} 
Oft ist es empfehlenswert die Funktionsargumente und die R\"uckgabewerte im
Debugger oder \ac{DBI} zu \"uberwachen. Zum Beispiel hat der Autor einmal
versucht die Bedeutung einer obskuren Funktion zu verstehen, die einen inkorrekten
Bubblesort-Algorithmus implementiert hatte\footnote{\url{https://yurichev.com/blog/weird_sort/}}
(Sie hat funktioniert, jedoch viel langsamer als normal). Die Eingaben und Ausgaben zur Laufzeit 
der Funktion zu \"uberwachen hilft sofort zu verstehen was die Funktion tut.

% TBT

% sections:
\input{digging_into_code/communication_win32_DE}
\input{digging_into_code/strings_DE}
\input{digging_into_code/assert_DE}
\input{digging_into_code/constants_DE}
\input{digging_into_code/instructions_DE}
\input{digging_into_code/suspicious_code_DE}
\input{digging_into_code/magic_numbers_tracing_DE}
\input{digging_into_code/loops_DE}
\input{digging_into_code/binary/main_DE}
\input{digging_into_code/snapshots_comparing_DE}
% TBT \input{digging_into_code/ISA_detect_DE}

\mysection{Andere Dinge}

\subsection{Die Idee}  

Ein Reverse Engineer sollte versuchen so oft wie m\"oglich in den Schuhen des
Programmierers zu laufen. Um ihren/seinen Standpunkt zu betrachten uns sich
selbst zu Fragen wie man einen Task in spezifischen F\"allen l\"osen w\"urde.

\subsection{Anordnung von Funktionen in Bin\"ar Code}  

S\"amtliche Funktionen die in einer einzelnen .c oder .cpp-Datei gefunden werden,
werden zu den entsprechenden Objekt Dateien (.o) kompiliert. Sp\"ater, f\"ugt
der Linker alle Objektdatein die er braucht zusammen, ohne die Reihenfolge oder
die Funktionen in Ihnen zu ver\"andern. Als eine Konsequenz, ergibt sich daraus
wenn man zwei oder mehr aufeinander folgende Funktionen sieht, bedeutet dass das
sie in der gleichen Source Code Datei platziert waren (Außer nat\"urlich man bewegt
sich an der Grenze zwischen zwei Dateien.). Das bedeutet das diese Funktionen etwas
gemeinsam haben, das sie aus dem gleichen \ac{API}-Level stammen oder aus der
gleichen Library, etc.

\subsection{kleine Funktionen} 

Sehr kleine oder leere Funktionen  (\myref{empty_func})
oder Funktionen die nur ``true'' (1) oder ``false'' (0) (\myref{ret_val_func}) sind weit verbreitet,
und fast jeder ordentlicher Compiler tendiert dazu nur solche Funktionen in den resultierenden ausf\"uhrbaren Code zu stecken,
sogar wenn es mehrere gleiche Funktionen im Source Code bereits gibt. 
Also, wann immer man solche kleinen Funktionen sieht die z.B nur aus \TT{mov eax, 1 / ret} bestehen und von mehreren 
Orten aus referenziert werden (und aufgerufen werden k\"onnen), und scheinbar keine Verbindung zu einander haben, dann 
ist das wahrscheinlich das Ergebnis einer Optimierung. 

\subsection{\Cpp}

\ac{RTTI}~(\myref{RTTI})-data ist vielleicht auch n\"utzlich f\"ur die \Cpp Klassen Identifikation.
}
\FR{\chapter{Trouver des choses importantes/intéressantes dans le code}

Le minimalisme n'est pas une caractéristique prépondérante des logiciels modernes.

\myindex{\Cpp!STL}

Pas parce que les programmeurs écrivent beaucoup, mais parce que de nombreuses bibliothèques
sont couramment liées statiquement aux fichiers exécutable.
Si toutes les bibliothèques externes étaient déplacées dans des fichiers DLL externes,
le monde serait différent. (Une autre raison pour C++ sont la \ac{STL} et autres
bibliothèques templates.)

\newcommand{\FOOTNOTEBOOST}{\footnote{\url{http://go.yurichev.com/17036}}}
\newcommand{\FOOTNOTELIBPNG}{\footnote{\url{http://go.yurichev.com/17037}}}

Ainsi, il est très important de déterminer l'origine de la fonction, si elle provient
d'une bibliothèque standard ou d'une bibliothèque bien connue (comme Boost\FOOTNOTEBOOST,
libpng\FOOTNOTELIBPNG), ou si elle est liée à ce que l'on essaye de trouver dans
le code.

Il est simplement absurde de tout récrire le code en \CCpp pour trouver ce que l'on
cherche.

Une des premières tâches d'un rétro-ingénieur est de trouver rapidement le code dont
il a besoin.

\myindex{\GrepUsage}

Le dés-assembleur \IDA nous permet de chercher parmi les chaînes de texte, les séquences
d'octets et les constantes.
Il est même possible d'exporter le code dans un fichier texte .lst ou .asm et d'utiliser
\TT{grep}, \TT{awk}, etc.

Lorsque vous essayez de comprendre ce que fait un certain code, ceci peut être facile
avec une bibliothèque open-source comme libpng.
Donc, lorsque vous voyez certaines constantes ou chaînes de texte qui vous semblent
familières, il vaut toujours la peine de les \emph{googler}.
Et si vous trouvez le projet open-source où elles sont utilisées, alors il suffit
de comparer les fonctions.
Ceci peut permettre de résoudre certaines parties du problème.

Par exemple, si un programme utilise des fichiers XML, la premières étape peut-être
de déterminer quelle bibliothèque XML est utilisée pour le traitement, puisque les
bibliothèques standards (ou bien connues) sont en général utilisées au lieu de code
fait maison.

\myindex{SAP}
\myindex{Windows!PDB}

Par exemple, l'auteur de ces lignes a essayé une fois de comprendre comment la compression/décompression
des paquets réseau fonctionne dans SAP 6.0.
C'est un logiciel gigantesque, mais un .\gls{PDB} détaillé avec des informations
de débogage est présent, et c'est pratique.
Il en est finalement arrivé à l'idée que l'une des fonctions, qui était appelée par
\emph{CsDecomprLZC}, effectuait la décompression des paquets réseau.
Immédiatement, il a essayé de googler le nom et a rapidement trouvé que la fonction
était utilisée dans MaxDB (c'est un projet open-source de SAP)
\footnote{Plus sur ce sujet dans la section concernée~(\myref{sec:SAPGUI})}.

\url{http://www.google.com/search?q=CsDecomprLZC}

Étonnament, les logiciels MaxDB et SAP 6.0 partagent du code comme ceci pour la compression/
décompression des paquets réseau.

% binary files might be also here

\mysection{Communication avec le mondes extérieur (niveau fonction)}
Il est souvent recommandé de suivre les arguments de la fonction et sa valeur de
retour dans un débogueur ou \ac{DBI}.
Par exemple, l'auteur a essayé une fois de comprendre la signification d'une fonction
obscure, qui s'est avérée être un tri à bulles mal implémenté\footnote{\url{https://yurichev.com/blog/weird_sort/}}.
(Il fonctionnait correctement, mais plus lentement.)
En même temps, regarder les entrées et sorties de cette fonction aide instantanément
à comprendre ce quelle fait.

Souvent, lorsque vous voyez une division par la multiplication (\myref{sec:divisionbymult}),
mais avez oublié tous les détails du mécanisme, vous pouvez seulement observer l'entrée
et la sortie, et trouver le diviseur rapidement.

% sections:
\input{digging_into_code/binary/main_FR}

\mysection{Autres choses}

\subsection{Idée générale}

Un rétro-ingénieur doit essayer se se mettre dans la peau d'un programmeur aussi
souvent que possible.
Pour prendre son point de vue et se demander comment il aurait résolu des taches
d'un cas spécifique.

\subsection{Ordre des fonctions dans le code binaire}

Toutes les fonctions situées dans un unique fichier .c ou .cpp sont compilées dans
le fichier objet (.o) correspondant.
Plus tard, l'éditeur de liens mets tous les fichiers dont il a besoin ensemble, sans
changer l'ordre ni les fonctions.
Par conséquent, si vous voyez deux ou plus fonctions consécutives, cela signifie
qu'elles étaient situées dans le même fichier source (à moins que vous ne soyez en
limite de deux fichiers objet, bien sûr).
Ceci signifie que ces fonctions ont quelque chose en commun, qu'elles sont des fonctions
du même niveau d'\ac{API}, de la même bibliothèque, etc.

\subsection{Fonctions minuscules}

Les fonctions minuscules comme les fonctions vides (\myref{empty_func})
ou les fonctions qui renvoient juste ``true'' (1) ou ``false'' (0) (\myref{ret_val_func})
sont très communes, et presque tous les compilateurs corrects tendent à ne mettre
qu'une seule fonction de ce genre dans le code de l'exécutable résultant, même si
il y avait plusieurs fonctions similaires dans le code source.
Donc, à chaque fois que vous voyez une fonction minuscule consistant seulement en
\TT{mov eax, 1 / ret} qui est référencée (et peut être appelée) dans plusieurs endroits
qui ne semblent pas reliés les uns au autres, ceci peut résulter d'une telle optimisation.%

\subsection{\Cpp}

Les données \ac{RTTI}~(\myref{RTTI})- peuvent être utiles pour l'identification des
classes \Cpp.

}

\chapter{\RU{Специфичное для ОС}\EN{OS-specific}%
\DE{Betriebssystem-spezifische Themen}%
\FR{Spécifique aux OS}}
\EN{\input{OS/calling_conventions/main_EN}}
\RU{\input{OS/calling_conventions/main_RU}}
\DE{\input{OS/calling_conventions/main_DE}}
\FR{\input{OS/calling_conventions/main_FR}}

\input{OS/TLS/main}

\EN{\input{OS/syscalls_EN}}
\RU{\input{OS/syscalls_RU}}
\DE{\input{OS/syscalls_DE}}
\FR{\input{OS/syscalls_FR}}

\mysection{Linux}
\EN{\input{OS/PIC_EN}}
\RU{\input{OS/PIC_RU}}
\DE{\input{OS/PIC_DE}}
\FR{\input{OS/PIC_FR}}

\EN{\input{OS/LD_PRELOAD/main_EN}}
\RU{\input{OS/LD_PRELOAD/main_RU}}
\DE{\input{OS/LD_PRELOAD/main_DE}}
\FR{\input{OS/LD_PRELOAD/main_FR}}
\IT{\input{OS/LD_PRELOAD/main_IT}}

\mysection{Windows NT}
\EN{\input{OS/win32_CRT/main_EN}}
\RU{\input{OS/win32_CRT/main_RU}}
\DE{\input{OS/win32_CRT/main_DE}}
\FR{\input{OS/win32_CRT/main_FR}}

\input{OS/PE/main}
\input{OS/SEH/main}

\EN{\input{OS/NT_critical_sections_EN}}
\RU{\input{OS/NT_critical_sections_RU}}
\DE{\input{OS/NT_critical_sections_DE}}
\FR{\input{OS/NT_critical_sections_FR}}

\EN{\input{tools_EN}}
\ES{\input{tools_ES}}
\RU{\input{tools_RU}}
\IT{\input{tools_IT}}
\FR{\input{tools_FR}}
\DE{\input{tools_DE}}
%\CN{\input{tools_CN}}
\JPN{\input{tools_JPN}}

\EN{\chapter{\RU{Примеры из практики}\EN{Case studies}\DE{Beispiele aus der Praxis}}

%\clearpage
\begin{center}
\vspace*{\fill}

\begin{figure}[H]
\centering
\myincludegraphics{cover4.jpg}
\end{figure}

\vspace*{\fill}
\end{center}

\clearpage

\EN{\input{examples/Knuth_EN}}\RU{\input{examples/Knuth_RU}}

% sections here
\EN{\input{examples/taskmgr/main_EN}}\RU{\input{examples/taskmgr/main_RU}}
\EN{\input{examples/lines/main_EN}}\RU{\input{examples/lines/main_RU}}
\EN{\input{examples/minesweeper/main_EN}}\RU{\input{examples/minesweeper/main_RU}}
\EN{\input{examples/timedate/main_EN}}\RU{\input{examples/timedate/main_RU}}
\input{examples/dongles/main}

% I never liked this part:
% \EN{\input{examples/qr9/qr9_EN}}\RU{\input{examples/qr9/qr9_RU}}

\EN{\input{examples/encrypted_DB1/main_EN}}\RU{\input{examples/encrypted_DB1/main_RU}}
\EN{\input{examples/bitcoin_miner/main_EN}}\RU{\input{examples/bitcoin_miner/main_RU}}
\EN{\input{examples/simple_exec_crypto/main_EN}}
\input{examples/SAP/main}
\input{examples/oracle/main}
\EN{\input{examples/handcoding/main_EN}}\RU{\input{examples/handcoding/main_RU}}
\input{examples/demos/main}
\RU{\input{examples/marriage/main_RU}}

\mysection{\EN{Other examples}\RU{Другие примеры}}

\RU{Здесь также был пример с Z3 и ручной декомпиляцией.}
\EN{An example about Z3 and manual decompilation was here.}
\RU{Он (временно) перемещен сюда:}
\EN{It is (temporarily) moved there:}
\url{http://yurichev.com/tmp/SAT_SMT_DRAFT.pdf}.


}\RU{\chapter{\RU{Примеры из практики}\EN{Case studies}\DE{Beispiele aus der Praxis}}

%\clearpage
\begin{center}
\vspace*{\fill}

\begin{figure}[H]
\centering
\myincludegraphics{cover4.jpg}
\end{figure}

\vspace*{\fill}
\end{center}

\clearpage

\EN{\input{examples/Knuth_EN}}\RU{\input{examples/Knuth_RU}}

% sections here
\EN{\input{examples/taskmgr/main_EN}}\RU{\input{examples/taskmgr/main_RU}}
\EN{\input{examples/lines/main_EN}}\RU{\input{examples/lines/main_RU}}
\EN{\input{examples/minesweeper/main_EN}}\RU{\input{examples/minesweeper/main_RU}}
\EN{\input{examples/timedate/main_EN}}\RU{\input{examples/timedate/main_RU}}
\input{examples/dongles/main}

% I never liked this part:
% \EN{\input{examples/qr9/qr9_EN}}\RU{\input{examples/qr9/qr9_RU}}

\EN{\input{examples/encrypted_DB1/main_EN}}\RU{\input{examples/encrypted_DB1/main_RU}}
\EN{\input{examples/bitcoin_miner/main_EN}}\RU{\input{examples/bitcoin_miner/main_RU}}
\EN{\input{examples/simple_exec_crypto/main_EN}}
\input{examples/SAP/main}
\input{examples/oracle/main}
\EN{\input{examples/handcoding/main_EN}}\RU{\input{examples/handcoding/main_RU}}
\input{examples/demos/main}
\RU{\input{examples/marriage/main_RU}}

\mysection{\EN{Other examples}\RU{Другие примеры}}

\RU{Здесь также был пример с Z3 и ручной декомпиляцией.}
\EN{An example about Z3 and manual decompilation was here.}
\RU{Он (временно) перемещен сюда:}
\EN{It is (temporarily) moved there:}
\url{http://yurichev.com/tmp/SAT_SMT_DRAFT.pdf}.


}\DE{\chapter{\RU{Примеры из практики}\EN{Case studies}\DE{Beispiele aus der Praxis}}

%\clearpage
\begin{center}
\vspace*{\fill}

\begin{figure}[H]
\centering
\myincludegraphics{cover4.jpg}
\end{figure}

\vspace*{\fill}
\end{center}

\clearpage

\EN{\input{examples/Knuth_EN}}\RU{\input{examples/Knuth_RU}}

% sections here
\EN{\input{examples/taskmgr/main_EN}}\RU{\input{examples/taskmgr/main_RU}}
\EN{\input{examples/lines/main_EN}}\RU{\input{examples/lines/main_RU}}
\EN{\input{examples/minesweeper/main_EN}}\RU{\input{examples/minesweeper/main_RU}}
\EN{\input{examples/timedate/main_EN}}\RU{\input{examples/timedate/main_RU}}
\input{examples/dongles/main}

% I never liked this part:
% \EN{\input{examples/qr9/qr9_EN}}\RU{\input{examples/qr9/qr9_RU}}

\EN{\input{examples/encrypted_DB1/main_EN}}\RU{\input{examples/encrypted_DB1/main_RU}}
\EN{\input{examples/bitcoin_miner/main_EN}}\RU{\input{examples/bitcoin_miner/main_RU}}
\EN{\input{examples/simple_exec_crypto/main_EN}}
\input{examples/SAP/main}
\input{examples/oracle/main}
\EN{\input{examples/handcoding/main_EN}}\RU{\input{examples/handcoding/main_RU}}
\input{examples/demos/main}
\RU{\input{examples/marriage/main_RU}}

\mysection{\EN{Other examples}\RU{Другие примеры}}

\RU{Здесь также был пример с Z3 и ручной декомпиляцией.}
\EN{An example about Z3 and manual decompilation was here.}
\RU{Он (временно) перемещен сюда:}
\EN{It is (temporarily) moved there:}
\url{http://yurichev.com/tmp/SAT_SMT_DRAFT.pdf}.


}
\RU{\chapter{Примеры разбора закрытых (proprietary) форматов файлов}}
\EN{\chapter{Examples of reversing proprietary file formats}}
\DE{\chapter{Beispiele für das Reverse Engineering proprietärer Dateiformate}}
\FR{\chapter{Exemples de Reverse Engineering de format de fichier propriétaire}}

% sections
\input{ff/XOR/main}

\EN{\input{ff/entropy/main_EN}}
\RU{\input{ff/entropy/main_RU}}

\EN{\input{ff/millenium/main_EN}}
\RU{\input{ff/millenium/main_RU}}
\FR{\input{ff/millenium/main_FR}}

\EN{\input{ff/fortune/main_EN}}
\RU{\input{ff/fortune/main_RU}}
\FR{\input{ff/fortune/main_FR}}

\EN{\input{ff/Oracle_SYM/main_EN}}
\RU{\input{ff/Oracle_SYM/main_RU}}

\EN{\input{ff/Oracle_MSB/main_EN}}
\RU{\input{ff/Oracle_MSB/main_RU}}

\EN{\input{ff/exercise_EN}}
\RU{\input{ff/exercise_RU}}
\FR{\input{ff/exercise_FR}}

\mysection{\EN{Further reading}\RU{Дальнейшее чтение}\FR{Pour aller plus loin}\DEph{}}

\href{https://yurichev.com/mirrors/SSTIC2016-Article-cryptanalyse_en_boite_noire_de_chiffrement_proprietaire-capillon.pdf}{Pierre Capillon -- Black-box cryptanalysis of home-made encryption algorithms: a practical case study}.

\href{https://alexhude.github.io/2019/01/24/hacking-leica-m240.html}{How to Hack an Expensive Camera and Not Get Killed by Your Wife}.


\EN{\input{DBI/main_EN}}
\DE{\input{DBI/main_DE}}
\FR{\input{DBI/main_FR}}
% TBT: RU, etc...


% another place for `unsorted' sections

\chapter{\RU{Прочее}\EN{Other things}\DE{Weitere Themen}\FR{Autres sujets}\JA{その他}}

% sections:
\EN{\input{other/patching_EN}}
\RU{\input{other/patching_RU}}
\DE{\input{other/patching_DE}}
\FR{\input{other/patching_FR}}

\EN{\input{other/args_stat_EN}}
\RU{\input{other/args_stat_RU}}
\DE{\input{other/args_stat_DE}}
\FR{\input{other/args_stat_FR}}

\EN{\input{other/compiler_intrinsic_EN}}
\RU{\input{other/compiler_intrinsic_RU}}
\DE{\input{other/compiler_intrinsic_DE}}
\FR{\input{other/compiler_intrinsic_FR}}

\EN{\input{other/compiler_anomalies_EN}}
\RU{\input{other/compiler_anomalies_RU}}
\DE{\input{other/compiler_anomalies_DE}}
\FR{\input{other/compiler_anomalies_FR}}

\EN{\input{other/itanium/main_EN}}
\RU{\input{other/itanium/main_RU}}
\DE{\input{other/itanium/main_DE}}
\FR{\input{other/itanium/main_FR}}

\EN{\input{other/8086mm_EN}}
\RU{\input{other/8086mm_RU}}
\DE{\input{other/8086mm_DE}}
\FR{\input{other/8086mm_FR}}

\EN{\input{other/bb_reordering_EN}}
\RU{\input{other/bb_reordering_RU}}
\DE{\input{other/bb_reordering_DE}}
\FR{\input{other/bb_reordering_FR}}

\EN{\input{other/hexrays_EN}}
\RU{\input{other/hexrays_RU}}



\EN{\subsubsection{Reading outside array bounds}

So, array indexing is just \emph{array\lbrack{}index\rbrack}.
If you study the generated code closely, you'll probably note the missing index bounds checking,
which could check \emph{if it is less than 20}.
What if the index is 20 or greater?
That's the one \CCpp feature it is often blamed for.

Here is a code that successfully compiles and works:

\lstinputlisting[style=customc]{patterns/13_arrays/2_BO/r.c}

Compilation results (MSVC 2008):

\lstinputlisting[caption=\NonOptimizing MSVC 2008,style=customasmx86]{patterns/13_arrays/2_BO/r_msvc.asm}

The code produced this result:

\lstinputlisting[caption=\olly: console output]{patterns/13_arrays/2_BO/console.txt}

It is just \emph{something} that has been lying in the stack near to the array, 80 bytes away from its first element.

\clearpage
\myindex{\olly}
Let's try to find out where did this value come from, using \olly.

Let's load and find the value located right after the last array element:

\begin{figure}[H]
\centering
\myincludegraphics{patterns/13_arrays/2_BO/olly_r1.png}
\caption{\olly: reading of the 20th element and execution of \printf}
\label{fig:array_BO_olly_r1}
\end{figure}

What is this? 
Judging by the stack layout,
this is the saved value of the EBP register.
\clearpage
Let's trace further and see how it gets restored:

\begin{figure}[H]
\centering
\myincludegraphics{patterns/13_arrays/2_BO/olly_r2.png}
\caption{\olly: restoring value of EBP}
\label{fig:array_BO_olly_r2}
\end{figure}

Indeed, how it could be different?
The compiler may generate some additional code to check the index value to be always
in the array's bounds (like in higher-level programming languages\footnote{Java, Python, etc.})
but this makes the code slower.

}
\ES{\input{reading_ES}}
\RU{% TODO sync with English version
\chapter{Что стоит почитать}

\mysection{Книги и прочие материалы}

\subsection{Reverse Engineering}

\input{RE_books}

Дмитрий Скляров --- ``Искусство защиты и взлома информации''.

Также, книги Криса Касперски.

\subsection{Windows}

\input{Win_reading}

\subsection{\CCpp}

\input{CCppBooks}

\subsection{x86 / x86-64}

\label{x86_manuals}
\begin{itemize}
\item Документация от Intel\footnote{\AlsoAvailableAs \url{http://www.intel.com/content/www/us/en/processors/architectures-software-developer-manuals.html}}

\item Документация от AMD\footnote{\AlsoAvailableAs \url{http://developer.amd.com/resources/developer-guides-manuals/}}

\item \AgnerFog{}\footnote{\AlsoAvailableAs \url{http://agner.org/optimize/microarchitecture.pdf}}

\item \AgnerFogCC{}\footnote{\AlsoAvailableAs \url{http://www.agner.org/optimize/calling_conventions.pdf}}

\item \IntelOptimization

\item \AMDOptimization
\end{itemize}

Немного устарело, но всё равно интересно почитать:

\MAbrash\footnote{\AlsoAvailableAs \url{https://github.com/jagregory/abrash-black-book}}
(он известен своей работой над низкоуровневой оптимизацией в таких проектах как Windows NT 3.1 и id Quake).

\subsection{ARM}

\begin{itemize}
\item Документация от ARM\footnote{\AlsoAvailableAs \url{http://infocenter.arm.com/help/index.jsp?topic=/com.arm.doc.subset.architecture.reference/index.html}}

\item \ARMSevenRef

\item \ARMSixFourRefURL

\item \ARMCookBook\footnote{\AlsoAvailableAs \url{http://go.yurichev.com/17273}}
\end{itemize}

\subsection{Язык ассемблера}

Richard Blum --- Professional Assembly Language.

\subsection{Java}

\JavaBook.

\subsection{UNIX}

\TAOUP

\subsection{Программирование}

\begin{itemize}

\item \RobPikePractice

\item Александр Шень\footnote{\url{http://imperium.lenin.ru/~verbit/Shen.dir/shen-progra.html}}

\item \HenryWarren.
Некоторые люди говорят, что трюки и хаки из этой книги уже не нужны, потому что годились только для \ac{RISC}-процессоров,
где инструкции перехода слишком дорогие.
Тем не менее, всё это здорово помогает лучше понять булеву алгебру и всю математику рядом.

\item (Для хард-корных гиков от информатики и математики) \TAOCP.
Некоторые люди спорят, стоит ли среднему программисту читать этот довольно тяжелый фундаментальный труд.
Но я бы сказал, что стоит просто пролистать, чтобы понимать, из чего вообще состоит информатика.

\end{itemize}

% subsection:
\input{crypto_reading}

\iffalse
\subsection{Посвящение}

Как написано на первой странице этой книги, ``Эта книга посвящается Роберту Джордейну, Джону Соухэ, Ральфу Брауну и Питеру Абелю''.
Это авторы хорошо известных книг и справочников по языку ассемблеру из 1980-х и 1990-х:

\begin{itemize}
\item Роберт Джордейн -- Справочник программиста персональных компьютеров типа IBM PC, XT и AT (1986, в бывшем СССР издался в 1992)

\item Питер Нортон и Джон Соухэ -- Персональный компьютер фирмы IBM и операционная система MS-DOS (в бывшем ССР издалась в 1991), Язык ассемблера для IBM PC (там же издалась в 1992)
На самом деле, Джон Соухэ настоящий автор этих книг, можно сказать, он был литературным негром.
Также он настоящий автор Norton Commander-а.

\item Ральф Браун был известен справочником по прерываниям MS-DOS и BIOS -- ``Ralf Brown's Interrupt List''\footnote{\url{http://www.ctyme.com/rbrown.htm}}.

\item Питер Абель -- Ассемблер и программирование для IBM PC (в бывшем СССР издалась в 1992)
\end{itemize}

Это всё устаревшие книги, конечно.
Но может быть кто-то вспомнит ``те времена''.
\fi

}
\FR{\chapter{Livres/blogs qui valent le détour}

\mysection{Livres et autres matériels}

\subsection{Rétro-ingénierie}

\input{RE_books}

Également, les livres de Kris Kaspersky.

\subsection{Windows}

\input{Win_reading}

\subsection{\CCpp}

\input{CCppBooks}

\subsection{Architecture x86 / x86-64}

\label{x86_manuals}
\begin{itemize}
\item Manuels Intel\footnote{\AlsoAvailableAs \url{http://www.intel.com/content/www/us/en/processors/architectures-software-developer-manuals.html}}

\item Manuels AMD\footnote{\AlsoAvailableAs \url{http://developer.amd.com/resources/developer-guides-manuals/}}

\item \AgnerFog{}\footnote{\AlsoAvailableAs \url{http://agner.org/optimize/microarchitecture.pdf}}

\item \AgnerFogCC{}\footnote{\AlsoAvailableAs \url{http://www.agner.org/optimize/calling_conventions.pdf}}

\item \IntelOptimization

\item \AMDOptimization
\end{itemize}

Quelque peu vieux, mais toujours intéressant à lire :

\MAbrash\footnote{\AlsoAvailableAs \url{https://github.com/jagregory/abrash-black-book}}
(il est connu pour son travail sur l'optimisation bas niveau pour des projets tels que Windows NT 3.1 et id Quake).

\subsection{ARM}

\begin{itemize}
\item Manuels ARM\footnote{\AlsoAvailableAs \url{http://infocenter.arm.com/help/index.jsp?topic=/com.arm.doc.subset.architecture.reference/index.html}}

\item \ARMSevenRef

\item \ARMSixFourRefURL

\item \ARMCookBook\footnote{\AlsoAvailableAs \url{http://go.yurichev.com/17273}}
\end{itemize}

\subsection{Langage d'assemblage}

Richard Blum --- Professional Assembly Language.

\subsection{Java}

\JavaBook.

\subsection{UNIX}

\TAOUP

\subsection{Programmation en général}

\begin{itemize}

\item \RobPikePractice

\item \HenryWarren
Certaines personnes disent ques les trucs et astuces de ce livre ne sont plus pertinents
aujourd'hui, car ils n'étaient valables que pour les \ac{CPU}s \ac{RISC}, où les instructions
de branchement sont coûteuses.
Néanmoins, ils peuvent énormément aider à comprendre l'algèbre booléenne et toutes les
mathématiques associées.

\item{(Pour les passionnés avec des connaissances en informatique et mathématiques) Donald E. Knuth, \emph{The Art of Computer Programming}}.

\end{itemize}

%subsection:
\input{crypto_reading}


}
\DE{\subsubsection{Lesezugriff außerhalb von Arraygrenzen}
Der indizierte Zugriff auf ein Array wird durch \emph{array\lbrack{}index\rbrack} realisiert.
Wenn man sich den erzeugten Code genau ansieht, bemerkt man, dass eine Prüfung der Indexgrenzen fehlt, welche die
Bedingung \emph{kleiner als 20} validiert.
Was also passiert, wenn der Index 20 oder größer ist? 
Hier haben wir es mit einem unschönen Feature von \CCpp zu tun

Hier ein Beipsielcode der erfolgreich kompiliert wurde und funktioniert:

\lstinputlisting[style=customc]{patterns/13_arrays/2_BO/r.c}

Ergebnis des Kompiliervorgangs (MSVC 2008):

\lstinputlisting[caption=\NonOptimizing MSVC 2008,style=customasmx86]{patterns/13_arrays/2_BO/r_msvc.asm}

Der Code produziert dieses Ergebnis:

\lstinputlisting[caption=\olly: console output]{patterns/13_arrays/2_BO/console.txt}
Es handelt sich um \emph{irgendetwas}, das auf dem Stack in der Nähe des Arrays gelegen hat, 80 Byte von dessen erstem
Element entfernt.

\clearpage
\myindex{\olly}
Versuchen wir mit \olly herauszufinden, woher dieser Wert kommt.

Laden und finden wir also den Wert, der sich direkt hinter dem letzten Arrayelement befindet:

\begin{figure}[H]
\centering
\myincludegraphics{patterns/13_arrays/2_BO/olly_r1.png}
\caption{\olly: das 20. Element lesen und \printf ausführen}
\label{fig:array_BO_olly_r1}
\end{figure}

Worum handelt es sich? 
Dem Stacklayout nach zu urteilen ist dies der gespeicherte Wert des EBP Registers.
\clearpage
Verfolgen wir das ganze weiter und schauen uns an, wie dieser wiederhergestellt wird:

\begin{figure}[H]
\centering
\myincludegraphics{patterns/13_arrays/2_BO/olly_r2.png}
\caption{\olly: Wert von EBP wiederherstellen}
\label{fig:array_BO_olly_r2}
\end{figure}
Wie könnte es anders gelöst werden?
Der Compiler könnte zusätzlichen Code erzeugen, der sicherstellt, dass der Index sich stets innerhalb der Arraygrenzen
befindet (wie in höheren Programmiersprachen\footnote{Java, Python, etc.}), aber das würde den Code langsamer machen.
}
\ITA{\chapter{Libri/blog da leggere}

\mysection{Libri ed altro materiale}

\subsection{Reverse Engineering}

\input{RE_books}

Inoltre, i libri di Kris Kaspersky.

\subsection{Windows}

\input{Win_reading}

\subsection{\CCpp}

\input{CCppBooks}

\subsection{x86 / x86-64}

\label{x86_manuals}
\begin{itemize}
\item Manuali Intel\footnote{\AlsoAvailableAs \url{http://www.intel.com/content/www/us/en/processors/architectures-software-developer-manuals.html}}

\item Manuali AMD\footnote{\AlsoAvailableAs \url{http://developer.amd.com/resources/developer-guides-manuals/}}

\item \AgnerFog{}\footnote{\AlsoAvailableAs \url{http://agner.org/optimize/microarchitecture.pdf}}

\item \AgnerFogCC{}\footnote{\AlsoAvailableAs \url{http://www.agner.org/optimize/calling_conventions.pdf}}

\item \IntelOptimization

\item \AMDOptimization
\end{itemize}

Un po' datati ma sempre interessanti:

\MAbrash\footnote{\AlsoAvailableAs \url{https://github.com/jagregory/abrash-black-book}}
(è conosciuto per i suoi lavori di ottimizzazione a basso livello su progetti come Windows NT 3.1 e id Quake).

\subsection{ARM}

\begin{itemize}
\item Manuali ARM\footnote{\AlsoAvailableAs \url{http://infocenter.arm.com/help/index.jsp?topic=/com.arm.doc.subset.architecture.reference/index.html}}

\item \ARMSevenRef

\item \ARMSixFourRefURL

\item \ARMCookBook\footnote{\AlsoAvailableAs \url{http://go.yurichev.com/17273}}
\end{itemize}

\subsection{Assembly}

Richard Blum --- Professional Assembly Language.

\subsection{Java}

\JavaBook.

\subsection{UNIX}

\TAOUP

\subsection{Programmazione in generale}

\begin{itemize}

\item \RobPikePractice

\item \HenryWarren.
Alcune persone sostengono che i trucchi e gli hack di questo libro non siano più attuali adesso perchè erano validi solo per le \ac{CPU} \ac{RISC},
dove le istruzioni di branching sono costose.
Ad ogni modo, possono aiutare enormemente a comprendere l'algebra booleana e tutta la matematica coinvolta.

\item (Per hard-core geek con background in computer science e matematica) \TAOCP.
Alcune persone si domandano se sia utile, per programmatori mediocri, il tentare di leggere questi testi abbastanza complessi sui fondamenti della programmazione.
Io direi che è utile sfogliarli, per comprendere di cosa tratta la \ac{CS}.

\end{itemize}

% subsection:
\input{crypto_reading}

\iffalse
\subsection{Dedica}

Come scritto nella prima pagina di questo libro, ``Questo libro è dedicato a Robert Jourdain, John Socha, Ralf Brown e Peter Abel''.
Si tratta di autori famosi di libri sul linguaggio assembly e di riferimento negli anni 1980 e 1990:

\begin{itemize}
\item Robert Jourdain -- Programmer's problem solver for the IBM PC, XT, \& AT (1986)

\item Peter Norton e John Socha -- The Peter Norton Programmer's Guide to the IBM PC (1985), Peter Norton's Assembly Language Book for the IBM PC (1989).
Di fatto, John Socha è un vero autore di questi libri, si può dire, era un ghostwriter.
Inoltre è autore del Norton Commander.

\item Ralph Brown era conosciuto per la ``Ralf Brown's Interrupt List''\footnote{\url{http://www.ctyme.com/rbrown.htm}}.

\item Peter Abel -- IBM PC assembly language and programming (1991)
\end{itemize}

Chiaramente si tratta di libri antiquati.
Ma magari qualcuno si ricorderà di ``quei tempi''.
\fi
}
\JPN{\subsubsection{配列の範囲外の読み込み}

配列のインデックス化は単に\emph{array\lbrack{}index\rbrack}です。
生成されたコードを詳しく研究したなら、\emph{20未満であるか}チェックするような
インデックスの境界チェックがないことに気づくでしょう。
もしインデックスが20以上だったらどうでしょうか。
これは \CCpp が批判される1つの特徴です。

コンパイルされて動作するコードがあります。

\lstinputlisting[style=customc]{patterns/13_arrays/2_BO/r.c}

コンパイル結果（MSVC 2008）

\lstinputlisting[caption=\NonOptimizing MSVC 2008,style=customasmx86]{patterns/13_arrays/2_BO/r_msvc.asm}

コードは次の結果を生成します。

\lstinputlisting[caption=\olly: console output]{patterns/13_arrays/2_BO/console.txt}

これは単に配列のそばのスタックにある \emph{何か} です。配列の最初の要素から80バイト離れています。

\clearpage
\myindex{\olly}
この値がどこから来るのか \olly を使って見つけてみましょう。

最後の配列の要素のすぐあとに配置された値をロードして見つけましょう。

\begin{figure}[H]
\centering
\myincludegraphics{patterns/13_arrays/2_BO/olly_r1.png}
\caption{\olly: 20番目の要素を読み込み、 \printf を実行する}
\label{fig:array_BO_olly_r1}
\end{figure}

これは何でしょうか？
スタックレイアウトで判断すると、
これは保存されたEBPレジスタの値です。
\clearpage
もっとトレースしてどのようにリストアされるか見てみましょう。

\begin{figure}[H]
\centering
\myincludegraphics{patterns/13_arrays/2_BO/olly_r2.png}
\caption{\olly: EBPの値をリストア}
\label{fig:array_BO_olly_r2}
\end{figure}

本当に、異なっていますか？
コンパイラはインデックス値が配列の境界内かを常にチェックする追加のコードを生成するかもしれません。
（高水準プログラミング言語\footnote{Java, Pythonなど}のように）
しかし、これはコードを遅くします。
}

\EN{\input{community_EN}}
\RU{\input{community_RU}}
\DE{\input{community_DE}}
\FR{\input{community_FR}}
\ITA{\chapter{Community}

Esistono due eccellenti subreddit riguardo il \ac{RE} su reddit.com:
\href{http://go.yurichev.com/17027}{reddit.com/r/ReverseEngineering/} e
\href{http://go.yurichev.com/17028}{reddit.com/r/remath}
(Sugli argomenti di intersezione fra \ac{RE} e matematica).

Esiste anche una sezione relativa a \ac{RE} sul sito Stack Exchange:
\href{http://go.yurichev.com/17029}{reverseengineering.stackexchange.com}.

In IRC c'è un canale \#\#re su
FreeNode\footnote{\href{http://go.yurichev.com/17030}{freenode.net}}.
}

\EN{\input{afterword_EN}}
\ES{\input{afterword_ES}}
\RU{\input{afterword_RU}}
\NL{\input{afterword_NL}}
\IT{\input{afterword_IT}}
\FR{\input{afterword_FR}}
\DE{\input{afterword_DE}}
%\CN{\input{afterword_CN}}
\JPN{\input{afterword_JPN}}
\PTBR{\input{afterword_PTBR}}




\EN{\part*{\RU{Приложение}\EN{Appendix}\DE{Anhang}\FR{Appendice}\IT{Appendice}}
\appendix
\addcontentsline{toc}{part}{\RU{Приложение}\EN{Appendix}\DE{Anhang}\FR{Appendice}\IT{Appendice}}

% chapters
\EN{\mysection{\MinesweeperWinXPExampleChapterName}
\label{minesweeper_winxp}
\myindex{Windows!Windows XP}

For those who are not very good at playing Minesweeper, we could try to reveal the hidden mines in the debugger.

\myindex{\CStandardLibrary!rand()}
\myindex{Windows!PDB}

As we know, Minesweeper places mines randomly, so there has to be some kind of random number generator or
a call to the standard \TT{rand()} C-function.

What is really cool about reversing Microsoft products is that there are \gls{PDB} 
file with symbols (function names, \etc{}).
When we load \TT{winmine.exe} into \IDA, it downloads the 
\gls{PDB} file exactly for this 
executable and shows all names.

So here it is, the only call to \TT{rand()} is this function:

\lstinputlisting[style=customasmx86]{examples/minesweeper/tmp1.lst}

\IDA named it so, and it was the name given to it by Minesweeper's developers.

The function is very simple:

\begin{lstlisting}[style=customc]
int Rnd(int limit)
{
    return rand() % limit;
};
\end{lstlisting}

(There is no \q{limit} name in the \gls{PDB} file; we manually named this argument like this.)

So it returns 
a random value from 0 to a specified limit.

\TT{Rnd()} is called only from one place, 
a function called \TT{StartGame()}, 
and as it seems, this is exactly 
the code which place the mines:

\begin{lstlisting}[style=customasmx86]
.text:010036C7                 push    _xBoxMac
.text:010036CD                 call    _Rnd@4          ; Rnd(x)
.text:010036D2                 push    _yBoxMac
.text:010036D8                 mov     esi, eax
.text:010036DA                 inc     esi
.text:010036DB                 call    _Rnd@4          ; Rnd(x)
.text:010036E0                 inc     eax
.text:010036E1                 mov     ecx, eax
.text:010036E3                 shl     ecx, 5          ; ECX=ECX*32
.text:010036E6                 test    _rgBlk[ecx+esi], 80h
.text:010036EE                 jnz     short loc_10036C7
.text:010036F0                 shl     eax, 5          ; EAX=EAX*32
.text:010036F3                 lea     eax, _rgBlk[eax+esi]
.text:010036FA                 or      byte ptr [eax], 80h
.text:010036FD                 dec     _cBombStart
.text:01003703                 jnz     short loc_10036C7
\end{lstlisting}

Minesweeper allows you to set the board size, so the X (xBoxMac) and Y (yBoxMac) of the board are global variables.
They are passed to \TT{Rnd()} and random 
coordinates are generated.
A mine is placed by the \TT{OR} instruction at \TT{0x010036FA}. 
And if it has been placed before 
(it's possible if the pair of \TT{Rnd()} 
generates a coordinates pair which has been already 
generated), 
then \TT{TEST} and \TT{JNZ} at \TT{0x010036E6} 
jumps to the generation routine again.

\TT{cBombStart} is the global variable containing total number of mines. So this is loop.

The width of the array is 32 
(we can conclude this by looking at the \TT{SHL} instruction, which multiplies one of the coordinates by 32).

The size of the \TT{rgBlk} 
global array can be easily determined by the difference 
between the \TT{rgBlk} 
label in the data segment and the next known one. 
It is 0x360 (864):

\begin{lstlisting}[style=customasmx86]
.data:01005340 _rgBlk          db 360h dup(?)          ; DATA XREF: MainWndProc(x,x,x,x)+574
.data:01005340                                         ; DisplayBlk(x,x)+23
.data:010056A0 _Preferences    dd ?                    ; DATA XREF: FixMenus()+2
...
\end{lstlisting}

$864/32=27$.

So the array size is $27*32$?
It is close to what we know: when we try to set board size to $100*100$ in Minesweeper settings, it fallbacks to a board of size $24*30$.
So this is the maximal board size here.
And the array has a fixed size for any board size.

So let's see all this in \olly.
We will ran Minesweeper, attaching \olly to it and now we can see the memory dump at the address of the \TT{rgBlk} array (\TT{0x01005340})
\footnote{All addresses here are for Minesweeper for Windows XP SP3 English. 
They may differ for other service packs.}.

So we got this memory dump of the array:

\lstinputlisting[style=customasmx86]{examples/minesweeper/1.lst}

\olly, like any other hexadecimal editor, shows 16 bytes per line.
So each 32-byte array row occupies exactly 2 lines here.

This is beginner level (9*9 board).

There is some square 
structure can be seen visually (0x10 bytes).

We will click \q{Run} in \olly to unfreeze the Minesweeper process, then we'll clicked randomly at the Minesweeper window 
and trapped into mine, but now all mines are visible:

\begin{figure}[H]
\centering
\myincludegraphicsSmall{examples/minesweeper/1.png}
\caption{Mines}
\label{fig:minesweeper1}
\end{figure}

By comparing the mine places and the dump, we can conclude that 0x10 stands for border, 0x0F---empty block, 0x8F---mine.
Perhaps, 0x10 is just a \emph{sentinel value}.

Now we'll add comments and also enclose all 0x8F bytes into square brackets:

\lstinputlisting[style=customasmx86]{examples/minesweeper/2.lst}

Now we'll remove all \emph{border bytes} (0x10) and what's beyond those:

\lstinputlisting[style=customasmx86]{examples/minesweeper/3.lst}

Yes, these are mines, now it can be clearly seen and compared with the screenshot.

\clearpage
What is interesting is that we can modify the array right in \olly.
We can remove all mines by changing all 0x8F bytes by 0x0F, and here is what we'll get in Minesweeper:

\begin{figure}[H]
\centering
\myincludegraphicsSmall{examples/minesweeper/3.png}
\caption{All mines are removed in debugger}
\label{fig:minesweeper3}
\end{figure}

We can also move all of them to the first line: 

\begin{figure}[H]
\centering
\myincludegraphicsSmall{examples/minesweeper/2.png}
\caption{Mines set in debugger}
\label{fig:minesweeper2}
\end{figure}

Well, the debugger is not very convenient for eavesdropping (which is our goal anyway), so we'll write a small utility
to dump the contents of the board:

\lstinputlisting[style=customc]{examples/minesweeper/minesweeper_cheater.c}

Just set the \ac{PID}
\footnote{PID it can be seen in Task Manager 
(enable it in \q{View $\rightarrow$ Select Columns})} 
and the address of the array (\TT{0x01005340} for Windows XP SP3 English) 
and it will dump it
\footnote{The compiled executable is here: 
\href{http://go.yurichev.com/17165}{beginners.re}}.

It attaches itself to a win32 process by \ac{PID} and just reads process memory at the address.

\subsection{Finding grid automatically}

This is kind of nuisance to set address each time when we run our utility.
Also, various Minesweeper versions may have the array on different address.
Knowing the fact that there is always a border (0x10 bytes), we can just find it in memory:

\lstinputlisting[style=customc]{examples/minesweeper/cheater2_fragment.c}

Full source code: \url{https://github.com/DennisYurichev/RE-for-beginners/blob/master/examples/minesweeper/minesweeper_cheater2.c}.

\subsection{\Exercises}

\begin{itemize}

\item 
Why do the \emph{border bytes} (or \emph{sentinel values}) (0x10) exist in the array?

What they are for if they are not visible in Minesweeper's interface?
How could it work without them?

\item 
As it turns out, there are more values possible (for open blocks, for flagged by user, \etc{}).
Try to find the meaning of each one.

\item 
Modify my utility so it can remove all mines or set them in a fixed pattern that you want in the Minesweeper
process currently running.

\end{itemize}
}
\RU{\mysection{Разгон майнера биткоинов Cointerra}
\index{Bitcoin}
\index{BeagleBone}

Был такой майнер биткоинов Cointerra, выглядящий так:

\begin{figure}[H]
\centering
\myincludegraphics{examples/bitcoin_miner/board.jpg}
\caption{Board}
\end{figure}

И была также (возможно утекшая) утилита\footnote{Можно скачать здесь: \url{https://github.com/DennisYurichev/RE-for-beginners/raw/master/examples/bitcoin_miner/files/cointool-overclock}}
которая могла выставлять тактовую частоту платы.
Она запускается на дополнительной плате BeagleBone на ARM с Linux (маленькая плата внизу фотографии).

И у автора (этих строк) однажды спросили, можно ли хакнуть эту утилиту и посмотреть, какие частоты можно выставлять, и какие нет.
И можно ли твикнуть её?

Утилиту нужно запускать так: \TT{./cointool-overclock 0 0 900}, где 900 это частота в МГц.
Если частота слишком большая, утилита выведет ошибку \q{Error with arguments} и закончит работу.

Вот фрагмент кода вокруг ссылки на текстовую строку \q{Error with arguments}:

\begin{lstlisting}[style=customasmARM]

...

.text:0000ABC4         STR      R3, [R11,#var_28]
.text:0000ABC8         MOV      R3, #optind
.text:0000ABD0         LDR      R3, [R3]
.text:0000ABD4         ADD      R3, R3, #1
.text:0000ABD8         MOV      R3, R3,LSL#2
.text:0000ABDC         LDR      R2, [R11,#argv]
.text:0000ABE0         ADD      R3, R2, R3
.text:0000ABE4         LDR      R3, [R3]
.text:0000ABE8         MOV      R0, R3  ; nptr
.text:0000ABEC         MOV      R1, #0  ; endptr
.text:0000ABF0         MOV      R2, #0  ; base
.text:0000ABF4         BL       strtoll
.text:0000ABF8         MOV      R2, R0
.text:0000ABFC         MOV      R3, R1
.text:0000AC00         MOV      R3, R2
.text:0000AC04         STR      R3, [R11,#var_2C]
.text:0000AC08         MOV      R3, #optind
.text:0000AC10         LDR      R3, [R3]
.text:0000AC14         ADD      R3, R3, #2
.text:0000AC18         MOV      R3, R3,LSL#2
.text:0000AC1C         LDR      R2, [R11,#argv]
.text:0000AC20         ADD      R3, R2, R3
.text:0000AC24         LDR      R3, [R3]
.text:0000AC28         MOV      R0, R3  ; nptr
.text:0000AC2C         MOV      R1, #0  ; endptr
.text:0000AC30         MOV      R2, #0  ; base
.text:0000AC34         BL       strtoll
.text:0000AC38         MOV      R2, R0
.text:0000AC3C         MOV      R3, R1
.text:0000AC40         MOV      R3, R2
.text:0000AC44         STR      R3, [R11,#third_argument]
.text:0000AC48         LDR      R3, [R11,#var_28]
.text:0000AC4C         CMP      R3, #0
.text:0000AC50         BLT      errors_with_arguments
.text:0000AC54         LDR      R3, [R11,#var_28]
.text:0000AC58         CMP      R3, #1
.text:0000AC5C         BGT      errors_with_arguments
.text:0000AC60         LDR      R3, [R11,#var_2C]
.text:0000AC64         CMP      R3, #0
.text:0000AC68         BLT      errors_with_arguments
.text:0000AC6C         LDR      R3, [R11,#var_2C]
.text:0000AC70         CMP      R3, #3
.text:0000AC74         BGT      errors_with_arguments
.text:0000AC78         LDR      R3, [R11,#third_argument]
.text:0000AC7C         CMP      R3, #0x31
.text:0000AC80         BLE      errors_with_arguments
.text:0000AC84         LDR      R2, [R11,#third_argument]
.text:0000AC88         MOV      R3, #950
.text:0000AC8C         CMP      R2, R3
.text:0000AC90         BGT      errors_with_arguments
.text:0000AC94         LDR      R2, [R11,#third_argument]
.text:0000AC98         MOV      R3, #0x51EB851F
.text:0000ACA0         SMULL    R1, R3, R3, R2
.text:0000ACA4         MOV      R1, R3,ASR#4
.text:0000ACA8         MOV      R3, R2,ASR#31
.text:0000ACAC         RSB      R3, R3, R1
.text:0000ACB0         MOV      R1, #50
.text:0000ACB4         MUL      R3, R1, R3
.text:0000ACB8         RSB      R3, R3, R2
.text:0000ACBC         CMP      R3, #0
.text:0000ACC0         BEQ      loc_ACEC
.text:0000ACC4
.text:0000ACC4 errors_with_arguments
.text:0000ACC4                                         
.text:0000ACC4         LDR      R3, [R11,#argv]
.text:0000ACC8         LDR      R3, [R3]
.text:0000ACCC         MOV      R0, R3  ; path
.text:0000ACD0         BL       __xpg_basename
.text:0000ACD4         MOV      R3, R0
.text:0000ACD8         MOV      R0, #aSErrorWithArgu ; format
.text:0000ACE0         MOV      R1, R3
.text:0000ACE4         BL       printf
.text:0000ACE8         B        loc_ADD4
.text:0000ACEC ; ------------------------------------------------------------
.text:0000ACEC
.text:0000ACEC loc_ACEC                 ; CODE XREF: main+66C
.text:0000ACEC         LDR      R2, [R11,#third_argument]
.text:0000ACF0         MOV      R3, #499
.text:0000ACF4         CMP      R2, R3
.text:0000ACF8         BGT      loc_AD08
.text:0000ACFC         MOV      R3, #0x64
.text:0000AD00         STR      R3, [R11,#unk_constant]
.text:0000AD04         B        jump_to_write_power
.text:0000AD08 ; ------------------------------------------------------------
.text:0000AD08
.text:0000AD08 loc_AD08                 ; CODE XREF: main+6A4
.text:0000AD08         LDR      R2, [R11,#third_argument]
.text:0000AD0C         MOV      R3, #799
.text:0000AD10         CMP      R2, R3
.text:0000AD14         BGT      loc_AD24
.text:0000AD18         MOV      R3, #0x5F
.text:0000AD1C         STR      R3, [R11,#unk_constant]
.text:0000AD20         B        jump_to_write_power
.text:0000AD24 ; ------------------------------------------------------------
.text:0000AD24
.text:0000AD24 loc_AD24                 ; CODE XREF: main+6C0
.text:0000AD24         LDR      R2, [R11,#third_argument]
.text:0000AD28         MOV      R3, #899
.text:0000AD2C         CMP      R2, R3
.text:0000AD30         BGT      loc_AD40
.text:0000AD34         MOV      R3, #0x5A
.text:0000AD38         STR      R3, [R11,#unk_constant]
.text:0000AD3C         B        jump_to_write_power
.text:0000AD40 ; ------------------------------------------------------------
.text:0000AD40
.text:0000AD40 loc_AD40                 ; CODE XREF: main+6DC
.text:0000AD40         LDR      R2, [R11,#third_argument]
.text:0000AD44         MOV      R3, #999
.text:0000AD48         CMP      R2, R3
.text:0000AD4C         BGT      loc_AD5C
.text:0000AD50         MOV      R3, #0x55
.text:0000AD54         STR      R3, [R11,#unk_constant]
.text:0000AD58         B        jump_to_write_power
.text:0000AD5C ; ------------------------------------------------------------
.text:0000AD5C
.text:0000AD5C loc_AD5C                 ; CODE XREF: main+6F8
.text:0000AD5C         LDR      R2, [R11,#third_argument]
.text:0000AD60         MOV      R3, #1099
.text:0000AD64         CMP      R2, R3
.text:0000AD68         BGT      jump_to_write_power
.text:0000AD6C         MOV      R3, #0x50
.text:0000AD70         STR      R3, [R11,#unk_constant]
.text:0000AD74
.text:0000AD74 jump_to_write_power                     ; CODE XREF: main+6B0
.text:0000AD74                                         ; main+6CC ...
.text:0000AD74         LDR      R3, [R11,#var_28]
.text:0000AD78         UXTB     R1, R3
.text:0000AD7C         LDR      R3, [R11,#var_2C]
.text:0000AD80         UXTB     R2, R3
.text:0000AD84         LDR      R3, [R11,#unk_constant]
.text:0000AD88         UXTB     R3, R3
.text:0000AD8C         LDR      R0, [R11,#third_argument]
.text:0000AD90         UXTH     R0, R0
.text:0000AD94         STR      R0, [SP,#0x44+var_44]
.text:0000AD98         LDR      R0, [R11,#var_24]
.text:0000AD9C         BL       write_power
.text:0000ADA0         LDR      R0, [R11,#var_24]
.text:0000ADA4         MOV      R1, #0x5A
.text:0000ADA8         BL       read_loop
.text:0000ADAC         B        loc_ADD4

...

.rodata:0000B378 aSErrorWithArgu DCB "%s: Error with arguments",0xA,0 ; DATA XREF: main+684

...

\end{lstlisting}

Имена ф-ций присутствовали в отладочной информации в оригинальном исполняемом файле,
такие как \TT{write\_power}, \TT{read\_loop}.
Но имена меткам внутри ф-ции дал я.

\myindex{UNIX!getopt}
\myindex{strtoll()}
Имя \TT{optind} звучит знакомо. Это библиотека \IT{getopt} из *NIX предназначенная для парсинга командной строки ---
и это то, что внутри и происходит.
Затем, третий аргумент (где передается значение частоты) конвертируется из строку в число используя вызов ф-ции \IT{strtoll()}.

Значение затем сравнивается с разными константами.
На 0xACEC есть проверка, меньше ли оно или равно 499, и если это так, то 0x64 будет передано в ф-цию
\TT{write\_power()} (которая посылает команду через USB используя \TT{send\_msg()}).
Если значение больше 499, происходит переход на 0xAD08.

На 0xAD08 есть проверка, меньше ли оно или равно 799. Если это так, то 0x5F передается в ф-цию \TT{write\_power()}.

Есть еще проверки: на 899 на 0xAD24, на 0x999 на 0xAD40, и наконец, на 1099 на 0xAD5C.
Если входная частота меньше или равна 1099, 0x50 (на 0xAD6C) будет передано в ф-цию \TT{write\_power()}.
И тут что-то вроде баги.
Если значение все еще больше 1099, само значение будет передано в ф-цию \TT{write\_power()}.
Но с другой стороны это не бага, потому что мы не можем попасть сюда: значение в начале проверяется с 950 на 0xAC88,
и если оно больше, выводится сообщение об ошибке и утилита заканчивает работу.

Вот таблица между частотами в МГц и значениями передаваемыми в ф-цию \TT{write\_power()}:

\begin{center}
\begin{longtable}{ | l | l | l | }
\hline
\HeaderColor МГц & \HeaderColor шестнадцатеричное представление & \HeaderColor десятичное \\
\hline
499MHz & 0x64 & 100 \\
\hline
799MHz & 0x5f & 95 \\
\hline
899MHz & 0x5a & 90 \\
\hline
999MHz & 0x55 & 85 \\
\hline
1099MHz & 0x50 & 80 \\
\hline
\end{longtable}
\end{center}

Как видно, значение передаваемое в плату постепенно уменьшается с ростом частоты.

Видно что значение в 950МГц это жесткий предел, по крайней мере в этой утилите. Можно ли её обмануть?

Вернемся к этому фрагменту кода:

\begin{lstlisting}[style=customasmARM]
.text:0000AC84      LDR     R2, [R11,#third_argument]
.text:0000AC88      MOV     R3, #950
.text:0000AC8C      CMP     R2, R3
.text:0000AC90      BGT     errors_with_arguments ; Я пропатчил здесь на 00 00 00 00
\end{lstlisting}

Нам нужно как-то запретить инструкцию перехода \INS{BGT} на 0xAC90. И это ARM в режиме ARM, потому что, как мы видим,
все адреса увеличиваются на 4, т.е., длина каждой инструкции это 4 байта.
Инструкция \TT{NOP} (нет операции) в режиме ARM это просто 4 нулевых байта: \TT{00 00 00 00}.
Так что, записывая 4 нуля по адресу 0xAC90 (или по физическому смещению в файле: 0x2C90) мы можем выключить
эту проверку.

Теперь можно выставлять частоты вплоть до 1050МГц. И даже больше, но из-за ошибки, если входное значение больше 1099,
значение в МГц, \IT{как есть}, будет передано в плату, что неправильно.

Дальше я не разбирался, но если бы продолжил, я бы уменьшал значение передаваемое в ф-цию \TT{write\_power()}.

Теперь страшный фрагмент кода, который я в начале пропустил:

\lstinputlisting[style=customasmARM]{examples/bitcoin_miner/tmp1.lst}

Здесь используется деление через умножение, и константа 0x51EB851F.
Я написал для себя простой программистский калькулятор\footnote{\url{https://github.com/DennisYurichev/progcalc}}.
И там есть возможность вычислять обратное число по модулю.

\begin{lstlisting}
modinv32(0x51EB851F)
Warning, result is not integer: 3.125000
(unsigned) dec: 3 hex: 0x3 bin: 11
\end{lstlisting}

Это значит что инструкция \INS{SMULL} на 0xACA0 просто делит 3-й аргумент на 3.125.
На самом деле, все что делает ф-ция \TT{modinv32()} в моем калькуляторе, это:

\[
\frac{1}{\frac{input}{2^{32}}} = \frac{2^{32}}{input}
\]

Потом там есть дополнительные сдвиги и теперь мы видим что 3-й аргумент просто делится на 50.
И затем умножается снова на 50.
Зачем?
Это простейшая проверка, можно ли делить входное значение на 50 без остатка.
Если значение этого выражения ненулевое, $x$ не может быть разделено на 50 без остатка:

\[
x-((\frac{x}{50}) \cdot 50)
\]

На самом деле, это простой способ вычисления остатка от деления.

И затем, если остаток ненулевой, выводится сообщение об ошибке.
Так что эта утилита берет значения частотв вроде 850, 900, 950, 1000, итд, но не 855 или 911.

Вот и всё! Если вы делаете что-то такое, имейте ввиду, что это может испортить вашу плату, как и в случае разгона
чипов вроде \ac{CPU}, \ac{GPU}, итд.
Если у вас есть плата Cointerra, делайте всё это на свой собственный риск!

}
\DE{\mysection{\Stack}
\label{sec:stack}
\myindex{\Stack}

Der Stack ist eine der fundamentalen Datenstrukturen in der Informatik.
\footnote{\href{http://go.yurichev.com/17119}{wikipedia.org/wiki/Call\_Stack}}.
\ac{AKA} \ac{LIFO}.

Technisch betrachtet ist es ein Stapelspeicher innerhalb des Prozessspeichers der zusammen mit den \ESP (x86), \RSP (x64) oder dem \ac{SP} (ARM) Register als ein Zeiger in diesem Speicherblock fungiert.

\myindex{ARM!\Instructions!PUSH}
\myindex{ARM!\Instructions!POP}
\myindex{x86!\Instructions!PUSH}
\myindex{x86!\Instructions!POP}

Die häufigsten Stack-Zugriffsinstruktionen sind die \PUSH- und \POP-Instruktionen (in beidem x86 und ARM Thumb-Modus). \PUSH subtrahiert vom \ESP/\RSP/\ac{SP} 4 Byte im 32-Bit Modus (oder 8 im 64-Bit Modus) und schreibt dann den Inhalt des Zeigers an die Adresse auf die von \ESP/\RSP/\ac{SP} gezeigt wird.

\POP ist die umgekehrte Operation: Die Daten des Zeigers für die Speicherregion auf die von \ac{SP}
gezeigt wird werden ausgelesen und die Inhalte in den Instruktionsoperanden geschreiben (oft ist das ein Register). Dann werden 4 (beziehungsweise 8) Byte zum \gls{stack pointer} addiert.

Nach der Stackallokation, zeigt der \gls{stack pointer} auf den Boden des Stacks.
\PUSH verringert den \gls{stack pointer} und \POP erhöht ihn.
Der Boden des Stacks ist eigentlich der Anfang der Speicherregion die für den Stack reserviert wurde.
Das wirkt zunächst seltsam, aber so funktioniert es.

ARM unterstützt beides, aufsteigende und absteigende Stacks.

\myindex{ARM!\Instructions!STMFD}
\myindex{ARM!\Instructions!LDMFD}
\myindex{ARM!\Instructions!STMED}
\myindex{ARM!\Instructions!LDMED}
\myindex{ARM!\Instructions!STMFA}
\myindex{ARM!\Instructions!LDMFA}
\myindex{ARM!\Instructions!STMEA}
\myindex{ARM!\Instructions!LDMEA}

Zum Beispiel die \ac{STMFD}/\ac{LDMFD} und \ac{STMED}/\ac{LDMED} Instruktionen sind alle dafür gedacht mit einem absteigendem Stack zu arbeiten ( wächst nach unten, fängt mit hohen Adressen an und entwickelt sich zu niedrigeren Adressen). Die \ac{STMFA}/\ac{LDMFA} und \ac{STMEA}/\ac{LDMEA} Instruktionen sind dazu gedacht mit einem aufsteigendem Stack zu arbeiten (wächst nach oben und fängt mit niedrigeren Adressen an und wächst nach oben).

% It might be worth mentioning that STMED and STMEA write first,
% and then move the pointer, and that LDMED and LDMEA move the pointer first, and then read.
% In other words, ARM not only lets the stack grow in a non-standard direction,
% but also in a non-standard order.
% Maybe this can be in the glossary, which would explain why E stands for "empty".

\subsection{Warum wächst der Stack nach unten?}
\label{stack_grow_backwards}

Intuitiv, würden man annehmen das der Stack nach oben wächst z.B Richtung höherer Adressen, so wie bei jeder anderen Datenstruktur.

Der Grund das der Stack rückwärts wächst ist wohl historisch bedingt. Als Computer so groß waren das sie einen ganzen Raum beansprucht haben war es einfach Speicher in zwei Sektionen zu unterteilen, einen Teil für den \gls{heap} und einen Teil für den Stack. Sicher war zu dieser Zeit nicht bekannt wie groß der \gls{heap} und der Stack wachsen würden, während der Programm Laufzeit, also war die Lösung die einfachste mögliche.

\input{patterns/02_stack/stack_and_heap}

In \RitchieThompsonUNIX können wir folgendes lesen:

\begin{framed}
\begin{quotation}
Der user-core eines Programm Images wird in drei logische Segmente unterteilt. Das Programm-Text Segment beginnt bei 0 im virtuellen Adress Speicher. Während der Ausführung wird das Segment als schreibgeschützt markiert und eine einzelne Kopie des Segments wird unter allen Prozessen geteilt die das Programm ausführen. An der ersten 8K grenze über dem Programm Text Segment im Virtuellen Speicher, fängt der ``nonshared'' Bereich an, der nach Bedarf von Syscalls erweitert werden kann. Beginnend bei der höchsten Adresse im Virtuellen Speicher ist das Stack Segment, das Automatisch nach unten wächst während der Hardware Stackpointer sich ändert.
\end{quotation}
\end{framed}

Das erinnert daran wie manche Schüler Notizen zu  zwei Vorträgen in einem Notebook dokumentieren:
Notizen für den ersten Vortrag werden normal notiert, und Notizen zur zum zweiten Vortrag werden 
ans Ende des Notizbuches geschrieben, indem man das Notizbuch umdreht. Die Notizen treffen sich irgendwann
im Notizbuch aufgrund des fehlenden Freien Platzes.

% I think if we want to expand on this analogy,
% one might remember that the line number increases as as you go down a page.
% So when you decrease the address when pushing to the stack, visually,
% the stack does grow upwards.
% Of course, the problem is that in most human languages,
% just as with computers,
% we write downwards, so this direction is what makes buffer overflows so messy.

\subsection{Für was wird der Stack benutzt?}

% subsections
\input{patterns/02_stack/01_saving_ret_addr}
\input{patterns/02_stack/02_args_passing}
\EN{\input{patterns/02_stack/03_local_vars_EN}}
\RU{\input{patterns/02_stack/03_local_vars_RU}}
\DE{\input{patterns/02_stack/03_local_vars_DE}}
\PTBR{\input{patterns/02_stack/03_local_vars_PTBR}}
\input{patterns/02_stack/04_alloca/main}
\input{patterns/02_stack/05_SEH}
\input{patterns/02_stack/06_BO_protection}

\subsubsection{Automatisches deallokieren der Daten auf dem Stack}

Vielleicht ist der Grund warum man lokale Variablen und SEH Einträge auf dem Stack speichert, weil sie beim 
verlassen der Funktion automatisch aufgeräumt werden. Man braucht dabei nur eine Instruktion um die Position
des Stackpointers zu korrigieren (oftmals ist es die \ADD Instruktion). Funktions Argumente, könnte man sagen 
werden auch am Ende der Funktion deallokiert. Im Kontrast dazu, alles was auf dem \emph{heap} gespeichert wird muss
explizit deallokiert werden. 

% sections
\EN{\input{patterns/02_stack/07_layout_EN}}
\RU{\input{patterns/02_stack/07_layout_RU}}
\DE{\input{patterns/02_stack/07_layout_DE}}
\PTBR{\input{patterns/02_stack/07_layout_PTBR}}
\input{patterns/02_stack/08_noise/main}
\input{patterns/02_stack/exercises}
}
\FR{\mysection{x86}

\subsection{Terminologie}

Commun en 16-bit (8086/80286), 32-bit (80386, etc.), 64-bit.

\myindex{IEEE 754}
\myindex{MS-DOS}
\begin{description}
	\item[octet] 8-bit.
		La directive d'assembleur DB est utilisée pour définir les variables et les tableaux d'octets.
		Les octets sont passés dans les parties 8-bit des registres: \TT{AL/BL/CL/DL/AH/BH/CH/DH/SIL/DIL/R*L}.
	\item[mot] 16-bit. 
		directive assembleur DW \dittoclosing.
		Les mots sont passés dans la partie 16-bit des registres:\\
			\TT{AX/BX/CX/DX/SI/DI/R*W}.
	\item[double mot] (\q{dword}) 32-bit.
		directive assembleur DD \dittoclosing.
		Les double mots sont passés dans les registres (x86) ou dans la partie 32-bit de registres (x64).
		Dans du code 16-bit, les double mots sont passés dans une paire de registres 16-bit.
	\item[quadruple mot] (\q{qword}) 64-bit.
		directive assembleur DQ \dittoclosing.
		En environnement 32-bit, les quadruple mots sont passés dans une paire de registres 32-bit.
	\item[tbyte] (10 bytes) 80-bit ou 10 octets (utilisé pour les registres FPU IEEE 754).
	\item[paragraph] (16 bytes)---le terme était répandu dans l'environnement MS-DOS.
\end{description}

\myindex{Windows!API}

Des types de données de même taille (BYTE, WORD, DWORD) existent aussi dans l'\ac{API} Windows.

\input{appendix/x86/registers_FR} % subsection
\input{appendix/x86/instructions_FR} % subsection
\input{appendix/x86/npad} % subsection

}
\IT{\mysection{x86}

\subsection{Terminologia}

Comune per 16-bit (8086/80286), 32-bit (80386, etc.), 64-bit.

\myindex{IEEE 754}
\myindex{MS-DOS}
\begin{description}
	\item[byte] 8-bit.
		La direttiva DB in assembly è utilizzata per definire variabili ed array di byte.
		I byte sono passati nella parte ad 8-bit dei registri: \TT{AL/BL/CL/DL/AH/BH/CH/DH/SIL/DIL/R*L}.
	\item[word] 16-bit.
		La direttiva DW in assembly \dittoclosing.
		Le Word sono passate nella parte a 16-bit dei registri:\\
			\TT{AX/BX/CX/DX/SI/DI/R*W}.
	\item[double word] (\q{dword}) 32-bit.
		La direttiva DD in assembly \dittoclosing.
		Le Double words sono passate nei registri (x86) o nelle parti a 32-bit dei registri (x64).
		Nel codice a 16-bit, le double words sono passate in coppie di registri a 16-bit.
	\item[quad word] (\q{qword}) 64-bit.
		La direttiva DQ in assembly \dittoclosing.
		Negli ambienti a 32-bit, le quad words sono passate in coppie di registri a 32-bit.
	\item[tbyte] (10 bytes) 80-bit o 10 bytes (usati per i registri FPU IEEE 754).
	\item[paragraph] (16 bytes)---termine popolare nell'ambiente MS-DOS. % TODO link to a part about 8086 memory model...
\end{description}

\myindex{Windows!API}

Tipi di dati della stessa dimensione (BYTE, WORD, DWORD) sono gli stessi nelle Windows \ac{API}.

%TBT
\input{appendix/x86/registers} % subsection
\input{appendix/x86/instructions} % subsection
\input{appendix/x86/npad} % subsection
}
\mysection{ARM}
\myindex{ARM}

\subsection{\RU{Терминология}\EN{Terminology}}

ARM \RU{изначально разрабатывался как 32-битный}\EN{was initially developed as 32-bit} \ac{CPU}, 
\RU{поэтому \emph{слово} здесь, в отличие от x86, 32-битное}\EN{so that's why a \emph{word} here, unlike x86, is 32-bit}.

\begin{description}
	\item[byte] 8-\bitENRU.
		\RU{Для определения переменных и массива байт используется директива ассемблера DCB}
		\EN{The DB assembly directive is used for defining variables and arrays of bytes}.
	\item[halfword] 16-\bitENRU. \RU{\dittoclosing директива ассемблера DCW}
					\EN{DCW assembly directive \dittoclosing}.	
	\item[word] 32-\bitENRU. \RU{\dittoclosing директива ассемблера DCD}
					\EN{DCD assembly directive \dittoclosing}.
	\item[doubleword] 64-\bitENRU.
	\item[quadword] 128-\bitENRU.
\end{description}

\subsection{\RU{Версии}\EN{Versions}}

\begin{itemize}
\item ARMv4: \RU{появился режим Thumb}\EN{Thumb mode introduced}.

\item ARMv6: \RU{использовался в}\EN{used in} iPhone 1st gen., iPhone 3G 
(Samsung 32-bit RISC ARM 1176JZ(F)-S \RU{поддерживающий}\EN{that supports} Thumb-2)

\item ARMv7: \RU{появился }Thumb-2\EN{ was added} (2003).
\RU{Использовался в}\EN{was used in} iPhone 3GS, iPhone 4, iPad 1st gen. (ARM Cortex-A8), iPad 2 (Cortex-A9),
iPad 3rd gen.

\item ARMv7s: \RU{Добавлены новые инструкции}\EN{New instructions added}.
\RU{Использовался в}\EN{Was used in} iPhone 5, iPhone 5c, iPad 4th gen. (Apple A6).

\item ARMv8: 64-\RU{битный процессор}\EN{bit CPU}, \ac{AKA} ARM64 \ac{AKA} AArch64.
\RU{Использовался в}\EN{Was used in} iPhone 5S, iPad Air (Apple A7).
\RU{В 64-битном режиме, режима Thumb больше нет, только режим ARM (4-байтные инструкции).}
\EN{There is no Thumb mode in 64-bit mode, only ARM (4-byte instructions).}
\end{itemize}

% sections
\input{appendix/ARM/registers}
\input{appendix/ARM/instructions}

\mysection{MIPS}

\subsection{\Registers}
\label{MIPS_registers_ref}

\myindex{MIPS!O32}
( \RU{Соглашение о вызовах O32}\EN{O32 calling convention} )

\subsubsection{\RU{Регистры общего пользования}\EN{General purpose registers} \ac{GPR}}

%\small
\begin{center}
\begin{tabular}{ | l | l | l | }
\hline
\HeaderColor \RU{Номер}\EN{Number} & 
\HeaderColor \RU{Псевдоимя}\EN{Pseudoname} & 
\HeaderColor \RU{Описание}\EN{Description} \\
\hline
\$0             & \$ZERO          & \RU{Всегда ноль. Запись в этот регистр это как}\EN{Always zero. Writing to this register is like} \ac{NOP}. \\
\hline
\$1             & \$AT            & \RU{Используется как временный регистр}\EN{Used as a temporary register} \\
                &                 & \RU{для ассемблерных макросов и псевдоинструкций}\EN{for assembly macros and pseudo instructions}. \\
\hline
\$2 \dots \$3   & \$V0 \dots \$V1 & \RU{Здесь возвращается результат функции}\EN{Function result is returned here}. \\
\hline
\$4 \dots \$7   & \$A0 \dots \$A3 & \RU{Аргументы функции}\EN{Function arguments}. \\
\hline
\$8 \dots \$15  & \$T0 \dots \$T7 & \RU{Используется для временных данных}\EN{Used for temporary data}. \\
\hline
\$16 \dots \$23 & \$S0 \dots \$S7 & \RU{Используется для временных данных}\EN{Used for temporary data}\AsteriskOne{}. \\
\hline
\$24 \dots \$25 & \$T8 \dots \$T9 & \RU{Используется для временных данных}\EN{Used for temporary data}. \\
\hline
\$26 \dots \$27 & \$K0 \dots \$K1 & \RU{Зарезервировано для ядра \ac{OS}}\EN{Reserved for \ac{OS} kernel}. \\
\hline
\$28            & \$GP            & \RU{Глобальный указатель}\EN{Global Pointer}\AsteriskTwo{}. \\
\hline
\$29            & \$SP            & \ac{SP}\AsteriskOne{}. \\
\hline
\$30            & \$FP            & \ac{FP}\AsteriskOne{}. \\
\hline
\$31            & \$RA            & \ac{RA}. \\
\hline
n/a             & PC              & \ac{PC}. \\
\hline
n/a             & HI              & \RU{старшие 32 бита результата умножения или остаток от деления}\EN{high 32 bit of multiplication or division remainder}\AsteriskThree{}. \\
\hline
n/a             & LO              & \RU{младшие 32 бита результата умножения или результат деления}\EN{low 32 bit of multiplication and division remainder}\AsteriskThree{}. \\
\hline
\end{tabular}
\end{center}
%\normalsize

\subsubsection{\RU{Регистры для работы с числами с плавающей точкой}\EN{Floating-point registers}}
\label{MIPS_FPU_registers}

\begin{center}
\begin{tabular}{ | l | l | l | }
\hline
\HeaderColor \RU{Название}\EN{Name} & \HeaderColor \RU{Описание}\EN{Description} \\
\hline
\$F0..\$F1   & \RU{Здесь возвращается результат функции}\EN{Function result returned here}. \\
\hline
\$F2..\$F3   & \RU{Не используется}\EN{Not used}. \\
\hline
\$F4..\$F11  & \RU{Используется для временных данных}\EN{Used for temporary data}. \\
\hline
\$F12..\$F15 & \RU{Первые два аргумента функции}\EN{First two function arguments}. \\
\hline
\$F16..\$F19 & \RU{Используется для временных данных}\EN{Used for temporary data}. \\
\hline
\$F20..\$F31 & \RU{Используется для временных данных}\EN{Used for temporary data}\AsteriskOne{}. \\
\hline
%fcr31 & Control/status register. \\
%\hline
\end{tabular}
\end{center}

\AsteriskOne{} --- \Gls{callee} \RU{должен сохранять}\EN{must preserve the value}.\\
\AsteriskTwo{} --- \Gls{callee} \RU{должен сохранять}\EN{must preserve the value} (\RU{кроме \ac{PIC}-кода}
\EN{except in \ac{PIC} code}).\\
\myindex{MIPS!\Instructions!MFLO}
\myindex{MIPS!\Instructions!MFHI}
\AsteriskThree{} --- \RU{доступны используя инструкции}\EN{accessible using the} \TT{MFHI} \AndENRU \TT{MFLO}\EN{ instructions}.\\

\subsection{\Instructions}

\RU{Есть три типа инструкций}\EN{There are 3 kinds of instructions}:

\begin{itemize}

\item \RU{Тип R: имеющие 3 регистра. R-инструкции обычно имеют такой вид:}
\EN{R-type: those which have 3 registers. R-instruction usually have the following form:}

\begin{lstlisting}
instruction destination, source1, source2
\end{lstlisting}

\RU{Важно помнить, что если первый и второй регистр один и тот же, IDA может показать инструкцию в сокращенной
форме:}
\EN{One important thing to keep in mind is that when the first and second register are the same, 
IDA may show the instruction in its shorter form:}

\begin{lstlisting}
instruction destination/source1, source2
\end{lstlisting}

\RU{Это немного напоминает Интеловский синтаксис ассемблера x86.}
\EN{That somewhat reminds us of the Intel syntax for x86 assembly language.}

\item \RU{Тип I: имеющие 2 регистра и 16-битное \q{immediate}-значение.}
\EN{I-type: those which have 2 registers and a 16-bit immediate value.}

\item \RU{Тип J: инструкции перехода, имеют 26 бит для кодирования смещения.}
\EN{J-type: jump/branch instructions, have 26 bits for encoding the offset.}

\end{itemize}

\subsubsection{\RU{Инструкции перехода}\EN{Jump instructions}}

\RU{Какая разница между инструкциями начинающихся с B- (BEQ, B, итд) и с J- (JAL, JALR, итд)?}
\EN{What is the difference between B- instructions (BEQ, B, etc.) and J- ones (JAL, JALR, etc.)?}

\RU{B-инструкции имеют тип I, так что, смещение в этих инструкциях кодируется как 16-битное значение.}
\EN{The B-instructions have an I-type, hence, the B-instructions' offset is encoded as a 16-bit immediate.}
\RU{Инструкции JR и JALR имеют тип R, и они делают переход по абсолютному адресу указанному в регистре.}
\EN{JR and JALR are R-type and jump to an absolute address specified in a register.}
\RU{J и JAL имеют тип J, так что смещение кодируется как 26-битное значение.}
\EN{J and JAL are J-type, hence the offset is encoded as a 26-bit immediate.}

\RU{Коротко говоря, в B-инструкциях можно кодировать условие 
(B на самом деле это псевдоинструкция для \TT{BEQ \$ZERO, \$ZERO, LABEL}),
а в J-инструкциях нельзя.}
\EN{In short, B-instructions can encode a condition 
(B is in fact pseudo instruction for \TT{BEQ \$ZERO, \$ZERO, LABEL}), 
while J-instructions can't.}

\mysection{\RU{Некоторые библиотечные функции GCC}\EN{Some GCC library functions}\DE{Einige GCC-Bibliotheks-Funktionen}%
\FR{Quelques fonctions de la bibliothèque de GCC}}
\myindex{GCC}
\label{sec:GCC_library_func}

%__ashldi3
%__ashrdi3
%__floatundidf
%__floatdisf
%__floatdixf
%__floatundidf
%__floatundisf
%__floatundixf
%__lshrdi3
%__muldi3

\begin{center}
\begin{tabular}{ | l | l | }
\hline
\HeaderColor \RU{имя}\EN{name}\DE{Name}\FR{nom} & \HeaderColor \RU{значение}\EN{meaning}\DE{Bedeutung}%
\FR{signification} \\
\hline \TT{\_\_divdi3} & \RU{знаковое деление}\EN{signed division}\DE{vorzeichenbehaftete Division}%
\FR{division signée} \\
\hline \TT{\_\_moddi3} & \RU{остаток от знакового деления}\EN{getting remainder (modulo) of signed division}%
\DE{Rest (Modulo) einer vorzeichenbehafteten Division}%
\FR{reste (modulo) d'une division signée} \\
\hline \TT{\_\_udivdi3} & \RU{беззнаковое деление}\EN{unsigned division}\DE{vorzeichenlose Division}%
\FR{division non signée} \\
\hline \TT{\_\_umoddi3} & \RU{остаток от беззнакового деления}\EN{getting remainder (modulo) of unsigned division}%
\DE{Rest (Modulo) einer vorzeichenlosen Division}%
\FR{reste (modulo) d'une division non signée} \\
\hline
\end{tabular}
\end{center}



\mysection{\RU{Некоторые библиотечные функции MSVC}\EN{Some MSVC library functions}\DE{Einige MSVC-Bibliotheks-Funktionen}%
\FR{Quelques fonctions de la bibliothèque MSVC}}
\myindex{MSVC}
\label{sec:MSVC_library_func}

\TT{ll} \RU{в имени функции означает}\EN{in function name stands for}\DE{in Funktionsnamen steht für}%
\FR{dans une fontion signifie} \q{long long}, \RU{т.е. 64-битный тип данных}
\EN{e.g., a 64-bit data type}\DE{z.B. einen 64-Bit-Datentyp}\FR{i.e., type de donées 64-bit}.

\begin{center}
\begin{tabular}{ | l | l | }
\hline
\HeaderColor \RU{имя}\EN{name}\DE{Name}\FR{nom} & \HeaderColor \RU{значение}\EN{meaning}\DE{Bedeutung}\FR{signification} \\
\hline \TT{\_\_alldiv} & \RU{знаковое деление}\EN{signed division}\DE{vorzeichenbehaftete Division}\FR{division signée} \\
\hline \TT{\_\_allmul} & \RU{умножение}\EN{multiplication}\DE{Multiplikation}\FR{multiplication} \\
\hline \TT{\_\_allrem} & \RU{остаток от знакового деления}\EN{remainder of signed division}\DE{Rest einer vorzeichenbehafteten Division}%
\FR{reste de la division signée} \\
\hline \TT{\_\_allshl} & \RU{сдвиг влево}\EN{shift left}\DE{Schiebe links}\FR{décalage à gauche} \\
\hline \TT{\_\_allshr} & \RU{знаковый сдвиг вправо}\EN{signed shift right}\DE{Schiebe links, vorzeichenbehaftet}%
\FR{décalage signé à droite} \\
\hline \TT{\_\_aulldiv} & \RU{беззнаковое деление}\EN{unsigned division}\DE{vorzeichenlose Division}%
\FR{division non signée} \\
\hline \TT{\_\_aullrem} & \RU{остаток от беззнакового деления}\EN{remainder of unsigned division}\DE{Rest (Modulo) einer vorzeichenlosen Division}%
\FR{reste de la division non signée} \\
\hline \TT{\_\_aullshr} & \RU{беззнаковый сдвиг вправо}\EN{unsigned shift right}\DE{Schiebe rechts, vorzeichenlos}%
\FR{décalage non signé à droite} \\
\hline
\end{tabular}
\end{center}

\RU{Процедуры умножения и сдвига влево, одни и те же и для знаковых чисел, и для беззнаковых,
поэтому здесь только одна функция для каждой операции}
\EN{Multiplication and shift left procedures are the same for both signed and unsigned numbers, hence there is only one function 
for each operation here}
\DE{Multiplikation und Links-Schiebebefehle sind sowohl für vorzeichenbehaftete als auch vorzeichenlose Zahlen,
da hier für jede Operation nur ein Befehl existiert}
\FR{La multiplication et le décalage à gauche sont similaire pour les nombres signés
et non signés, donc il n'y a qu'une seule fonction ici}. \\
\\
\RU{Исходные коды этих функций можно найти в установленной \ac{MSVS}, в}\EN{The source code of these function
can be found in the installed \ac{MSVS}, in}%
\DE{Der Quellcode dieser Funktionen kann im Pfad des installierten \ac{MSVS}, gefunden werden: }%
\FR{Le code source des ces fonctions peut être trouvé dans l'installation de \ac{MSVS},
dans} \TT{VC/crt/src/intel/*.asm}.


\input{appendix/cheatsheets}
}
\RU{\part*{\RU{Приложение}\EN{Appendix}\DE{Anhang}\FR{Appendice}\IT{Appendice}}
\appendix
\addcontentsline{toc}{part}{\RU{Приложение}\EN{Appendix}\DE{Anhang}\FR{Appendice}\IT{Appendice}}

% chapters
\EN{\mysection{\MinesweeperWinXPExampleChapterName}
\label{minesweeper_winxp}
\myindex{Windows!Windows XP}

For those who are not very good at playing Minesweeper, we could try to reveal the hidden mines in the debugger.

\myindex{\CStandardLibrary!rand()}
\myindex{Windows!PDB}

As we know, Minesweeper places mines randomly, so there has to be some kind of random number generator or
a call to the standard \TT{rand()} C-function.

What is really cool about reversing Microsoft products is that there are \gls{PDB} 
file with symbols (function names, \etc{}).
When we load \TT{winmine.exe} into \IDA, it downloads the 
\gls{PDB} file exactly for this 
executable and shows all names.

So here it is, the only call to \TT{rand()} is this function:

\lstinputlisting[style=customasmx86]{examples/minesweeper/tmp1.lst}

\IDA named it so, and it was the name given to it by Minesweeper's developers.

The function is very simple:

\begin{lstlisting}[style=customc]
int Rnd(int limit)
{
    return rand() % limit;
};
\end{lstlisting}

(There is no \q{limit} name in the \gls{PDB} file; we manually named this argument like this.)

So it returns 
a random value from 0 to a specified limit.

\TT{Rnd()} is called only from one place, 
a function called \TT{StartGame()}, 
and as it seems, this is exactly 
the code which place the mines:

\begin{lstlisting}[style=customasmx86]
.text:010036C7                 push    _xBoxMac
.text:010036CD                 call    _Rnd@4          ; Rnd(x)
.text:010036D2                 push    _yBoxMac
.text:010036D8                 mov     esi, eax
.text:010036DA                 inc     esi
.text:010036DB                 call    _Rnd@4          ; Rnd(x)
.text:010036E0                 inc     eax
.text:010036E1                 mov     ecx, eax
.text:010036E3                 shl     ecx, 5          ; ECX=ECX*32
.text:010036E6                 test    _rgBlk[ecx+esi], 80h
.text:010036EE                 jnz     short loc_10036C7
.text:010036F0                 shl     eax, 5          ; EAX=EAX*32
.text:010036F3                 lea     eax, _rgBlk[eax+esi]
.text:010036FA                 or      byte ptr [eax], 80h
.text:010036FD                 dec     _cBombStart
.text:01003703                 jnz     short loc_10036C7
\end{lstlisting}

Minesweeper allows you to set the board size, so the X (xBoxMac) and Y (yBoxMac) of the board are global variables.
They are passed to \TT{Rnd()} and random 
coordinates are generated.
A mine is placed by the \TT{OR} instruction at \TT{0x010036FA}. 
And if it has been placed before 
(it's possible if the pair of \TT{Rnd()} 
generates a coordinates pair which has been already 
generated), 
then \TT{TEST} and \TT{JNZ} at \TT{0x010036E6} 
jumps to the generation routine again.

\TT{cBombStart} is the global variable containing total number of mines. So this is loop.

The width of the array is 32 
(we can conclude this by looking at the \TT{SHL} instruction, which multiplies one of the coordinates by 32).

The size of the \TT{rgBlk} 
global array can be easily determined by the difference 
between the \TT{rgBlk} 
label in the data segment and the next known one. 
It is 0x360 (864):

\begin{lstlisting}[style=customasmx86]
.data:01005340 _rgBlk          db 360h dup(?)          ; DATA XREF: MainWndProc(x,x,x,x)+574
.data:01005340                                         ; DisplayBlk(x,x)+23
.data:010056A0 _Preferences    dd ?                    ; DATA XREF: FixMenus()+2
...
\end{lstlisting}

$864/32=27$.

So the array size is $27*32$?
It is close to what we know: when we try to set board size to $100*100$ in Minesweeper settings, it fallbacks to a board of size $24*30$.
So this is the maximal board size here.
And the array has a fixed size for any board size.

So let's see all this in \olly.
We will ran Minesweeper, attaching \olly to it and now we can see the memory dump at the address of the \TT{rgBlk} array (\TT{0x01005340})
\footnote{All addresses here are for Minesweeper for Windows XP SP3 English. 
They may differ for other service packs.}.

So we got this memory dump of the array:

\lstinputlisting[style=customasmx86]{examples/minesweeper/1.lst}

\olly, like any other hexadecimal editor, shows 16 bytes per line.
So each 32-byte array row occupies exactly 2 lines here.

This is beginner level (9*9 board).

There is some square 
structure can be seen visually (0x10 bytes).

We will click \q{Run} in \olly to unfreeze the Minesweeper process, then we'll clicked randomly at the Minesweeper window 
and trapped into mine, but now all mines are visible:

\begin{figure}[H]
\centering
\myincludegraphicsSmall{examples/minesweeper/1.png}
\caption{Mines}
\label{fig:minesweeper1}
\end{figure}

By comparing the mine places and the dump, we can conclude that 0x10 stands for border, 0x0F---empty block, 0x8F---mine.
Perhaps, 0x10 is just a \emph{sentinel value}.

Now we'll add comments and also enclose all 0x8F bytes into square brackets:

\lstinputlisting[style=customasmx86]{examples/minesweeper/2.lst}

Now we'll remove all \emph{border bytes} (0x10) and what's beyond those:

\lstinputlisting[style=customasmx86]{examples/minesweeper/3.lst}

Yes, these are mines, now it can be clearly seen and compared with the screenshot.

\clearpage
What is interesting is that we can modify the array right in \olly.
We can remove all mines by changing all 0x8F bytes by 0x0F, and here is what we'll get in Minesweeper:

\begin{figure}[H]
\centering
\myincludegraphicsSmall{examples/minesweeper/3.png}
\caption{All mines are removed in debugger}
\label{fig:minesweeper3}
\end{figure}

We can also move all of them to the first line: 

\begin{figure}[H]
\centering
\myincludegraphicsSmall{examples/minesweeper/2.png}
\caption{Mines set in debugger}
\label{fig:minesweeper2}
\end{figure}

Well, the debugger is not very convenient for eavesdropping (which is our goal anyway), so we'll write a small utility
to dump the contents of the board:

\lstinputlisting[style=customc]{examples/minesweeper/minesweeper_cheater.c}

Just set the \ac{PID}
\footnote{PID it can be seen in Task Manager 
(enable it in \q{View $\rightarrow$ Select Columns})} 
and the address of the array (\TT{0x01005340} for Windows XP SP3 English) 
and it will dump it
\footnote{The compiled executable is here: 
\href{http://go.yurichev.com/17165}{beginners.re}}.

It attaches itself to a win32 process by \ac{PID} and just reads process memory at the address.

\subsection{Finding grid automatically}

This is kind of nuisance to set address each time when we run our utility.
Also, various Minesweeper versions may have the array on different address.
Knowing the fact that there is always a border (0x10 bytes), we can just find it in memory:

\lstinputlisting[style=customc]{examples/minesweeper/cheater2_fragment.c}

Full source code: \url{https://github.com/DennisYurichev/RE-for-beginners/blob/master/examples/minesweeper/minesweeper_cheater2.c}.

\subsection{\Exercises}

\begin{itemize}

\item 
Why do the \emph{border bytes} (or \emph{sentinel values}) (0x10) exist in the array?

What they are for if they are not visible in Minesweeper's interface?
How could it work without them?

\item 
As it turns out, there are more values possible (for open blocks, for flagged by user, \etc{}).
Try to find the meaning of each one.

\item 
Modify my utility so it can remove all mines or set them in a fixed pattern that you want in the Minesweeper
process currently running.

\end{itemize}
}
\RU{\mysection{Разгон майнера биткоинов Cointerra}
\index{Bitcoin}
\index{BeagleBone}

Был такой майнер биткоинов Cointerra, выглядящий так:

\begin{figure}[H]
\centering
\myincludegraphics{examples/bitcoin_miner/board.jpg}
\caption{Board}
\end{figure}

И была также (возможно утекшая) утилита\footnote{Можно скачать здесь: \url{https://github.com/DennisYurichev/RE-for-beginners/raw/master/examples/bitcoin_miner/files/cointool-overclock}}
которая могла выставлять тактовую частоту платы.
Она запускается на дополнительной плате BeagleBone на ARM с Linux (маленькая плата внизу фотографии).

И у автора (этих строк) однажды спросили, можно ли хакнуть эту утилиту и посмотреть, какие частоты можно выставлять, и какие нет.
И можно ли твикнуть её?

Утилиту нужно запускать так: \TT{./cointool-overclock 0 0 900}, где 900 это частота в МГц.
Если частота слишком большая, утилита выведет ошибку \q{Error with arguments} и закончит работу.

Вот фрагмент кода вокруг ссылки на текстовую строку \q{Error with arguments}:

\begin{lstlisting}[style=customasmARM]

...

.text:0000ABC4         STR      R3, [R11,#var_28]
.text:0000ABC8         MOV      R3, #optind
.text:0000ABD0         LDR      R3, [R3]
.text:0000ABD4         ADD      R3, R3, #1
.text:0000ABD8         MOV      R3, R3,LSL#2
.text:0000ABDC         LDR      R2, [R11,#argv]
.text:0000ABE0         ADD      R3, R2, R3
.text:0000ABE4         LDR      R3, [R3]
.text:0000ABE8         MOV      R0, R3  ; nptr
.text:0000ABEC         MOV      R1, #0  ; endptr
.text:0000ABF0         MOV      R2, #0  ; base
.text:0000ABF4         BL       strtoll
.text:0000ABF8         MOV      R2, R0
.text:0000ABFC         MOV      R3, R1
.text:0000AC00         MOV      R3, R2
.text:0000AC04         STR      R3, [R11,#var_2C]
.text:0000AC08         MOV      R3, #optind
.text:0000AC10         LDR      R3, [R3]
.text:0000AC14         ADD      R3, R3, #2
.text:0000AC18         MOV      R3, R3,LSL#2
.text:0000AC1C         LDR      R2, [R11,#argv]
.text:0000AC20         ADD      R3, R2, R3
.text:0000AC24         LDR      R3, [R3]
.text:0000AC28         MOV      R0, R3  ; nptr
.text:0000AC2C         MOV      R1, #0  ; endptr
.text:0000AC30         MOV      R2, #0  ; base
.text:0000AC34         BL       strtoll
.text:0000AC38         MOV      R2, R0
.text:0000AC3C         MOV      R3, R1
.text:0000AC40         MOV      R3, R2
.text:0000AC44         STR      R3, [R11,#third_argument]
.text:0000AC48         LDR      R3, [R11,#var_28]
.text:0000AC4C         CMP      R3, #0
.text:0000AC50         BLT      errors_with_arguments
.text:0000AC54         LDR      R3, [R11,#var_28]
.text:0000AC58         CMP      R3, #1
.text:0000AC5C         BGT      errors_with_arguments
.text:0000AC60         LDR      R3, [R11,#var_2C]
.text:0000AC64         CMP      R3, #0
.text:0000AC68         BLT      errors_with_arguments
.text:0000AC6C         LDR      R3, [R11,#var_2C]
.text:0000AC70         CMP      R3, #3
.text:0000AC74         BGT      errors_with_arguments
.text:0000AC78         LDR      R3, [R11,#third_argument]
.text:0000AC7C         CMP      R3, #0x31
.text:0000AC80         BLE      errors_with_arguments
.text:0000AC84         LDR      R2, [R11,#third_argument]
.text:0000AC88         MOV      R3, #950
.text:0000AC8C         CMP      R2, R3
.text:0000AC90         BGT      errors_with_arguments
.text:0000AC94         LDR      R2, [R11,#third_argument]
.text:0000AC98         MOV      R3, #0x51EB851F
.text:0000ACA0         SMULL    R1, R3, R3, R2
.text:0000ACA4         MOV      R1, R3,ASR#4
.text:0000ACA8         MOV      R3, R2,ASR#31
.text:0000ACAC         RSB      R3, R3, R1
.text:0000ACB0         MOV      R1, #50
.text:0000ACB4         MUL      R3, R1, R3
.text:0000ACB8         RSB      R3, R3, R2
.text:0000ACBC         CMP      R3, #0
.text:0000ACC0         BEQ      loc_ACEC
.text:0000ACC4
.text:0000ACC4 errors_with_arguments
.text:0000ACC4                                         
.text:0000ACC4         LDR      R3, [R11,#argv]
.text:0000ACC8         LDR      R3, [R3]
.text:0000ACCC         MOV      R0, R3  ; path
.text:0000ACD0         BL       __xpg_basename
.text:0000ACD4         MOV      R3, R0
.text:0000ACD8         MOV      R0, #aSErrorWithArgu ; format
.text:0000ACE0         MOV      R1, R3
.text:0000ACE4         BL       printf
.text:0000ACE8         B        loc_ADD4
.text:0000ACEC ; ------------------------------------------------------------
.text:0000ACEC
.text:0000ACEC loc_ACEC                 ; CODE XREF: main+66C
.text:0000ACEC         LDR      R2, [R11,#third_argument]
.text:0000ACF0         MOV      R3, #499
.text:0000ACF4         CMP      R2, R3
.text:0000ACF8         BGT      loc_AD08
.text:0000ACFC         MOV      R3, #0x64
.text:0000AD00         STR      R3, [R11,#unk_constant]
.text:0000AD04         B        jump_to_write_power
.text:0000AD08 ; ------------------------------------------------------------
.text:0000AD08
.text:0000AD08 loc_AD08                 ; CODE XREF: main+6A4
.text:0000AD08         LDR      R2, [R11,#third_argument]
.text:0000AD0C         MOV      R3, #799
.text:0000AD10         CMP      R2, R3
.text:0000AD14         BGT      loc_AD24
.text:0000AD18         MOV      R3, #0x5F
.text:0000AD1C         STR      R3, [R11,#unk_constant]
.text:0000AD20         B        jump_to_write_power
.text:0000AD24 ; ------------------------------------------------------------
.text:0000AD24
.text:0000AD24 loc_AD24                 ; CODE XREF: main+6C0
.text:0000AD24         LDR      R2, [R11,#third_argument]
.text:0000AD28         MOV      R3, #899
.text:0000AD2C         CMP      R2, R3
.text:0000AD30         BGT      loc_AD40
.text:0000AD34         MOV      R3, #0x5A
.text:0000AD38         STR      R3, [R11,#unk_constant]
.text:0000AD3C         B        jump_to_write_power
.text:0000AD40 ; ------------------------------------------------------------
.text:0000AD40
.text:0000AD40 loc_AD40                 ; CODE XREF: main+6DC
.text:0000AD40         LDR      R2, [R11,#third_argument]
.text:0000AD44         MOV      R3, #999
.text:0000AD48         CMP      R2, R3
.text:0000AD4C         BGT      loc_AD5C
.text:0000AD50         MOV      R3, #0x55
.text:0000AD54         STR      R3, [R11,#unk_constant]
.text:0000AD58         B        jump_to_write_power
.text:0000AD5C ; ------------------------------------------------------------
.text:0000AD5C
.text:0000AD5C loc_AD5C                 ; CODE XREF: main+6F8
.text:0000AD5C         LDR      R2, [R11,#third_argument]
.text:0000AD60         MOV      R3, #1099
.text:0000AD64         CMP      R2, R3
.text:0000AD68         BGT      jump_to_write_power
.text:0000AD6C         MOV      R3, #0x50
.text:0000AD70         STR      R3, [R11,#unk_constant]
.text:0000AD74
.text:0000AD74 jump_to_write_power                     ; CODE XREF: main+6B0
.text:0000AD74                                         ; main+6CC ...
.text:0000AD74         LDR      R3, [R11,#var_28]
.text:0000AD78         UXTB     R1, R3
.text:0000AD7C         LDR      R3, [R11,#var_2C]
.text:0000AD80         UXTB     R2, R3
.text:0000AD84         LDR      R3, [R11,#unk_constant]
.text:0000AD88         UXTB     R3, R3
.text:0000AD8C         LDR      R0, [R11,#third_argument]
.text:0000AD90         UXTH     R0, R0
.text:0000AD94         STR      R0, [SP,#0x44+var_44]
.text:0000AD98         LDR      R0, [R11,#var_24]
.text:0000AD9C         BL       write_power
.text:0000ADA0         LDR      R0, [R11,#var_24]
.text:0000ADA4         MOV      R1, #0x5A
.text:0000ADA8         BL       read_loop
.text:0000ADAC         B        loc_ADD4

...

.rodata:0000B378 aSErrorWithArgu DCB "%s: Error with arguments",0xA,0 ; DATA XREF: main+684

...

\end{lstlisting}

Имена ф-ций присутствовали в отладочной информации в оригинальном исполняемом файле,
такие как \TT{write\_power}, \TT{read\_loop}.
Но имена меткам внутри ф-ции дал я.

\myindex{UNIX!getopt}
\myindex{strtoll()}
Имя \TT{optind} звучит знакомо. Это библиотека \IT{getopt} из *NIX предназначенная для парсинга командной строки ---
и это то, что внутри и происходит.
Затем, третий аргумент (где передается значение частоты) конвертируется из строку в число используя вызов ф-ции \IT{strtoll()}.

Значение затем сравнивается с разными константами.
На 0xACEC есть проверка, меньше ли оно или равно 499, и если это так, то 0x64 будет передано в ф-цию
\TT{write\_power()} (которая посылает команду через USB используя \TT{send\_msg()}).
Если значение больше 499, происходит переход на 0xAD08.

На 0xAD08 есть проверка, меньше ли оно или равно 799. Если это так, то 0x5F передается в ф-цию \TT{write\_power()}.

Есть еще проверки: на 899 на 0xAD24, на 0x999 на 0xAD40, и наконец, на 1099 на 0xAD5C.
Если входная частота меньше или равна 1099, 0x50 (на 0xAD6C) будет передано в ф-цию \TT{write\_power()}.
И тут что-то вроде баги.
Если значение все еще больше 1099, само значение будет передано в ф-цию \TT{write\_power()}.
Но с другой стороны это не бага, потому что мы не можем попасть сюда: значение в начале проверяется с 950 на 0xAC88,
и если оно больше, выводится сообщение об ошибке и утилита заканчивает работу.

Вот таблица между частотами в МГц и значениями передаваемыми в ф-цию \TT{write\_power()}:

\begin{center}
\begin{longtable}{ | l | l | l | }
\hline
\HeaderColor МГц & \HeaderColor шестнадцатеричное представление & \HeaderColor десятичное \\
\hline
499MHz & 0x64 & 100 \\
\hline
799MHz & 0x5f & 95 \\
\hline
899MHz & 0x5a & 90 \\
\hline
999MHz & 0x55 & 85 \\
\hline
1099MHz & 0x50 & 80 \\
\hline
\end{longtable}
\end{center}

Как видно, значение передаваемое в плату постепенно уменьшается с ростом частоты.

Видно что значение в 950МГц это жесткий предел, по крайней мере в этой утилите. Можно ли её обмануть?

Вернемся к этому фрагменту кода:

\begin{lstlisting}[style=customasmARM]
.text:0000AC84      LDR     R2, [R11,#third_argument]
.text:0000AC88      MOV     R3, #950
.text:0000AC8C      CMP     R2, R3
.text:0000AC90      BGT     errors_with_arguments ; Я пропатчил здесь на 00 00 00 00
\end{lstlisting}

Нам нужно как-то запретить инструкцию перехода \INS{BGT} на 0xAC90. И это ARM в режиме ARM, потому что, как мы видим,
все адреса увеличиваются на 4, т.е., длина каждой инструкции это 4 байта.
Инструкция \TT{NOP} (нет операции) в режиме ARM это просто 4 нулевых байта: \TT{00 00 00 00}.
Так что, записывая 4 нуля по адресу 0xAC90 (или по физическому смещению в файле: 0x2C90) мы можем выключить
эту проверку.

Теперь можно выставлять частоты вплоть до 1050МГц. И даже больше, но из-за ошибки, если входное значение больше 1099,
значение в МГц, \IT{как есть}, будет передано в плату, что неправильно.

Дальше я не разбирался, но если бы продолжил, я бы уменьшал значение передаваемое в ф-цию \TT{write\_power()}.

Теперь страшный фрагмент кода, который я в начале пропустил:

\lstinputlisting[style=customasmARM]{examples/bitcoin_miner/tmp1.lst}

Здесь используется деление через умножение, и константа 0x51EB851F.
Я написал для себя простой программистский калькулятор\footnote{\url{https://github.com/DennisYurichev/progcalc}}.
И там есть возможность вычислять обратное число по модулю.

\begin{lstlisting}
modinv32(0x51EB851F)
Warning, result is not integer: 3.125000
(unsigned) dec: 3 hex: 0x3 bin: 11
\end{lstlisting}

Это значит что инструкция \INS{SMULL} на 0xACA0 просто делит 3-й аргумент на 3.125.
На самом деле, все что делает ф-ция \TT{modinv32()} в моем калькуляторе, это:

\[
\frac{1}{\frac{input}{2^{32}}} = \frac{2^{32}}{input}
\]

Потом там есть дополнительные сдвиги и теперь мы видим что 3-й аргумент просто делится на 50.
И затем умножается снова на 50.
Зачем?
Это простейшая проверка, можно ли делить входное значение на 50 без остатка.
Если значение этого выражения ненулевое, $x$ не может быть разделено на 50 без остатка:

\[
x-((\frac{x}{50}) \cdot 50)
\]

На самом деле, это простой способ вычисления остатка от деления.

И затем, если остаток ненулевой, выводится сообщение об ошибке.
Так что эта утилита берет значения частотв вроде 850, 900, 950, 1000, итд, но не 855 или 911.

Вот и всё! Если вы делаете что-то такое, имейте ввиду, что это может испортить вашу плату, как и в случае разгона
чипов вроде \ac{CPU}, \ac{GPU}, итд.
Если у вас есть плата Cointerra, делайте всё это на свой собственный риск!

}
\DE{\mysection{\Stack}
\label{sec:stack}
\myindex{\Stack}

Der Stack ist eine der fundamentalen Datenstrukturen in der Informatik.
\footnote{\href{http://go.yurichev.com/17119}{wikipedia.org/wiki/Call\_Stack}}.
\ac{AKA} \ac{LIFO}.

Technisch betrachtet ist es ein Stapelspeicher innerhalb des Prozessspeichers der zusammen mit den \ESP (x86), \RSP (x64) oder dem \ac{SP} (ARM) Register als ein Zeiger in diesem Speicherblock fungiert.

\myindex{ARM!\Instructions!PUSH}
\myindex{ARM!\Instructions!POP}
\myindex{x86!\Instructions!PUSH}
\myindex{x86!\Instructions!POP}

Die häufigsten Stack-Zugriffsinstruktionen sind die \PUSH- und \POP-Instruktionen (in beidem x86 und ARM Thumb-Modus). \PUSH subtrahiert vom \ESP/\RSP/\ac{SP} 4 Byte im 32-Bit Modus (oder 8 im 64-Bit Modus) und schreibt dann den Inhalt des Zeigers an die Adresse auf die von \ESP/\RSP/\ac{SP} gezeigt wird.

\POP ist die umgekehrte Operation: Die Daten des Zeigers für die Speicherregion auf die von \ac{SP}
gezeigt wird werden ausgelesen und die Inhalte in den Instruktionsoperanden geschreiben (oft ist das ein Register). Dann werden 4 (beziehungsweise 8) Byte zum \gls{stack pointer} addiert.

Nach der Stackallokation, zeigt der \gls{stack pointer} auf den Boden des Stacks.
\PUSH verringert den \gls{stack pointer} und \POP erhöht ihn.
Der Boden des Stacks ist eigentlich der Anfang der Speicherregion die für den Stack reserviert wurde.
Das wirkt zunächst seltsam, aber so funktioniert es.

ARM unterstützt beides, aufsteigende und absteigende Stacks.

\myindex{ARM!\Instructions!STMFD}
\myindex{ARM!\Instructions!LDMFD}
\myindex{ARM!\Instructions!STMED}
\myindex{ARM!\Instructions!LDMED}
\myindex{ARM!\Instructions!STMFA}
\myindex{ARM!\Instructions!LDMFA}
\myindex{ARM!\Instructions!STMEA}
\myindex{ARM!\Instructions!LDMEA}

Zum Beispiel die \ac{STMFD}/\ac{LDMFD} und \ac{STMED}/\ac{LDMED} Instruktionen sind alle dafür gedacht mit einem absteigendem Stack zu arbeiten ( wächst nach unten, fängt mit hohen Adressen an und entwickelt sich zu niedrigeren Adressen). Die \ac{STMFA}/\ac{LDMFA} und \ac{STMEA}/\ac{LDMEA} Instruktionen sind dazu gedacht mit einem aufsteigendem Stack zu arbeiten (wächst nach oben und fängt mit niedrigeren Adressen an und wächst nach oben).

% It might be worth mentioning that STMED and STMEA write first,
% and then move the pointer, and that LDMED and LDMEA move the pointer first, and then read.
% In other words, ARM not only lets the stack grow in a non-standard direction,
% but also in a non-standard order.
% Maybe this can be in the glossary, which would explain why E stands for "empty".

\subsection{Warum wächst der Stack nach unten?}
\label{stack_grow_backwards}

Intuitiv, würden man annehmen das der Stack nach oben wächst z.B Richtung höherer Adressen, so wie bei jeder anderen Datenstruktur.

Der Grund das der Stack rückwärts wächst ist wohl historisch bedingt. Als Computer so groß waren das sie einen ganzen Raum beansprucht haben war es einfach Speicher in zwei Sektionen zu unterteilen, einen Teil für den \gls{heap} und einen Teil für den Stack. Sicher war zu dieser Zeit nicht bekannt wie groß der \gls{heap} und der Stack wachsen würden, während der Programm Laufzeit, also war die Lösung die einfachste mögliche.

\input{patterns/02_stack/stack_and_heap}

In \RitchieThompsonUNIX können wir folgendes lesen:

\begin{framed}
\begin{quotation}
Der user-core eines Programm Images wird in drei logische Segmente unterteilt. Das Programm-Text Segment beginnt bei 0 im virtuellen Adress Speicher. Während der Ausführung wird das Segment als schreibgeschützt markiert und eine einzelne Kopie des Segments wird unter allen Prozessen geteilt die das Programm ausführen. An der ersten 8K grenze über dem Programm Text Segment im Virtuellen Speicher, fängt der ``nonshared'' Bereich an, der nach Bedarf von Syscalls erweitert werden kann. Beginnend bei der höchsten Adresse im Virtuellen Speicher ist das Stack Segment, das Automatisch nach unten wächst während der Hardware Stackpointer sich ändert.
\end{quotation}
\end{framed}

Das erinnert daran wie manche Schüler Notizen zu  zwei Vorträgen in einem Notebook dokumentieren:
Notizen für den ersten Vortrag werden normal notiert, und Notizen zur zum zweiten Vortrag werden 
ans Ende des Notizbuches geschrieben, indem man das Notizbuch umdreht. Die Notizen treffen sich irgendwann
im Notizbuch aufgrund des fehlenden Freien Platzes.

% I think if we want to expand on this analogy,
% one might remember that the line number increases as as you go down a page.
% So when you decrease the address when pushing to the stack, visually,
% the stack does grow upwards.
% Of course, the problem is that in most human languages,
% just as with computers,
% we write downwards, so this direction is what makes buffer overflows so messy.

\subsection{Für was wird der Stack benutzt?}

% subsections
\input{patterns/02_stack/01_saving_ret_addr}
\input{patterns/02_stack/02_args_passing}
\EN{\input{patterns/02_stack/03_local_vars_EN}}
\RU{\input{patterns/02_stack/03_local_vars_RU}}
\DE{\input{patterns/02_stack/03_local_vars_DE}}
\PTBR{\input{patterns/02_stack/03_local_vars_PTBR}}
\input{patterns/02_stack/04_alloca/main}
\input{patterns/02_stack/05_SEH}
\input{patterns/02_stack/06_BO_protection}

\subsubsection{Automatisches deallokieren der Daten auf dem Stack}

Vielleicht ist der Grund warum man lokale Variablen und SEH Einträge auf dem Stack speichert, weil sie beim 
verlassen der Funktion automatisch aufgeräumt werden. Man braucht dabei nur eine Instruktion um die Position
des Stackpointers zu korrigieren (oftmals ist es die \ADD Instruktion). Funktions Argumente, könnte man sagen 
werden auch am Ende der Funktion deallokiert. Im Kontrast dazu, alles was auf dem \emph{heap} gespeichert wird muss
explizit deallokiert werden. 

% sections
\EN{\input{patterns/02_stack/07_layout_EN}}
\RU{\input{patterns/02_stack/07_layout_RU}}
\DE{\input{patterns/02_stack/07_layout_DE}}
\PTBR{\input{patterns/02_stack/07_layout_PTBR}}
\input{patterns/02_stack/08_noise/main}
\input{patterns/02_stack/exercises}
}
\FR{\mysection{x86}

\subsection{Terminologie}

Commun en 16-bit (8086/80286), 32-bit (80386, etc.), 64-bit.

\myindex{IEEE 754}
\myindex{MS-DOS}
\begin{description}
	\item[octet] 8-bit.
		La directive d'assembleur DB est utilisée pour définir les variables et les tableaux d'octets.
		Les octets sont passés dans les parties 8-bit des registres: \TT{AL/BL/CL/DL/AH/BH/CH/DH/SIL/DIL/R*L}.
	\item[mot] 16-bit. 
		directive assembleur DW \dittoclosing.
		Les mots sont passés dans la partie 16-bit des registres:\\
			\TT{AX/BX/CX/DX/SI/DI/R*W}.
	\item[double mot] (\q{dword}) 32-bit.
		directive assembleur DD \dittoclosing.
		Les double mots sont passés dans les registres (x86) ou dans la partie 32-bit de registres (x64).
		Dans du code 16-bit, les double mots sont passés dans une paire de registres 16-bit.
	\item[quadruple mot] (\q{qword}) 64-bit.
		directive assembleur DQ \dittoclosing.
		En environnement 32-bit, les quadruple mots sont passés dans une paire de registres 32-bit.
	\item[tbyte] (10 bytes) 80-bit ou 10 octets (utilisé pour les registres FPU IEEE 754).
	\item[paragraph] (16 bytes)---le terme était répandu dans l'environnement MS-DOS.
\end{description}

\myindex{Windows!API}

Des types de données de même taille (BYTE, WORD, DWORD) existent aussi dans l'\ac{API} Windows.

\input{appendix/x86/registers_FR} % subsection
\input{appendix/x86/instructions_FR} % subsection
\input{appendix/x86/npad} % subsection

}
\IT{\mysection{x86}

\subsection{Terminologia}

Comune per 16-bit (8086/80286), 32-bit (80386, etc.), 64-bit.

\myindex{IEEE 754}
\myindex{MS-DOS}
\begin{description}
	\item[byte] 8-bit.
		La direttiva DB in assembly è utilizzata per definire variabili ed array di byte.
		I byte sono passati nella parte ad 8-bit dei registri: \TT{AL/BL/CL/DL/AH/BH/CH/DH/SIL/DIL/R*L}.
	\item[word] 16-bit.
		La direttiva DW in assembly \dittoclosing.
		Le Word sono passate nella parte a 16-bit dei registri:\\
			\TT{AX/BX/CX/DX/SI/DI/R*W}.
	\item[double word] (\q{dword}) 32-bit.
		La direttiva DD in assembly \dittoclosing.
		Le Double words sono passate nei registri (x86) o nelle parti a 32-bit dei registri (x64).
		Nel codice a 16-bit, le double words sono passate in coppie di registri a 16-bit.
	\item[quad word] (\q{qword}) 64-bit.
		La direttiva DQ in assembly \dittoclosing.
		Negli ambienti a 32-bit, le quad words sono passate in coppie di registri a 32-bit.
	\item[tbyte] (10 bytes) 80-bit o 10 bytes (usati per i registri FPU IEEE 754).
	\item[paragraph] (16 bytes)---termine popolare nell'ambiente MS-DOS. % TODO link to a part about 8086 memory model...
\end{description}

\myindex{Windows!API}

Tipi di dati della stessa dimensione (BYTE, WORD, DWORD) sono gli stessi nelle Windows \ac{API}.

%TBT
\input{appendix/x86/registers} % subsection
\input{appendix/x86/instructions} % subsection
\input{appendix/x86/npad} % subsection
}
\mysection{ARM}
\myindex{ARM}

\subsection{\RU{Терминология}\EN{Terminology}}

ARM \RU{изначально разрабатывался как 32-битный}\EN{was initially developed as 32-bit} \ac{CPU}, 
\RU{поэтому \emph{слово} здесь, в отличие от x86, 32-битное}\EN{so that's why a \emph{word} here, unlike x86, is 32-bit}.

\begin{description}
	\item[byte] 8-\bitENRU.
		\RU{Для определения переменных и массива байт используется директива ассемблера DCB}
		\EN{The DB assembly directive is used for defining variables and arrays of bytes}.
	\item[halfword] 16-\bitENRU. \RU{\dittoclosing директива ассемблера DCW}
					\EN{DCW assembly directive \dittoclosing}.	
	\item[word] 32-\bitENRU. \RU{\dittoclosing директива ассемблера DCD}
					\EN{DCD assembly directive \dittoclosing}.
	\item[doubleword] 64-\bitENRU.
	\item[quadword] 128-\bitENRU.
\end{description}

\subsection{\RU{Версии}\EN{Versions}}

\begin{itemize}
\item ARMv4: \RU{появился режим Thumb}\EN{Thumb mode introduced}.

\item ARMv6: \RU{использовался в}\EN{used in} iPhone 1st gen., iPhone 3G 
(Samsung 32-bit RISC ARM 1176JZ(F)-S \RU{поддерживающий}\EN{that supports} Thumb-2)

\item ARMv7: \RU{появился }Thumb-2\EN{ was added} (2003).
\RU{Использовался в}\EN{was used in} iPhone 3GS, iPhone 4, iPad 1st gen. (ARM Cortex-A8), iPad 2 (Cortex-A9),
iPad 3rd gen.

\item ARMv7s: \RU{Добавлены новые инструкции}\EN{New instructions added}.
\RU{Использовался в}\EN{Was used in} iPhone 5, iPhone 5c, iPad 4th gen. (Apple A6).

\item ARMv8: 64-\RU{битный процессор}\EN{bit CPU}, \ac{AKA} ARM64 \ac{AKA} AArch64.
\RU{Использовался в}\EN{Was used in} iPhone 5S, iPad Air (Apple A7).
\RU{В 64-битном режиме, режима Thumb больше нет, только режим ARM (4-байтные инструкции).}
\EN{There is no Thumb mode in 64-bit mode, only ARM (4-byte instructions).}
\end{itemize}

% sections
\input{appendix/ARM/registers}
\input{appendix/ARM/instructions}

\mysection{MIPS}

\subsection{\Registers}
\label{MIPS_registers_ref}

\myindex{MIPS!O32}
( \RU{Соглашение о вызовах O32}\EN{O32 calling convention} )

\subsubsection{\RU{Регистры общего пользования}\EN{General purpose registers} \ac{GPR}}

%\small
\begin{center}
\begin{tabular}{ | l | l | l | }
\hline
\HeaderColor \RU{Номер}\EN{Number} & 
\HeaderColor \RU{Псевдоимя}\EN{Pseudoname} & 
\HeaderColor \RU{Описание}\EN{Description} \\
\hline
\$0             & \$ZERO          & \RU{Всегда ноль. Запись в этот регистр это как}\EN{Always zero. Writing to this register is like} \ac{NOP}. \\
\hline
\$1             & \$AT            & \RU{Используется как временный регистр}\EN{Used as a temporary register} \\
                &                 & \RU{для ассемблерных макросов и псевдоинструкций}\EN{for assembly macros and pseudo instructions}. \\
\hline
\$2 \dots \$3   & \$V0 \dots \$V1 & \RU{Здесь возвращается результат функции}\EN{Function result is returned here}. \\
\hline
\$4 \dots \$7   & \$A0 \dots \$A3 & \RU{Аргументы функции}\EN{Function arguments}. \\
\hline
\$8 \dots \$15  & \$T0 \dots \$T7 & \RU{Используется для временных данных}\EN{Used for temporary data}. \\
\hline
\$16 \dots \$23 & \$S0 \dots \$S7 & \RU{Используется для временных данных}\EN{Used for temporary data}\AsteriskOne{}. \\
\hline
\$24 \dots \$25 & \$T8 \dots \$T9 & \RU{Используется для временных данных}\EN{Used for temporary data}. \\
\hline
\$26 \dots \$27 & \$K0 \dots \$K1 & \RU{Зарезервировано для ядра \ac{OS}}\EN{Reserved for \ac{OS} kernel}. \\
\hline
\$28            & \$GP            & \RU{Глобальный указатель}\EN{Global Pointer}\AsteriskTwo{}. \\
\hline
\$29            & \$SP            & \ac{SP}\AsteriskOne{}. \\
\hline
\$30            & \$FP            & \ac{FP}\AsteriskOne{}. \\
\hline
\$31            & \$RA            & \ac{RA}. \\
\hline
n/a             & PC              & \ac{PC}. \\
\hline
n/a             & HI              & \RU{старшие 32 бита результата умножения или остаток от деления}\EN{high 32 bit of multiplication or division remainder}\AsteriskThree{}. \\
\hline
n/a             & LO              & \RU{младшие 32 бита результата умножения или результат деления}\EN{low 32 bit of multiplication and division remainder}\AsteriskThree{}. \\
\hline
\end{tabular}
\end{center}
%\normalsize

\subsubsection{\RU{Регистры для работы с числами с плавающей точкой}\EN{Floating-point registers}}
\label{MIPS_FPU_registers}

\begin{center}
\begin{tabular}{ | l | l | l | }
\hline
\HeaderColor \RU{Название}\EN{Name} & \HeaderColor \RU{Описание}\EN{Description} \\
\hline
\$F0..\$F1   & \RU{Здесь возвращается результат функции}\EN{Function result returned here}. \\
\hline
\$F2..\$F3   & \RU{Не используется}\EN{Not used}. \\
\hline
\$F4..\$F11  & \RU{Используется для временных данных}\EN{Used for temporary data}. \\
\hline
\$F12..\$F15 & \RU{Первые два аргумента функции}\EN{First two function arguments}. \\
\hline
\$F16..\$F19 & \RU{Используется для временных данных}\EN{Used for temporary data}. \\
\hline
\$F20..\$F31 & \RU{Используется для временных данных}\EN{Used for temporary data}\AsteriskOne{}. \\
\hline
%fcr31 & Control/status register. \\
%\hline
\end{tabular}
\end{center}

\AsteriskOne{} --- \Gls{callee} \RU{должен сохранять}\EN{must preserve the value}.\\
\AsteriskTwo{} --- \Gls{callee} \RU{должен сохранять}\EN{must preserve the value} (\RU{кроме \ac{PIC}-кода}
\EN{except in \ac{PIC} code}).\\
\myindex{MIPS!\Instructions!MFLO}
\myindex{MIPS!\Instructions!MFHI}
\AsteriskThree{} --- \RU{доступны используя инструкции}\EN{accessible using the} \TT{MFHI} \AndENRU \TT{MFLO}\EN{ instructions}.\\

\subsection{\Instructions}

\RU{Есть три типа инструкций}\EN{There are 3 kinds of instructions}:

\begin{itemize}

\item \RU{Тип R: имеющие 3 регистра. R-инструкции обычно имеют такой вид:}
\EN{R-type: those which have 3 registers. R-instruction usually have the following form:}

\begin{lstlisting}
instruction destination, source1, source2
\end{lstlisting}

\RU{Важно помнить, что если первый и второй регистр один и тот же, IDA может показать инструкцию в сокращенной
форме:}
\EN{One important thing to keep in mind is that when the first and second register are the same, 
IDA may show the instruction in its shorter form:}

\begin{lstlisting}
instruction destination/source1, source2
\end{lstlisting}

\RU{Это немного напоминает Интеловский синтаксис ассемблера x86.}
\EN{That somewhat reminds us of the Intel syntax for x86 assembly language.}

\item \RU{Тип I: имеющие 2 регистра и 16-битное \q{immediate}-значение.}
\EN{I-type: those which have 2 registers and a 16-bit immediate value.}

\item \RU{Тип J: инструкции перехода, имеют 26 бит для кодирования смещения.}
\EN{J-type: jump/branch instructions, have 26 bits for encoding the offset.}

\end{itemize}

\subsubsection{\RU{Инструкции перехода}\EN{Jump instructions}}

\RU{Какая разница между инструкциями начинающихся с B- (BEQ, B, итд) и с J- (JAL, JALR, итд)?}
\EN{What is the difference between B- instructions (BEQ, B, etc.) and J- ones (JAL, JALR, etc.)?}

\RU{B-инструкции имеют тип I, так что, смещение в этих инструкциях кодируется как 16-битное значение.}
\EN{The B-instructions have an I-type, hence, the B-instructions' offset is encoded as a 16-bit immediate.}
\RU{Инструкции JR и JALR имеют тип R, и они делают переход по абсолютному адресу указанному в регистре.}
\EN{JR and JALR are R-type and jump to an absolute address specified in a register.}
\RU{J и JAL имеют тип J, так что смещение кодируется как 26-битное значение.}
\EN{J and JAL are J-type, hence the offset is encoded as a 26-bit immediate.}

\RU{Коротко говоря, в B-инструкциях можно кодировать условие 
(B на самом деле это псевдоинструкция для \TT{BEQ \$ZERO, \$ZERO, LABEL}),
а в J-инструкциях нельзя.}
\EN{In short, B-instructions can encode a condition 
(B is in fact pseudo instruction for \TT{BEQ \$ZERO, \$ZERO, LABEL}), 
while J-instructions can't.}

\mysection{\RU{Некоторые библиотечные функции GCC}\EN{Some GCC library functions}\DE{Einige GCC-Bibliotheks-Funktionen}%
\FR{Quelques fonctions de la bibliothèque de GCC}}
\myindex{GCC}
\label{sec:GCC_library_func}

%__ashldi3
%__ashrdi3
%__floatundidf
%__floatdisf
%__floatdixf
%__floatundidf
%__floatundisf
%__floatundixf
%__lshrdi3
%__muldi3

\begin{center}
\begin{tabular}{ | l | l | }
\hline
\HeaderColor \RU{имя}\EN{name}\DE{Name}\FR{nom} & \HeaderColor \RU{значение}\EN{meaning}\DE{Bedeutung}%
\FR{signification} \\
\hline \TT{\_\_divdi3} & \RU{знаковое деление}\EN{signed division}\DE{vorzeichenbehaftete Division}%
\FR{division signée} \\
\hline \TT{\_\_moddi3} & \RU{остаток от знакового деления}\EN{getting remainder (modulo) of signed division}%
\DE{Rest (Modulo) einer vorzeichenbehafteten Division}%
\FR{reste (modulo) d'une division signée} \\
\hline \TT{\_\_udivdi3} & \RU{беззнаковое деление}\EN{unsigned division}\DE{vorzeichenlose Division}%
\FR{division non signée} \\
\hline \TT{\_\_umoddi3} & \RU{остаток от беззнакового деления}\EN{getting remainder (modulo) of unsigned division}%
\DE{Rest (Modulo) einer vorzeichenlosen Division}%
\FR{reste (modulo) d'une division non signée} \\
\hline
\end{tabular}
\end{center}



\mysection{\RU{Некоторые библиотечные функции MSVC}\EN{Some MSVC library functions}\DE{Einige MSVC-Bibliotheks-Funktionen}%
\FR{Quelques fonctions de la bibliothèque MSVC}}
\myindex{MSVC}
\label{sec:MSVC_library_func}

\TT{ll} \RU{в имени функции означает}\EN{in function name stands for}\DE{in Funktionsnamen steht für}%
\FR{dans une fontion signifie} \q{long long}, \RU{т.е. 64-битный тип данных}
\EN{e.g., a 64-bit data type}\DE{z.B. einen 64-Bit-Datentyp}\FR{i.e., type de donées 64-bit}.

\begin{center}
\begin{tabular}{ | l | l | }
\hline
\HeaderColor \RU{имя}\EN{name}\DE{Name}\FR{nom} & \HeaderColor \RU{значение}\EN{meaning}\DE{Bedeutung}\FR{signification} \\
\hline \TT{\_\_alldiv} & \RU{знаковое деление}\EN{signed division}\DE{vorzeichenbehaftete Division}\FR{division signée} \\
\hline \TT{\_\_allmul} & \RU{умножение}\EN{multiplication}\DE{Multiplikation}\FR{multiplication} \\
\hline \TT{\_\_allrem} & \RU{остаток от знакового деления}\EN{remainder of signed division}\DE{Rest einer vorzeichenbehafteten Division}%
\FR{reste de la division signée} \\
\hline \TT{\_\_allshl} & \RU{сдвиг влево}\EN{shift left}\DE{Schiebe links}\FR{décalage à gauche} \\
\hline \TT{\_\_allshr} & \RU{знаковый сдвиг вправо}\EN{signed shift right}\DE{Schiebe links, vorzeichenbehaftet}%
\FR{décalage signé à droite} \\
\hline \TT{\_\_aulldiv} & \RU{беззнаковое деление}\EN{unsigned division}\DE{vorzeichenlose Division}%
\FR{division non signée} \\
\hline \TT{\_\_aullrem} & \RU{остаток от беззнакового деления}\EN{remainder of unsigned division}\DE{Rest (Modulo) einer vorzeichenlosen Division}%
\FR{reste de la division non signée} \\
\hline \TT{\_\_aullshr} & \RU{беззнаковый сдвиг вправо}\EN{unsigned shift right}\DE{Schiebe rechts, vorzeichenlos}%
\FR{décalage non signé à droite} \\
\hline
\end{tabular}
\end{center}

\RU{Процедуры умножения и сдвига влево, одни и те же и для знаковых чисел, и для беззнаковых,
поэтому здесь только одна функция для каждой операции}
\EN{Multiplication and shift left procedures are the same for both signed and unsigned numbers, hence there is only one function 
for each operation here}
\DE{Multiplikation und Links-Schiebebefehle sind sowohl für vorzeichenbehaftete als auch vorzeichenlose Zahlen,
da hier für jede Operation nur ein Befehl existiert}
\FR{La multiplication et le décalage à gauche sont similaire pour les nombres signés
et non signés, donc il n'y a qu'une seule fonction ici}. \\
\\
\RU{Исходные коды этих функций можно найти в установленной \ac{MSVS}, в}\EN{The source code of these function
can be found in the installed \ac{MSVS}, in}%
\DE{Der Quellcode dieser Funktionen kann im Pfad des installierten \ac{MSVS}, gefunden werden: }%
\FR{Le code source des ces fonctions peut être trouvé dans l'installation de \ac{MSVS},
dans} \TT{VC/crt/src/intel/*.asm}.


\input{appendix/cheatsheets}
}
\DE{\part*{\RU{Приложение}\EN{Appendix}\DE{Anhang}\FR{Appendice}\IT{Appendice}}
\appendix
\addcontentsline{toc}{part}{\RU{Приложение}\EN{Appendix}\DE{Anhang}\FR{Appendice}\IT{Appendice}}

% chapters
\EN{\mysection{\MinesweeperWinXPExampleChapterName}
\label{minesweeper_winxp}
\myindex{Windows!Windows XP}

For those who are not very good at playing Minesweeper, we could try to reveal the hidden mines in the debugger.

\myindex{\CStandardLibrary!rand()}
\myindex{Windows!PDB}

As we know, Minesweeper places mines randomly, so there has to be some kind of random number generator or
a call to the standard \TT{rand()} C-function.

What is really cool about reversing Microsoft products is that there are \gls{PDB} 
file with symbols (function names, \etc{}).
When we load \TT{winmine.exe} into \IDA, it downloads the 
\gls{PDB} file exactly for this 
executable and shows all names.

So here it is, the only call to \TT{rand()} is this function:

\lstinputlisting[style=customasmx86]{examples/minesweeper/tmp1.lst}

\IDA named it so, and it was the name given to it by Minesweeper's developers.

The function is very simple:

\begin{lstlisting}[style=customc]
int Rnd(int limit)
{
    return rand() % limit;
};
\end{lstlisting}

(There is no \q{limit} name in the \gls{PDB} file; we manually named this argument like this.)

So it returns 
a random value from 0 to a specified limit.

\TT{Rnd()} is called only from one place, 
a function called \TT{StartGame()}, 
and as it seems, this is exactly 
the code which place the mines:

\begin{lstlisting}[style=customasmx86]
.text:010036C7                 push    _xBoxMac
.text:010036CD                 call    _Rnd@4          ; Rnd(x)
.text:010036D2                 push    _yBoxMac
.text:010036D8                 mov     esi, eax
.text:010036DA                 inc     esi
.text:010036DB                 call    _Rnd@4          ; Rnd(x)
.text:010036E0                 inc     eax
.text:010036E1                 mov     ecx, eax
.text:010036E3                 shl     ecx, 5          ; ECX=ECX*32
.text:010036E6                 test    _rgBlk[ecx+esi], 80h
.text:010036EE                 jnz     short loc_10036C7
.text:010036F0                 shl     eax, 5          ; EAX=EAX*32
.text:010036F3                 lea     eax, _rgBlk[eax+esi]
.text:010036FA                 or      byte ptr [eax], 80h
.text:010036FD                 dec     _cBombStart
.text:01003703                 jnz     short loc_10036C7
\end{lstlisting}

Minesweeper allows you to set the board size, so the X (xBoxMac) and Y (yBoxMac) of the board are global variables.
They are passed to \TT{Rnd()} and random 
coordinates are generated.
A mine is placed by the \TT{OR} instruction at \TT{0x010036FA}. 
And if it has been placed before 
(it's possible if the pair of \TT{Rnd()} 
generates a coordinates pair which has been already 
generated), 
then \TT{TEST} and \TT{JNZ} at \TT{0x010036E6} 
jumps to the generation routine again.

\TT{cBombStart} is the global variable containing total number of mines. So this is loop.

The width of the array is 32 
(we can conclude this by looking at the \TT{SHL} instruction, which multiplies one of the coordinates by 32).

The size of the \TT{rgBlk} 
global array can be easily determined by the difference 
between the \TT{rgBlk} 
label in the data segment and the next known one. 
It is 0x360 (864):

\begin{lstlisting}[style=customasmx86]
.data:01005340 _rgBlk          db 360h dup(?)          ; DATA XREF: MainWndProc(x,x,x,x)+574
.data:01005340                                         ; DisplayBlk(x,x)+23
.data:010056A0 _Preferences    dd ?                    ; DATA XREF: FixMenus()+2
...
\end{lstlisting}

$864/32=27$.

So the array size is $27*32$?
It is close to what we know: when we try to set board size to $100*100$ in Minesweeper settings, it fallbacks to a board of size $24*30$.
So this is the maximal board size here.
And the array has a fixed size for any board size.

So let's see all this in \olly.
We will ran Minesweeper, attaching \olly to it and now we can see the memory dump at the address of the \TT{rgBlk} array (\TT{0x01005340})
\footnote{All addresses here are for Minesweeper for Windows XP SP3 English. 
They may differ for other service packs.}.

So we got this memory dump of the array:

\lstinputlisting[style=customasmx86]{examples/minesweeper/1.lst}

\olly, like any other hexadecimal editor, shows 16 bytes per line.
So each 32-byte array row occupies exactly 2 lines here.

This is beginner level (9*9 board).

There is some square 
structure can be seen visually (0x10 bytes).

We will click \q{Run} in \olly to unfreeze the Minesweeper process, then we'll clicked randomly at the Minesweeper window 
and trapped into mine, but now all mines are visible:

\begin{figure}[H]
\centering
\myincludegraphicsSmall{examples/minesweeper/1.png}
\caption{Mines}
\label{fig:minesweeper1}
\end{figure}

By comparing the mine places and the dump, we can conclude that 0x10 stands for border, 0x0F---empty block, 0x8F---mine.
Perhaps, 0x10 is just a \emph{sentinel value}.

Now we'll add comments and also enclose all 0x8F bytes into square brackets:

\lstinputlisting[style=customasmx86]{examples/minesweeper/2.lst}

Now we'll remove all \emph{border bytes} (0x10) and what's beyond those:

\lstinputlisting[style=customasmx86]{examples/minesweeper/3.lst}

Yes, these are mines, now it can be clearly seen and compared with the screenshot.

\clearpage
What is interesting is that we can modify the array right in \olly.
We can remove all mines by changing all 0x8F bytes by 0x0F, and here is what we'll get in Minesweeper:

\begin{figure}[H]
\centering
\myincludegraphicsSmall{examples/minesweeper/3.png}
\caption{All mines are removed in debugger}
\label{fig:minesweeper3}
\end{figure}

We can also move all of them to the first line: 

\begin{figure}[H]
\centering
\myincludegraphicsSmall{examples/minesweeper/2.png}
\caption{Mines set in debugger}
\label{fig:minesweeper2}
\end{figure}

Well, the debugger is not very convenient for eavesdropping (which is our goal anyway), so we'll write a small utility
to dump the contents of the board:

\lstinputlisting[style=customc]{examples/minesweeper/minesweeper_cheater.c}

Just set the \ac{PID}
\footnote{PID it can be seen in Task Manager 
(enable it in \q{View $\rightarrow$ Select Columns})} 
and the address of the array (\TT{0x01005340} for Windows XP SP3 English) 
and it will dump it
\footnote{The compiled executable is here: 
\href{http://go.yurichev.com/17165}{beginners.re}}.

It attaches itself to a win32 process by \ac{PID} and just reads process memory at the address.

\subsection{Finding grid automatically}

This is kind of nuisance to set address each time when we run our utility.
Also, various Minesweeper versions may have the array on different address.
Knowing the fact that there is always a border (0x10 bytes), we can just find it in memory:

\lstinputlisting[style=customc]{examples/minesweeper/cheater2_fragment.c}

Full source code: \url{https://github.com/DennisYurichev/RE-for-beginners/blob/master/examples/minesweeper/minesweeper_cheater2.c}.

\subsection{\Exercises}

\begin{itemize}

\item 
Why do the \emph{border bytes} (or \emph{sentinel values}) (0x10) exist in the array?

What they are for if they are not visible in Minesweeper's interface?
How could it work without them?

\item 
As it turns out, there are more values possible (for open blocks, for flagged by user, \etc{}).
Try to find the meaning of each one.

\item 
Modify my utility so it can remove all mines or set them in a fixed pattern that you want in the Minesweeper
process currently running.

\end{itemize}
}
\RU{\mysection{Разгон майнера биткоинов Cointerra}
\index{Bitcoin}
\index{BeagleBone}

Был такой майнер биткоинов Cointerra, выглядящий так:

\begin{figure}[H]
\centering
\myincludegraphics{examples/bitcoin_miner/board.jpg}
\caption{Board}
\end{figure}

И была также (возможно утекшая) утилита\footnote{Можно скачать здесь: \url{https://github.com/DennisYurichev/RE-for-beginners/raw/master/examples/bitcoin_miner/files/cointool-overclock}}
которая могла выставлять тактовую частоту платы.
Она запускается на дополнительной плате BeagleBone на ARM с Linux (маленькая плата внизу фотографии).

И у автора (этих строк) однажды спросили, можно ли хакнуть эту утилиту и посмотреть, какие частоты можно выставлять, и какие нет.
И можно ли твикнуть её?

Утилиту нужно запускать так: \TT{./cointool-overclock 0 0 900}, где 900 это частота в МГц.
Если частота слишком большая, утилита выведет ошибку \q{Error with arguments} и закончит работу.

Вот фрагмент кода вокруг ссылки на текстовую строку \q{Error with arguments}:

\begin{lstlisting}[style=customasmARM]

...

.text:0000ABC4         STR      R3, [R11,#var_28]
.text:0000ABC8         MOV      R3, #optind
.text:0000ABD0         LDR      R3, [R3]
.text:0000ABD4         ADD      R3, R3, #1
.text:0000ABD8         MOV      R3, R3,LSL#2
.text:0000ABDC         LDR      R2, [R11,#argv]
.text:0000ABE0         ADD      R3, R2, R3
.text:0000ABE4         LDR      R3, [R3]
.text:0000ABE8         MOV      R0, R3  ; nptr
.text:0000ABEC         MOV      R1, #0  ; endptr
.text:0000ABF0         MOV      R2, #0  ; base
.text:0000ABF4         BL       strtoll
.text:0000ABF8         MOV      R2, R0
.text:0000ABFC         MOV      R3, R1
.text:0000AC00         MOV      R3, R2
.text:0000AC04         STR      R3, [R11,#var_2C]
.text:0000AC08         MOV      R3, #optind
.text:0000AC10         LDR      R3, [R3]
.text:0000AC14         ADD      R3, R3, #2
.text:0000AC18         MOV      R3, R3,LSL#2
.text:0000AC1C         LDR      R2, [R11,#argv]
.text:0000AC20         ADD      R3, R2, R3
.text:0000AC24         LDR      R3, [R3]
.text:0000AC28         MOV      R0, R3  ; nptr
.text:0000AC2C         MOV      R1, #0  ; endptr
.text:0000AC30         MOV      R2, #0  ; base
.text:0000AC34         BL       strtoll
.text:0000AC38         MOV      R2, R0
.text:0000AC3C         MOV      R3, R1
.text:0000AC40         MOV      R3, R2
.text:0000AC44         STR      R3, [R11,#third_argument]
.text:0000AC48         LDR      R3, [R11,#var_28]
.text:0000AC4C         CMP      R3, #0
.text:0000AC50         BLT      errors_with_arguments
.text:0000AC54         LDR      R3, [R11,#var_28]
.text:0000AC58         CMP      R3, #1
.text:0000AC5C         BGT      errors_with_arguments
.text:0000AC60         LDR      R3, [R11,#var_2C]
.text:0000AC64         CMP      R3, #0
.text:0000AC68         BLT      errors_with_arguments
.text:0000AC6C         LDR      R3, [R11,#var_2C]
.text:0000AC70         CMP      R3, #3
.text:0000AC74         BGT      errors_with_arguments
.text:0000AC78         LDR      R3, [R11,#third_argument]
.text:0000AC7C         CMP      R3, #0x31
.text:0000AC80         BLE      errors_with_arguments
.text:0000AC84         LDR      R2, [R11,#third_argument]
.text:0000AC88         MOV      R3, #950
.text:0000AC8C         CMP      R2, R3
.text:0000AC90         BGT      errors_with_arguments
.text:0000AC94         LDR      R2, [R11,#third_argument]
.text:0000AC98         MOV      R3, #0x51EB851F
.text:0000ACA0         SMULL    R1, R3, R3, R2
.text:0000ACA4         MOV      R1, R3,ASR#4
.text:0000ACA8         MOV      R3, R2,ASR#31
.text:0000ACAC         RSB      R3, R3, R1
.text:0000ACB0         MOV      R1, #50
.text:0000ACB4         MUL      R3, R1, R3
.text:0000ACB8         RSB      R3, R3, R2
.text:0000ACBC         CMP      R3, #0
.text:0000ACC0         BEQ      loc_ACEC
.text:0000ACC4
.text:0000ACC4 errors_with_arguments
.text:0000ACC4                                         
.text:0000ACC4         LDR      R3, [R11,#argv]
.text:0000ACC8         LDR      R3, [R3]
.text:0000ACCC         MOV      R0, R3  ; path
.text:0000ACD0         BL       __xpg_basename
.text:0000ACD4         MOV      R3, R0
.text:0000ACD8         MOV      R0, #aSErrorWithArgu ; format
.text:0000ACE0         MOV      R1, R3
.text:0000ACE4         BL       printf
.text:0000ACE8         B        loc_ADD4
.text:0000ACEC ; ------------------------------------------------------------
.text:0000ACEC
.text:0000ACEC loc_ACEC                 ; CODE XREF: main+66C
.text:0000ACEC         LDR      R2, [R11,#third_argument]
.text:0000ACF0         MOV      R3, #499
.text:0000ACF4         CMP      R2, R3
.text:0000ACF8         BGT      loc_AD08
.text:0000ACFC         MOV      R3, #0x64
.text:0000AD00         STR      R3, [R11,#unk_constant]
.text:0000AD04         B        jump_to_write_power
.text:0000AD08 ; ------------------------------------------------------------
.text:0000AD08
.text:0000AD08 loc_AD08                 ; CODE XREF: main+6A4
.text:0000AD08         LDR      R2, [R11,#third_argument]
.text:0000AD0C         MOV      R3, #799
.text:0000AD10         CMP      R2, R3
.text:0000AD14         BGT      loc_AD24
.text:0000AD18         MOV      R3, #0x5F
.text:0000AD1C         STR      R3, [R11,#unk_constant]
.text:0000AD20         B        jump_to_write_power
.text:0000AD24 ; ------------------------------------------------------------
.text:0000AD24
.text:0000AD24 loc_AD24                 ; CODE XREF: main+6C0
.text:0000AD24         LDR      R2, [R11,#third_argument]
.text:0000AD28         MOV      R3, #899
.text:0000AD2C         CMP      R2, R3
.text:0000AD30         BGT      loc_AD40
.text:0000AD34         MOV      R3, #0x5A
.text:0000AD38         STR      R3, [R11,#unk_constant]
.text:0000AD3C         B        jump_to_write_power
.text:0000AD40 ; ------------------------------------------------------------
.text:0000AD40
.text:0000AD40 loc_AD40                 ; CODE XREF: main+6DC
.text:0000AD40         LDR      R2, [R11,#third_argument]
.text:0000AD44         MOV      R3, #999
.text:0000AD48         CMP      R2, R3
.text:0000AD4C         BGT      loc_AD5C
.text:0000AD50         MOV      R3, #0x55
.text:0000AD54         STR      R3, [R11,#unk_constant]
.text:0000AD58         B        jump_to_write_power
.text:0000AD5C ; ------------------------------------------------------------
.text:0000AD5C
.text:0000AD5C loc_AD5C                 ; CODE XREF: main+6F8
.text:0000AD5C         LDR      R2, [R11,#third_argument]
.text:0000AD60         MOV      R3, #1099
.text:0000AD64         CMP      R2, R3
.text:0000AD68         BGT      jump_to_write_power
.text:0000AD6C         MOV      R3, #0x50
.text:0000AD70         STR      R3, [R11,#unk_constant]
.text:0000AD74
.text:0000AD74 jump_to_write_power                     ; CODE XREF: main+6B0
.text:0000AD74                                         ; main+6CC ...
.text:0000AD74         LDR      R3, [R11,#var_28]
.text:0000AD78         UXTB     R1, R3
.text:0000AD7C         LDR      R3, [R11,#var_2C]
.text:0000AD80         UXTB     R2, R3
.text:0000AD84         LDR      R3, [R11,#unk_constant]
.text:0000AD88         UXTB     R3, R3
.text:0000AD8C         LDR      R0, [R11,#third_argument]
.text:0000AD90         UXTH     R0, R0
.text:0000AD94         STR      R0, [SP,#0x44+var_44]
.text:0000AD98         LDR      R0, [R11,#var_24]
.text:0000AD9C         BL       write_power
.text:0000ADA0         LDR      R0, [R11,#var_24]
.text:0000ADA4         MOV      R1, #0x5A
.text:0000ADA8         BL       read_loop
.text:0000ADAC         B        loc_ADD4

...

.rodata:0000B378 aSErrorWithArgu DCB "%s: Error with arguments",0xA,0 ; DATA XREF: main+684

...

\end{lstlisting}

Имена ф-ций присутствовали в отладочной информации в оригинальном исполняемом файле,
такие как \TT{write\_power}, \TT{read\_loop}.
Но имена меткам внутри ф-ции дал я.

\myindex{UNIX!getopt}
\myindex{strtoll()}
Имя \TT{optind} звучит знакомо. Это библиотека \IT{getopt} из *NIX предназначенная для парсинга командной строки ---
и это то, что внутри и происходит.
Затем, третий аргумент (где передается значение частоты) конвертируется из строку в число используя вызов ф-ции \IT{strtoll()}.

Значение затем сравнивается с разными константами.
На 0xACEC есть проверка, меньше ли оно или равно 499, и если это так, то 0x64 будет передано в ф-цию
\TT{write\_power()} (которая посылает команду через USB используя \TT{send\_msg()}).
Если значение больше 499, происходит переход на 0xAD08.

На 0xAD08 есть проверка, меньше ли оно или равно 799. Если это так, то 0x5F передается в ф-цию \TT{write\_power()}.

Есть еще проверки: на 899 на 0xAD24, на 0x999 на 0xAD40, и наконец, на 1099 на 0xAD5C.
Если входная частота меньше или равна 1099, 0x50 (на 0xAD6C) будет передано в ф-цию \TT{write\_power()}.
И тут что-то вроде баги.
Если значение все еще больше 1099, само значение будет передано в ф-цию \TT{write\_power()}.
Но с другой стороны это не бага, потому что мы не можем попасть сюда: значение в начале проверяется с 950 на 0xAC88,
и если оно больше, выводится сообщение об ошибке и утилита заканчивает работу.

Вот таблица между частотами в МГц и значениями передаваемыми в ф-цию \TT{write\_power()}:

\begin{center}
\begin{longtable}{ | l | l | l | }
\hline
\HeaderColor МГц & \HeaderColor шестнадцатеричное представление & \HeaderColor десятичное \\
\hline
499MHz & 0x64 & 100 \\
\hline
799MHz & 0x5f & 95 \\
\hline
899MHz & 0x5a & 90 \\
\hline
999MHz & 0x55 & 85 \\
\hline
1099MHz & 0x50 & 80 \\
\hline
\end{longtable}
\end{center}

Как видно, значение передаваемое в плату постепенно уменьшается с ростом частоты.

Видно что значение в 950МГц это жесткий предел, по крайней мере в этой утилите. Можно ли её обмануть?

Вернемся к этому фрагменту кода:

\begin{lstlisting}[style=customasmARM]
.text:0000AC84      LDR     R2, [R11,#third_argument]
.text:0000AC88      MOV     R3, #950
.text:0000AC8C      CMP     R2, R3
.text:0000AC90      BGT     errors_with_arguments ; Я пропатчил здесь на 00 00 00 00
\end{lstlisting}

Нам нужно как-то запретить инструкцию перехода \INS{BGT} на 0xAC90. И это ARM в режиме ARM, потому что, как мы видим,
все адреса увеличиваются на 4, т.е., длина каждой инструкции это 4 байта.
Инструкция \TT{NOP} (нет операции) в режиме ARM это просто 4 нулевых байта: \TT{00 00 00 00}.
Так что, записывая 4 нуля по адресу 0xAC90 (или по физическому смещению в файле: 0x2C90) мы можем выключить
эту проверку.

Теперь можно выставлять частоты вплоть до 1050МГц. И даже больше, но из-за ошибки, если входное значение больше 1099,
значение в МГц, \IT{как есть}, будет передано в плату, что неправильно.

Дальше я не разбирался, но если бы продолжил, я бы уменьшал значение передаваемое в ф-цию \TT{write\_power()}.

Теперь страшный фрагмент кода, который я в начале пропустил:

\lstinputlisting[style=customasmARM]{examples/bitcoin_miner/tmp1.lst}

Здесь используется деление через умножение, и константа 0x51EB851F.
Я написал для себя простой программистский калькулятор\footnote{\url{https://github.com/DennisYurichev/progcalc}}.
И там есть возможность вычислять обратное число по модулю.

\begin{lstlisting}
modinv32(0x51EB851F)
Warning, result is not integer: 3.125000
(unsigned) dec: 3 hex: 0x3 bin: 11
\end{lstlisting}

Это значит что инструкция \INS{SMULL} на 0xACA0 просто делит 3-й аргумент на 3.125.
На самом деле, все что делает ф-ция \TT{modinv32()} в моем калькуляторе, это:

\[
\frac{1}{\frac{input}{2^{32}}} = \frac{2^{32}}{input}
\]

Потом там есть дополнительные сдвиги и теперь мы видим что 3-й аргумент просто делится на 50.
И затем умножается снова на 50.
Зачем?
Это простейшая проверка, можно ли делить входное значение на 50 без остатка.
Если значение этого выражения ненулевое, $x$ не может быть разделено на 50 без остатка:

\[
x-((\frac{x}{50}) \cdot 50)
\]

На самом деле, это простой способ вычисления остатка от деления.

И затем, если остаток ненулевой, выводится сообщение об ошибке.
Так что эта утилита берет значения частотв вроде 850, 900, 950, 1000, итд, но не 855 или 911.

Вот и всё! Если вы делаете что-то такое, имейте ввиду, что это может испортить вашу плату, как и в случае разгона
чипов вроде \ac{CPU}, \ac{GPU}, итд.
Если у вас есть плата Cointerra, делайте всё это на свой собственный риск!

}
\DE{\mysection{\Stack}
\label{sec:stack}
\myindex{\Stack}

Der Stack ist eine der fundamentalen Datenstrukturen in der Informatik.
\footnote{\href{http://go.yurichev.com/17119}{wikipedia.org/wiki/Call\_Stack}}.
\ac{AKA} \ac{LIFO}.

Technisch betrachtet ist es ein Stapelspeicher innerhalb des Prozessspeichers der zusammen mit den \ESP (x86), \RSP (x64) oder dem \ac{SP} (ARM) Register als ein Zeiger in diesem Speicherblock fungiert.

\myindex{ARM!\Instructions!PUSH}
\myindex{ARM!\Instructions!POP}
\myindex{x86!\Instructions!PUSH}
\myindex{x86!\Instructions!POP}

Die häufigsten Stack-Zugriffsinstruktionen sind die \PUSH- und \POP-Instruktionen (in beidem x86 und ARM Thumb-Modus). \PUSH subtrahiert vom \ESP/\RSP/\ac{SP} 4 Byte im 32-Bit Modus (oder 8 im 64-Bit Modus) und schreibt dann den Inhalt des Zeigers an die Adresse auf die von \ESP/\RSP/\ac{SP} gezeigt wird.

\POP ist die umgekehrte Operation: Die Daten des Zeigers für die Speicherregion auf die von \ac{SP}
gezeigt wird werden ausgelesen und die Inhalte in den Instruktionsoperanden geschreiben (oft ist das ein Register). Dann werden 4 (beziehungsweise 8) Byte zum \gls{stack pointer} addiert.

Nach der Stackallokation, zeigt der \gls{stack pointer} auf den Boden des Stacks.
\PUSH verringert den \gls{stack pointer} und \POP erhöht ihn.
Der Boden des Stacks ist eigentlich der Anfang der Speicherregion die für den Stack reserviert wurde.
Das wirkt zunächst seltsam, aber so funktioniert es.

ARM unterstützt beides, aufsteigende und absteigende Stacks.

\myindex{ARM!\Instructions!STMFD}
\myindex{ARM!\Instructions!LDMFD}
\myindex{ARM!\Instructions!STMED}
\myindex{ARM!\Instructions!LDMED}
\myindex{ARM!\Instructions!STMFA}
\myindex{ARM!\Instructions!LDMFA}
\myindex{ARM!\Instructions!STMEA}
\myindex{ARM!\Instructions!LDMEA}

Zum Beispiel die \ac{STMFD}/\ac{LDMFD} und \ac{STMED}/\ac{LDMED} Instruktionen sind alle dafür gedacht mit einem absteigendem Stack zu arbeiten ( wächst nach unten, fängt mit hohen Adressen an und entwickelt sich zu niedrigeren Adressen). Die \ac{STMFA}/\ac{LDMFA} und \ac{STMEA}/\ac{LDMEA} Instruktionen sind dazu gedacht mit einem aufsteigendem Stack zu arbeiten (wächst nach oben und fängt mit niedrigeren Adressen an und wächst nach oben).

% It might be worth mentioning that STMED and STMEA write first,
% and then move the pointer, and that LDMED and LDMEA move the pointer first, and then read.
% In other words, ARM not only lets the stack grow in a non-standard direction,
% but also in a non-standard order.
% Maybe this can be in the glossary, which would explain why E stands for "empty".

\subsection{Warum wächst der Stack nach unten?}
\label{stack_grow_backwards}

Intuitiv, würden man annehmen das der Stack nach oben wächst z.B Richtung höherer Adressen, so wie bei jeder anderen Datenstruktur.

Der Grund das der Stack rückwärts wächst ist wohl historisch bedingt. Als Computer so groß waren das sie einen ganzen Raum beansprucht haben war es einfach Speicher in zwei Sektionen zu unterteilen, einen Teil für den \gls{heap} und einen Teil für den Stack. Sicher war zu dieser Zeit nicht bekannt wie groß der \gls{heap} und der Stack wachsen würden, während der Programm Laufzeit, also war die Lösung die einfachste mögliche.

\input{patterns/02_stack/stack_and_heap}

In \RitchieThompsonUNIX können wir folgendes lesen:

\begin{framed}
\begin{quotation}
Der user-core eines Programm Images wird in drei logische Segmente unterteilt. Das Programm-Text Segment beginnt bei 0 im virtuellen Adress Speicher. Während der Ausführung wird das Segment als schreibgeschützt markiert und eine einzelne Kopie des Segments wird unter allen Prozessen geteilt die das Programm ausführen. An der ersten 8K grenze über dem Programm Text Segment im Virtuellen Speicher, fängt der ``nonshared'' Bereich an, der nach Bedarf von Syscalls erweitert werden kann. Beginnend bei der höchsten Adresse im Virtuellen Speicher ist das Stack Segment, das Automatisch nach unten wächst während der Hardware Stackpointer sich ändert.
\end{quotation}
\end{framed}

Das erinnert daran wie manche Schüler Notizen zu  zwei Vorträgen in einem Notebook dokumentieren:
Notizen für den ersten Vortrag werden normal notiert, und Notizen zur zum zweiten Vortrag werden 
ans Ende des Notizbuches geschrieben, indem man das Notizbuch umdreht. Die Notizen treffen sich irgendwann
im Notizbuch aufgrund des fehlenden Freien Platzes.

% I think if we want to expand on this analogy,
% one might remember that the line number increases as as you go down a page.
% So when you decrease the address when pushing to the stack, visually,
% the stack does grow upwards.
% Of course, the problem is that in most human languages,
% just as with computers,
% we write downwards, so this direction is what makes buffer overflows so messy.

\subsection{Für was wird der Stack benutzt?}

% subsections
\input{patterns/02_stack/01_saving_ret_addr}
\input{patterns/02_stack/02_args_passing}
\EN{\input{patterns/02_stack/03_local_vars_EN}}
\RU{\input{patterns/02_stack/03_local_vars_RU}}
\DE{\input{patterns/02_stack/03_local_vars_DE}}
\PTBR{\input{patterns/02_stack/03_local_vars_PTBR}}
\input{patterns/02_stack/04_alloca/main}
\input{patterns/02_stack/05_SEH}
\input{patterns/02_stack/06_BO_protection}

\subsubsection{Automatisches deallokieren der Daten auf dem Stack}

Vielleicht ist der Grund warum man lokale Variablen und SEH Einträge auf dem Stack speichert, weil sie beim 
verlassen der Funktion automatisch aufgeräumt werden. Man braucht dabei nur eine Instruktion um die Position
des Stackpointers zu korrigieren (oftmals ist es die \ADD Instruktion). Funktions Argumente, könnte man sagen 
werden auch am Ende der Funktion deallokiert. Im Kontrast dazu, alles was auf dem \emph{heap} gespeichert wird muss
explizit deallokiert werden. 

% sections
\EN{\input{patterns/02_stack/07_layout_EN}}
\RU{\input{patterns/02_stack/07_layout_RU}}
\DE{\input{patterns/02_stack/07_layout_DE}}
\PTBR{\input{patterns/02_stack/07_layout_PTBR}}
\input{patterns/02_stack/08_noise/main}
\input{patterns/02_stack/exercises}
}
\FR{\mysection{x86}

\subsection{Terminologie}

Commun en 16-bit (8086/80286), 32-bit (80386, etc.), 64-bit.

\myindex{IEEE 754}
\myindex{MS-DOS}
\begin{description}
	\item[octet] 8-bit.
		La directive d'assembleur DB est utilisée pour définir les variables et les tableaux d'octets.
		Les octets sont passés dans les parties 8-bit des registres: \TT{AL/BL/CL/DL/AH/BH/CH/DH/SIL/DIL/R*L}.
	\item[mot] 16-bit. 
		directive assembleur DW \dittoclosing.
		Les mots sont passés dans la partie 16-bit des registres:\\
			\TT{AX/BX/CX/DX/SI/DI/R*W}.
	\item[double mot] (\q{dword}) 32-bit.
		directive assembleur DD \dittoclosing.
		Les double mots sont passés dans les registres (x86) ou dans la partie 32-bit de registres (x64).
		Dans du code 16-bit, les double mots sont passés dans une paire de registres 16-bit.
	\item[quadruple mot] (\q{qword}) 64-bit.
		directive assembleur DQ \dittoclosing.
		En environnement 32-bit, les quadruple mots sont passés dans une paire de registres 32-bit.
	\item[tbyte] (10 bytes) 80-bit ou 10 octets (utilisé pour les registres FPU IEEE 754).
	\item[paragraph] (16 bytes)---le terme était répandu dans l'environnement MS-DOS.
\end{description}

\myindex{Windows!API}

Des types de données de même taille (BYTE, WORD, DWORD) existent aussi dans l'\ac{API} Windows.

\input{appendix/x86/registers_FR} % subsection
\input{appendix/x86/instructions_FR} % subsection
\input{appendix/x86/npad} % subsection

}
\IT{\mysection{x86}

\subsection{Terminologia}

Comune per 16-bit (8086/80286), 32-bit (80386, etc.), 64-bit.

\myindex{IEEE 754}
\myindex{MS-DOS}
\begin{description}
	\item[byte] 8-bit.
		La direttiva DB in assembly è utilizzata per definire variabili ed array di byte.
		I byte sono passati nella parte ad 8-bit dei registri: \TT{AL/BL/CL/DL/AH/BH/CH/DH/SIL/DIL/R*L}.
	\item[word] 16-bit.
		La direttiva DW in assembly \dittoclosing.
		Le Word sono passate nella parte a 16-bit dei registri:\\
			\TT{AX/BX/CX/DX/SI/DI/R*W}.
	\item[double word] (\q{dword}) 32-bit.
		La direttiva DD in assembly \dittoclosing.
		Le Double words sono passate nei registri (x86) o nelle parti a 32-bit dei registri (x64).
		Nel codice a 16-bit, le double words sono passate in coppie di registri a 16-bit.
	\item[quad word] (\q{qword}) 64-bit.
		La direttiva DQ in assembly \dittoclosing.
		Negli ambienti a 32-bit, le quad words sono passate in coppie di registri a 32-bit.
	\item[tbyte] (10 bytes) 80-bit o 10 bytes (usati per i registri FPU IEEE 754).
	\item[paragraph] (16 bytes)---termine popolare nell'ambiente MS-DOS. % TODO link to a part about 8086 memory model...
\end{description}

\myindex{Windows!API}

Tipi di dati della stessa dimensione (BYTE, WORD, DWORD) sono gli stessi nelle Windows \ac{API}.

%TBT
\input{appendix/x86/registers} % subsection
\input{appendix/x86/instructions} % subsection
\input{appendix/x86/npad} % subsection
}
\mysection{ARM}
\myindex{ARM}

\subsection{\RU{Терминология}\EN{Terminology}}

ARM \RU{изначально разрабатывался как 32-битный}\EN{was initially developed as 32-bit} \ac{CPU}, 
\RU{поэтому \emph{слово} здесь, в отличие от x86, 32-битное}\EN{so that's why a \emph{word} here, unlike x86, is 32-bit}.

\begin{description}
	\item[byte] 8-\bitENRU.
		\RU{Для определения переменных и массива байт используется директива ассемблера DCB}
		\EN{The DB assembly directive is used for defining variables and arrays of bytes}.
	\item[halfword] 16-\bitENRU. \RU{\dittoclosing директива ассемблера DCW}
					\EN{DCW assembly directive \dittoclosing}.	
	\item[word] 32-\bitENRU. \RU{\dittoclosing директива ассемблера DCD}
					\EN{DCD assembly directive \dittoclosing}.
	\item[doubleword] 64-\bitENRU.
	\item[quadword] 128-\bitENRU.
\end{description}

\subsection{\RU{Версии}\EN{Versions}}

\begin{itemize}
\item ARMv4: \RU{появился режим Thumb}\EN{Thumb mode introduced}.

\item ARMv6: \RU{использовался в}\EN{used in} iPhone 1st gen., iPhone 3G 
(Samsung 32-bit RISC ARM 1176JZ(F)-S \RU{поддерживающий}\EN{that supports} Thumb-2)

\item ARMv7: \RU{появился }Thumb-2\EN{ was added} (2003).
\RU{Использовался в}\EN{was used in} iPhone 3GS, iPhone 4, iPad 1st gen. (ARM Cortex-A8), iPad 2 (Cortex-A9),
iPad 3rd gen.

\item ARMv7s: \RU{Добавлены новые инструкции}\EN{New instructions added}.
\RU{Использовался в}\EN{Was used in} iPhone 5, iPhone 5c, iPad 4th gen. (Apple A6).

\item ARMv8: 64-\RU{битный процессор}\EN{bit CPU}, \ac{AKA} ARM64 \ac{AKA} AArch64.
\RU{Использовался в}\EN{Was used in} iPhone 5S, iPad Air (Apple A7).
\RU{В 64-битном режиме, режима Thumb больше нет, только режим ARM (4-байтные инструкции).}
\EN{There is no Thumb mode in 64-bit mode, only ARM (4-byte instructions).}
\end{itemize}

% sections
\input{appendix/ARM/registers}
\input{appendix/ARM/instructions}

\mysection{MIPS}

\subsection{\Registers}
\label{MIPS_registers_ref}

\myindex{MIPS!O32}
( \RU{Соглашение о вызовах O32}\EN{O32 calling convention} )

\subsubsection{\RU{Регистры общего пользования}\EN{General purpose registers} \ac{GPR}}

%\small
\begin{center}
\begin{tabular}{ | l | l | l | }
\hline
\HeaderColor \RU{Номер}\EN{Number} & 
\HeaderColor \RU{Псевдоимя}\EN{Pseudoname} & 
\HeaderColor \RU{Описание}\EN{Description} \\
\hline
\$0             & \$ZERO          & \RU{Всегда ноль. Запись в этот регистр это как}\EN{Always zero. Writing to this register is like} \ac{NOP}. \\
\hline
\$1             & \$AT            & \RU{Используется как временный регистр}\EN{Used as a temporary register} \\
                &                 & \RU{для ассемблерных макросов и псевдоинструкций}\EN{for assembly macros and pseudo instructions}. \\
\hline
\$2 \dots \$3   & \$V0 \dots \$V1 & \RU{Здесь возвращается результат функции}\EN{Function result is returned here}. \\
\hline
\$4 \dots \$7   & \$A0 \dots \$A3 & \RU{Аргументы функции}\EN{Function arguments}. \\
\hline
\$8 \dots \$15  & \$T0 \dots \$T7 & \RU{Используется для временных данных}\EN{Used for temporary data}. \\
\hline
\$16 \dots \$23 & \$S0 \dots \$S7 & \RU{Используется для временных данных}\EN{Used for temporary data}\AsteriskOne{}. \\
\hline
\$24 \dots \$25 & \$T8 \dots \$T9 & \RU{Используется для временных данных}\EN{Used for temporary data}. \\
\hline
\$26 \dots \$27 & \$K0 \dots \$K1 & \RU{Зарезервировано для ядра \ac{OS}}\EN{Reserved for \ac{OS} kernel}. \\
\hline
\$28            & \$GP            & \RU{Глобальный указатель}\EN{Global Pointer}\AsteriskTwo{}. \\
\hline
\$29            & \$SP            & \ac{SP}\AsteriskOne{}. \\
\hline
\$30            & \$FP            & \ac{FP}\AsteriskOne{}. \\
\hline
\$31            & \$RA            & \ac{RA}. \\
\hline
n/a             & PC              & \ac{PC}. \\
\hline
n/a             & HI              & \RU{старшие 32 бита результата умножения или остаток от деления}\EN{high 32 bit of multiplication or division remainder}\AsteriskThree{}. \\
\hline
n/a             & LO              & \RU{младшие 32 бита результата умножения или результат деления}\EN{low 32 bit of multiplication and division remainder}\AsteriskThree{}. \\
\hline
\end{tabular}
\end{center}
%\normalsize

\subsubsection{\RU{Регистры для работы с числами с плавающей точкой}\EN{Floating-point registers}}
\label{MIPS_FPU_registers}

\begin{center}
\begin{tabular}{ | l | l | l | }
\hline
\HeaderColor \RU{Название}\EN{Name} & \HeaderColor \RU{Описание}\EN{Description} \\
\hline
\$F0..\$F1   & \RU{Здесь возвращается результат функции}\EN{Function result returned here}. \\
\hline
\$F2..\$F3   & \RU{Не используется}\EN{Not used}. \\
\hline
\$F4..\$F11  & \RU{Используется для временных данных}\EN{Used for temporary data}. \\
\hline
\$F12..\$F15 & \RU{Первые два аргумента функции}\EN{First two function arguments}. \\
\hline
\$F16..\$F19 & \RU{Используется для временных данных}\EN{Used for temporary data}. \\
\hline
\$F20..\$F31 & \RU{Используется для временных данных}\EN{Used for temporary data}\AsteriskOne{}. \\
\hline
%fcr31 & Control/status register. \\
%\hline
\end{tabular}
\end{center}

\AsteriskOne{} --- \Gls{callee} \RU{должен сохранять}\EN{must preserve the value}.\\
\AsteriskTwo{} --- \Gls{callee} \RU{должен сохранять}\EN{must preserve the value} (\RU{кроме \ac{PIC}-кода}
\EN{except in \ac{PIC} code}).\\
\myindex{MIPS!\Instructions!MFLO}
\myindex{MIPS!\Instructions!MFHI}
\AsteriskThree{} --- \RU{доступны используя инструкции}\EN{accessible using the} \TT{MFHI} \AndENRU \TT{MFLO}\EN{ instructions}.\\

\subsection{\Instructions}

\RU{Есть три типа инструкций}\EN{There are 3 kinds of instructions}:

\begin{itemize}

\item \RU{Тип R: имеющие 3 регистра. R-инструкции обычно имеют такой вид:}
\EN{R-type: those which have 3 registers. R-instruction usually have the following form:}

\begin{lstlisting}
instruction destination, source1, source2
\end{lstlisting}

\RU{Важно помнить, что если первый и второй регистр один и тот же, IDA может показать инструкцию в сокращенной
форме:}
\EN{One important thing to keep in mind is that when the first and second register are the same, 
IDA may show the instruction in its shorter form:}

\begin{lstlisting}
instruction destination/source1, source2
\end{lstlisting}

\RU{Это немного напоминает Интеловский синтаксис ассемблера x86.}
\EN{That somewhat reminds us of the Intel syntax for x86 assembly language.}

\item \RU{Тип I: имеющие 2 регистра и 16-битное \q{immediate}-значение.}
\EN{I-type: those which have 2 registers and a 16-bit immediate value.}

\item \RU{Тип J: инструкции перехода, имеют 26 бит для кодирования смещения.}
\EN{J-type: jump/branch instructions, have 26 bits for encoding the offset.}

\end{itemize}

\subsubsection{\RU{Инструкции перехода}\EN{Jump instructions}}

\RU{Какая разница между инструкциями начинающихся с B- (BEQ, B, итд) и с J- (JAL, JALR, итд)?}
\EN{What is the difference between B- instructions (BEQ, B, etc.) and J- ones (JAL, JALR, etc.)?}

\RU{B-инструкции имеют тип I, так что, смещение в этих инструкциях кодируется как 16-битное значение.}
\EN{The B-instructions have an I-type, hence, the B-instructions' offset is encoded as a 16-bit immediate.}
\RU{Инструкции JR и JALR имеют тип R, и они делают переход по абсолютному адресу указанному в регистре.}
\EN{JR and JALR are R-type and jump to an absolute address specified in a register.}
\RU{J и JAL имеют тип J, так что смещение кодируется как 26-битное значение.}
\EN{J and JAL are J-type, hence the offset is encoded as a 26-bit immediate.}

\RU{Коротко говоря, в B-инструкциях можно кодировать условие 
(B на самом деле это псевдоинструкция для \TT{BEQ \$ZERO, \$ZERO, LABEL}),
а в J-инструкциях нельзя.}
\EN{In short, B-instructions can encode a condition 
(B is in fact pseudo instruction for \TT{BEQ \$ZERO, \$ZERO, LABEL}), 
while J-instructions can't.}

\mysection{\RU{Некоторые библиотечные функции GCC}\EN{Some GCC library functions}\DE{Einige GCC-Bibliotheks-Funktionen}%
\FR{Quelques fonctions de la bibliothèque de GCC}}
\myindex{GCC}
\label{sec:GCC_library_func}

%__ashldi3
%__ashrdi3
%__floatundidf
%__floatdisf
%__floatdixf
%__floatundidf
%__floatundisf
%__floatundixf
%__lshrdi3
%__muldi3

\begin{center}
\begin{tabular}{ | l | l | }
\hline
\HeaderColor \RU{имя}\EN{name}\DE{Name}\FR{nom} & \HeaderColor \RU{значение}\EN{meaning}\DE{Bedeutung}%
\FR{signification} \\
\hline \TT{\_\_divdi3} & \RU{знаковое деление}\EN{signed division}\DE{vorzeichenbehaftete Division}%
\FR{division signée} \\
\hline \TT{\_\_moddi3} & \RU{остаток от знакового деления}\EN{getting remainder (modulo) of signed division}%
\DE{Rest (Modulo) einer vorzeichenbehafteten Division}%
\FR{reste (modulo) d'une division signée} \\
\hline \TT{\_\_udivdi3} & \RU{беззнаковое деление}\EN{unsigned division}\DE{vorzeichenlose Division}%
\FR{division non signée} \\
\hline \TT{\_\_umoddi3} & \RU{остаток от беззнакового деления}\EN{getting remainder (modulo) of unsigned division}%
\DE{Rest (Modulo) einer vorzeichenlosen Division}%
\FR{reste (modulo) d'une division non signée} \\
\hline
\end{tabular}
\end{center}



\mysection{\RU{Некоторые библиотечные функции MSVC}\EN{Some MSVC library functions}\DE{Einige MSVC-Bibliotheks-Funktionen}%
\FR{Quelques fonctions de la bibliothèque MSVC}}
\myindex{MSVC}
\label{sec:MSVC_library_func}

\TT{ll} \RU{в имени функции означает}\EN{in function name stands for}\DE{in Funktionsnamen steht für}%
\FR{dans une fontion signifie} \q{long long}, \RU{т.е. 64-битный тип данных}
\EN{e.g., a 64-bit data type}\DE{z.B. einen 64-Bit-Datentyp}\FR{i.e., type de donées 64-bit}.

\begin{center}
\begin{tabular}{ | l | l | }
\hline
\HeaderColor \RU{имя}\EN{name}\DE{Name}\FR{nom} & \HeaderColor \RU{значение}\EN{meaning}\DE{Bedeutung}\FR{signification} \\
\hline \TT{\_\_alldiv} & \RU{знаковое деление}\EN{signed division}\DE{vorzeichenbehaftete Division}\FR{division signée} \\
\hline \TT{\_\_allmul} & \RU{умножение}\EN{multiplication}\DE{Multiplikation}\FR{multiplication} \\
\hline \TT{\_\_allrem} & \RU{остаток от знакового деления}\EN{remainder of signed division}\DE{Rest einer vorzeichenbehafteten Division}%
\FR{reste de la division signée} \\
\hline \TT{\_\_allshl} & \RU{сдвиг влево}\EN{shift left}\DE{Schiebe links}\FR{décalage à gauche} \\
\hline \TT{\_\_allshr} & \RU{знаковый сдвиг вправо}\EN{signed shift right}\DE{Schiebe links, vorzeichenbehaftet}%
\FR{décalage signé à droite} \\
\hline \TT{\_\_aulldiv} & \RU{беззнаковое деление}\EN{unsigned division}\DE{vorzeichenlose Division}%
\FR{division non signée} \\
\hline \TT{\_\_aullrem} & \RU{остаток от беззнакового деления}\EN{remainder of unsigned division}\DE{Rest (Modulo) einer vorzeichenlosen Division}%
\FR{reste de la division non signée} \\
\hline \TT{\_\_aullshr} & \RU{беззнаковый сдвиг вправо}\EN{unsigned shift right}\DE{Schiebe rechts, vorzeichenlos}%
\FR{décalage non signé à droite} \\
\hline
\end{tabular}
\end{center}

\RU{Процедуры умножения и сдвига влево, одни и те же и для знаковых чисел, и для беззнаковых,
поэтому здесь только одна функция для каждой операции}
\EN{Multiplication and shift left procedures are the same for both signed and unsigned numbers, hence there is only one function 
for each operation here}
\DE{Multiplikation und Links-Schiebebefehle sind sowohl für vorzeichenbehaftete als auch vorzeichenlose Zahlen,
da hier für jede Operation nur ein Befehl existiert}
\FR{La multiplication et le décalage à gauche sont similaire pour les nombres signés
et non signés, donc il n'y a qu'une seule fonction ici}. \\
\\
\RU{Исходные коды этих функций можно найти в установленной \ac{MSVS}, в}\EN{The source code of these function
can be found in the installed \ac{MSVS}, in}%
\DE{Der Quellcode dieser Funktionen kann im Pfad des installierten \ac{MSVS}, gefunden werden: }%
\FR{Le code source des ces fonctions peut être trouvé dans l'installation de \ac{MSVS},
dans} \TT{VC/crt/src/intel/*.asm}.


\input{appendix/cheatsheets}
}
\FR{\part*{\RU{Приложение}\EN{Appendix}\DE{Anhang}\FR{Appendice}\IT{Appendice}}
\appendix
\addcontentsline{toc}{part}{\RU{Приложение}\EN{Appendix}\DE{Anhang}\FR{Appendice}\IT{Appendice}}

% chapters
\EN{\mysection{\MinesweeperWinXPExampleChapterName}
\label{minesweeper_winxp}
\myindex{Windows!Windows XP}

For those who are not very good at playing Minesweeper, we could try to reveal the hidden mines in the debugger.

\myindex{\CStandardLibrary!rand()}
\myindex{Windows!PDB}

As we know, Minesweeper places mines randomly, so there has to be some kind of random number generator or
a call to the standard \TT{rand()} C-function.

What is really cool about reversing Microsoft products is that there are \gls{PDB} 
file with symbols (function names, \etc{}).
When we load \TT{winmine.exe} into \IDA, it downloads the 
\gls{PDB} file exactly for this 
executable and shows all names.

So here it is, the only call to \TT{rand()} is this function:

\lstinputlisting[style=customasmx86]{examples/minesweeper/tmp1.lst}

\IDA named it so, and it was the name given to it by Minesweeper's developers.

The function is very simple:

\begin{lstlisting}[style=customc]
int Rnd(int limit)
{
    return rand() % limit;
};
\end{lstlisting}

(There is no \q{limit} name in the \gls{PDB} file; we manually named this argument like this.)

So it returns 
a random value from 0 to a specified limit.

\TT{Rnd()} is called only from one place, 
a function called \TT{StartGame()}, 
and as it seems, this is exactly 
the code which place the mines:

\begin{lstlisting}[style=customasmx86]
.text:010036C7                 push    _xBoxMac
.text:010036CD                 call    _Rnd@4          ; Rnd(x)
.text:010036D2                 push    _yBoxMac
.text:010036D8                 mov     esi, eax
.text:010036DA                 inc     esi
.text:010036DB                 call    _Rnd@4          ; Rnd(x)
.text:010036E0                 inc     eax
.text:010036E1                 mov     ecx, eax
.text:010036E3                 shl     ecx, 5          ; ECX=ECX*32
.text:010036E6                 test    _rgBlk[ecx+esi], 80h
.text:010036EE                 jnz     short loc_10036C7
.text:010036F0                 shl     eax, 5          ; EAX=EAX*32
.text:010036F3                 lea     eax, _rgBlk[eax+esi]
.text:010036FA                 or      byte ptr [eax], 80h
.text:010036FD                 dec     _cBombStart
.text:01003703                 jnz     short loc_10036C7
\end{lstlisting}

Minesweeper allows you to set the board size, so the X (xBoxMac) and Y (yBoxMac) of the board are global variables.
They are passed to \TT{Rnd()} and random 
coordinates are generated.
A mine is placed by the \TT{OR} instruction at \TT{0x010036FA}. 
And if it has been placed before 
(it's possible if the pair of \TT{Rnd()} 
generates a coordinates pair which has been already 
generated), 
then \TT{TEST} and \TT{JNZ} at \TT{0x010036E6} 
jumps to the generation routine again.

\TT{cBombStart} is the global variable containing total number of mines. So this is loop.

The width of the array is 32 
(we can conclude this by looking at the \TT{SHL} instruction, which multiplies one of the coordinates by 32).

The size of the \TT{rgBlk} 
global array can be easily determined by the difference 
between the \TT{rgBlk} 
label in the data segment and the next known one. 
It is 0x360 (864):

\begin{lstlisting}[style=customasmx86]
.data:01005340 _rgBlk          db 360h dup(?)          ; DATA XREF: MainWndProc(x,x,x,x)+574
.data:01005340                                         ; DisplayBlk(x,x)+23
.data:010056A0 _Preferences    dd ?                    ; DATA XREF: FixMenus()+2
...
\end{lstlisting}

$864/32=27$.

So the array size is $27*32$?
It is close to what we know: when we try to set board size to $100*100$ in Minesweeper settings, it fallbacks to a board of size $24*30$.
So this is the maximal board size here.
And the array has a fixed size for any board size.

So let's see all this in \olly.
We will ran Minesweeper, attaching \olly to it and now we can see the memory dump at the address of the \TT{rgBlk} array (\TT{0x01005340})
\footnote{All addresses here are for Minesweeper for Windows XP SP3 English. 
They may differ for other service packs.}.

So we got this memory dump of the array:

\lstinputlisting[style=customasmx86]{examples/minesweeper/1.lst}

\olly, like any other hexadecimal editor, shows 16 bytes per line.
So each 32-byte array row occupies exactly 2 lines here.

This is beginner level (9*9 board).

There is some square 
structure can be seen visually (0x10 bytes).

We will click \q{Run} in \olly to unfreeze the Minesweeper process, then we'll clicked randomly at the Minesweeper window 
and trapped into mine, but now all mines are visible:

\begin{figure}[H]
\centering
\myincludegraphicsSmall{examples/minesweeper/1.png}
\caption{Mines}
\label{fig:minesweeper1}
\end{figure}

By comparing the mine places and the dump, we can conclude that 0x10 stands for border, 0x0F---empty block, 0x8F---mine.
Perhaps, 0x10 is just a \emph{sentinel value}.

Now we'll add comments and also enclose all 0x8F bytes into square brackets:

\lstinputlisting[style=customasmx86]{examples/minesweeper/2.lst}

Now we'll remove all \emph{border bytes} (0x10) and what's beyond those:

\lstinputlisting[style=customasmx86]{examples/minesweeper/3.lst}

Yes, these are mines, now it can be clearly seen and compared with the screenshot.

\clearpage
What is interesting is that we can modify the array right in \olly.
We can remove all mines by changing all 0x8F bytes by 0x0F, and here is what we'll get in Minesweeper:

\begin{figure}[H]
\centering
\myincludegraphicsSmall{examples/minesweeper/3.png}
\caption{All mines are removed in debugger}
\label{fig:minesweeper3}
\end{figure}

We can also move all of them to the first line: 

\begin{figure}[H]
\centering
\myincludegraphicsSmall{examples/minesweeper/2.png}
\caption{Mines set in debugger}
\label{fig:minesweeper2}
\end{figure}

Well, the debugger is not very convenient for eavesdropping (which is our goal anyway), so we'll write a small utility
to dump the contents of the board:

\lstinputlisting[style=customc]{examples/minesweeper/minesweeper_cheater.c}

Just set the \ac{PID}
\footnote{PID it can be seen in Task Manager 
(enable it in \q{View $\rightarrow$ Select Columns})} 
and the address of the array (\TT{0x01005340} for Windows XP SP3 English) 
and it will dump it
\footnote{The compiled executable is here: 
\href{http://go.yurichev.com/17165}{beginners.re}}.

It attaches itself to a win32 process by \ac{PID} and just reads process memory at the address.

\subsection{Finding grid automatically}

This is kind of nuisance to set address each time when we run our utility.
Also, various Minesweeper versions may have the array on different address.
Knowing the fact that there is always a border (0x10 bytes), we can just find it in memory:

\lstinputlisting[style=customc]{examples/minesweeper/cheater2_fragment.c}

Full source code: \url{https://github.com/DennisYurichev/RE-for-beginners/blob/master/examples/minesweeper/minesweeper_cheater2.c}.

\subsection{\Exercises}

\begin{itemize}

\item 
Why do the \emph{border bytes} (or \emph{sentinel values}) (0x10) exist in the array?

What they are for if they are not visible in Minesweeper's interface?
How could it work without them?

\item 
As it turns out, there are more values possible (for open blocks, for flagged by user, \etc{}).
Try to find the meaning of each one.

\item 
Modify my utility so it can remove all mines or set them in a fixed pattern that you want in the Minesweeper
process currently running.

\end{itemize}
}
\RU{\mysection{Разгон майнера биткоинов Cointerra}
\index{Bitcoin}
\index{BeagleBone}

Был такой майнер биткоинов Cointerra, выглядящий так:

\begin{figure}[H]
\centering
\myincludegraphics{examples/bitcoin_miner/board.jpg}
\caption{Board}
\end{figure}

И была также (возможно утекшая) утилита\footnote{Можно скачать здесь: \url{https://github.com/DennisYurichev/RE-for-beginners/raw/master/examples/bitcoin_miner/files/cointool-overclock}}
которая могла выставлять тактовую частоту платы.
Она запускается на дополнительной плате BeagleBone на ARM с Linux (маленькая плата внизу фотографии).

И у автора (этих строк) однажды спросили, можно ли хакнуть эту утилиту и посмотреть, какие частоты можно выставлять, и какие нет.
И можно ли твикнуть её?

Утилиту нужно запускать так: \TT{./cointool-overclock 0 0 900}, где 900 это частота в МГц.
Если частота слишком большая, утилита выведет ошибку \q{Error with arguments} и закончит работу.

Вот фрагмент кода вокруг ссылки на текстовую строку \q{Error with arguments}:

\begin{lstlisting}[style=customasmARM]

...

.text:0000ABC4         STR      R3, [R11,#var_28]
.text:0000ABC8         MOV      R3, #optind
.text:0000ABD0         LDR      R3, [R3]
.text:0000ABD4         ADD      R3, R3, #1
.text:0000ABD8         MOV      R3, R3,LSL#2
.text:0000ABDC         LDR      R2, [R11,#argv]
.text:0000ABE0         ADD      R3, R2, R3
.text:0000ABE4         LDR      R3, [R3]
.text:0000ABE8         MOV      R0, R3  ; nptr
.text:0000ABEC         MOV      R1, #0  ; endptr
.text:0000ABF0         MOV      R2, #0  ; base
.text:0000ABF4         BL       strtoll
.text:0000ABF8         MOV      R2, R0
.text:0000ABFC         MOV      R3, R1
.text:0000AC00         MOV      R3, R2
.text:0000AC04         STR      R3, [R11,#var_2C]
.text:0000AC08         MOV      R3, #optind
.text:0000AC10         LDR      R3, [R3]
.text:0000AC14         ADD      R3, R3, #2
.text:0000AC18         MOV      R3, R3,LSL#2
.text:0000AC1C         LDR      R2, [R11,#argv]
.text:0000AC20         ADD      R3, R2, R3
.text:0000AC24         LDR      R3, [R3]
.text:0000AC28         MOV      R0, R3  ; nptr
.text:0000AC2C         MOV      R1, #0  ; endptr
.text:0000AC30         MOV      R2, #0  ; base
.text:0000AC34         BL       strtoll
.text:0000AC38         MOV      R2, R0
.text:0000AC3C         MOV      R3, R1
.text:0000AC40         MOV      R3, R2
.text:0000AC44         STR      R3, [R11,#third_argument]
.text:0000AC48         LDR      R3, [R11,#var_28]
.text:0000AC4C         CMP      R3, #0
.text:0000AC50         BLT      errors_with_arguments
.text:0000AC54         LDR      R3, [R11,#var_28]
.text:0000AC58         CMP      R3, #1
.text:0000AC5C         BGT      errors_with_arguments
.text:0000AC60         LDR      R3, [R11,#var_2C]
.text:0000AC64         CMP      R3, #0
.text:0000AC68         BLT      errors_with_arguments
.text:0000AC6C         LDR      R3, [R11,#var_2C]
.text:0000AC70         CMP      R3, #3
.text:0000AC74         BGT      errors_with_arguments
.text:0000AC78         LDR      R3, [R11,#third_argument]
.text:0000AC7C         CMP      R3, #0x31
.text:0000AC80         BLE      errors_with_arguments
.text:0000AC84         LDR      R2, [R11,#third_argument]
.text:0000AC88         MOV      R3, #950
.text:0000AC8C         CMP      R2, R3
.text:0000AC90         BGT      errors_with_arguments
.text:0000AC94         LDR      R2, [R11,#third_argument]
.text:0000AC98         MOV      R3, #0x51EB851F
.text:0000ACA0         SMULL    R1, R3, R3, R2
.text:0000ACA4         MOV      R1, R3,ASR#4
.text:0000ACA8         MOV      R3, R2,ASR#31
.text:0000ACAC         RSB      R3, R3, R1
.text:0000ACB0         MOV      R1, #50
.text:0000ACB4         MUL      R3, R1, R3
.text:0000ACB8         RSB      R3, R3, R2
.text:0000ACBC         CMP      R3, #0
.text:0000ACC0         BEQ      loc_ACEC
.text:0000ACC4
.text:0000ACC4 errors_with_arguments
.text:0000ACC4                                         
.text:0000ACC4         LDR      R3, [R11,#argv]
.text:0000ACC8         LDR      R3, [R3]
.text:0000ACCC         MOV      R0, R3  ; path
.text:0000ACD0         BL       __xpg_basename
.text:0000ACD4         MOV      R3, R0
.text:0000ACD8         MOV      R0, #aSErrorWithArgu ; format
.text:0000ACE0         MOV      R1, R3
.text:0000ACE4         BL       printf
.text:0000ACE8         B        loc_ADD4
.text:0000ACEC ; ------------------------------------------------------------
.text:0000ACEC
.text:0000ACEC loc_ACEC                 ; CODE XREF: main+66C
.text:0000ACEC         LDR      R2, [R11,#third_argument]
.text:0000ACF0         MOV      R3, #499
.text:0000ACF4         CMP      R2, R3
.text:0000ACF8         BGT      loc_AD08
.text:0000ACFC         MOV      R3, #0x64
.text:0000AD00         STR      R3, [R11,#unk_constant]
.text:0000AD04         B        jump_to_write_power
.text:0000AD08 ; ------------------------------------------------------------
.text:0000AD08
.text:0000AD08 loc_AD08                 ; CODE XREF: main+6A4
.text:0000AD08         LDR      R2, [R11,#third_argument]
.text:0000AD0C         MOV      R3, #799
.text:0000AD10         CMP      R2, R3
.text:0000AD14         BGT      loc_AD24
.text:0000AD18         MOV      R3, #0x5F
.text:0000AD1C         STR      R3, [R11,#unk_constant]
.text:0000AD20         B        jump_to_write_power
.text:0000AD24 ; ------------------------------------------------------------
.text:0000AD24
.text:0000AD24 loc_AD24                 ; CODE XREF: main+6C0
.text:0000AD24         LDR      R2, [R11,#third_argument]
.text:0000AD28         MOV      R3, #899
.text:0000AD2C         CMP      R2, R3
.text:0000AD30         BGT      loc_AD40
.text:0000AD34         MOV      R3, #0x5A
.text:0000AD38         STR      R3, [R11,#unk_constant]
.text:0000AD3C         B        jump_to_write_power
.text:0000AD40 ; ------------------------------------------------------------
.text:0000AD40
.text:0000AD40 loc_AD40                 ; CODE XREF: main+6DC
.text:0000AD40         LDR      R2, [R11,#third_argument]
.text:0000AD44         MOV      R3, #999
.text:0000AD48         CMP      R2, R3
.text:0000AD4C         BGT      loc_AD5C
.text:0000AD50         MOV      R3, #0x55
.text:0000AD54         STR      R3, [R11,#unk_constant]
.text:0000AD58         B        jump_to_write_power
.text:0000AD5C ; ------------------------------------------------------------
.text:0000AD5C
.text:0000AD5C loc_AD5C                 ; CODE XREF: main+6F8
.text:0000AD5C         LDR      R2, [R11,#third_argument]
.text:0000AD60         MOV      R3, #1099
.text:0000AD64         CMP      R2, R3
.text:0000AD68         BGT      jump_to_write_power
.text:0000AD6C         MOV      R3, #0x50
.text:0000AD70         STR      R3, [R11,#unk_constant]
.text:0000AD74
.text:0000AD74 jump_to_write_power                     ; CODE XREF: main+6B0
.text:0000AD74                                         ; main+6CC ...
.text:0000AD74         LDR      R3, [R11,#var_28]
.text:0000AD78         UXTB     R1, R3
.text:0000AD7C         LDR      R3, [R11,#var_2C]
.text:0000AD80         UXTB     R2, R3
.text:0000AD84         LDR      R3, [R11,#unk_constant]
.text:0000AD88         UXTB     R3, R3
.text:0000AD8C         LDR      R0, [R11,#third_argument]
.text:0000AD90         UXTH     R0, R0
.text:0000AD94         STR      R0, [SP,#0x44+var_44]
.text:0000AD98         LDR      R0, [R11,#var_24]
.text:0000AD9C         BL       write_power
.text:0000ADA0         LDR      R0, [R11,#var_24]
.text:0000ADA4         MOV      R1, #0x5A
.text:0000ADA8         BL       read_loop
.text:0000ADAC         B        loc_ADD4

...

.rodata:0000B378 aSErrorWithArgu DCB "%s: Error with arguments",0xA,0 ; DATA XREF: main+684

...

\end{lstlisting}

Имена ф-ций присутствовали в отладочной информации в оригинальном исполняемом файле,
такие как \TT{write\_power}, \TT{read\_loop}.
Но имена меткам внутри ф-ции дал я.

\myindex{UNIX!getopt}
\myindex{strtoll()}
Имя \TT{optind} звучит знакомо. Это библиотека \IT{getopt} из *NIX предназначенная для парсинга командной строки ---
и это то, что внутри и происходит.
Затем, третий аргумент (где передается значение частоты) конвертируется из строку в число используя вызов ф-ции \IT{strtoll()}.

Значение затем сравнивается с разными константами.
На 0xACEC есть проверка, меньше ли оно или равно 499, и если это так, то 0x64 будет передано в ф-цию
\TT{write\_power()} (которая посылает команду через USB используя \TT{send\_msg()}).
Если значение больше 499, происходит переход на 0xAD08.

На 0xAD08 есть проверка, меньше ли оно или равно 799. Если это так, то 0x5F передается в ф-цию \TT{write\_power()}.

Есть еще проверки: на 899 на 0xAD24, на 0x999 на 0xAD40, и наконец, на 1099 на 0xAD5C.
Если входная частота меньше или равна 1099, 0x50 (на 0xAD6C) будет передано в ф-цию \TT{write\_power()}.
И тут что-то вроде баги.
Если значение все еще больше 1099, само значение будет передано в ф-цию \TT{write\_power()}.
Но с другой стороны это не бага, потому что мы не можем попасть сюда: значение в начале проверяется с 950 на 0xAC88,
и если оно больше, выводится сообщение об ошибке и утилита заканчивает работу.

Вот таблица между частотами в МГц и значениями передаваемыми в ф-цию \TT{write\_power()}:

\begin{center}
\begin{longtable}{ | l | l | l | }
\hline
\HeaderColor МГц & \HeaderColor шестнадцатеричное представление & \HeaderColor десятичное \\
\hline
499MHz & 0x64 & 100 \\
\hline
799MHz & 0x5f & 95 \\
\hline
899MHz & 0x5a & 90 \\
\hline
999MHz & 0x55 & 85 \\
\hline
1099MHz & 0x50 & 80 \\
\hline
\end{longtable}
\end{center}

Как видно, значение передаваемое в плату постепенно уменьшается с ростом частоты.

Видно что значение в 950МГц это жесткий предел, по крайней мере в этой утилите. Можно ли её обмануть?

Вернемся к этому фрагменту кода:

\begin{lstlisting}[style=customasmARM]
.text:0000AC84      LDR     R2, [R11,#third_argument]
.text:0000AC88      MOV     R3, #950
.text:0000AC8C      CMP     R2, R3
.text:0000AC90      BGT     errors_with_arguments ; Я пропатчил здесь на 00 00 00 00
\end{lstlisting}

Нам нужно как-то запретить инструкцию перехода \INS{BGT} на 0xAC90. И это ARM в режиме ARM, потому что, как мы видим,
все адреса увеличиваются на 4, т.е., длина каждой инструкции это 4 байта.
Инструкция \TT{NOP} (нет операции) в режиме ARM это просто 4 нулевых байта: \TT{00 00 00 00}.
Так что, записывая 4 нуля по адресу 0xAC90 (или по физическому смещению в файле: 0x2C90) мы можем выключить
эту проверку.

Теперь можно выставлять частоты вплоть до 1050МГц. И даже больше, но из-за ошибки, если входное значение больше 1099,
значение в МГц, \IT{как есть}, будет передано в плату, что неправильно.

Дальше я не разбирался, но если бы продолжил, я бы уменьшал значение передаваемое в ф-цию \TT{write\_power()}.

Теперь страшный фрагмент кода, который я в начале пропустил:

\lstinputlisting[style=customasmARM]{examples/bitcoin_miner/tmp1.lst}

Здесь используется деление через умножение, и константа 0x51EB851F.
Я написал для себя простой программистский калькулятор\footnote{\url{https://github.com/DennisYurichev/progcalc}}.
И там есть возможность вычислять обратное число по модулю.

\begin{lstlisting}
modinv32(0x51EB851F)
Warning, result is not integer: 3.125000
(unsigned) dec: 3 hex: 0x3 bin: 11
\end{lstlisting}

Это значит что инструкция \INS{SMULL} на 0xACA0 просто делит 3-й аргумент на 3.125.
На самом деле, все что делает ф-ция \TT{modinv32()} в моем калькуляторе, это:

\[
\frac{1}{\frac{input}{2^{32}}} = \frac{2^{32}}{input}
\]

Потом там есть дополнительные сдвиги и теперь мы видим что 3-й аргумент просто делится на 50.
И затем умножается снова на 50.
Зачем?
Это простейшая проверка, можно ли делить входное значение на 50 без остатка.
Если значение этого выражения ненулевое, $x$ не может быть разделено на 50 без остатка:

\[
x-((\frac{x}{50}) \cdot 50)
\]

На самом деле, это простой способ вычисления остатка от деления.

И затем, если остаток ненулевой, выводится сообщение об ошибке.
Так что эта утилита берет значения частотв вроде 850, 900, 950, 1000, итд, но не 855 или 911.

Вот и всё! Если вы делаете что-то такое, имейте ввиду, что это может испортить вашу плату, как и в случае разгона
чипов вроде \ac{CPU}, \ac{GPU}, итд.
Если у вас есть плата Cointerra, делайте всё это на свой собственный риск!

}
\DE{\mysection{\Stack}
\label{sec:stack}
\myindex{\Stack}

Der Stack ist eine der fundamentalen Datenstrukturen in der Informatik.
\footnote{\href{http://go.yurichev.com/17119}{wikipedia.org/wiki/Call\_Stack}}.
\ac{AKA} \ac{LIFO}.

Technisch betrachtet ist es ein Stapelspeicher innerhalb des Prozessspeichers der zusammen mit den \ESP (x86), \RSP (x64) oder dem \ac{SP} (ARM) Register als ein Zeiger in diesem Speicherblock fungiert.

\myindex{ARM!\Instructions!PUSH}
\myindex{ARM!\Instructions!POP}
\myindex{x86!\Instructions!PUSH}
\myindex{x86!\Instructions!POP}

Die häufigsten Stack-Zugriffsinstruktionen sind die \PUSH- und \POP-Instruktionen (in beidem x86 und ARM Thumb-Modus). \PUSH subtrahiert vom \ESP/\RSP/\ac{SP} 4 Byte im 32-Bit Modus (oder 8 im 64-Bit Modus) und schreibt dann den Inhalt des Zeigers an die Adresse auf die von \ESP/\RSP/\ac{SP} gezeigt wird.

\POP ist die umgekehrte Operation: Die Daten des Zeigers für die Speicherregion auf die von \ac{SP}
gezeigt wird werden ausgelesen und die Inhalte in den Instruktionsoperanden geschreiben (oft ist das ein Register). Dann werden 4 (beziehungsweise 8) Byte zum \gls{stack pointer} addiert.

Nach der Stackallokation, zeigt der \gls{stack pointer} auf den Boden des Stacks.
\PUSH verringert den \gls{stack pointer} und \POP erhöht ihn.
Der Boden des Stacks ist eigentlich der Anfang der Speicherregion die für den Stack reserviert wurde.
Das wirkt zunächst seltsam, aber so funktioniert es.

ARM unterstützt beides, aufsteigende und absteigende Stacks.

\myindex{ARM!\Instructions!STMFD}
\myindex{ARM!\Instructions!LDMFD}
\myindex{ARM!\Instructions!STMED}
\myindex{ARM!\Instructions!LDMED}
\myindex{ARM!\Instructions!STMFA}
\myindex{ARM!\Instructions!LDMFA}
\myindex{ARM!\Instructions!STMEA}
\myindex{ARM!\Instructions!LDMEA}

Zum Beispiel die \ac{STMFD}/\ac{LDMFD} und \ac{STMED}/\ac{LDMED} Instruktionen sind alle dafür gedacht mit einem absteigendem Stack zu arbeiten ( wächst nach unten, fängt mit hohen Adressen an und entwickelt sich zu niedrigeren Adressen). Die \ac{STMFA}/\ac{LDMFA} und \ac{STMEA}/\ac{LDMEA} Instruktionen sind dazu gedacht mit einem aufsteigendem Stack zu arbeiten (wächst nach oben und fängt mit niedrigeren Adressen an und wächst nach oben).

% It might be worth mentioning that STMED and STMEA write first,
% and then move the pointer, and that LDMED and LDMEA move the pointer first, and then read.
% In other words, ARM not only lets the stack grow in a non-standard direction,
% but also in a non-standard order.
% Maybe this can be in the glossary, which would explain why E stands for "empty".

\subsection{Warum wächst der Stack nach unten?}
\label{stack_grow_backwards}

Intuitiv, würden man annehmen das der Stack nach oben wächst z.B Richtung höherer Adressen, so wie bei jeder anderen Datenstruktur.

Der Grund das der Stack rückwärts wächst ist wohl historisch bedingt. Als Computer so groß waren das sie einen ganzen Raum beansprucht haben war es einfach Speicher in zwei Sektionen zu unterteilen, einen Teil für den \gls{heap} und einen Teil für den Stack. Sicher war zu dieser Zeit nicht bekannt wie groß der \gls{heap} und der Stack wachsen würden, während der Programm Laufzeit, also war die Lösung die einfachste mögliche.

\input{patterns/02_stack/stack_and_heap}

In \RitchieThompsonUNIX können wir folgendes lesen:

\begin{framed}
\begin{quotation}
Der user-core eines Programm Images wird in drei logische Segmente unterteilt. Das Programm-Text Segment beginnt bei 0 im virtuellen Adress Speicher. Während der Ausführung wird das Segment als schreibgeschützt markiert und eine einzelne Kopie des Segments wird unter allen Prozessen geteilt die das Programm ausführen. An der ersten 8K grenze über dem Programm Text Segment im Virtuellen Speicher, fängt der ``nonshared'' Bereich an, der nach Bedarf von Syscalls erweitert werden kann. Beginnend bei der höchsten Adresse im Virtuellen Speicher ist das Stack Segment, das Automatisch nach unten wächst während der Hardware Stackpointer sich ändert.
\end{quotation}
\end{framed}

Das erinnert daran wie manche Schüler Notizen zu  zwei Vorträgen in einem Notebook dokumentieren:
Notizen für den ersten Vortrag werden normal notiert, und Notizen zur zum zweiten Vortrag werden 
ans Ende des Notizbuches geschrieben, indem man das Notizbuch umdreht. Die Notizen treffen sich irgendwann
im Notizbuch aufgrund des fehlenden Freien Platzes.

% I think if we want to expand on this analogy,
% one might remember that the line number increases as as you go down a page.
% So when you decrease the address when pushing to the stack, visually,
% the stack does grow upwards.
% Of course, the problem is that in most human languages,
% just as with computers,
% we write downwards, so this direction is what makes buffer overflows so messy.

\subsection{Für was wird der Stack benutzt?}

% subsections
\input{patterns/02_stack/01_saving_ret_addr}
\input{patterns/02_stack/02_args_passing}
\EN{\input{patterns/02_stack/03_local_vars_EN}}
\RU{\input{patterns/02_stack/03_local_vars_RU}}
\DE{\input{patterns/02_stack/03_local_vars_DE}}
\PTBR{\input{patterns/02_stack/03_local_vars_PTBR}}
\input{patterns/02_stack/04_alloca/main}
\input{patterns/02_stack/05_SEH}
\input{patterns/02_stack/06_BO_protection}

\subsubsection{Automatisches deallokieren der Daten auf dem Stack}

Vielleicht ist der Grund warum man lokale Variablen und SEH Einträge auf dem Stack speichert, weil sie beim 
verlassen der Funktion automatisch aufgeräumt werden. Man braucht dabei nur eine Instruktion um die Position
des Stackpointers zu korrigieren (oftmals ist es die \ADD Instruktion). Funktions Argumente, könnte man sagen 
werden auch am Ende der Funktion deallokiert. Im Kontrast dazu, alles was auf dem \emph{heap} gespeichert wird muss
explizit deallokiert werden. 

% sections
\EN{\input{patterns/02_stack/07_layout_EN}}
\RU{\input{patterns/02_stack/07_layout_RU}}
\DE{\input{patterns/02_stack/07_layout_DE}}
\PTBR{\input{patterns/02_stack/07_layout_PTBR}}
\input{patterns/02_stack/08_noise/main}
\input{patterns/02_stack/exercises}
}
\FR{\mysection{x86}

\subsection{Terminologie}

Commun en 16-bit (8086/80286), 32-bit (80386, etc.), 64-bit.

\myindex{IEEE 754}
\myindex{MS-DOS}
\begin{description}
	\item[octet] 8-bit.
		La directive d'assembleur DB est utilisée pour définir les variables et les tableaux d'octets.
		Les octets sont passés dans les parties 8-bit des registres: \TT{AL/BL/CL/DL/AH/BH/CH/DH/SIL/DIL/R*L}.
	\item[mot] 16-bit. 
		directive assembleur DW \dittoclosing.
		Les mots sont passés dans la partie 16-bit des registres:\\
			\TT{AX/BX/CX/DX/SI/DI/R*W}.
	\item[double mot] (\q{dword}) 32-bit.
		directive assembleur DD \dittoclosing.
		Les double mots sont passés dans les registres (x86) ou dans la partie 32-bit de registres (x64).
		Dans du code 16-bit, les double mots sont passés dans une paire de registres 16-bit.
	\item[quadruple mot] (\q{qword}) 64-bit.
		directive assembleur DQ \dittoclosing.
		En environnement 32-bit, les quadruple mots sont passés dans une paire de registres 32-bit.
	\item[tbyte] (10 bytes) 80-bit ou 10 octets (utilisé pour les registres FPU IEEE 754).
	\item[paragraph] (16 bytes)---le terme était répandu dans l'environnement MS-DOS.
\end{description}

\myindex{Windows!API}

Des types de données de même taille (BYTE, WORD, DWORD) existent aussi dans l'\ac{API} Windows.

\input{appendix/x86/registers_FR} % subsection
\input{appendix/x86/instructions_FR} % subsection
\input{appendix/x86/npad} % subsection

}
\IT{\mysection{x86}

\subsection{Terminologia}

Comune per 16-bit (8086/80286), 32-bit (80386, etc.), 64-bit.

\myindex{IEEE 754}
\myindex{MS-DOS}
\begin{description}
	\item[byte] 8-bit.
		La direttiva DB in assembly è utilizzata per definire variabili ed array di byte.
		I byte sono passati nella parte ad 8-bit dei registri: \TT{AL/BL/CL/DL/AH/BH/CH/DH/SIL/DIL/R*L}.
	\item[word] 16-bit.
		La direttiva DW in assembly \dittoclosing.
		Le Word sono passate nella parte a 16-bit dei registri:\\
			\TT{AX/BX/CX/DX/SI/DI/R*W}.
	\item[double word] (\q{dword}) 32-bit.
		La direttiva DD in assembly \dittoclosing.
		Le Double words sono passate nei registri (x86) o nelle parti a 32-bit dei registri (x64).
		Nel codice a 16-bit, le double words sono passate in coppie di registri a 16-bit.
	\item[quad word] (\q{qword}) 64-bit.
		La direttiva DQ in assembly \dittoclosing.
		Negli ambienti a 32-bit, le quad words sono passate in coppie di registri a 32-bit.
	\item[tbyte] (10 bytes) 80-bit o 10 bytes (usati per i registri FPU IEEE 754).
	\item[paragraph] (16 bytes)---termine popolare nell'ambiente MS-DOS. % TODO link to a part about 8086 memory model...
\end{description}

\myindex{Windows!API}

Tipi di dati della stessa dimensione (BYTE, WORD, DWORD) sono gli stessi nelle Windows \ac{API}.

%TBT
\input{appendix/x86/registers} % subsection
\input{appendix/x86/instructions} % subsection
\input{appendix/x86/npad} % subsection
}
\mysection{ARM}
\myindex{ARM}

\subsection{\RU{Терминология}\EN{Terminology}}

ARM \RU{изначально разрабатывался как 32-битный}\EN{was initially developed as 32-bit} \ac{CPU}, 
\RU{поэтому \emph{слово} здесь, в отличие от x86, 32-битное}\EN{so that's why a \emph{word} here, unlike x86, is 32-bit}.

\begin{description}
	\item[byte] 8-\bitENRU.
		\RU{Для определения переменных и массива байт используется директива ассемблера DCB}
		\EN{The DB assembly directive is used for defining variables and arrays of bytes}.
	\item[halfword] 16-\bitENRU. \RU{\dittoclosing директива ассемблера DCW}
					\EN{DCW assembly directive \dittoclosing}.	
	\item[word] 32-\bitENRU. \RU{\dittoclosing директива ассемблера DCD}
					\EN{DCD assembly directive \dittoclosing}.
	\item[doubleword] 64-\bitENRU.
	\item[quadword] 128-\bitENRU.
\end{description}

\subsection{\RU{Версии}\EN{Versions}}

\begin{itemize}
\item ARMv4: \RU{появился режим Thumb}\EN{Thumb mode introduced}.

\item ARMv6: \RU{использовался в}\EN{used in} iPhone 1st gen., iPhone 3G 
(Samsung 32-bit RISC ARM 1176JZ(F)-S \RU{поддерживающий}\EN{that supports} Thumb-2)

\item ARMv7: \RU{появился }Thumb-2\EN{ was added} (2003).
\RU{Использовался в}\EN{was used in} iPhone 3GS, iPhone 4, iPad 1st gen. (ARM Cortex-A8), iPad 2 (Cortex-A9),
iPad 3rd gen.

\item ARMv7s: \RU{Добавлены новые инструкции}\EN{New instructions added}.
\RU{Использовался в}\EN{Was used in} iPhone 5, iPhone 5c, iPad 4th gen. (Apple A6).

\item ARMv8: 64-\RU{битный процессор}\EN{bit CPU}, \ac{AKA} ARM64 \ac{AKA} AArch64.
\RU{Использовался в}\EN{Was used in} iPhone 5S, iPad Air (Apple A7).
\RU{В 64-битном режиме, режима Thumb больше нет, только режим ARM (4-байтные инструкции).}
\EN{There is no Thumb mode in 64-bit mode, only ARM (4-byte instructions).}
\end{itemize}

% sections
\input{appendix/ARM/registers}
\input{appendix/ARM/instructions}

\mysection{MIPS}

\subsection{\Registers}
\label{MIPS_registers_ref}

\myindex{MIPS!O32}
( \RU{Соглашение о вызовах O32}\EN{O32 calling convention} )

\subsubsection{\RU{Регистры общего пользования}\EN{General purpose registers} \ac{GPR}}

%\small
\begin{center}
\begin{tabular}{ | l | l | l | }
\hline
\HeaderColor \RU{Номер}\EN{Number} & 
\HeaderColor \RU{Псевдоимя}\EN{Pseudoname} & 
\HeaderColor \RU{Описание}\EN{Description} \\
\hline
\$0             & \$ZERO          & \RU{Всегда ноль. Запись в этот регистр это как}\EN{Always zero. Writing to this register is like} \ac{NOP}. \\
\hline
\$1             & \$AT            & \RU{Используется как временный регистр}\EN{Used as a temporary register} \\
                &                 & \RU{для ассемблерных макросов и псевдоинструкций}\EN{for assembly macros and pseudo instructions}. \\
\hline
\$2 \dots \$3   & \$V0 \dots \$V1 & \RU{Здесь возвращается результат функции}\EN{Function result is returned here}. \\
\hline
\$4 \dots \$7   & \$A0 \dots \$A3 & \RU{Аргументы функции}\EN{Function arguments}. \\
\hline
\$8 \dots \$15  & \$T0 \dots \$T7 & \RU{Используется для временных данных}\EN{Used for temporary data}. \\
\hline
\$16 \dots \$23 & \$S0 \dots \$S7 & \RU{Используется для временных данных}\EN{Used for temporary data}\AsteriskOne{}. \\
\hline
\$24 \dots \$25 & \$T8 \dots \$T9 & \RU{Используется для временных данных}\EN{Used for temporary data}. \\
\hline
\$26 \dots \$27 & \$K0 \dots \$K1 & \RU{Зарезервировано для ядра \ac{OS}}\EN{Reserved for \ac{OS} kernel}. \\
\hline
\$28            & \$GP            & \RU{Глобальный указатель}\EN{Global Pointer}\AsteriskTwo{}. \\
\hline
\$29            & \$SP            & \ac{SP}\AsteriskOne{}. \\
\hline
\$30            & \$FP            & \ac{FP}\AsteriskOne{}. \\
\hline
\$31            & \$RA            & \ac{RA}. \\
\hline
n/a             & PC              & \ac{PC}. \\
\hline
n/a             & HI              & \RU{старшие 32 бита результата умножения или остаток от деления}\EN{high 32 bit of multiplication or division remainder}\AsteriskThree{}. \\
\hline
n/a             & LO              & \RU{младшие 32 бита результата умножения или результат деления}\EN{low 32 bit of multiplication and division remainder}\AsteriskThree{}. \\
\hline
\end{tabular}
\end{center}
%\normalsize

\subsubsection{\RU{Регистры для работы с числами с плавающей точкой}\EN{Floating-point registers}}
\label{MIPS_FPU_registers}

\begin{center}
\begin{tabular}{ | l | l | l | }
\hline
\HeaderColor \RU{Название}\EN{Name} & \HeaderColor \RU{Описание}\EN{Description} \\
\hline
\$F0..\$F1   & \RU{Здесь возвращается результат функции}\EN{Function result returned here}. \\
\hline
\$F2..\$F3   & \RU{Не используется}\EN{Not used}. \\
\hline
\$F4..\$F11  & \RU{Используется для временных данных}\EN{Used for temporary data}. \\
\hline
\$F12..\$F15 & \RU{Первые два аргумента функции}\EN{First two function arguments}. \\
\hline
\$F16..\$F19 & \RU{Используется для временных данных}\EN{Used for temporary data}. \\
\hline
\$F20..\$F31 & \RU{Используется для временных данных}\EN{Used for temporary data}\AsteriskOne{}. \\
\hline
%fcr31 & Control/status register. \\
%\hline
\end{tabular}
\end{center}

\AsteriskOne{} --- \Gls{callee} \RU{должен сохранять}\EN{must preserve the value}.\\
\AsteriskTwo{} --- \Gls{callee} \RU{должен сохранять}\EN{must preserve the value} (\RU{кроме \ac{PIC}-кода}
\EN{except in \ac{PIC} code}).\\
\myindex{MIPS!\Instructions!MFLO}
\myindex{MIPS!\Instructions!MFHI}
\AsteriskThree{} --- \RU{доступны используя инструкции}\EN{accessible using the} \TT{MFHI} \AndENRU \TT{MFLO}\EN{ instructions}.\\

\subsection{\Instructions}

\RU{Есть три типа инструкций}\EN{There are 3 kinds of instructions}:

\begin{itemize}

\item \RU{Тип R: имеющие 3 регистра. R-инструкции обычно имеют такой вид:}
\EN{R-type: those which have 3 registers. R-instruction usually have the following form:}

\begin{lstlisting}
instruction destination, source1, source2
\end{lstlisting}

\RU{Важно помнить, что если первый и второй регистр один и тот же, IDA может показать инструкцию в сокращенной
форме:}
\EN{One important thing to keep in mind is that when the first and second register are the same, 
IDA may show the instruction in its shorter form:}

\begin{lstlisting}
instruction destination/source1, source2
\end{lstlisting}

\RU{Это немного напоминает Интеловский синтаксис ассемблера x86.}
\EN{That somewhat reminds us of the Intel syntax for x86 assembly language.}

\item \RU{Тип I: имеющие 2 регистра и 16-битное \q{immediate}-значение.}
\EN{I-type: those which have 2 registers and a 16-bit immediate value.}

\item \RU{Тип J: инструкции перехода, имеют 26 бит для кодирования смещения.}
\EN{J-type: jump/branch instructions, have 26 bits for encoding the offset.}

\end{itemize}

\subsubsection{\RU{Инструкции перехода}\EN{Jump instructions}}

\RU{Какая разница между инструкциями начинающихся с B- (BEQ, B, итд) и с J- (JAL, JALR, итд)?}
\EN{What is the difference between B- instructions (BEQ, B, etc.) and J- ones (JAL, JALR, etc.)?}

\RU{B-инструкции имеют тип I, так что, смещение в этих инструкциях кодируется как 16-битное значение.}
\EN{The B-instructions have an I-type, hence, the B-instructions' offset is encoded as a 16-bit immediate.}
\RU{Инструкции JR и JALR имеют тип R, и они делают переход по абсолютному адресу указанному в регистре.}
\EN{JR and JALR are R-type and jump to an absolute address specified in a register.}
\RU{J и JAL имеют тип J, так что смещение кодируется как 26-битное значение.}
\EN{J and JAL are J-type, hence the offset is encoded as a 26-bit immediate.}

\RU{Коротко говоря, в B-инструкциях можно кодировать условие 
(B на самом деле это псевдоинструкция для \TT{BEQ \$ZERO, \$ZERO, LABEL}),
а в J-инструкциях нельзя.}
\EN{In short, B-instructions can encode a condition 
(B is in fact pseudo instruction for \TT{BEQ \$ZERO, \$ZERO, LABEL}), 
while J-instructions can't.}

\mysection{\RU{Некоторые библиотечные функции GCC}\EN{Some GCC library functions}\DE{Einige GCC-Bibliotheks-Funktionen}%
\FR{Quelques fonctions de la bibliothèque de GCC}}
\myindex{GCC}
\label{sec:GCC_library_func}

%__ashldi3
%__ashrdi3
%__floatundidf
%__floatdisf
%__floatdixf
%__floatundidf
%__floatundisf
%__floatundixf
%__lshrdi3
%__muldi3

\begin{center}
\begin{tabular}{ | l | l | }
\hline
\HeaderColor \RU{имя}\EN{name}\DE{Name}\FR{nom} & \HeaderColor \RU{значение}\EN{meaning}\DE{Bedeutung}%
\FR{signification} \\
\hline \TT{\_\_divdi3} & \RU{знаковое деление}\EN{signed division}\DE{vorzeichenbehaftete Division}%
\FR{division signée} \\
\hline \TT{\_\_moddi3} & \RU{остаток от знакового деления}\EN{getting remainder (modulo) of signed division}%
\DE{Rest (Modulo) einer vorzeichenbehafteten Division}%
\FR{reste (modulo) d'une division signée} \\
\hline \TT{\_\_udivdi3} & \RU{беззнаковое деление}\EN{unsigned division}\DE{vorzeichenlose Division}%
\FR{division non signée} \\
\hline \TT{\_\_umoddi3} & \RU{остаток от беззнакового деления}\EN{getting remainder (modulo) of unsigned division}%
\DE{Rest (Modulo) einer vorzeichenlosen Division}%
\FR{reste (modulo) d'une division non signée} \\
\hline
\end{tabular}
\end{center}



\mysection{\RU{Некоторые библиотечные функции MSVC}\EN{Some MSVC library functions}\DE{Einige MSVC-Bibliotheks-Funktionen}%
\FR{Quelques fonctions de la bibliothèque MSVC}}
\myindex{MSVC}
\label{sec:MSVC_library_func}

\TT{ll} \RU{в имени функции означает}\EN{in function name stands for}\DE{in Funktionsnamen steht für}%
\FR{dans une fontion signifie} \q{long long}, \RU{т.е. 64-битный тип данных}
\EN{e.g., a 64-bit data type}\DE{z.B. einen 64-Bit-Datentyp}\FR{i.e., type de donées 64-bit}.

\begin{center}
\begin{tabular}{ | l | l | }
\hline
\HeaderColor \RU{имя}\EN{name}\DE{Name}\FR{nom} & \HeaderColor \RU{значение}\EN{meaning}\DE{Bedeutung}\FR{signification} \\
\hline \TT{\_\_alldiv} & \RU{знаковое деление}\EN{signed division}\DE{vorzeichenbehaftete Division}\FR{division signée} \\
\hline \TT{\_\_allmul} & \RU{умножение}\EN{multiplication}\DE{Multiplikation}\FR{multiplication} \\
\hline \TT{\_\_allrem} & \RU{остаток от знакового деления}\EN{remainder of signed division}\DE{Rest einer vorzeichenbehafteten Division}%
\FR{reste de la division signée} \\
\hline \TT{\_\_allshl} & \RU{сдвиг влево}\EN{shift left}\DE{Schiebe links}\FR{décalage à gauche} \\
\hline \TT{\_\_allshr} & \RU{знаковый сдвиг вправо}\EN{signed shift right}\DE{Schiebe links, vorzeichenbehaftet}%
\FR{décalage signé à droite} \\
\hline \TT{\_\_aulldiv} & \RU{беззнаковое деление}\EN{unsigned division}\DE{vorzeichenlose Division}%
\FR{division non signée} \\
\hline \TT{\_\_aullrem} & \RU{остаток от беззнакового деления}\EN{remainder of unsigned division}\DE{Rest (Modulo) einer vorzeichenlosen Division}%
\FR{reste de la division non signée} \\
\hline \TT{\_\_aullshr} & \RU{беззнаковый сдвиг вправо}\EN{unsigned shift right}\DE{Schiebe rechts, vorzeichenlos}%
\FR{décalage non signé à droite} \\
\hline
\end{tabular}
\end{center}

\RU{Процедуры умножения и сдвига влево, одни и те же и для знаковых чисел, и для беззнаковых,
поэтому здесь только одна функция для каждой операции}
\EN{Multiplication and shift left procedures are the same for both signed and unsigned numbers, hence there is only one function 
for each operation here}
\DE{Multiplikation und Links-Schiebebefehle sind sowohl für vorzeichenbehaftete als auch vorzeichenlose Zahlen,
da hier für jede Operation nur ein Befehl existiert}
\FR{La multiplication et le décalage à gauche sont similaire pour les nombres signés
et non signés, donc il n'y a qu'une seule fonction ici}. \\
\\
\RU{Исходные коды этих функций можно найти в установленной \ac{MSVS}, в}\EN{The source code of these function
can be found in the installed \ac{MSVS}, in}%
\DE{Der Quellcode dieser Funktionen kann im Pfad des installierten \ac{MSVS}, gefunden werden: }%
\FR{Le code source des ces fonctions peut être trouvé dans l'installation de \ac{MSVS},
dans} \TT{VC/crt/src/intel/*.asm}.


\input{appendix/cheatsheets}
}
\IT{\part*{\RU{Приложение}\EN{Appendix}\DE{Anhang}\FR{Appendice}\IT{Appendice}}
\appendix
\addcontentsline{toc}{part}{\RU{Приложение}\EN{Appendix}\DE{Anhang}\FR{Appendice}\IT{Appendice}}

% chapters
\EN{\mysection{\MinesweeperWinXPExampleChapterName}
\label{minesweeper_winxp}
\myindex{Windows!Windows XP}

For those who are not very good at playing Minesweeper, we could try to reveal the hidden mines in the debugger.

\myindex{\CStandardLibrary!rand()}
\myindex{Windows!PDB}

As we know, Minesweeper places mines randomly, so there has to be some kind of random number generator or
a call to the standard \TT{rand()} C-function.

What is really cool about reversing Microsoft products is that there are \gls{PDB} 
file with symbols (function names, \etc{}).
When we load \TT{winmine.exe} into \IDA, it downloads the 
\gls{PDB} file exactly for this 
executable and shows all names.

So here it is, the only call to \TT{rand()} is this function:

\lstinputlisting[style=customasmx86]{examples/minesweeper/tmp1.lst}

\IDA named it so, and it was the name given to it by Minesweeper's developers.

The function is very simple:

\begin{lstlisting}[style=customc]
int Rnd(int limit)
{
    return rand() % limit;
};
\end{lstlisting}

(There is no \q{limit} name in the \gls{PDB} file; we manually named this argument like this.)

So it returns 
a random value from 0 to a specified limit.

\TT{Rnd()} is called only from one place, 
a function called \TT{StartGame()}, 
and as it seems, this is exactly 
the code which place the mines:

\begin{lstlisting}[style=customasmx86]
.text:010036C7                 push    _xBoxMac
.text:010036CD                 call    _Rnd@4          ; Rnd(x)
.text:010036D2                 push    _yBoxMac
.text:010036D8                 mov     esi, eax
.text:010036DA                 inc     esi
.text:010036DB                 call    _Rnd@4          ; Rnd(x)
.text:010036E0                 inc     eax
.text:010036E1                 mov     ecx, eax
.text:010036E3                 shl     ecx, 5          ; ECX=ECX*32
.text:010036E6                 test    _rgBlk[ecx+esi], 80h
.text:010036EE                 jnz     short loc_10036C7
.text:010036F0                 shl     eax, 5          ; EAX=EAX*32
.text:010036F3                 lea     eax, _rgBlk[eax+esi]
.text:010036FA                 or      byte ptr [eax], 80h
.text:010036FD                 dec     _cBombStart
.text:01003703                 jnz     short loc_10036C7
\end{lstlisting}

Minesweeper allows you to set the board size, so the X (xBoxMac) and Y (yBoxMac) of the board are global variables.
They are passed to \TT{Rnd()} and random 
coordinates are generated.
A mine is placed by the \TT{OR} instruction at \TT{0x010036FA}. 
And if it has been placed before 
(it's possible if the pair of \TT{Rnd()} 
generates a coordinates pair which has been already 
generated), 
then \TT{TEST} and \TT{JNZ} at \TT{0x010036E6} 
jumps to the generation routine again.

\TT{cBombStart} is the global variable containing total number of mines. So this is loop.

The width of the array is 32 
(we can conclude this by looking at the \TT{SHL} instruction, which multiplies one of the coordinates by 32).

The size of the \TT{rgBlk} 
global array can be easily determined by the difference 
between the \TT{rgBlk} 
label in the data segment and the next known one. 
It is 0x360 (864):

\begin{lstlisting}[style=customasmx86]
.data:01005340 _rgBlk          db 360h dup(?)          ; DATA XREF: MainWndProc(x,x,x,x)+574
.data:01005340                                         ; DisplayBlk(x,x)+23
.data:010056A0 _Preferences    dd ?                    ; DATA XREF: FixMenus()+2
...
\end{lstlisting}

$864/32=27$.

So the array size is $27*32$?
It is close to what we know: when we try to set board size to $100*100$ in Minesweeper settings, it fallbacks to a board of size $24*30$.
So this is the maximal board size here.
And the array has a fixed size for any board size.

So let's see all this in \olly.
We will ran Minesweeper, attaching \olly to it and now we can see the memory dump at the address of the \TT{rgBlk} array (\TT{0x01005340})
\footnote{All addresses here are for Minesweeper for Windows XP SP3 English. 
They may differ for other service packs.}.

So we got this memory dump of the array:

\lstinputlisting[style=customasmx86]{examples/minesweeper/1.lst}

\olly, like any other hexadecimal editor, shows 16 bytes per line.
So each 32-byte array row occupies exactly 2 lines here.

This is beginner level (9*9 board).

There is some square 
structure can be seen visually (0x10 bytes).

We will click \q{Run} in \olly to unfreeze the Minesweeper process, then we'll clicked randomly at the Minesweeper window 
and trapped into mine, but now all mines are visible:

\begin{figure}[H]
\centering
\myincludegraphicsSmall{examples/minesweeper/1.png}
\caption{Mines}
\label{fig:minesweeper1}
\end{figure}

By comparing the mine places and the dump, we can conclude that 0x10 stands for border, 0x0F---empty block, 0x8F---mine.
Perhaps, 0x10 is just a \emph{sentinel value}.

Now we'll add comments and also enclose all 0x8F bytes into square brackets:

\lstinputlisting[style=customasmx86]{examples/minesweeper/2.lst}

Now we'll remove all \emph{border bytes} (0x10) and what's beyond those:

\lstinputlisting[style=customasmx86]{examples/minesweeper/3.lst}

Yes, these are mines, now it can be clearly seen and compared with the screenshot.

\clearpage
What is interesting is that we can modify the array right in \olly.
We can remove all mines by changing all 0x8F bytes by 0x0F, and here is what we'll get in Minesweeper:

\begin{figure}[H]
\centering
\myincludegraphicsSmall{examples/minesweeper/3.png}
\caption{All mines are removed in debugger}
\label{fig:minesweeper3}
\end{figure}

We can also move all of them to the first line: 

\begin{figure}[H]
\centering
\myincludegraphicsSmall{examples/minesweeper/2.png}
\caption{Mines set in debugger}
\label{fig:minesweeper2}
\end{figure}

Well, the debugger is not very convenient for eavesdropping (which is our goal anyway), so we'll write a small utility
to dump the contents of the board:

\lstinputlisting[style=customc]{examples/minesweeper/minesweeper_cheater.c}

Just set the \ac{PID}
\footnote{PID it can be seen in Task Manager 
(enable it in \q{View $\rightarrow$ Select Columns})} 
and the address of the array (\TT{0x01005340} for Windows XP SP3 English) 
and it will dump it
\footnote{The compiled executable is here: 
\href{http://go.yurichev.com/17165}{beginners.re}}.

It attaches itself to a win32 process by \ac{PID} and just reads process memory at the address.

\subsection{Finding grid automatically}

This is kind of nuisance to set address each time when we run our utility.
Also, various Minesweeper versions may have the array on different address.
Knowing the fact that there is always a border (0x10 bytes), we can just find it in memory:

\lstinputlisting[style=customc]{examples/minesweeper/cheater2_fragment.c}

Full source code: \url{https://github.com/DennisYurichev/RE-for-beginners/blob/master/examples/minesweeper/minesweeper_cheater2.c}.

\subsection{\Exercises}

\begin{itemize}

\item 
Why do the \emph{border bytes} (or \emph{sentinel values}) (0x10) exist in the array?

What they are for if they are not visible in Minesweeper's interface?
How could it work without them?

\item 
As it turns out, there are more values possible (for open blocks, for flagged by user, \etc{}).
Try to find the meaning of each one.

\item 
Modify my utility so it can remove all mines or set them in a fixed pattern that you want in the Minesweeper
process currently running.

\end{itemize}
}
\RU{\mysection{Разгон майнера биткоинов Cointerra}
\index{Bitcoin}
\index{BeagleBone}

Был такой майнер биткоинов Cointerra, выглядящий так:

\begin{figure}[H]
\centering
\myincludegraphics{examples/bitcoin_miner/board.jpg}
\caption{Board}
\end{figure}

И была также (возможно утекшая) утилита\footnote{Можно скачать здесь: \url{https://github.com/DennisYurichev/RE-for-beginners/raw/master/examples/bitcoin_miner/files/cointool-overclock}}
которая могла выставлять тактовую частоту платы.
Она запускается на дополнительной плате BeagleBone на ARM с Linux (маленькая плата внизу фотографии).

И у автора (этих строк) однажды спросили, можно ли хакнуть эту утилиту и посмотреть, какие частоты можно выставлять, и какие нет.
И можно ли твикнуть её?

Утилиту нужно запускать так: \TT{./cointool-overclock 0 0 900}, где 900 это частота в МГц.
Если частота слишком большая, утилита выведет ошибку \q{Error with arguments} и закончит работу.

Вот фрагмент кода вокруг ссылки на текстовую строку \q{Error with arguments}:

\begin{lstlisting}[style=customasmARM]

...

.text:0000ABC4         STR      R3, [R11,#var_28]
.text:0000ABC8         MOV      R3, #optind
.text:0000ABD0         LDR      R3, [R3]
.text:0000ABD4         ADD      R3, R3, #1
.text:0000ABD8         MOV      R3, R3,LSL#2
.text:0000ABDC         LDR      R2, [R11,#argv]
.text:0000ABE0         ADD      R3, R2, R3
.text:0000ABE4         LDR      R3, [R3]
.text:0000ABE8         MOV      R0, R3  ; nptr
.text:0000ABEC         MOV      R1, #0  ; endptr
.text:0000ABF0         MOV      R2, #0  ; base
.text:0000ABF4         BL       strtoll
.text:0000ABF8         MOV      R2, R0
.text:0000ABFC         MOV      R3, R1
.text:0000AC00         MOV      R3, R2
.text:0000AC04         STR      R3, [R11,#var_2C]
.text:0000AC08         MOV      R3, #optind
.text:0000AC10         LDR      R3, [R3]
.text:0000AC14         ADD      R3, R3, #2
.text:0000AC18         MOV      R3, R3,LSL#2
.text:0000AC1C         LDR      R2, [R11,#argv]
.text:0000AC20         ADD      R3, R2, R3
.text:0000AC24         LDR      R3, [R3]
.text:0000AC28         MOV      R0, R3  ; nptr
.text:0000AC2C         MOV      R1, #0  ; endptr
.text:0000AC30         MOV      R2, #0  ; base
.text:0000AC34         BL       strtoll
.text:0000AC38         MOV      R2, R0
.text:0000AC3C         MOV      R3, R1
.text:0000AC40         MOV      R3, R2
.text:0000AC44         STR      R3, [R11,#third_argument]
.text:0000AC48         LDR      R3, [R11,#var_28]
.text:0000AC4C         CMP      R3, #0
.text:0000AC50         BLT      errors_with_arguments
.text:0000AC54         LDR      R3, [R11,#var_28]
.text:0000AC58         CMP      R3, #1
.text:0000AC5C         BGT      errors_with_arguments
.text:0000AC60         LDR      R3, [R11,#var_2C]
.text:0000AC64         CMP      R3, #0
.text:0000AC68         BLT      errors_with_arguments
.text:0000AC6C         LDR      R3, [R11,#var_2C]
.text:0000AC70         CMP      R3, #3
.text:0000AC74         BGT      errors_with_arguments
.text:0000AC78         LDR      R3, [R11,#third_argument]
.text:0000AC7C         CMP      R3, #0x31
.text:0000AC80         BLE      errors_with_arguments
.text:0000AC84         LDR      R2, [R11,#third_argument]
.text:0000AC88         MOV      R3, #950
.text:0000AC8C         CMP      R2, R3
.text:0000AC90         BGT      errors_with_arguments
.text:0000AC94         LDR      R2, [R11,#third_argument]
.text:0000AC98         MOV      R3, #0x51EB851F
.text:0000ACA0         SMULL    R1, R3, R3, R2
.text:0000ACA4         MOV      R1, R3,ASR#4
.text:0000ACA8         MOV      R3, R2,ASR#31
.text:0000ACAC         RSB      R3, R3, R1
.text:0000ACB0         MOV      R1, #50
.text:0000ACB4         MUL      R3, R1, R3
.text:0000ACB8         RSB      R3, R3, R2
.text:0000ACBC         CMP      R3, #0
.text:0000ACC0         BEQ      loc_ACEC
.text:0000ACC4
.text:0000ACC4 errors_with_arguments
.text:0000ACC4                                         
.text:0000ACC4         LDR      R3, [R11,#argv]
.text:0000ACC8         LDR      R3, [R3]
.text:0000ACCC         MOV      R0, R3  ; path
.text:0000ACD0         BL       __xpg_basename
.text:0000ACD4         MOV      R3, R0
.text:0000ACD8         MOV      R0, #aSErrorWithArgu ; format
.text:0000ACE0         MOV      R1, R3
.text:0000ACE4         BL       printf
.text:0000ACE8         B        loc_ADD4
.text:0000ACEC ; ------------------------------------------------------------
.text:0000ACEC
.text:0000ACEC loc_ACEC                 ; CODE XREF: main+66C
.text:0000ACEC         LDR      R2, [R11,#third_argument]
.text:0000ACF0         MOV      R3, #499
.text:0000ACF4         CMP      R2, R3
.text:0000ACF8         BGT      loc_AD08
.text:0000ACFC         MOV      R3, #0x64
.text:0000AD00         STR      R3, [R11,#unk_constant]
.text:0000AD04         B        jump_to_write_power
.text:0000AD08 ; ------------------------------------------------------------
.text:0000AD08
.text:0000AD08 loc_AD08                 ; CODE XREF: main+6A4
.text:0000AD08         LDR      R2, [R11,#third_argument]
.text:0000AD0C         MOV      R3, #799
.text:0000AD10         CMP      R2, R3
.text:0000AD14         BGT      loc_AD24
.text:0000AD18         MOV      R3, #0x5F
.text:0000AD1C         STR      R3, [R11,#unk_constant]
.text:0000AD20         B        jump_to_write_power
.text:0000AD24 ; ------------------------------------------------------------
.text:0000AD24
.text:0000AD24 loc_AD24                 ; CODE XREF: main+6C0
.text:0000AD24         LDR      R2, [R11,#third_argument]
.text:0000AD28         MOV      R3, #899
.text:0000AD2C         CMP      R2, R3
.text:0000AD30         BGT      loc_AD40
.text:0000AD34         MOV      R3, #0x5A
.text:0000AD38         STR      R3, [R11,#unk_constant]
.text:0000AD3C         B        jump_to_write_power
.text:0000AD40 ; ------------------------------------------------------------
.text:0000AD40
.text:0000AD40 loc_AD40                 ; CODE XREF: main+6DC
.text:0000AD40         LDR      R2, [R11,#third_argument]
.text:0000AD44         MOV      R3, #999
.text:0000AD48         CMP      R2, R3
.text:0000AD4C         BGT      loc_AD5C
.text:0000AD50         MOV      R3, #0x55
.text:0000AD54         STR      R3, [R11,#unk_constant]
.text:0000AD58         B        jump_to_write_power
.text:0000AD5C ; ------------------------------------------------------------
.text:0000AD5C
.text:0000AD5C loc_AD5C                 ; CODE XREF: main+6F8
.text:0000AD5C         LDR      R2, [R11,#third_argument]
.text:0000AD60         MOV      R3, #1099
.text:0000AD64         CMP      R2, R3
.text:0000AD68         BGT      jump_to_write_power
.text:0000AD6C         MOV      R3, #0x50
.text:0000AD70         STR      R3, [R11,#unk_constant]
.text:0000AD74
.text:0000AD74 jump_to_write_power                     ; CODE XREF: main+6B0
.text:0000AD74                                         ; main+6CC ...
.text:0000AD74         LDR      R3, [R11,#var_28]
.text:0000AD78         UXTB     R1, R3
.text:0000AD7C         LDR      R3, [R11,#var_2C]
.text:0000AD80         UXTB     R2, R3
.text:0000AD84         LDR      R3, [R11,#unk_constant]
.text:0000AD88         UXTB     R3, R3
.text:0000AD8C         LDR      R0, [R11,#third_argument]
.text:0000AD90         UXTH     R0, R0
.text:0000AD94         STR      R0, [SP,#0x44+var_44]
.text:0000AD98         LDR      R0, [R11,#var_24]
.text:0000AD9C         BL       write_power
.text:0000ADA0         LDR      R0, [R11,#var_24]
.text:0000ADA4         MOV      R1, #0x5A
.text:0000ADA8         BL       read_loop
.text:0000ADAC         B        loc_ADD4

...

.rodata:0000B378 aSErrorWithArgu DCB "%s: Error with arguments",0xA,0 ; DATA XREF: main+684

...

\end{lstlisting}

Имена ф-ций присутствовали в отладочной информации в оригинальном исполняемом файле,
такие как \TT{write\_power}, \TT{read\_loop}.
Но имена меткам внутри ф-ции дал я.

\myindex{UNIX!getopt}
\myindex{strtoll()}
Имя \TT{optind} звучит знакомо. Это библиотека \IT{getopt} из *NIX предназначенная для парсинга командной строки ---
и это то, что внутри и происходит.
Затем, третий аргумент (где передается значение частоты) конвертируется из строку в число используя вызов ф-ции \IT{strtoll()}.

Значение затем сравнивается с разными константами.
На 0xACEC есть проверка, меньше ли оно или равно 499, и если это так, то 0x64 будет передано в ф-цию
\TT{write\_power()} (которая посылает команду через USB используя \TT{send\_msg()}).
Если значение больше 499, происходит переход на 0xAD08.

На 0xAD08 есть проверка, меньше ли оно или равно 799. Если это так, то 0x5F передается в ф-цию \TT{write\_power()}.

Есть еще проверки: на 899 на 0xAD24, на 0x999 на 0xAD40, и наконец, на 1099 на 0xAD5C.
Если входная частота меньше или равна 1099, 0x50 (на 0xAD6C) будет передано в ф-цию \TT{write\_power()}.
И тут что-то вроде баги.
Если значение все еще больше 1099, само значение будет передано в ф-цию \TT{write\_power()}.
Но с другой стороны это не бага, потому что мы не можем попасть сюда: значение в начале проверяется с 950 на 0xAC88,
и если оно больше, выводится сообщение об ошибке и утилита заканчивает работу.

Вот таблица между частотами в МГц и значениями передаваемыми в ф-цию \TT{write\_power()}:

\begin{center}
\begin{longtable}{ | l | l | l | }
\hline
\HeaderColor МГц & \HeaderColor шестнадцатеричное представление & \HeaderColor десятичное \\
\hline
499MHz & 0x64 & 100 \\
\hline
799MHz & 0x5f & 95 \\
\hline
899MHz & 0x5a & 90 \\
\hline
999MHz & 0x55 & 85 \\
\hline
1099MHz & 0x50 & 80 \\
\hline
\end{longtable}
\end{center}

Как видно, значение передаваемое в плату постепенно уменьшается с ростом частоты.

Видно что значение в 950МГц это жесткий предел, по крайней мере в этой утилите. Можно ли её обмануть?

Вернемся к этому фрагменту кода:

\begin{lstlisting}[style=customasmARM]
.text:0000AC84      LDR     R2, [R11,#third_argument]
.text:0000AC88      MOV     R3, #950
.text:0000AC8C      CMP     R2, R3
.text:0000AC90      BGT     errors_with_arguments ; Я пропатчил здесь на 00 00 00 00
\end{lstlisting}

Нам нужно как-то запретить инструкцию перехода \INS{BGT} на 0xAC90. И это ARM в режиме ARM, потому что, как мы видим,
все адреса увеличиваются на 4, т.е., длина каждой инструкции это 4 байта.
Инструкция \TT{NOP} (нет операции) в режиме ARM это просто 4 нулевых байта: \TT{00 00 00 00}.
Так что, записывая 4 нуля по адресу 0xAC90 (или по физическому смещению в файле: 0x2C90) мы можем выключить
эту проверку.

Теперь можно выставлять частоты вплоть до 1050МГц. И даже больше, но из-за ошибки, если входное значение больше 1099,
значение в МГц, \IT{как есть}, будет передано в плату, что неправильно.

Дальше я не разбирался, но если бы продолжил, я бы уменьшал значение передаваемое в ф-цию \TT{write\_power()}.

Теперь страшный фрагмент кода, который я в начале пропустил:

\lstinputlisting[style=customasmARM]{examples/bitcoin_miner/tmp1.lst}

Здесь используется деление через умножение, и константа 0x51EB851F.
Я написал для себя простой программистский калькулятор\footnote{\url{https://github.com/DennisYurichev/progcalc}}.
И там есть возможность вычислять обратное число по модулю.

\begin{lstlisting}
modinv32(0x51EB851F)
Warning, result is not integer: 3.125000
(unsigned) dec: 3 hex: 0x3 bin: 11
\end{lstlisting}

Это значит что инструкция \INS{SMULL} на 0xACA0 просто делит 3-й аргумент на 3.125.
На самом деле, все что делает ф-ция \TT{modinv32()} в моем калькуляторе, это:

\[
\frac{1}{\frac{input}{2^{32}}} = \frac{2^{32}}{input}
\]

Потом там есть дополнительные сдвиги и теперь мы видим что 3-й аргумент просто делится на 50.
И затем умножается снова на 50.
Зачем?
Это простейшая проверка, можно ли делить входное значение на 50 без остатка.
Если значение этого выражения ненулевое, $x$ не может быть разделено на 50 без остатка:

\[
x-((\frac{x}{50}) \cdot 50)
\]

На самом деле, это простой способ вычисления остатка от деления.

И затем, если остаток ненулевой, выводится сообщение об ошибке.
Так что эта утилита берет значения частотв вроде 850, 900, 950, 1000, итд, но не 855 или 911.

Вот и всё! Если вы делаете что-то такое, имейте ввиду, что это может испортить вашу плату, как и в случае разгона
чипов вроде \ac{CPU}, \ac{GPU}, итд.
Если у вас есть плата Cointerra, делайте всё это на свой собственный риск!

}
\DE{\mysection{\Stack}
\label{sec:stack}
\myindex{\Stack}

Der Stack ist eine der fundamentalen Datenstrukturen in der Informatik.
\footnote{\href{http://go.yurichev.com/17119}{wikipedia.org/wiki/Call\_Stack}}.
\ac{AKA} \ac{LIFO}.

Technisch betrachtet ist es ein Stapelspeicher innerhalb des Prozessspeichers der zusammen mit den \ESP (x86), \RSP (x64) oder dem \ac{SP} (ARM) Register als ein Zeiger in diesem Speicherblock fungiert.

\myindex{ARM!\Instructions!PUSH}
\myindex{ARM!\Instructions!POP}
\myindex{x86!\Instructions!PUSH}
\myindex{x86!\Instructions!POP}

Die häufigsten Stack-Zugriffsinstruktionen sind die \PUSH- und \POP-Instruktionen (in beidem x86 und ARM Thumb-Modus). \PUSH subtrahiert vom \ESP/\RSP/\ac{SP} 4 Byte im 32-Bit Modus (oder 8 im 64-Bit Modus) und schreibt dann den Inhalt des Zeigers an die Adresse auf die von \ESP/\RSP/\ac{SP} gezeigt wird.

\POP ist die umgekehrte Operation: Die Daten des Zeigers für die Speicherregion auf die von \ac{SP}
gezeigt wird werden ausgelesen und die Inhalte in den Instruktionsoperanden geschreiben (oft ist das ein Register). Dann werden 4 (beziehungsweise 8) Byte zum \gls{stack pointer} addiert.

Nach der Stackallokation, zeigt der \gls{stack pointer} auf den Boden des Stacks.
\PUSH verringert den \gls{stack pointer} und \POP erhöht ihn.
Der Boden des Stacks ist eigentlich der Anfang der Speicherregion die für den Stack reserviert wurde.
Das wirkt zunächst seltsam, aber so funktioniert es.

ARM unterstützt beides, aufsteigende und absteigende Stacks.

\myindex{ARM!\Instructions!STMFD}
\myindex{ARM!\Instructions!LDMFD}
\myindex{ARM!\Instructions!STMED}
\myindex{ARM!\Instructions!LDMED}
\myindex{ARM!\Instructions!STMFA}
\myindex{ARM!\Instructions!LDMFA}
\myindex{ARM!\Instructions!STMEA}
\myindex{ARM!\Instructions!LDMEA}

Zum Beispiel die \ac{STMFD}/\ac{LDMFD} und \ac{STMED}/\ac{LDMED} Instruktionen sind alle dafür gedacht mit einem absteigendem Stack zu arbeiten ( wächst nach unten, fängt mit hohen Adressen an und entwickelt sich zu niedrigeren Adressen). Die \ac{STMFA}/\ac{LDMFA} und \ac{STMEA}/\ac{LDMEA} Instruktionen sind dazu gedacht mit einem aufsteigendem Stack zu arbeiten (wächst nach oben und fängt mit niedrigeren Adressen an und wächst nach oben).

% It might be worth mentioning that STMED and STMEA write first,
% and then move the pointer, and that LDMED and LDMEA move the pointer first, and then read.
% In other words, ARM not only lets the stack grow in a non-standard direction,
% but also in a non-standard order.
% Maybe this can be in the glossary, which would explain why E stands for "empty".

\subsection{Warum wächst der Stack nach unten?}
\label{stack_grow_backwards}

Intuitiv, würden man annehmen das der Stack nach oben wächst z.B Richtung höherer Adressen, so wie bei jeder anderen Datenstruktur.

Der Grund das der Stack rückwärts wächst ist wohl historisch bedingt. Als Computer so groß waren das sie einen ganzen Raum beansprucht haben war es einfach Speicher in zwei Sektionen zu unterteilen, einen Teil für den \gls{heap} und einen Teil für den Stack. Sicher war zu dieser Zeit nicht bekannt wie groß der \gls{heap} und der Stack wachsen würden, während der Programm Laufzeit, also war die Lösung die einfachste mögliche.

\input{patterns/02_stack/stack_and_heap}

In \RitchieThompsonUNIX können wir folgendes lesen:

\begin{framed}
\begin{quotation}
Der user-core eines Programm Images wird in drei logische Segmente unterteilt. Das Programm-Text Segment beginnt bei 0 im virtuellen Adress Speicher. Während der Ausführung wird das Segment als schreibgeschützt markiert und eine einzelne Kopie des Segments wird unter allen Prozessen geteilt die das Programm ausführen. An der ersten 8K grenze über dem Programm Text Segment im Virtuellen Speicher, fängt der ``nonshared'' Bereich an, der nach Bedarf von Syscalls erweitert werden kann. Beginnend bei der höchsten Adresse im Virtuellen Speicher ist das Stack Segment, das Automatisch nach unten wächst während der Hardware Stackpointer sich ändert.
\end{quotation}
\end{framed}

Das erinnert daran wie manche Schüler Notizen zu  zwei Vorträgen in einem Notebook dokumentieren:
Notizen für den ersten Vortrag werden normal notiert, und Notizen zur zum zweiten Vortrag werden 
ans Ende des Notizbuches geschrieben, indem man das Notizbuch umdreht. Die Notizen treffen sich irgendwann
im Notizbuch aufgrund des fehlenden Freien Platzes.

% I think if we want to expand on this analogy,
% one might remember that the line number increases as as you go down a page.
% So when you decrease the address when pushing to the stack, visually,
% the stack does grow upwards.
% Of course, the problem is that in most human languages,
% just as with computers,
% we write downwards, so this direction is what makes buffer overflows so messy.

\subsection{Für was wird der Stack benutzt?}

% subsections
\input{patterns/02_stack/01_saving_ret_addr}
\input{patterns/02_stack/02_args_passing}
\EN{\input{patterns/02_stack/03_local_vars_EN}}
\RU{\input{patterns/02_stack/03_local_vars_RU}}
\DE{\input{patterns/02_stack/03_local_vars_DE}}
\PTBR{\input{patterns/02_stack/03_local_vars_PTBR}}
\input{patterns/02_stack/04_alloca/main}
\input{patterns/02_stack/05_SEH}
\input{patterns/02_stack/06_BO_protection}

\subsubsection{Automatisches deallokieren der Daten auf dem Stack}

Vielleicht ist der Grund warum man lokale Variablen und SEH Einträge auf dem Stack speichert, weil sie beim 
verlassen der Funktion automatisch aufgeräumt werden. Man braucht dabei nur eine Instruktion um die Position
des Stackpointers zu korrigieren (oftmals ist es die \ADD Instruktion). Funktions Argumente, könnte man sagen 
werden auch am Ende der Funktion deallokiert. Im Kontrast dazu, alles was auf dem \emph{heap} gespeichert wird muss
explizit deallokiert werden. 

% sections
\EN{\input{patterns/02_stack/07_layout_EN}}
\RU{\input{patterns/02_stack/07_layout_RU}}
\DE{\input{patterns/02_stack/07_layout_DE}}
\PTBR{\input{patterns/02_stack/07_layout_PTBR}}
\input{patterns/02_stack/08_noise/main}
\input{patterns/02_stack/exercises}
}
\FR{\mysection{x86}

\subsection{Terminologie}

Commun en 16-bit (8086/80286), 32-bit (80386, etc.), 64-bit.

\myindex{IEEE 754}
\myindex{MS-DOS}
\begin{description}
	\item[octet] 8-bit.
		La directive d'assembleur DB est utilisée pour définir les variables et les tableaux d'octets.
		Les octets sont passés dans les parties 8-bit des registres: \TT{AL/BL/CL/DL/AH/BH/CH/DH/SIL/DIL/R*L}.
	\item[mot] 16-bit. 
		directive assembleur DW \dittoclosing.
		Les mots sont passés dans la partie 16-bit des registres:\\
			\TT{AX/BX/CX/DX/SI/DI/R*W}.
	\item[double mot] (\q{dword}) 32-bit.
		directive assembleur DD \dittoclosing.
		Les double mots sont passés dans les registres (x86) ou dans la partie 32-bit de registres (x64).
		Dans du code 16-bit, les double mots sont passés dans une paire de registres 16-bit.
	\item[quadruple mot] (\q{qword}) 64-bit.
		directive assembleur DQ \dittoclosing.
		En environnement 32-bit, les quadruple mots sont passés dans une paire de registres 32-bit.
	\item[tbyte] (10 bytes) 80-bit ou 10 octets (utilisé pour les registres FPU IEEE 754).
	\item[paragraph] (16 bytes)---le terme était répandu dans l'environnement MS-DOS.
\end{description}

\myindex{Windows!API}

Des types de données de même taille (BYTE, WORD, DWORD) existent aussi dans l'\ac{API} Windows.

\input{appendix/x86/registers_FR} % subsection
\input{appendix/x86/instructions_FR} % subsection
\input{appendix/x86/npad} % subsection

}
\IT{\mysection{x86}

\subsection{Terminologia}

Comune per 16-bit (8086/80286), 32-bit (80386, etc.), 64-bit.

\myindex{IEEE 754}
\myindex{MS-DOS}
\begin{description}
	\item[byte] 8-bit.
		La direttiva DB in assembly è utilizzata per definire variabili ed array di byte.
		I byte sono passati nella parte ad 8-bit dei registri: \TT{AL/BL/CL/DL/AH/BH/CH/DH/SIL/DIL/R*L}.
	\item[word] 16-bit.
		La direttiva DW in assembly \dittoclosing.
		Le Word sono passate nella parte a 16-bit dei registri:\\
			\TT{AX/BX/CX/DX/SI/DI/R*W}.
	\item[double word] (\q{dword}) 32-bit.
		La direttiva DD in assembly \dittoclosing.
		Le Double words sono passate nei registri (x86) o nelle parti a 32-bit dei registri (x64).
		Nel codice a 16-bit, le double words sono passate in coppie di registri a 16-bit.
	\item[quad word] (\q{qword}) 64-bit.
		La direttiva DQ in assembly \dittoclosing.
		Negli ambienti a 32-bit, le quad words sono passate in coppie di registri a 32-bit.
	\item[tbyte] (10 bytes) 80-bit o 10 bytes (usati per i registri FPU IEEE 754).
	\item[paragraph] (16 bytes)---termine popolare nell'ambiente MS-DOS. % TODO link to a part about 8086 memory model...
\end{description}

\myindex{Windows!API}

Tipi di dati della stessa dimensione (BYTE, WORD, DWORD) sono gli stessi nelle Windows \ac{API}.

%TBT
\input{appendix/x86/registers} % subsection
\input{appendix/x86/instructions} % subsection
\input{appendix/x86/npad} % subsection
}
\mysection{ARM}
\myindex{ARM}

\subsection{\RU{Терминология}\EN{Terminology}}

ARM \RU{изначально разрабатывался как 32-битный}\EN{was initially developed as 32-bit} \ac{CPU}, 
\RU{поэтому \emph{слово} здесь, в отличие от x86, 32-битное}\EN{so that's why a \emph{word} here, unlike x86, is 32-bit}.

\begin{description}
	\item[byte] 8-\bitENRU.
		\RU{Для определения переменных и массива байт используется директива ассемблера DCB}
		\EN{The DB assembly directive is used for defining variables and arrays of bytes}.
	\item[halfword] 16-\bitENRU. \RU{\dittoclosing директива ассемблера DCW}
					\EN{DCW assembly directive \dittoclosing}.	
	\item[word] 32-\bitENRU. \RU{\dittoclosing директива ассемблера DCD}
					\EN{DCD assembly directive \dittoclosing}.
	\item[doubleword] 64-\bitENRU.
	\item[quadword] 128-\bitENRU.
\end{description}

\subsection{\RU{Версии}\EN{Versions}}

\begin{itemize}
\item ARMv4: \RU{появился режим Thumb}\EN{Thumb mode introduced}.

\item ARMv6: \RU{использовался в}\EN{used in} iPhone 1st gen., iPhone 3G 
(Samsung 32-bit RISC ARM 1176JZ(F)-S \RU{поддерживающий}\EN{that supports} Thumb-2)

\item ARMv7: \RU{появился }Thumb-2\EN{ was added} (2003).
\RU{Использовался в}\EN{was used in} iPhone 3GS, iPhone 4, iPad 1st gen. (ARM Cortex-A8), iPad 2 (Cortex-A9),
iPad 3rd gen.

\item ARMv7s: \RU{Добавлены новые инструкции}\EN{New instructions added}.
\RU{Использовался в}\EN{Was used in} iPhone 5, iPhone 5c, iPad 4th gen. (Apple A6).

\item ARMv8: 64-\RU{битный процессор}\EN{bit CPU}, \ac{AKA} ARM64 \ac{AKA} AArch64.
\RU{Использовался в}\EN{Was used in} iPhone 5S, iPad Air (Apple A7).
\RU{В 64-битном режиме, режима Thumb больше нет, только режим ARM (4-байтные инструкции).}
\EN{There is no Thumb mode in 64-bit mode, only ARM (4-byte instructions).}
\end{itemize}

% sections
\input{appendix/ARM/registers}
\input{appendix/ARM/instructions}

\mysection{MIPS}

\subsection{\Registers}
\label{MIPS_registers_ref}

\myindex{MIPS!O32}
( \RU{Соглашение о вызовах O32}\EN{O32 calling convention} )

\subsubsection{\RU{Регистры общего пользования}\EN{General purpose registers} \ac{GPR}}

%\small
\begin{center}
\begin{tabular}{ | l | l | l | }
\hline
\HeaderColor \RU{Номер}\EN{Number} & 
\HeaderColor \RU{Псевдоимя}\EN{Pseudoname} & 
\HeaderColor \RU{Описание}\EN{Description} \\
\hline
\$0             & \$ZERO          & \RU{Всегда ноль. Запись в этот регистр это как}\EN{Always zero. Writing to this register is like} \ac{NOP}. \\
\hline
\$1             & \$AT            & \RU{Используется как временный регистр}\EN{Used as a temporary register} \\
                &                 & \RU{для ассемблерных макросов и псевдоинструкций}\EN{for assembly macros and pseudo instructions}. \\
\hline
\$2 \dots \$3   & \$V0 \dots \$V1 & \RU{Здесь возвращается результат функции}\EN{Function result is returned here}. \\
\hline
\$4 \dots \$7   & \$A0 \dots \$A3 & \RU{Аргументы функции}\EN{Function arguments}. \\
\hline
\$8 \dots \$15  & \$T0 \dots \$T7 & \RU{Используется для временных данных}\EN{Used for temporary data}. \\
\hline
\$16 \dots \$23 & \$S0 \dots \$S7 & \RU{Используется для временных данных}\EN{Used for temporary data}\AsteriskOne{}. \\
\hline
\$24 \dots \$25 & \$T8 \dots \$T9 & \RU{Используется для временных данных}\EN{Used for temporary data}. \\
\hline
\$26 \dots \$27 & \$K0 \dots \$K1 & \RU{Зарезервировано для ядра \ac{OS}}\EN{Reserved for \ac{OS} kernel}. \\
\hline
\$28            & \$GP            & \RU{Глобальный указатель}\EN{Global Pointer}\AsteriskTwo{}. \\
\hline
\$29            & \$SP            & \ac{SP}\AsteriskOne{}. \\
\hline
\$30            & \$FP            & \ac{FP}\AsteriskOne{}. \\
\hline
\$31            & \$RA            & \ac{RA}. \\
\hline
n/a             & PC              & \ac{PC}. \\
\hline
n/a             & HI              & \RU{старшие 32 бита результата умножения или остаток от деления}\EN{high 32 bit of multiplication or division remainder}\AsteriskThree{}. \\
\hline
n/a             & LO              & \RU{младшие 32 бита результата умножения или результат деления}\EN{low 32 bit of multiplication and division remainder}\AsteriskThree{}. \\
\hline
\end{tabular}
\end{center}
%\normalsize

\subsubsection{\RU{Регистры для работы с числами с плавающей точкой}\EN{Floating-point registers}}
\label{MIPS_FPU_registers}

\begin{center}
\begin{tabular}{ | l | l | l | }
\hline
\HeaderColor \RU{Название}\EN{Name} & \HeaderColor \RU{Описание}\EN{Description} \\
\hline
\$F0..\$F1   & \RU{Здесь возвращается результат функции}\EN{Function result returned here}. \\
\hline
\$F2..\$F3   & \RU{Не используется}\EN{Not used}. \\
\hline
\$F4..\$F11  & \RU{Используется для временных данных}\EN{Used for temporary data}. \\
\hline
\$F12..\$F15 & \RU{Первые два аргумента функции}\EN{First two function arguments}. \\
\hline
\$F16..\$F19 & \RU{Используется для временных данных}\EN{Used for temporary data}. \\
\hline
\$F20..\$F31 & \RU{Используется для временных данных}\EN{Used for temporary data}\AsteriskOne{}. \\
\hline
%fcr31 & Control/status register. \\
%\hline
\end{tabular}
\end{center}

\AsteriskOne{} --- \Gls{callee} \RU{должен сохранять}\EN{must preserve the value}.\\
\AsteriskTwo{} --- \Gls{callee} \RU{должен сохранять}\EN{must preserve the value} (\RU{кроме \ac{PIC}-кода}
\EN{except in \ac{PIC} code}).\\
\myindex{MIPS!\Instructions!MFLO}
\myindex{MIPS!\Instructions!MFHI}
\AsteriskThree{} --- \RU{доступны используя инструкции}\EN{accessible using the} \TT{MFHI} \AndENRU \TT{MFLO}\EN{ instructions}.\\

\subsection{\Instructions}

\RU{Есть три типа инструкций}\EN{There are 3 kinds of instructions}:

\begin{itemize}

\item \RU{Тип R: имеющие 3 регистра. R-инструкции обычно имеют такой вид:}
\EN{R-type: those which have 3 registers. R-instruction usually have the following form:}

\begin{lstlisting}
instruction destination, source1, source2
\end{lstlisting}

\RU{Важно помнить, что если первый и второй регистр один и тот же, IDA может показать инструкцию в сокращенной
форме:}
\EN{One important thing to keep in mind is that when the first and second register are the same, 
IDA may show the instruction in its shorter form:}

\begin{lstlisting}
instruction destination/source1, source2
\end{lstlisting}

\RU{Это немного напоминает Интеловский синтаксис ассемблера x86.}
\EN{That somewhat reminds us of the Intel syntax for x86 assembly language.}

\item \RU{Тип I: имеющие 2 регистра и 16-битное \q{immediate}-значение.}
\EN{I-type: those which have 2 registers and a 16-bit immediate value.}

\item \RU{Тип J: инструкции перехода, имеют 26 бит для кодирования смещения.}
\EN{J-type: jump/branch instructions, have 26 bits for encoding the offset.}

\end{itemize}

\subsubsection{\RU{Инструкции перехода}\EN{Jump instructions}}

\RU{Какая разница между инструкциями начинающихся с B- (BEQ, B, итд) и с J- (JAL, JALR, итд)?}
\EN{What is the difference between B- instructions (BEQ, B, etc.) and J- ones (JAL, JALR, etc.)?}

\RU{B-инструкции имеют тип I, так что, смещение в этих инструкциях кодируется как 16-битное значение.}
\EN{The B-instructions have an I-type, hence, the B-instructions' offset is encoded as a 16-bit immediate.}
\RU{Инструкции JR и JALR имеют тип R, и они делают переход по абсолютному адресу указанному в регистре.}
\EN{JR and JALR are R-type and jump to an absolute address specified in a register.}
\RU{J и JAL имеют тип J, так что смещение кодируется как 26-битное значение.}
\EN{J and JAL are J-type, hence the offset is encoded as a 26-bit immediate.}

\RU{Коротко говоря, в B-инструкциях можно кодировать условие 
(B на самом деле это псевдоинструкция для \TT{BEQ \$ZERO, \$ZERO, LABEL}),
а в J-инструкциях нельзя.}
\EN{In short, B-instructions can encode a condition 
(B is in fact pseudo instruction for \TT{BEQ \$ZERO, \$ZERO, LABEL}), 
while J-instructions can't.}

\mysection{\RU{Некоторые библиотечные функции GCC}\EN{Some GCC library functions}\DE{Einige GCC-Bibliotheks-Funktionen}%
\FR{Quelques fonctions de la bibliothèque de GCC}}
\myindex{GCC}
\label{sec:GCC_library_func}

%__ashldi3
%__ashrdi3
%__floatundidf
%__floatdisf
%__floatdixf
%__floatundidf
%__floatundisf
%__floatundixf
%__lshrdi3
%__muldi3

\begin{center}
\begin{tabular}{ | l | l | }
\hline
\HeaderColor \RU{имя}\EN{name}\DE{Name}\FR{nom} & \HeaderColor \RU{значение}\EN{meaning}\DE{Bedeutung}%
\FR{signification} \\
\hline \TT{\_\_divdi3} & \RU{знаковое деление}\EN{signed division}\DE{vorzeichenbehaftete Division}%
\FR{division signée} \\
\hline \TT{\_\_moddi3} & \RU{остаток от знакового деления}\EN{getting remainder (modulo) of signed division}%
\DE{Rest (Modulo) einer vorzeichenbehafteten Division}%
\FR{reste (modulo) d'une division signée} \\
\hline \TT{\_\_udivdi3} & \RU{беззнаковое деление}\EN{unsigned division}\DE{vorzeichenlose Division}%
\FR{division non signée} \\
\hline \TT{\_\_umoddi3} & \RU{остаток от беззнакового деления}\EN{getting remainder (modulo) of unsigned division}%
\DE{Rest (Modulo) einer vorzeichenlosen Division}%
\FR{reste (modulo) d'une division non signée} \\
\hline
\end{tabular}
\end{center}



\mysection{\RU{Некоторые библиотечные функции MSVC}\EN{Some MSVC library functions}\DE{Einige MSVC-Bibliotheks-Funktionen}%
\FR{Quelques fonctions de la bibliothèque MSVC}}
\myindex{MSVC}
\label{sec:MSVC_library_func}

\TT{ll} \RU{в имени функции означает}\EN{in function name stands for}\DE{in Funktionsnamen steht für}%
\FR{dans une fontion signifie} \q{long long}, \RU{т.е. 64-битный тип данных}
\EN{e.g., a 64-bit data type}\DE{z.B. einen 64-Bit-Datentyp}\FR{i.e., type de donées 64-bit}.

\begin{center}
\begin{tabular}{ | l | l | }
\hline
\HeaderColor \RU{имя}\EN{name}\DE{Name}\FR{nom} & \HeaderColor \RU{значение}\EN{meaning}\DE{Bedeutung}\FR{signification} \\
\hline \TT{\_\_alldiv} & \RU{знаковое деление}\EN{signed division}\DE{vorzeichenbehaftete Division}\FR{division signée} \\
\hline \TT{\_\_allmul} & \RU{умножение}\EN{multiplication}\DE{Multiplikation}\FR{multiplication} \\
\hline \TT{\_\_allrem} & \RU{остаток от знакового деления}\EN{remainder of signed division}\DE{Rest einer vorzeichenbehafteten Division}%
\FR{reste de la division signée} \\
\hline \TT{\_\_allshl} & \RU{сдвиг влево}\EN{shift left}\DE{Schiebe links}\FR{décalage à gauche} \\
\hline \TT{\_\_allshr} & \RU{знаковый сдвиг вправо}\EN{signed shift right}\DE{Schiebe links, vorzeichenbehaftet}%
\FR{décalage signé à droite} \\
\hline \TT{\_\_aulldiv} & \RU{беззнаковое деление}\EN{unsigned division}\DE{vorzeichenlose Division}%
\FR{division non signée} \\
\hline \TT{\_\_aullrem} & \RU{остаток от беззнакового деления}\EN{remainder of unsigned division}\DE{Rest (Modulo) einer vorzeichenlosen Division}%
\FR{reste de la division non signée} \\
\hline \TT{\_\_aullshr} & \RU{беззнаковый сдвиг вправо}\EN{unsigned shift right}\DE{Schiebe rechts, vorzeichenlos}%
\FR{décalage non signé à droite} \\
\hline
\end{tabular}
\end{center}

\RU{Процедуры умножения и сдвига влево, одни и те же и для знаковых чисел, и для беззнаковых,
поэтому здесь только одна функция для каждой операции}
\EN{Multiplication and shift left procedures are the same for both signed and unsigned numbers, hence there is only one function 
for each operation here}
\DE{Multiplikation und Links-Schiebebefehle sind sowohl für vorzeichenbehaftete als auch vorzeichenlose Zahlen,
da hier für jede Operation nur ein Befehl existiert}
\FR{La multiplication et le décalage à gauche sont similaire pour les nombres signés
et non signés, donc il n'y a qu'une seule fonction ici}. \\
\\
\RU{Исходные коды этих функций можно найти в установленной \ac{MSVS}, в}\EN{The source code of these function
can be found in the installed \ac{MSVS}, in}%
\DE{Der Quellcode dieser Funktionen kann im Pfad des installierten \ac{MSVS}, gefunden werden: }%
\FR{Le code source des ces fonctions peut être trouvé dans l'installation de \ac{MSVS},
dans} \TT{VC/crt/src/intel/*.asm}.


\input{appendix/cheatsheets}
}
% TODO split
\part*{\AcronymsUsed}

\addcontentsline{toc}{part}{\AcronymsUsed}

\begin{acronym}

\RU{	\acro{OS}[ОС]{Операционная Система}
	\acro{FAQ}[ЧаВО]{Часто задаваемые вопросы}
	\acro{OOP}[ООП]{Объектно-Ориентированное Программирование}
	\acro{PL}[ЯП]{Язык Программирования}
	\acro{PRNG}[ГПСЧ]{Генератор псевдослучайных чисел}
	\acro{ROM}[ПЗУ]{Постоянное запоминающее устройство}
	\acro{ALU}[АЛУ]{Арифметико-логическое устройство}
	\acro{PID}{ID программы/процесса}
	\acro{LF}{Line feed (подача строки) (10 или '\textbackslash{}n' в \CCpp)}
	\acro{CR}{Carriage return (возврат каретки) (13 или '\textbackslash{}r' в \CCpp)}
	\acro{LIFO}{Last In First Out (последним вошел, первым вышел)}
	\acro{MSB}{Most significant bit (самый старший бит)} % NOT BYTE!
	\acro{LSB}{Least significant bit (самый младший бит)} % NOT BYTE!
	\acro{NSA}[АНБ]{Агентство национальной безопасности}
	\acro{CFB}{Режим обратной связи по шифротексту (Cipher Feedback)}
	\acro{CSPRNG}{Криптографически стойкий генератор псевдослучайных чисел (cryptographically secure pseudorandom number generator)}
	\acro{SICP}{Структура и интерпретация компьютерных программ (Structure and Interpretation of Computer Programs)}
}
\EN{	\acro{OS}{Operating System}
	\acro{FAQ}{Frequently Asked Questions}
	\acro{OOP}{Object-Oriented Programming}
	\acro{PL}{Programming Language}
	\acro{PRNG}{Pseudorandom Number Generator}
	\acro{ROM}{Read-Only Memory}
	\acro{ALU}{Arithmetic Logic Unit}
	\acro{PID}{Program/process ID}
	\acro{LF}{Line Feed (10 or '\textbackslash{}n' in \CCpp)}
	\acro{CR}{Carriage Return (13 or '\textbackslash{}r' in \CCpp)}
	\acro{LIFO}{Last In First Out}
	\acro{MSB}{Most Significant Bit} % NOT BYTE!
	\acro{LSB}{Least Significant Bit} % NOT BYTE!
	\acro{NSA}{National Security Agency}
	\acro{CFB}{Cipher Feedback}
	\acro{CSPRNG}{Cryptographically Secure Pseudorandom Number Generator}
	\acro{SICP}{Structure and Interpretation of Computer Programs}
	\acro{ABI}{Application Binary Interface}
}
\ES{	\acro{OS}[SO]{Sistema Operativo}
	\acro{FAQ}{Preguntas Frecuentes}
	\acro{OOP}[POO]{Programaci\'on Orientada a Objetos}
	\acro{PL}[LP]{Lenguaje de Programaci\'on}
	\acro{PRNG}[GPAN]{Generador Pseudo-Aleatorio de N\'umeros}
	\acro{ROM}{Memoria de Solo Lectura}
	\acro{ALU}{Unidad Aritm\'etica L\'ogica}
}
\DE{	\acro{OS}[BS]{Betriebssystem}
	\acro{FAQ}{Häufig gestellte Fragen}
	\acro{OOP}{Objektorientierte Programmierung}
	\acro{PL}[PS]{Programmiersprache}
	\acro{PRNG}{Pseudozufallszahlen-Generator}
	\acro{ALU}[ALE]{Arithmetisch-logische Einheit}

}
\IT{	\acro{OS}{Sistema Operativo (Operating System)}
	\acro{FAQ}{Domande Frequente (Frequently Asked Questions)}
	\acro{OOP}{Programmazione ad oggetti (Object-Oriented Programming)}
	\acro{PL}{Linguaggio di programmazione (Programming Language)}
	\acro{PRNG}{Generatore di numeri pseudo-casuali (Pseudorandom Number Generator)}
	\acro{ROM}{Read-Only Memory}
	\acro{ALU}{Arithmetic Logic Unit}
	\acro{PID}{Program/process ID}
	\acro{LF}{Line Feed (10 o '\textbackslash{}n' in \CCpp)}
	\acro{CR}{Carriage Return (13 o '\textbackslash{}r' in \CCpp)}
	\acro{LIFO}{Last In First Out}
	\acro{MSB}{Most Significant Bit} % NOT BYTE!
	\acro{LSB}{Least Significant Bit} % NOT BYTE!
	\acro{NSA}{National Security Agency}
	\acro{CFB}{Cipher Feedback}
	\acro{CSPRNG}{Cryptographically Secure Pseudorandom Number Generator}
	\acro{SICP}{Structure and Interpretation of Computer Programs}
	\acro{ABI}{Application Binary Interface}
}
\NL{	\acro{OS}{\NL{Operating System}}
	\acro{FAQ}{\NL{Veelvoorkomende vragen}}
	\acro{OOP}{\NL{Object-Oriented Programmeren}}
	\acro{PL}[PT]{\NL{Programmeertaal}}
	\acro{PRNG}{\NL{Pseudorandom number generator}}
	\acro{ROM}{\NL{Read-only memory}}
	\acro{ALU}{\NL{Arithmetic logic unit}}
}
\FR{	\acro{OS}[OS]{Système d'exploitation (Operating System)}
	\acro{FAQ}{Foire Aux Questions}
	\acro{OOP}[POO]{Programmation orientée objet}
	\acro{PL}[LP]{Langage de programmation}
	\acro{PRNG}{Nombre généré pseudo-aléatoirement}
	\acro{ROM}{Mémoire morte}
	\acro{ALU}[UAL]{Unité arithmétique et logique}
	\acro{PID}{ID d'un processus}
	\acro{LF}{Line feed (10 ou '\textbackslash{}n' en \CCpp)}
	\acro{CR}{Carriage return (13 ou '\textbackslash{}r' en \CCpp)}
	\acro{LIFO}{Dernier entré, premier sorti}
	\acro{MSB}{Bit le plus significatif} % NOT BYTE!
	\acro{LSB}{Bit le moins significatif} % NOT BYTE!
	\acro{NSA}{National Security Agency (Agence Nationale de la Sécurité)} % translation not used in French
	\acro{CSPRNG}{Cryptographically Secure Pseudorandom Number Generator (générateur de nombres pseudo-aléatoire cryptographiquement sûr)}
	\acro{SICP}{Structure and Interpretation of Computer Programs}
	\acro{ABI}{Application Binary Interface}
}
\JA{	\acro{OS}{オペレーティングシステム}
	\acro{FAQ}{よくある質問}
	\acro{OOP}{オブジェクト指向プログラミング}
	\acro{PL}{プログラミング言語}
	\acro{PRNG}{擬似乱数生成器}
	\acro{ROM}{読み取り専用メモリ}
	\acro{ALU}{算術論理ユニット}
	\acro{PID}{プログラム／プロセスID}
	\acro{LF}{ラインフィード (\CCpp で10 または '\textbackslash{}n')}
	\acro{CR}{キャリッジリターン (\CCpp で13 または '\textbackslash{}r')}
	\acro{LIFO}{後入れ先出し}
	\acro{MSB}{最上位ビット} % NOT BYTE!
	\acro{LSB}{最下位ビット} % NOT BYTE!
	\acro{NSA}{国家安全保障局}
	\acro{CFB}{暗号フィードバック}
	\acro{CSPRNG}{暗号論的擬似乱数生成器}
	\acro{SICP}{計算機プログラムの構造と解釈}
	\acro{ABI}{アプリケーション・バイナリー・インタフェース}
}

\acro{RA}{\ReturnAddress}
\acro{PE}{Portable Executable}
\acro{SP}{\gls{stack pointer}. SP/ESP/RSP \InENRU x86/x64. SP \InENRU ARM.}
\acro{DLL}{Dynamic-Link Library}
\acro{PC}{Program Counter. IP/EIP/RIP \InENRU x86/64. PC \InENRU ARM.}
\acro{LR}{Link Register}
\acro{IDA}{
	\RU{Интерактивный дизассемблер и отладчик, разработан \href{https://hex-rays.com/}{Hex-Rays}}%
	\EN{Interactive Disassembler and Debugger developed by \href{https://hex-rays.com/}{Hex-Rays}}%
	\IT{\ac{TBT} by \href{https://hex-rays.com/}{Hex-Rays}}%
	\ES{Desensamblador Interactivo y depurador desarrollado por \href{https://hex-rays.com/}{Hex-Rays}}%
	\NL{Interactive Disassembler en debugger ontwikkeld door \href{https://hex-rays.com}{Hex-Rays}}
	\DE{Interaktiver Disassembler und Debugger entwickelt von \href{https://hex-rays.com/}{Hex-Rays}}%
	\FR{Désassembleur interactif et débogueur développé par \href{https://hex-rays.com/}{Hex-Rays}}%
	\JA{\href{https://hex-rays.com/}{Hex-Rays} によって開発されたインタラクティブなディスアセンブラ・デバッガ}%
}
\acro{IAT}{Import Address Table}
\acro{INT}{Import Name Table}
\acro{RVA}{Relative Virtual Address}
\acro{VA}{Virtual Address}
\acro{OEP}{Original Entry Point}
\acro{MSVC}{Microsoft Visual C++}
\acro{MSVS}{Microsoft Visual Studio}
\acro{ASLR}{Address Space Layout Randomization}
\acro{MFC}{Microsoft Foundation Classes}
\acro{TLS}{Thread Local Storage}
\acro{AKA}{
        \EN{Also Known As}%
	\FR{Aussi connu sous le nom de}%
	\RU{ - (Также известный как)}%
	\ES{ - (Tambi\'en Conocido Como)}%
	\NL{ - (Ook gekend als)}%
	\ITph{}
	\JA{別名}%
	\DEph{}
}
\acro{CRT}{C Runtime library}
\acro{CPU}{Central Processing Unit}
\acro{GPU}{Graphics Processing Unit}
\acro{FPU}{Floating-Point Unit}
\acro{CISC}{Complex Instruction Set Computing}
\acro{RISC}{Reduced Instruction Set Computing}
\acro{GUI}{Graphical User Interface}
\acro{RTTI}{Run-Time Type Information}
\acro{BSS}{Block Started by Symbol}
\acro{SIMD}{Single Instruction, Multiple Data}
\acro{BSOD}{Blue Screen of Death}
\acro{DBMS}{Database Management Systems}
\acro{ISA}{Instruction Set Architecture\RU{ (Архитектура набора команд)}}
\acro{CGI}{Common Gateway Interface}
\acro{HPC}{High-Performance Computing}
\acro{SOC}{System on Chip}
\acro{SEH}{Structured Exception Handling}
\acro{ELF}{\RU{Формат исполняемых файлов, использующийся в Linux и некоторых других *NIX}
\EN{Executable File format widely used in *NIX systems including Linux}
\FR{Format de fichier exécutable couramment utilisé sur les systèmes *NIX, Linux inclus}
\JA{Linuxを含め*NIXシステムで広く使用される実行ファイルフォーマット}
\ITph{}
\DEph{}}
\acro{TIB}{Thread Information Block}
\acro{TEA}{Tiny Encryption Algorithm}
\acro{PIC}{Position Independent Code}
\acro{NAN}{Not a Number}
\acro{NOP}{No Operation}
\acro{BEQ}{(PowerPC, ARM) Branch if Equal}
\acro{BNE}{(PowerPC, ARM) Branch if Not Equal}
\acro{BLR}{(PowerPC) Branch to Link Register}
\acro{XOR}{eXclusive OR\RU{ (исключающее \q{ИЛИ})}\FR{ (OU exclusif)}}
\acro{MCU}{Microcontroller Unit}
\acro{RAM}{Random-Access Memory}
\acro{GCC}{GNU Compiler Collection}
\acro{EGA}{Enhanced Graphics Adapter}
\acro{VGA}{Video Graphics Array}
\acro{API}{Application Programming Interface}
\acro{ASCII}{American Standard Code for Information Interchange}
\acro{ASCIIZ}{ASCII Zero (\RU{ASCII-строка заканчивающаяся нулем}\EN{null-terminated ASCII string}
\FR{chaîne ASCII terminée par un octet nul (à zéro)}\JA{ヌル終端文字列})}
\acro{IA64}{Intel Architecture 64 (Itanium)}
\acro{EPIC}{Explicitly Parallel Instruction Computing}
\acro{OOE}{Out-of-Order Execution}
\acro{MSDN}{Microsoft Developer Network}
\acro{STL}{(\Cpp) Standard Template Library}
\acro{PODT}{(\Cpp) Plain Old Data Type}
\acro{HDD}{Hard Disk Drive}
\acro{VM}{Virtual Memory\RU{ (виртуальная память)}\FR{ (mémoire virtuelle)}}
\acro{WRK}{Windows Research Kernel}
\acro{GPR}{General Purpose Registers\RU{ (регистры общего пользования)}}
\acro{SSDT}{System Service Dispatch Table}
\acro{RE}{Reverse Engineering}
\acro{RAID}{Redundant Array of Independent Disks}
\acro{SSE}{Streaming SIMD Extensions}
\acro{BCD}{Binary-Coded Decimal}
\acro{BOM}{Byte Order Mark}
\acro{GDB}{GNU Debugger}
\acro{FP}{Frame Pointer}
\acro{MBR}{Master Boot Record}
\acro{JPE}{Jump Parity Even (\RU{инструкция x86}\EN{x86 instruction}\FR{instruction x86}\JA{x86命令}\DEph{})}
\acro{CIDR}{Classless Inter-Domain Routing}
\acro{STMFD}{Store Multiple Full Descending (\DEph{}\RU{инструкция ARM}\EN{ARM instruction}\FR{instruction ARM}\JA{ARM命令})}
\acro{LDMFD}{Load Multiple Full Descending (\DEph{}\RU{инструкция ARM}\EN{ARM instruction}\FR{instruction ARM}\JA{ARM命令})}
\acro{STMED}{Store Multiple Empty Descending (\DEph{}\RU{инструкция ARM}\EN{ARM instruction}\FR{instruction ARM}\JA{ARM命令})}
\acro{LDMED}{Load Multiple Empty Descending (\DEph{}\RU{инструкция ARM}\EN{ARM instruction}\FR{instruction ARM}\JA{ARM命令})}
\acro{STMFA}{Store Multiple Full Ascending (\DEph{}\RU{инструкция ARM}\EN{ARM instruction}\FR{instruction ARM}\JA{ARM命令})}
\acro{LDMFA}{Load Multiple Full Ascending (\DEph{}\RU{инструкция ARM}\EN{ARM instruction}\FR{instruction ARM}\JA{ARM命令})}
\acro{STMEA}{Store Multiple Empty Ascending (\DEph{}\RU{инструкция ARM}\EN{ARM instruction}\FR{instruction ARM}\JA{ARM命令})}
\acro{LDMEA}{Load Multiple Empty Ascending (\DEph{}\RU{инструкция ARM}\EN{ARM instruction}\FR{instruction ARM}\JA{ARM命令})}
\acro{APSR}{(ARM) Application Program Status Register}
\acro{FPSCR}{(ARM) Floating-Point Status and Control Register}
\acro{RFC}{Request for Comments}
\acro{TOS}{Top of Stack\RU{ (вершина стека)}}
\acro{LVA}{(Java) Local Variable Array\RU{ (массив локальных переменных)}}
\acro{JVM}{Java Virtual Machine}
\acro{JIT}{Just-In-Time compilation}
\acro{CDFS}{Compact Disc File System}
\acro{CD}{Compact Disc}
\acro{ADC}{Analog-to-Digital Converter}
\acro{EOF}{End of File\DEph{}\RU{ (конец файла)}\FR{ (fin de fichier)}\JA{（ファイル終端）}}
\acro{TBT}{To be Translated. The presence of this acronym in this place means that the English version has some new/modified content which is to be translated and placed right here.}
\acro{DIY}{Do It Yourself}
\acro{MMU}{Memory Management Unit}
\acro{DES}{Data Encryption Standard}
\acro{MIME}{Multipurpose Internet Mail Extensions}
\acro{DBI}{Dynamic Binary Instrumentation}
\acro{XML}{Extensible Markup Language}
\acro{JSON}{JavaScript Object Notation}
\acro{URL}{Uniform Resource Locator}
\acro{ISP}{Internet Service Provider}
\acro{IV}{Initialization Vector}
\acro{RSA}{Rivest Shamir Adleman}
\acro{CPRNG}{Cryptographically secure PseudoRandom Number Generator}
\acro{GiB}{Gibibyte}
\acro{CRC}{Cyclic redundancy check}
\acro{AES}{Advanced Encryption Standard}
\acro{GC}{Garbage Collector}

\end{acronym}


\bookmarksetup{startatroot}

\clearpage
\phantomsection
\addcontentsline{toc}{chapter}{%
    \RU{Глоссарий}%
    \EN{Glossary}%
    \ES{Glosario}%
    \PTBRph{}%
    \DE{Glossar}%
    \PLph{}%
    \IT{Glossario}%
    \THAph{}\NLph{}%
    \FR{Glossaire}%
    \JA{用語}
    \TR{Bolum}
}
\printglossaries

\clearpage
\phantomsection
\printindex

\end{document}

\subsection{ARM}

\EN{\subsubsection{ARM: 3 arguments}

ARM's traditional scheme for passing arguments (calling convention) behaves as follows:
the first 4 arguments are passed through the \Reg{0}-\Reg{3} registers; the remaining arguments via the stack.
This resembles the arguments passing scheme in 
fastcall~(\myref{fastcall}) or win64~(\myref{sec:callingconventions_win64}).

\myparagraph{32-bit ARM}

\mysubparagraph{\NonOptimizingKeilVI (\ARMMode)}

\begin{lstlisting}[caption=\NonOptimizingKeilVI (\ARMMode),style=customasmARM]
.text:00000000 main
.text:00000000 10 40 2D E9   STMFD   SP!, {R4,LR}
.text:00000004 03 30 A0 E3   MOV     R3, #3
.text:00000008 02 20 A0 E3   MOV     R2, #2
.text:0000000C 01 10 A0 E3   MOV     R1, #1
.text:00000010 08 00 8F E2   ADR     R0, aADBDCD     ; "a=\%d; b=\%d; c=\%d"
.text:00000014 06 00 00 EB   BL      __2printf
.text:00000018 00 00 A0 E3   MOV     R0, #0          ; return 0
.text:0000001C 10 80 BD E8   LDMFD   SP!, {R4,PC}
\end{lstlisting}

So, the first 4 arguments are passed via the \Reg{0}-\Reg{3} registers in this order:
a pointer to the \printf format string in 
\Reg{0}, then 1 in \Reg{1}, 2 in \Reg{2} and 3 in \Reg{3}.
The instruction at \GTT{0x18} writes 0 to \Reg{0}---this is \IT{return 0} C-statement.
There is nothing unusual so far.

\OptimizingKeilVI generates the same code.

\mysubparagraph{\OptimizingKeilVI (\ThumbMode)}

\begin{lstlisting}[caption=\OptimizingKeilVI (\ThumbMode),style=customasmARM]
.text:00000000 main
.text:00000000 10 B5        PUSH    {R4,LR}
.text:00000002 03 23        MOVS    R3, #3
.text:00000004 02 22        MOVS    R2, #2
.text:00000006 01 21        MOVS    R1, #1
.text:00000008 02 A0        ADR     R0, aADBDCD     ; "a=\%d; b=\%d; c=\%d"
.text:0000000A 00 F0 0D F8  BL      __2printf
.text:0000000E 00 20        MOVS    R0, #0
.text:00000010 10 BD        POP     {R4,PC}
\end{lstlisting}

There is no significant difference from the non-optimized code for ARM mode.

\mysubparagraph{\OptimizingKeilVI (\ARMMode) + let's remove return}
\label{ARM_B_to_printf}

Let's rework example slightly by removing \IT{return 0}:

\begin{lstlisting}[style=customc]
#include <stdio.h>

void main()
{
	printf("a=%d; b=%d; c=%d", 1, 2, 3);
};
\end{lstlisting}

The result is somewhat unusual:

\lstinputlisting[caption=\OptimizingKeilVI (\ARMMode),style=customasmARM]{patterns/03_printf/ARM/tmp3.asm}

\myindex{ARM!\Registers!Link Register}
\myindex{ARM!\Instructions!B}
\myindex{Function epilogue}
This is the optimized (\Othree) version for ARM mode and this time we see \INS{B} as the last instruction instead of the familiar \INS{BL}.
Another difference between this optimized version and the previous one (compiled without optimization)
is the lack of function prologue and epilogue (instructions preserving the \Reg{0} and \ac{LR} registers values).
\myindex{x86!\Instructions!JMP}
The \INS{B} instruction just jumps to another address, without any manipulation of the \ac{LR} register,
similar to \JMP in x86.
Why does it work? Because this code is, in fact, effectively equivalent to the previous.
There are two main reasons: 1) neither the stack nor \ac{SP} (the \gls{stack pointer}) is modified;
2) the call to \printf is the last instruction, so there is nothing going on afterwards.
On completion, the \printf function simply returns the control to the address 
stored in \ac{LR}.
Since the \ac{LR} currently stores the address of the point from where our function
has been called then the control from \printf will be returned to that point.
Therefore we do not have to save \ac{LR} because we do not have necessity to modify \ac{LR}.
And we do not have necessity to modify \ac{LR} because there are no other function calls except \printf. Furthermore,
after this call we do not to do anything else!
That is the reason such optimization is possible.

This optimization is often used in functions where the last statement is a call to another function.
A similar example is presented here:
\myref{jump_to_last_printf}.

\myparagraph{ARM64}

\mysubparagraph{\NonOptimizing GCC (Linaro) 4.9}

\lstinputlisting[caption=\NonOptimizing GCC (Linaro) 4.9,style=customasmARM]{patterns/03_printf/ARM/ARM3_O0_EN.lst}

\myindex{ARM!\Instructions!STP}

The first instruction \INS{STP} (\IT{Store Pair}) saves \ac{FP} (X29) and \ac{LR} (X30) in the stack.
The second \INS{ADD X29, SP, 0} instruction forms the stack frame.
It is just writing the value of \ac{SP} into X29.

\myindex{ARM!\Instructions!ADRP/ADD pair}
Next, we see the familiar \INS{ADRP}/\ADD instruction pair, which forms a pointer to the string.
\myindex{ARM64!lo12}
\IT{lo12} meaning low 12 bits, i.e., linker will write low 12 bits of LC1 address into the opcode of \ADD instruction.
\GTT{\%d} in \printf string format is a 32-bit \Tint, so the 1, 2 and 3 are loaded into 32-bit register parts.

\Optimizing GCC (Linaro) 4.9 generates the same code.

}
\RU{\subsubsection{ARM: 3 аргумента}

В ARM традиционно принята такая схема передачи аргументов в функцию: 
4 первых аргумента через регистры \Reg{0}-\Reg{3}; а остальные~--- через стек.
Это немного похоже на то, как аргументы передаются в 
fastcall~(\myref{fastcall}) или win64~(\myref{sec:callingconventions_win64}).

\myparagraph{32-битный ARM}

\mysubparagraph{\NonOptimizingKeilVI (\ARMMode)}

\begin{lstlisting}[caption=\NonOptimizingKeilVI (\ARMMode),style=customasmARM]
.text:00000000 main
.text:00000000 10 40 2D E9   STMFD   SP!, {R4,LR}
.text:00000004 03 30 A0 E3   MOV     R3, #3
.text:00000008 02 20 A0 E3   MOV     R2, #2
.text:0000000C 01 10 A0 E3   MOV     R1, #1
.text:00000010 08 00 8F E2   ADR     R0, aADBDCD     ; "a=\%d; b=\%d; c=\%d"
.text:00000014 06 00 00 EB   BL      __2printf
.text:00000018 00 00 A0 E3   MOV     R0, #0          ; return 0
.text:0000001C 10 80 BD E8   LDMFD   SP!, {R4,PC}
\end{lstlisting}

Итак, первые 4 аргумента передаются через регистры \Reg{0}-\Reg{3}, по порядку: 
указатель на формат-строку для \printf
в \Reg{0}, затем 1 в \Reg{1}, 2 в \Reg{2} и 3 в \Reg{3}.

Инструкция на \GTT{0x18} записывает 0 в \Reg{0} --- это выражение в Си \emph{return 0}.
Пока что здесь нет ничего необычного.
\OptimizingKeilVI генерирует точно такой же код.

\mysubparagraph{\OptimizingKeilVI (\ThumbMode)}

\begin{lstlisting}[caption=\OptimizingKeilVI (\ThumbMode),style=customasmARM]
.text:00000000 main
.text:00000000 10 B5        PUSH    {R4,LR}
.text:00000002 03 23        MOVS    R3, #3
.text:00000004 02 22        MOVS    R2, #2
.text:00000006 01 21        MOVS    R1, #1
.text:00000008 02 A0        ADR     R0, aADBDCD     ; "a=\%d; b=\%d; c=\%d"
.text:0000000A 00 F0 0D F8  BL      __2printf
.text:0000000E 00 20        MOVS    R0, #0
.text:00000010 10 BD        POP     {R4,PC}
\end{lstlisting}

Здесь нет особых отличий от неоптимизированного варианта для режима ARM.
\mysubparagraph{\OptimizingKeilVI (\ARMMode) + убираем return}
\label{ARM_B_to_printf}

Немного переделаем пример, убрав \emph{return 0}:

\begin{lstlisting}[style=customc]
#include <stdio.h>

void main()
{
	printf("a=%d; b=%d; c=%d", 1, 2, 3);
};
\end{lstlisting}

Результат получится необычным:

\lstinputlisting[caption=\OptimizingKeilVI (\ARMMode),style=customasmARM]{patterns/03_printf/ARM/tmp3.asm}

\myindex{ARM!\Registers!Link Register}
\myindex{ARM!\Instructions!B}
\myindex{Function epilogue}
Это оптимизированная версия (\Othree) для режима ARM, и здесь мы видим последнюю инструкцию 
\INS{B} вместо привычной нам \INS{BL}.
Отличия между этой оптимизированной версией и предыдущей, скомпилированной без оптимизации, 
ещё и в том, что здесь нет пролога и эпилога функции (инструкций, сохраняющих состояние регистров \Reg{0} и \ac{LR}).
\myindex{x86!\Instructions!JMP}
Инструкция \INS{B} просто переходит на другой адрес, без манипуляций с регистром \ac{LR}, то есть это аналог \JMP в x86.
Почему это работает нормально? Потому что этот код эквивалентен предыдущему.

Основных причин две: 1) стек не модифицируется, как и \glslink{stack pointer}{указатель стека} \ac{SP}; 2) вызов функции \printf последний, после него ничего не происходит.
Функция \printf, отработав, просто возвращает управление по адресу, записанному в \ac{LR}.

Но в \ac{LR} находится адрес места, откуда была вызвана наша функция!
А следовательно, управление из \printf вернется сразу туда.

Значит нет нужды сохранять \ac{LR}, потому что нет нужны модифицировать \ac{LR}.
А нет нужды модифицировать \ac{LR}, потому что нет иных вызовов функций, кроме \printf, к тому же, после этого вызова не нужно ничего здесь больше делать!
Поэтому такая оптимизация возможна.

Эта оптимизация часто используется в функциях, где последнее выражение~--- это вызов другой функции.

Ещё один похожий пример описан здесь:
\myref{jump_to_last_printf}.

\myparagraph{ARM64}

\mysubparagraph{\NonOptimizing GCC (Linaro) 4.9}

\lstinputlisting[caption=\NonOptimizing GCC (Linaro) 4.9,style=customasmARM]{patterns/03_printf/ARM/ARM3_O0_RU.lst}

\myindex{ARM!\Instructions!STP}
Итак, первая инструкция \INS{STP} (\emph{Store Pair}) сохраняет \ac{FP} (X29) и \ac{LR} (X30) в стеке.
Вторая инструкция \INS{ADD X29, SP, 0} формирует стековый фрейм.
Это просто запись значения \ac{SP} в X29.

\myindex{ARM!\Instructions!ADRP/ADD pair}
Далее уже знакомая пара инструкций \INS{ADRP}/\ADD формирует указатель на строку.

\myindex{ARM64!lo12}
\emph{lo12} означает младшие 12 бит, т.е., линкер запишет младшие 12 бит адреса метки LC1 в опкод инструкции \ADD.
\GTT{\%d} в формате \printf это 32-битный \Tint, так что 1, 2 и 3 заносятся в 32-битные части регистров.
\Optimizing GCC (Linaro) 4.9 генерирует почти такой же код.

}
\FR{\subsubsection{ARM: 3 arguments}

Le schéma ARM traditionnel pour passer des arguments (convention d'appel) se
comporte de cette façon:
les 4 premiers arguments sont passés par les registres \Reg{0}-\Reg{3}; les autres
par la pile.
Cela ressemble au schéma de passage des arguments dans
fastcall~(\myref{fastcall}) ou win64~(\myref{sec:callingconventions_win64}).

\myparagraph{ARM 32-bit}

\mysubparagraph{\NonOptimizingKeilVI (\ARMMode)}

\begin{lstlisting}[caption=\NonOptimizingKeilVI (\ARMMode),style=customasmARM]
.text:00000000 main
.text:00000000 10 40 2D E9   STMFD   SP!, {R4,LR}
.text:00000004 03 30 A0 E3   MOV     R3, #3
.text:00000008 02 20 A0 E3   MOV     R2, #2
.text:0000000C 01 10 A0 E3   MOV     R1, #1
.text:00000010 08 00 8F E2   ADR     R0, aADBDCD     ; "a=\%d; b=\%d; c=\%d"
.text:00000014 06 00 00 EB   BL      __2printf
.text:00000018 00 00 A0 E3   MOV     R0, #0          ; renvoyer 0
.text:0000001C 10 80 BD E8   LDMFD   SP!, {R4,PC}
\end{lstlisting}

Donc, les 4 premiers arguments sont passés par les registres \Reg{0}-\Reg{3} dans
cet ordre:
un pointeur sur la chaîne de format de \printf dans \Reg{0}, puis 1 dans \Reg{1},
2 dans \Reg{2} et 3 dans \Reg{3}.
L'instruction en \GTT{0x18} écrit 0 dans \Reg{0}---c'est la déclaration C
de \emph{return 0}.

\OptimizingKeilVI génère le même code.

\mysubparagraph{\OptimizingKeilVI (\ThumbMode)}

\begin{lstlisting}[caption=\OptimizingKeilVI (\ThumbMode),style=customasmARM]
.text:00000000 main
.text:00000000 10 B5        PUSH    {R4,LR}
.text:00000002 03 23        MOVS    R3, #3
.text:00000004 02 22        MOVS    R2, #2
.text:00000006 01 21        MOVS    R1, #1
.text:00000008 02 A0        ADR     R0, aADBDCD     ; "a=\%d; b=\%d; c=\%d"
.text:0000000A 00 F0 0D F8  BL      __2printf
.text:0000000E 00 20        MOVS    R0, #0
.text:00000010 10 BD        POP     {R4,PC}
\end{lstlisting}

Il n'y a pas de différence significative avec le code non optimisé pour le mode ARM.

\mysubparagraph{\OptimizingKeilVI (\ARMMode) + supprimons le retour}
\label{ARM_B_to_printf}

Retravaillons légèrement l'exemple en supprimant \emph{return 0}:

\begin{lstlisting}[style=customc]
#include <stdio.h>

void main()
{
	printf("a=%d; b=%d; c=%d", 1, 2, 3);
};
\end{lstlisting}

Le résultat est quelque peu inhabituel:

\lstinputlisting[caption=\OptimizingKeilVI (\ARMMode),style=customasmARM]{patterns/03_printf/ARM/tmp3.asm}

\myindex{ARM!\Registers!Link Register}
\myindex{ARM!\Instructions!B}
\myindex{Epilogue de fonction}
C'est la version optimisée (\Othree) pour le mode ARM et cette fois nous voyons
\INS{B} comme dernière instruction au lieu du \INS{BL} habituel.
Une autre différence entre cette version optimisée et la précédente (compilée
sans optimisation) est l'absence de fonctions prologue et épilogue (les instructions
qui préservent les valeurs des registres \Reg{0} et \ac{LR}).
\myindex{x86!\Instructions!JMP}
L'instruction \INS{B} saute simplement à une autre adresse, sans manipuler le registre
\ac{LR}, de façon similaire au \JMP en x86.
Pourquoi est-ce que fonctionne? Parce ce code est en fait bien équivalent au précédent.
Il y a deux raisons principales: 1) Ni la pile ni \ac{SP} (\glslink{stack pointer}{pointeur de pile})
ne sont modifiés;
2) l'appel à \printf est la dernière instruction, donc il ne se passe rien après.
A la fin, la fonction \printf rend simplement le contrôle à l'adresse stockée
dans \ac{LR}.
Puisque \ac{LR} contient actuellement l'adresse du point depuis lequel notre fonction
a été appelée alors le contrôle après \printf sera redonné à ce point.
Par conséquent, nous n'avons pas besoin de sauver \ac{LR} car il ne nous est pas
nécessaire de le modifier.
Et il ne nous est non plus pas nécessaire de modifier \ac{LR} car il n'y a pas d'autre
appel de fonction excepté \printf. Par ailleurs, après cet appel nous ne faisons
rien d'autre!
C'est la raison pour laquelle une telle optimisation est possible.

Cette optimisation est souvent utilisée dans les fonctions où la dernière déclaration
est un appel à une autre fonction.
Un exemple similaire est présenté ici:
\myref{jump_to_last_printf}.

\myparagraph{ARM64}

\mysubparagraph{GCC (Linaro) 4.9 \NonOptimizing}

\lstinputlisting[caption=GCC (Linaro) 4.9 \NonOptimizing,style=customasmARM]{patterns/03_printf/ARM/ARM3_O0_FR.lst}

\myindex{ARM!\Instructions!STP}

La première instruction \INS{STP} (\emph{Store Pair}) sauve \ac{FP} (X29) et \ac{LR} (X30) sur la pile.

La seconde instruction, \INS{ADD X29, SP, 0} crée la pile.
Elle écrit simplement la valeur de \ac{SP} dans X29.

\myindex{ARM!\Instructions!ADRP/ADD pair}
Ensuite nous voyons la paire d'instructions habituelle \INS{ADRP}/\ADD, qui crée
le pointeur sur la chaîne.
\myindex{ARM64!lo12}
\emph{lo12} signifie les 12 bits de poids faible, i.e., le linker va écrire les 12 bits
de poids faible de l'adresse LC1 dans l'opcode de l'instruction \ADD.
\GTT{\%d} dans la chaîne de format de \printf est un \Tint 32-bit, les 1, 2 et
3 sont chargés dans les parties 32-bit des registres.

GCC (Linaro) 4.9 \Optimizing génère le même code.

}
\IT{\subsubsection{ARM: 3 arguments}

Lo schema tradizionale per il passaggio di argomenti (calling convention) di ARM si comporta in questo modo:
i primi 4 argomenti vengono passati attraverso i registri \Reg{0}-\Reg{3} , i restanti attraverso lo stack.
Cio' ricorda molto il metodo per il passaggio di argomenti in 
fastcall~(\myref{fastcall}) o win64~(\myref{sec:callingconventions_win64}).

\myparagraph{32-bit ARM}

\mysubparagraph{\NonOptimizingKeilVI (\ARMMode)}

\begin{lstlisting}[caption=\NonOptimizingKeilVI (\ARMMode),style=customasmARM]
.text:00000000 main
.text:00000000 10 40 2D E9   STMFD   SP!, {R4,LR}
.text:00000004 03 30 A0 E3   MOV     R3, #3
.text:00000008 02 20 A0 E3   MOV     R2, #2
.text:0000000C 01 10 A0 E3   MOV     R1, #1
.text:00000010 08 00 8F E2   ADR     R0, aADBDCD     ; "a=\%d; b=\%d; c=\%d"
.text:00000014 06 00 00 EB   BL      __2printf
.text:00000018 00 00 A0 E3   MOV     R0, #0          ; return 0
.text:0000001C 10 80 BD E8   LDMFD   SP!, {R4,PC}
\end{lstlisting}

I primi 4 argomenti sono quindi passati attraverso i registri \Reg{0}-\Reg{3} nel seguente ordine:
un puntatore alla format string di \printf in 
\Reg{0}, 1 in \Reg{1}, 2 in \Reg{2} e 3 in \Reg{3}.
L'istruzione a \GTT{0x18} scrive 0 in \Reg{0}---questo equivale allo statement C \emph{return 0}.
Niente di insolito fino a qui.

\OptimizingKeilVI genera lo stesso codice.

\mysubparagraph{\OptimizingKeilVI (\ThumbMode)}

\begin{lstlisting}[caption=\OptimizingKeilVI (\ThumbMode),style=customasmARM]
.text:00000000 main
.text:00000000 10 B5        PUSH    {R4,LR}
.text:00000002 03 23        MOVS    R3, #3
.text:00000004 02 22        MOVS    R2, #2
.text:00000006 01 21        MOVS    R1, #1
.text:00000008 02 A0        ADR     R0, aADBDCD     ; "a=\%d; b=\%d; c=\%d"
.text:0000000A 00 F0 0D F8  BL      __2printf
.text:0000000E 00 20        MOVS    R0, #0
.text:00000010 10 BD        POP     {R4,PC}
\end{lstlisting}

Non c'e' nessuna differenza significativa nel codice non ottimizzato per modo ARM.

\mysubparagraph{\OptimizingKeilVI (\ARMMode) + rimozione di return}
\label{ARM_B_to_printf}

Modifichiamo leggermente l'esempio rimuovendo \emph{return 0}:

\begin{lstlisting}[style=customc]
#include <stdio.h>

void main()
{
	printf("a=%d; b=%d; c=%d", 1, 2, 3);
};
\end{lstlisting}

Il risultato e' alquanto insolito:

\lstinputlisting[caption=\OptimizingKeilVI (\ARMMode),style=customasmARM]{patterns/03_printf/ARM/tmp3.asm}

\myindex{ARM!\Registers!Link Register}
\myindex{ARM!\Instructions!B}
\myindex{Function epilogue}
Questa e' la versione ottimizzata (\Othree) per ARM mode e stavolta notiamo \INS{B} come ultima istruzione, al posto della familiare \INS{BL}.
Un'altra differenza tra questa versione ottimizzata e la precedente (compilata senza ottimizzazione)
e' la mancanza di prologo ed epilogo della funzione (le istruzioni che preservano i valori dei registri \Reg{0} e \ac{LR}).
\myindex{x86!\Instructions!JMP}
L'istruzione \INS{B} salta semplicemente ad un altro indirizzo, senza alcuna manipolazione del registro \ac{LR}, in modo simile
a \JMP in x86.
Perche' funziona? Perche' questo codice e' infatti equivalente al precedente.
Principalmente per due motivi: 1) ne' lo stack ne' \ac{SP} (lo \gls{stack pointer}) vengono modificati;
2) la chiamata a \printf e' l'ultima istruzionie, quindi non succede niente dopo di essa.
Al completamento, \printf restituisce semplicemente il controllo all'indirizzo memorizzato in \ac{LR}.
Poiche' \ac{LR} attualmente contiene l'indirizzo del punto da cui la nostra funzione era stata chiamata, 
il controllo verra' restituito da \printf a quello stesso punto.
Pertanto non c'e' alcun bisogno di salvare \ac{LR} in quanto non abbiamo necessita' di modificare \ac{LR}.
E non vogliamo affatto modificare \ac{LR} poiche' non ci sono altre chiamate a funzione ad eccezzione di \printf. Inoltre, dopo
questa chiamata non abbiamo nient'altro da fare!
Queste le ragione per cui una simile ottimizzazione e' possibile.

Questa ottimizzazione e' spesso usata in funzioni in cui l'ultimo statement e' una chiamata ad un'altra funzione.
Un esempio simile e' fornito di seguito:
\myref{jump_to_last_printf}.

\myparagraph{ARM64}

\mysubparagraph{\NonOptimizing GCC (Linaro) 4.9}

% TODO translate to Italian:
\lstinputlisting[caption=\NonOptimizing GCC (Linaro) 4.9,style=customasmARM]{patterns/03_printf/ARM/ARM3_O0_EN.lst}

\myindex{ARM!\Instructions!STP}

La prima istruzione \INS{STP} (\emph{Store Pair}) salva \ac{FP} (X29) e \ac{LR} (X30) nello stack.
La seconda istruzione \INS{ADD X29, SP, 0} forma lo stack frame.
Scrive semplicemente il valore di \ac{SP} in X29.

\myindex{ARM!\Instructions!ADRP/ADD pair}
Successivamente vediamo la familiare coppia di istruzioni \INS{ADRP}/\ADD , che forma un puntatore alla stringa.
\myindex{ARM64!lo12}
\emph{lo12} indica 12 bit bassi (low 12 bits), ovvero, il linker scrivera' i 12 bit bassi dell'indirizzo di LC1 nell'opcode dell'istruzione \ADD .
\GTT{\%d} nella format string di \printf e' un \Tint a 32-bit, quindi u valori 1, 2 e 3 sono caricati nelle parti a 32-bit dei registri.

\Optimizing GCC (Linaro) 4.9 genera lo stesso codice.
}
\JA{\subsubsection{ARM: 3つの引数}

引数を渡すためのARMの伝統的なスキーム（呼び出し規約）は、次のように動作します。
最初の4つの引数は\Reg{0}-\Reg{3}レジスタに渡されます。 残りの引数はスタックを介して。 
これはfastcall~(\myref{fastcall})またはwin64~(\myref{sec:callingconventions_win64})の引数渡しスキームに似ています。

\myparagraph{32ビットARM}

\mysubparagraph{\NonOptimizingKeilVI (\ARMMode)}

\begin{lstlisting}[caption=\NonOptimizingKeilVI (\ARMMode),style=customasmARM]
.text:00000000 main
.text:00000000 10 40 2D E9   STMFD   SP!, {R4,LR}
.text:00000004 03 30 A0 E3   MOV     R3, #3
.text:00000008 02 20 A0 E3   MOV     R2, #2
.text:0000000C 01 10 A0 E3   MOV     R1, #1
.text:00000010 08 00 8F E2   ADR     R0, aADBDCD     ; "a=\%d; b=\%d; c=\%d"
.text:00000014 06 00 00 EB   BL      __2printf
.text:00000018 00 00 A0 E3   MOV     R0, #0          ; return 0
.text:0000001C 10 80 BD E8   LDMFD   SP!, {R4,PC}
\end{lstlisting}

したがって、最初の4つの引数は、\Reg{0}-\Reg{3}レジスタをこの順序で渡します。
\Reg{0}の \printf 形式文字列へのポインタ、\Reg{1}の1、\Reg{2}の2、\Reg{3}の3の順です。 
\GTT{0x18}の命令は\Reg{0}に0を書き込みます。 これは\emph{return 0} となるCの命令文です。 
珍しいことは何もありません。

\OptimizingKeilVI は同じコードを生成します。

\mysubparagraph{\OptimizingKeilVI (\ThumbMode)}

\begin{lstlisting}[caption=\OptimizingKeilVI (\ThumbMode),style=customasmARM]
.text:00000000 main
.text:00000000 10 B5        PUSH    {R4,LR}
.text:00000002 03 23        MOVS    R3, #3
.text:00000004 02 22        MOVS    R2, #2
.text:00000006 01 21        MOVS    R1, #1
.text:00000008 02 A0        ADR     R0, aADBDCD     ; "a=\%d; b=\%d; c=\%d"
.text:0000000A 00 F0 0D F8  BL      __2printf
.text:0000000E 00 20        MOVS    R0, #0
.text:00000010 10 BD        POP     {R4,PC}
\end{lstlisting}

ARMモードの最適化されていないコードとの大きな違いはありません。

\mysubparagraph{\OptimizingKeilVI (\ARMMode) + returnを削除してみる}
\label{ARM_B_to_printf}

\emph{return 0}を取り除いて例を少し修正しましょう：

\begin{lstlisting}[style=customc]
#include <stdio.h>

void main()
{
	printf("a=%d; b=%d; c=%d", 1, 2, 3);
};
\end{lstlisting}

結果はちょっと珍しくなりました。

\lstinputlisting[caption=\OptimizingKeilVI (\ARMMode),style=customasmARM]{patterns/03_printf/ARM/tmp3.asm}

\myindex{ARM!\Registers!Link Register}
\myindex{ARM!\Instructions!B}
\myindex{Function epilogue}
これはARMモード用に最適化された（\Othree）バージョンであり、今回は\INS{B}を使い慣れた\INS{BL}ではなく最後の命令と見なします。
この最適化されたバージョンと前のバージョン（最適化なしでコンパイルされたもの）との別の違いは、
関数のプロローグとエピローグ（\Reg{0} と \ac{LR}レジスタの値を保持する命令）の欠如です。 
\INS{B}命令は、x86の \JMP と同様に、\ac{LR}レジスタの操作なしで別のアドレスにジャンプするだけです。
それはなぜ機能するのでしょうか？実際、このコードは以前のコードと事実上同等です。
主な理由は2つあります。1）スタックも\ac{SP}（\gls{stack pointer}）も変更されていません。 
2） \printf の呼び出しが最後の命令なので、その後は何も起こりません。
完了すると、 \printf 関数は\ac{LR}に格納されているアドレスにコントロールを返します。 
\ac{LR}は現在、関数が呼び出されたポイントのアドレスを格納しているので、 \printf からの制御はそのポイントに返されます。
したがって、\ac{LR}を変更する必要がないため、\ac{LR}を節約する必要はありません。 
\printf 以外の関数呼び出しがないため、\ac{LR}を変更する必要はありません。
さらに、この呼び出しの後、私たちは何もしません！これがそのような最適化が可能な理由です。

この最適化は、最後のステートメントが別の関数の呼び出しである関数でよく使用されます。
同様の例をここに示します：\myref{jump_to_last_printf}

\myparagraph{ARM64}

\mysubparagraph{\NonOptimizing GCC (Linaro) 4.9}

\lstinputlisting[caption=\NonOptimizing GCC (Linaro) 4.9,style=customasmARM]{patterns/03_printf/ARM/ARM3_O0_JA.lst}

\myindex{ARM!\Instructions!STP}

第1の命令\INS{STP} (\emph{ストアペア})は、\ac{FP}（X29）および\ac{LR}（X30）をスタックに保存します。 
2番目の\INS{ADD X29, SP, 0}命令がスタックフレームを形成します。 
\ac{SP}の値をX29に書き込むだけです。

\myindex{ARM!\Instructions!ADRP/ADD pair}
次に、使い慣れた\INS{ADRP}/\ADD 命令ペアを参照します。これは、文字列へのポインタを形成します。 
\myindex{ARM64!lo12}
すなわち、\emph{lo12}は、LC1アドレスの下位12ビットを \ADD 命令のオペコードに書き込みます。
\printf の文字列書式の\GTT{\%d}は32ビット \Tint なので、1,2および3は32ビットのレジスタ部分にロードされます。

\Optimizing GCC (Linaro) 4.9 は同じコードを生成します。
}

\EN{\subsubsection{ARM: 8 arguments}

Let's use again the example with 9 arguments from the previous section: \myref{example_printf8_x64}.

\lstinputlisting[style=customc]{patterns/03_printf/2.c}

\myparagraph{\OptimizingKeilVI: \ARMMode}

\lstinputlisting[style=customasmARM]{patterns/03_printf/ARM/tmp1.asm}

This code can be divided into several parts:

\myindex{Function prologue}
\begin{itemize}
\item Function prologue:

\myindex{ARM!\Instructions!STR}
The very first \INS{STR LR, [SP,\#var\_4]!} instruction saves \ac{LR} on the stack, because we are going to use this register for the \printf call.
Exclamation mark at the end indicates \IT{pre-index}.

This implies that \ac{SP} is to be decreased by 4 first, and then \ac{LR} will be saved at the address stored in \ac{SP}.
This is similar to \PUSH in x86.
Read more about it at: \myref{ARM_postindex_vs_preindex}.

\myindex{ARM!\Instructions!SUB}
The second \INS{SUB SP, SP, \#0x14} instruction decreases \ac{SP} (the \gls{stack pointer}) in order to allocate \GTT{0x14} (20) bytes on the stack.
Indeed, we have to pass 5 32-bit values via the stack to the \printf function, and each one occupies 4 bytes, which is exactly $5*4=20$.
The other 4 32-bit values are to be passed through registers.

\item Passing 5, 6, 7 and 8 via the stack: they are stored in the \Reg{0}, \Reg{1}, \Reg{2} and \Reg{3} registers respectively.\\
Then, the \INS{ADD R12, SP, \#0x18+var\_14} instruction writes the stack address where these 4 variables are to be stored, into the \Reg{12} register.
\myindex{IDA!var\_?}
\IT{var\_14} is an assembly macro, equal to -0x14, created by \IDA to conveniently display the code accessing the stack.
The \IT{var\_?} macros generated by \IDA reflect local variables in the stack.

So, \GTT{SP+4} is to be stored into the \Reg{12} register. \\
\myindex{ARM!\Instructions!STMIA}
The next \INS{STMIA R12, {R0-R3}} instruction writes registers \Reg{0}-\Reg{3} contents to the memory pointed by \Reg{12}.
\INS{STMIA} abbreviates \IT{Store Multiple Increment After}. 
\IT{\q{Increment After}} implies that \Reg{12} is to be increased by 4 after each register value is written.

\item Passing 4 via the stack: 4 is stored in \Reg{0} and then this value, with the help of the \\
\INS{STR R0, [SP,\#0x18+var\_18]} instruction is saved on the stack.
\IT{var\_18} is -0x18, so the offset is to be 0, thus the value from the \Reg{0} register (4) is to be written to the address written in \ac{SP}.

\item Passing 1, 2 and 3 via registers:
The values of the first 3 numbers (a, b, c) (1, 2, 3 respectively) are passed through the 
\Reg{1}, \Reg{2} and \Reg{3}
registers right before the \printf call, and the other
5 values are passed via the stack:

\item \printf call.

\myindex{Function epilogue}
\item Function epilogue:

The \INS{ADD SP, SP, \#0x14} instruction restores the \ac{SP} pointer back to its former value,
thus annulling everything what has been stored on the stack.
Of course, what has been stored on the stack will stay there, but it will all be rewritten during the execution of subsequent functions.

\myindex{ARM!\Instructions!LDR}
The \INS{LDR PC, [SP+4+var\_4],\#4} instruction loads the saved \ac{LR} value from the stack into the \ac{PC} register, thus causing the function to exit.
There is no exclamation mark---indeed, \ac{PC} is loaded first from the address stored in \ac{SP} ($4+var\_4=4+(-4)=0$, so this instruction is analogous to \INS{LDR PC, [SP],\#4}), and then \ac{SP} is increased by 4.
This is referred as \IT{post-index}\footnote{Read more about it: \myref{ARM_postindex_vs_preindex}.}.
Why does \IDA display the instruction like that?
Because it wants to illustrate the stack layout and the fact that \GTT{var\_4} is allocated for saving the \ac{LR} value in the local stack.
This instruction is somewhat similar to \INS{POP PC} in x86\footnote{It is impossible to set \GTT{IP/EIP/RIP} value using \POP in x86, but anyway, you got the analogy right.}.

\end{itemize}

\myparagraph{\OptimizingKeilVI: \ThumbMode}

\lstinputlisting[style=customasmARM]{patterns/03_printf/ARM/tmp2.asm}

The output is almost like in the previous example. However, this is Thumb code and the values are packed into stack differently: 
8 goes first, then 5, 6, 7, and 4 goes third.

\myparagraph{\OptimizingXcodeIV: \ARMMode}

\begin{lstlisting}[style=customasmARM]
__text:0000290C             _printf_main2
__text:0000290C
__text:0000290C             var_1C = -0x1C
__text:0000290C             var_C  = -0xC
__text:0000290C
__text:0000290C 80 40 2D E9   STMFD  SP!, {R7,LR}
__text:00002910 0D 70 A0 E1   MOV    R7, SP
__text:00002914 14 D0 4D E2   SUB    SP, SP, #0x14
__text:00002918 70 05 01 E3   MOV    R0, #0x1570
__text:0000291C 07 C0 A0 E3   MOV    R12, #7
__text:00002920 00 00 40 E3   MOVT   R0, #0
__text:00002924 04 20 A0 E3   MOV    R2, #4
__text:00002928 00 00 8F E0   ADD    R0, PC, R0
__text:0000292C 06 30 A0 E3   MOV    R3, #6
__text:00002930 05 10 A0 E3   MOV    R1, #5
__text:00002934 00 20 8D E5   STR    R2, [SP,#0x1C+var_1C]
__text:00002938 0A 10 8D E9   STMFA  SP, {R1,R3,R12}
__text:0000293C 08 90 A0 E3   MOV    R9, #8
__text:00002940 01 10 A0 E3   MOV    R1, #1
__text:00002944 02 20 A0 E3   MOV    R2, #2
__text:00002948 03 30 A0 E3   MOV    R3, #3
__text:0000294C 10 90 8D E5   STR    R9, [SP,#0x1C+var_C]
__text:00002950 A4 05 00 EB   BL     _printf
__text:00002954 07 D0 A0 E1   MOV    SP, R7
__text:00002958 80 80 BD E8   LDMFD  SP!, {R7,PC}
\end{lstlisting}

\myindex{ARM!\Instructions!STMFA}
\myindex{ARM!\Instructions!STMIB}
Almost the same as what we have already seen, with the
exception of \INS{STMFA} (Store Multiple Full Ascending) instruction, which is a synonym of \INS{STMIB} (Store Multiple Increment Before) instruction. 
This instruction increases the value in the \ac{SP} register and only then writes the next register value into the memory, rather than performing those two actions in the opposite order.

Another thing that catches the eye is that the instructions are arranged seemingly random.
For example, the value in the \Reg{0} register is manipulated in three
places, at addresses \GTT{0x2918}, \GTT{0x2920} and \GTT{0x2928}, when it would be possible to do it in one point.

However, the optimizing compiler may have its own reasons on how to order the instructions so to achieve higher efficiency during the execution.

Usually, the processor attempts to simultaneously execute instructions located side-by-side.\\
For example, instructions like \INS{MOVT R0, \#0} and
\INS{ADD R0, PC, R0} cannot be executed simultaneously since they both modify the \Reg{0} register. 
On the other hand, \INS{MOVT R0, \#0} and \INS{MOV R2, \#4} 
instructions can be executed
simultaneously since the effects of their execution are not conflicting with each other.
Presumably, the compiler tries to generate code in such a manner (wherever it is possible).
 
\myparagraph{\OptimizingXcodeIV: \ThumbTwoMode}

\begin{lstlisting}[style=customasmARM]
__text:00002BA0               _printf_main2
__text:00002BA0
__text:00002BA0               var_1C = -0x1C
__text:00002BA0               var_18 = -0x18
__text:00002BA0               var_C  = -0xC
__text:00002BA0
__text:00002BA0 80 B5          PUSH     {R7,LR}
__text:00002BA2 6F 46          MOV      R7, SP
__text:00002BA4 85 B0          SUB      SP, SP, #0x14
__text:00002BA6 41 F2 D8 20    MOVW     R0, #0x12D8
__text:00002BAA 4F F0 07 0C    MOV.W    R12, #7
__text:00002BAE C0 F2 00 00    MOVT.W   R0, #0
__text:00002BB2 04 22          MOVS     R2, #4
__text:00002BB4 78 44          ADD      R0, PC  ; char *
__text:00002BB6 06 23          MOVS     R3, #6
__text:00002BB8 05 21          MOVS     R1, #5
__text:00002BBA 0D F1 04 0E    ADD.W    LR, SP, #0x1C+var_18
__text:00002BBE 00 92          STR      R2, [SP,#0x1C+var_1C]
__text:00002BC0 4F F0 08 09    MOV.W    R9, #8
__text:00002BC4 8E E8 0A 10    STMIA.W  LR, {R1,R3,R12}
__text:00002BC8 01 21          MOVS     R1, #1
__text:00002BCA 02 22          MOVS     R2, #2
__text:00002BCC 03 23          MOVS     R3, #3
__text:00002BCE CD F8 10 90    STR.W    R9, [SP,#0x1C+var_C]
__text:00002BD2 01 F0 0A EA    BLX      _printf
__text:00002BD6 05 B0          ADD      SP, SP, #0x14
__text:00002BD8 80 BD          POP      {R7,PC}
\end{lstlisting}

The output is almost the same as in the previous example, with the exception that Thumb-instructions are used instead.
% FIXME: also STMIA is used instead of STMIB,
% which is why it uses LR, which is 4 bytes ahead of SP

\myparagraph{ARM64}

\mysubparagraph{\NonOptimizing GCC (Linaro) 4.9}

\lstinputlisting[caption=\NonOptimizing GCC (Linaro) 4.9,style=customasmARM]{patterns/03_printf/ARM/ARM8_O0_EN.lst}

The first 8 arguments are passed in X- or W-registers: \ARMPCS.
A string pointer requires a 64-bit register, so it's passed in \RegX{0}.
All other values have a \Tint 32-bit type, so they are stored in the 32-bit part of the registers (W-).
The 9th argument (8) is passed via the stack.
Indeed: it's not possible to pass large number of arguments through registers, because the number of registers is limited.

\Optimizing GCC (Linaro) 4.9 generates the same code.
}
\RU{\subsubsection{ARM: 8 аргументов}

Снова воспользуемся примером с 9-ю аргументами из предыдущей секции: \myref{example_printf8_x64}.

\lstinputlisting[style=customc]{patterns/03_printf/2.c}

\myparagraph{\OptimizingKeilVI: \ARMMode}

\begin{lstlisting}[style=customasmARM]
.text:00000028             main
.text:00000028
.text:00000028             var_18 = -0x18
.text:00000028             var_14 = -0x14
.text:00000028             var_4  = -4
.text:00000028
.text:00000028 04 E0 2D E5  STR    LR, [SP,#var_4]!
.text:0000002C 14 D0 4D E2  SUB    SP, SP, #0x14
.text:00000030 08 30 A0 E3  MOV    R3, #8
.text:00000034 07 20 A0 E3  MOV    R2, #7
.text:00000038 06 10 A0 E3  MOV    R1, #6
.text:0000003C 05 00 A0 E3  MOV    R0, #5
.text:00000040 04 C0 8D E2  ADD    R12, SP, #0x18+var_14
.text:00000044 0F 00 8C E8  STMIA  R12, {R0-R3}
.text:00000048 04 00 A0 E3  MOV    R0, #4
.text:0000004C 00 00 8D E5  STR    R0, [SP,#0x18+var_18]
.text:00000050 03 30 A0 E3  MOV    R3, #3
.text:00000054 02 20 A0 E3  MOV    R2, #2
.text:00000058 01 10 A0 E3  MOV    R1, #1
.text:0000005C 6E 0F 8F E2  ADR    R0, aADBDCDDDEDFDGD ; "a=\%d; b=\%d; c=\%d; d=\%d; e=\%d; f=\%d; g=\%"...
.text:00000060 BC 18 00 EB  BL     __2printf
.text:00000064 14 D0 8D E2  ADD    SP, SP, #0x14
.text:00000068 04 F0 9D E4  LDR    PC, [SP+4+var_4],#4
\end{lstlisting}

Этот код можно условно разделить на несколько частей:

\myindex{Function prologue}
\begin{itemize}
\item Пролог функции:

\myindex{ARM!\Instructions!STR}
Самая первая инструкция \INS{STR LR, [SP,\#var\_4]!} 
сохраняет в стеке \ac{LR}, ведь нам придется использовать этот регистр для вызова \printf.
Восклицательный знак в конце означает \IT{pre-index}.
Это значит, что в начале \ac{SP} должно быть уменьшено на 4, затем по адресу в \ac{SP} должно быть записано значение \ac{LR}.

Это аналог знакомой в x86 инструкции \PUSH. Читайте больше об этом: \myref{ARM_postindex_vs_preindex}.

\myindex{ARM!\Instructions!SUB}
Вторая инструкция \INS{SUB SP, SP, \#0x14} уменьшает \glslink{stack pointer}{указатель стека} \ac{SP}, но, на самом деле, эта процедура нужна для выделения в локальном стеке места размером \GTT{0x14} (20) байт.
Действительно, нам нужно передать 5 32-битных значений через стек в \printf. Каждое значение занимает 4 байта, все вместе~--- $5*4=20$.
Остальные 4 32-битных значения будут переданы через регистры.

\item Передача 5, 6, 7 и 8 через стек:
они записываются в регистры \Reg{0}, \Reg{1}, \Reg{2} и \Reg{3} соответственно.\\
Затем инструкция \INS{ADD R12, SP, \#0x18+var\_14} 
записывает в регистр \Reg{12} адрес места в стеке, куда будут помещены эти 4 значения.
\myindex{IDA!var\_?}
\IT{var\_14}~--- это макрос ассемблера, равный -0x14.
Такие макросы создает \IDA, чтобы удобнее было показывать, как код обращается к стеку.

Макросы \IT{var\_?}, создаваемые \IDA, отражают локальные переменные в стеке. Так что в \Reg{12} будет записано \GTT{SP+4}.

\myindex{ARM!\Instructions!STMIA}
Следующая инструкция \INS{STMIA R12, {R0-R3}} записывает содержимое регистров \Reg{0}-\Reg{3} по адресу в памяти, на который указывает \Reg{12}.

Инструкция \INS{STMIA} означает \IT{Store Multiple Increment After}.\\
\IT{Increment After} означает, что \Reg{12} будет увеличиваться на 4 после записи каждого значения регистра.

\item Передача 4 через стек:
4 записывается в \Reg{0}, затем инструкция \INS{STR R0, [SP,\#0x18+var\_18]} записывает его в стек.
\IT{var\_18} равен -0x18, смещение будет 0, 
так что значение из регистра \Reg{0} (4) запишется туда, куда указывает \ac{SP}.

\item Передача 1, 2 и 3 через регистры:

Значения для первых трех чисел (a, b, c) (1, 2, 3 соответственно) передаются в регистрах 
\Reg{1}, \Reg{2} и \Reg{3} перед самим вызовом \printf, а остальные 5 значений передаются через стек, и вот как:

\item Вызов \printf.

\myindex{Function epilogue}
\item Эпилог функции:

Инструкция \INS{ADD SP, SP, \#0x14} возвращает \ac{SP} на прежнее место, 
аннулируя таким образом всё, что было записано в стек.
Конечно, то что было записано в стек, там пока и останется, но всё это будет многократно 
перезаписано во время исполнения последующих функций.

\myindex{ARM!\Instructions!LDR}
Инструкция \INS{LDR PC, [SP+4+var\_4],\#4} загружает в \ac{PC} 
сохраненное значение \ac{LR} из стека, обеспечивая таким образом выход из функции.

Здесь нет восклицательного знака~--- действительно, сначала \ac{PC} загружается из места, куда указывает \ac{SP}
($4+var\_4=4+(-4)=0$, так что эта инструкция аналогична \INS{LDR PC, [SP],\#4}), затем \ac{SP} увеличивается 
на 4.
Это называется \IT{post-index}\footnote{Читайте больше об этом: \myref{ARM_postindex_vs_preindex}.}.
Почему \IDA показывает инструкцию именно так?
Потому что она хочет показать разметку стека и тот факт, что \GTT{var\_4} выделена в локальном стеке именно для сохраненного
значения \ac{LR}.
Эта инструкция в каком-то смысле аналогична \INS{POP PC} в x86
\footnote{В x86 невозможно установить значение \GTT{IP/EIP/RIP} используя \POP, но будем надеяться, вы поняли аналогию.}.

\end{itemize}

\myparagraph{\OptimizingKeilVI: \ThumbMode}

\begin{lstlisting}[style=customasmARM]
.text:0000001C             printf_main2
.text:0000001C
.text:0000001C             var_18 = -0x18
.text:0000001C             var_14 = -0x14
.text:0000001C             var_8  = -8
.text:0000001C
.text:0000001C 00 B5        PUSH    {LR}
.text:0000001E 08 23        MOVS    R3, #8
.text:00000020 85 B0        SUB     SP, SP, #0x14
.text:00000022 04 93        STR     R3, [SP,#0x18+var_8]
.text:00000024 07 22        MOVS    R2, #7
.text:00000026 06 21        MOVS    R1, #6
.text:00000028 05 20        MOVS    R0, #5
.text:0000002A 01 AB        ADD     R3, SP, #0x18+var_14
.text:0000002C 07 C3        STMIA   R3!, {R0-R2}
.text:0000002E 04 20        MOVS    R0, #4
.text:00000030 00 90        STR     R0, [SP,#0x18+var_18]
.text:00000032 03 23        MOVS    R3, #3
.text:00000034 02 22        MOVS    R2, #2
.text:00000036 01 21        MOVS    R1, #1
.text:00000038 A0 A0        ADR     R0, aADBDCDDDEDFDGD ; "a=\%d; b=\%d; c=\%d; d=\%d; e=\%d; f=\%d; g=\%"...
.text:0000003A 06 F0 D9 F8  BL      __2printf
.text:0000003E
.text:0000003E             loc_3E   ; CODE XREF: example13\_f+16
.text:0000003E 05 B0        ADD     SP, SP, #0x14
.text:00000040 00 BD        POP     {PC}
\end{lstlisting}

Это почти то же самое что и в предыдущем примере, только код для Thumb и значения помещаются в 
стек немного иначе: сначала 8 за первый раз, затем 5, 6, 7 за второй раз и 4 за третий раз.

\myparagraph{\OptimizingXcodeIV: \ARMMode}

\begin{lstlisting}[style=customasmARM]
__text:0000290C             _printf_main2
__text:0000290C
__text:0000290C             var_1C = -0x1C
__text:0000290C             var_C  = -0xC
__text:0000290C
__text:0000290C 80 40 2D E9   STMFD  SP!, {R7,LR}
__text:00002910 0D 70 A0 E1   MOV    R7, SP
__text:00002914 14 D0 4D E2   SUB    SP, SP, #0x14
__text:00002918 70 05 01 E3   MOV    R0, #0x1570
__text:0000291C 07 C0 A0 E3   MOV    R12, #7
__text:00002920 00 00 40 E3   MOVT   R0, #0
__text:00002924 04 20 A0 E3   MOV    R2, #4
__text:00002928 00 00 8F E0   ADD    R0, PC, R0
__text:0000292C 06 30 A0 E3   MOV    R3, #6
__text:00002930 05 10 A0 E3   MOV    R1, #5
__text:00002934 00 20 8D E5   STR    R2, [SP,#0x1C+var_1C]
__text:00002938 0A 10 8D E9   STMFA  SP, {R1,R3,R12}
__text:0000293C 08 90 A0 E3   MOV    R9, #8
__text:00002940 01 10 A0 E3   MOV    R1, #1
__text:00002944 02 20 A0 E3   MOV    R2, #2
__text:00002948 03 30 A0 E3   MOV    R3, #3
__text:0000294C 10 90 8D E5   STR    R9, [SP,#0x1C+var_C]
__text:00002950 A4 05 00 EB   BL     _printf
__text:00002954 07 D0 A0 E1   MOV    SP, R7
__text:00002958 80 80 BD E8   LDMFD  SP!, {R7,PC}
\end{lstlisting}

\myindex{ARM!\Instructions!STMFA}
\myindex{ARM!\Instructions!STMIB}
Почти то же самое, что мы уже видели, за исключением того, что \INS{STMFA} (Store Multiple Full Ascending)~--- 
это синоним инструкции \INS{STMIB} (Store Multiple Increment Before). 
Эта инструкция увеличивает \ac{SP} и только затем записывает в память значение очередного регистра, но не наоборот.

Далее бросается в глаза то, что инструкции как будто бы расположены случайно.
Например, значение в регистре \Reg{0} подготавливается в трех местах, по адресам \GTT{0x2918}, \GTT{0x2920} и \GTT{0x2928}, 
когда это можно было бы сделать в одном месте.
Однако, у оптимизирующего компилятора могут быть свои доводы о том, как лучше составлять инструкции 
друг с другом для лучшей эффективности исполнения.
Процессор обычно пытается исполнять одновременно идущие друг за другом инструкции.
К примеру, инструкции \INS{MOVT R0, \#0} и \INS{ADD R0, PC, R0} не могут быть исполнены одновременно,
потому что обе инструкции модифицируют регистр \Reg{0}. 
А вот инструкции \INS{MOVT R0, \#0} и \INS{MOV R2, \#4} легко можно исполнить одновременно, 
потому что эффекты от их исполнения никак не конфликтуют друг с другом.
Вероятно, компилятор старается генерировать код именно таким образом там, где это возможно.
 
\myparagraph{\OptimizingXcodeIV: \ThumbTwoMode}

\begin{lstlisting}[style=customasmARM]
__text:00002BA0               _printf_main2
__text:00002BA0
__text:00002BA0               var_1C = -0x1C
__text:00002BA0               var_18 = -0x18
__text:00002BA0               var_C  = -0xC
__text:00002BA0
__text:00002BA0 80 B5          PUSH     {R7,LR}
__text:00002BA2 6F 46          MOV      R7, SP
__text:00002BA4 85 B0          SUB      SP, SP, #0x14
__text:00002BA6 41 F2 D8 20    MOVW     R0, #0x12D8
__text:00002BAA 4F F0 07 0C    MOV.W    R12, #7
__text:00002BAE C0 F2 00 00    MOVT.W   R0, #0
__text:00002BB2 04 22          MOVS     R2, #4
__text:00002BB4 78 44          ADD      R0, PC  ; char *
__text:00002BB6 06 23          MOVS     R3, #6
__text:00002BB8 05 21          MOVS     R1, #5
__text:00002BBA 0D F1 04 0E    ADD.W    LR, SP, #0x1C+var_18
__text:00002BBE 00 92          STR      R2, [SP,#0x1C+var_1C]
__text:00002BC0 4F F0 08 09    MOV.W    R9, #8
__text:00002BC4 8E E8 0A 10    STMIA.W  LR, {R1,R3,R12}
__text:00002BC8 01 21          MOVS     R1, #1
__text:00002BCA 02 22          MOVS     R2, #2
__text:00002BCC 03 23          MOVS     R3, #3
__text:00002BCE CD F8 10 90    STR.W    R9, [SP,#0x1C+var_C]
__text:00002BD2 01 F0 0A EA    BLX      _printf
__text:00002BD6 05 B0          ADD      SP, SP, #0x14
__text:00002BD8 80 BD          POP      {R7,PC}
\end{lstlisting}

Почти то же самое, что и в предыдущем примере,
лишь за тем исключением, что здесь используются Thumb-инструкции.

% FIXME: also STMIA is used instead of STMIB,
% which is why it uses LR, which is 4 bytes ahead of SP

\myparagraph{ARM64}

\mysubparagraph{\NonOptimizing GCC (Linaro) 4.9}

\lstinputlisting[caption=\NonOptimizing GCC (Linaro) 4.9,style=customasmARM]{patterns/03_printf/ARM/ARM8_O0_RU.lst}

Первые 8 аргументов передаются в X- или W-регистрах: \ARMPCS.
Указатель на строку требует 64-битного регистра, так что он передается в \RegX{0}.
Все остальные значения имеют 32-битный тип \Tint, так что они записываются в 32-битные части регистров (W-).
Девятый аргумент (8) передается через стек.
Действительно, невозможно передать большое количество аргументов в регистрах, потому что количество регистров ограничено.

\Optimizing GCC (Linaro) 4.9 генерирует почти такой же код.
}
\FR{\subsubsection{ARM: 8 arguments}

Utilisons de nouveau l'exemple avec 9 arguments de la section précédente: \myref{example_printf8_x64}.

\lstinputlisting[style=customc]{patterns/03_printf/2.c}

\myparagraph{\OptimizingKeilVI: \ARMMode}

\lstinputlisting[style=customasmARM]{patterns/03_printf/ARM/tmp1.asm}

Ce code peut être divisé en plusieurs parties:

\myindex{Prologue de fonction}
\begin{itemize}
\item Prologue de la fonction:

\myindex{ARM!\Instructions!STR}
La toute première instruction \INS{STR LR, [SP,\#var\_4]!} sauve \ac{LR} sur la pile,
car nous allons utiliser ce registre pour l'appel à \printf.
Le point d'exclamation à la fin indique un \IT{pré-index}.

Cela signifie que \ac{SP} est d'abord décrémenté de 4, et qu'ensuite \ac{LR}
va être sauvé à l'adresse stockée dans \ac{SP}.
C'est similaire à \PUSH en x86.
Lire aussi à ce propos: \myref{ARM_postindex_vs_preindex}.

\myindex{ARM!\Instructions!SUB}
La seconde instruction \INS{SUB SP, SP, \#0x14} décrémente \ac{SP} (le \glslink{stack pointer}{pointeur de pile})
afin d'allouer \GTT{0x14} (20) octets sur la pile.
En effet, nous devons passer 5 valeurs de 32-bit par la pile à la fonction \printf,
et chacune occupe 4 octets, ce qui fait exactement $5*4=20$.
Les 4 autres valeurs de 32-bit sont passées par les registres.

\item Passer 5, 6, 7 et 8 par la pile: ils sont stockés dans les registres \Reg{0},
\Reg{1}, \Reg{2} et \Reg{3} respectivement.\\
Ensuite, l'instruction \INS{ADD R12, SP, \#0x18+var\_14} écrit l'adresse de la pile
où ces 4 variables doivent être stockées dans le registre \Reg{12}.
\myindex{IDA!var\_?}
\IT{var\_14} est une macro d'assemblage, égal à -0x14, créée par \IDA pour afficher
commodément le code accédant à la pile.
Les macros \IT{var\_?} générée par \IDA reflètent les variables locales dans la pile.

Donc, \GTT{SP+4} doit être stocké dans le registre \Reg{12}. \\
\myindex{ARM!\Instructions!STMIA}
L'instruction suivante \INS{STMIA R12, {R0-R3}} écrit le contenu des registres
\Reg{0}-\Reg{3} dans la mémoire pointée par \Reg{12}.
\INS{STMIA} est l'abréviation de \IT{Store Multiple Increment After} (stocker plusieurs incrémenter après).
\IT{\q{Increment After}} signifie que \Reg{12} doit être incrémenté de 4 après
l'écriture de chaque valeur d'un registre.

\item Passer 4 par la pile: 4 est stocké dans \Reg{0} et ensuite, cette valeur, avec l'aide de \\
l'instruction \INS{STR R0, [SP,\#0x18+var\_18]} est sauvée dans la pile.
\IT{var\_18} est -0x18, donc l'offset est 0, donc la valeur du registre \Reg{0}
(4) est écrite à l'adresse écrite dans \ac{SP}.

\item Passer 1, 2 et 3 par des registres:
Les valeurs des 3 premiers nombres (a,b,c) (respectivement 1, 2, 3) sont passées
par les registres \Reg{1}, \Reg{2} et \Reg{3} juste avant l'appel de \printf, et
les 5 autres valeurs sont passées par la pile:

\item appel de \printf

\myindex{Epilogue de fonction}
\item Épilogue de fonction:

L'instruction \INS{ADD SP, SP, \#0x14} restaure le pointeur \ac{SP} à sa valeur
précédente, nettoyant ainsi la pile.
Bien sûr, ce qui a été stocké sur la pile y reste, mais sera récrit lors de
l'exécution ultérieure de fonctions.

\myindex{ARM!\Instructions!LDR}
L'instruction \INS{LDR PC, [SP+4+var\_4],\#4} charge la valeur sauvée de \ac{LR}
depuis la pile dans le registre \ac{PC}, provoquant ainsi la sortie de la fonction.
Il n'y a pas de point d'exclamation---effectivement, \ac{PC} est d'abord chargé
depuis l'adresse stockées dans \ac{SP} ($4+var\_4=4+(-4)=0$, donc cette instruction
est analogue à INS{LDR PC, [SP],\#4}), et ensuite \ac{SP} est incrémenté de 4.
Il s'agit de \IT{post-index}\footnote{Lire à ce propos: \myref{ARM_postindex_vs_preindex}.}.
Pourquoi est-ce qu'\IDA affiche l'instruction comme ça?
Parce qu'il veut illustrer la disposition de la pile et le fait que \GTT{var\_4}
est alloué pour sauver la valeur de \ac{LR} dans la pile locale.
Cette instruction est quelque peu similaire à \INS{POP PC} en x86\footnote{Il est
impossible de définir la valeur de \GTT{IP/EIP/RIP} en utilisant \POP en x86, mais
de toutes façons, vous avez le droit de faire l'analogie.}.

\end{itemize}

\myparagraph{\OptimizingKeilVI: \ThumbMode}

\lstinputlisting[style=customasmARM]{patterns/03_printf/ARM/tmp2.asm}

La sortie est presque comme dans les exemples précédents. Toutefois, c'est du code
Thumb et les valeurs sont arrangées différemment dans la pile:
8 vient en premier, puis 5, 6, 7 et 4 vient en troisième.

\myparagraph{\OptimizingXcodeIV: \ARMMode}

\begin{lstlisting}[style=customasmARM]
__text:0000290C             _printf_main2
__text:0000290C
__text:0000290C             var_1C = -0x1C
__text:0000290C             var_C  = -0xC
__text:0000290C
__text:0000290C 80 40 2D E9   STMFD  SP!, {R7,LR}
__text:00002910 0D 70 A0 E1   MOV    R7, SP
__text:00002914 14 D0 4D E2   SUB    SP, SP, #0x14
__text:00002918 70 05 01 E3   MOV    R0, #0x1570
__text:0000291C 07 C0 A0 E3   MOV    R12, #7
__text:00002920 00 00 40 E3   MOVT   R0, #0
__text:00002924 04 20 A0 E3   MOV    R2, #4
__text:00002928 00 00 8F E0   ADD    R0, PC, R0
__text:0000292C 06 30 A0 E3   MOV    R3, #6
__text:00002930 05 10 A0 E3   MOV    R1, #5
__text:00002934 00 20 8D E5   STR    R2, [SP,#0x1C+var_1C]
__text:00002938 0A 10 8D E9   STMFA  SP, {R1,R3,R12}
__text:0000293C 08 90 A0 E3   MOV    R9, #8
__text:00002940 01 10 A0 E3   MOV    R1, #1
__text:00002944 02 20 A0 E3   MOV    R2, #2
__text:00002948 03 30 A0 E3   MOV    R3, #3
__text:0000294C 10 90 8D E5   STR    R9, [SP,#0x1C+var_C]
__text:00002950 A4 05 00 EB   BL     _printf
__text:00002954 07 D0 A0 E1   MOV    SP, R7
__text:00002958 80 80 BD E8   LDMFD  SP!, {R7,PC}
\end{lstlisting}

\myindex{ARM!\Instructions!STMFA}
\myindex{ARM!\Instructions!STMIB}
Presque la même chose que ce que nous avons déjà vu, avec l'exception de l'instruction
\INS{STMFA} (Store Multiple Full Ascending), qui est un synonyme de l'instruction
\INS{STMIB} (Store Multiple Increment Before).
Cette instruction incrémente la valeur du registre \ac{SP} et écrit seulement
après la valeur registre suivant dans la mémoire, plutôt que d'effectuer ces
deux actions dans l'ordre inverse.

Une autre chose qui accroche le regard est que les instructions semblent être arrangées
de manière aléatoire.
Par exemple, la valeur dans le registre \Reg{0} est manipulée en trois endroits,
aux adresses \GTT{0x2918}, \GTT{0x2920} et \GTT{0x2928}, alors qu'il serait possible
de le faire en un seul endroit.

Toutefois, le compilateur qui optimise doit avoir ses propres raisons d'ordonner
les instructions pour avoir une plus grande efficacité à l'exécution.

D'habitude, le processeur essaye d'exécuter simultanément des instructions situées côte à côte.\\
Par exemple, des instructions comme \INS{MOVT R0, \#0} et \INS{ADD R0, PC, R0} ne
peuvent pas être exécutées simultanément puisqu'elles modifient toutes deux le
registre \Reg{0}.
D'un autre côté, les instructions \INS{MOVT R0, \#0} et \INS{MOV R2, \#4} peuvent
être exécutées simultanément puisque leurs effets n'interfèrent pas l'un avec l'autre
lors de leurs exécution.
Probablement que le compilateur essaye de générer du code arrangé de cette façon
(lorsque c'est possible).

\myparagraph{\OptimizingXcodeIV: \ThumbTwoMode}

\begin{lstlisting}[style=customasmARM]
__text:00002BA0               _printf_main2
__text:00002BA0
__text:00002BA0               var_1C = -0x1C
__text:00002BA0               var_18 = -0x18
__text:00002BA0               var_C  = -0xC
__text:00002BA0
__text:00002BA0 80 B5          PUSH     {R7,LR}
__text:00002BA2 6F 46          MOV      R7, SP
__text:00002BA4 85 B0          SUB      SP, SP, #0x14
__text:00002BA6 41 F2 D8 20    MOVW     R0, #0x12D8
__text:00002BAA 4F F0 07 0C    MOV.W    R12, #7
__text:00002BAE C0 F2 00 00    MOVT.W   R0, #0
__text:00002BB2 04 22          MOVS     R2, #4
__text:00002BB4 78 44          ADD      R0, PC  ; char *
__text:00002BB6 06 23          MOVS     R3, #6
__text:00002BB8 05 21          MOVS     R1, #5
__text:00002BBA 0D F1 04 0E    ADD.W    LR, SP, #0x1C+var_18
__text:00002BBE 00 92          STR      R2, [SP,#0x1C+var_1C]
__text:00002BC0 4F F0 08 09    MOV.W    R9, #8
__text:00002BC4 8E E8 0A 10    STMIA.W  LR, {R1,R3,R12}
__text:00002BC8 01 21          MOVS     R1, #1
__text:00002BCA 02 22          MOVS     R2, #2
__text:00002BCC 03 23          MOVS     R3, #3
__text:00002BCE CD F8 10 90    STR.W    R9, [SP,#0x1C+var_C]
__text:00002BD2 01 F0 0A EA    BLX      _printf
__text:00002BD6 05 B0          ADD      SP, SP, #0x14
__text:00002BD8 80 BD          POP      {R7,PC}
\end{lstlisting}

La sortie est presque la même que dans l'exemple précédent, avec l'exception que
des instructions Thumb sont utilisées à la place.
% FIXME: also STMIA is used instead of STMIB,
% which is why it uses LR, which is 4 bytes ahead of SP

\myparagraph{ARM64}

\mysubparagraph{GCC (Linaro) 4.9 \NonOptimizing}

\lstinputlisting[caption=GCC (Linaro) 4.9 \NonOptimizing,style=customasmARM]{patterns/03_printf/ARM/ARM8_O0_FR.lst}

Les 8 premiers arguments sont passés dans des registres X- ou W-: \ARMPCS.
Un pointeur de chaîne nécessite un registre 64-bit, donc il est passé dans \RegX{0}.
Toutes les autres valeurs ont un type \Tint 32-bit, donc elles sont stockées dans
la partie 32-bit des registres (W-).
Le 9ème argument (8) est passé par la pile.
En effet: il n'est pas possible de passer un grand nombre d'arguments par les registres,
car le nombre de registres est limité.

GCC (Linaro) 4.9 \Optimizing génère le même code.
}
\IT{\subsubsection{ARM: 8 arguments}

Usiamo nuovamente l'esempio con 9 argomenti della sezione precedente: \myref{example_printf8_x64}.

\lstinputlisting[style=customc]{patterns/03_printf/2.c}

\myparagraph{\OptimizingKeilVI: \ARMMode}

\lstinputlisting[style=customasmARM]{patterns/03_printf/ARM/tmp1.asm}

Il codice puo' essre diviso in piu' parti

\myindex{Function prologue}
\begin{itemize}
\item Preambolo della funzione:

\myindex{ARM!\Instructions!STR}
La prima istruzione \INS{STR LR, [SP,\#var\_4]!} salva \ac{LR} sullo stack, poiche' questo registro sara' usato per la chiamata a \printf.
Il punto esclamativo all fine indica il \emph{pre-index}.

Questo implica che \ac{SP} deve essere prima decrementato di 4, e successivamente \ac{LR} sara' salvato all'indirizzo memorizzato in \ac{SP}.
Tutto cio' e' simile a \PUSH in x86.
Maggiori informazioni qui: \myref{ARM_postindex_vs_preindex}.

\myindex{ARM!\Instructions!SUB}
La seconda istruzione \INS{SUB SP, SP, \#0x14} decrementa \ac{SP} (lo \gls{stack pointer}) per allocare \GTT{0x14} (20) byte sullo stack.
Infatti dobbiamo passare 5 valori a 32-bit tramite lo stack per la funzione \printf, e ciascuno di essi occupa 4 byte, che e' esattamente $5*4=20$.
Gli altri 4 valori a 32-bit saranno passati tramite registri.

\item Passaggio di 5, 6, 7 e 8 tramite lo stack: sono memorizzati nei registri \Reg{0}, \Reg{1}, \Reg{2} e \Reg{3}, rispettivamente.\\
Successivamente l'istruzione \INS{ADD R12, SP, \#0x18+var\_14} scrive l'indirizzo dello stack, dove queste 4 variabili saranno memorizzate,
nel registri \Reg{12}.
\myindex{IDA!var\_?}
\emph{var\_14} e' una macro assembly, uguale a -0x14, creata da \IDA per visualizzare in maniera conveniente il codice che accede allo stack.
Le macro \emph{var\_?} generate da \IDA riflettono le variabili locali nello stack.

Quindi, \GTT{SP+4} sara' memorizzato nel registro \Reg{12}.
\myindex{ARM!\Instructions!STMIA}
L'istruzione successiva \INS{STMIA R12, {R0-R3}} scrive il contenuto dei registri \Reg{0}-\Reg{3} alla memoria puntata da \Reg{12}.
\INS{STMIA} e' abbreviazione per \emph{Store Multiple Increment After}. 
\emph{Increment After} (incrementa dopo) implica che \Reg{12} deve essere incrementato di 4 dopo ciascuna scrittura di un valore nei registri.

\item Passaggio di 4 tramite lo stack: 4 e' memorizzato in \Reg{0} e questo valore, con l'aiuto dell'istruzione \INS{STR R0, [SP,\#0x18+var\_18]}, viene salvato sullo stack.
\emph{var\_18} e' -0x18, quindi l'offset deve essere 0, da cui il valore dal registro \Reg{0} (4) sara' scritto all'indirizzo memorizzato in \ac{SP}.

\item Passaggio di 1, 2 e 3 tramite registri:
I valori dei primi 3 numeri (a, b, c) (1, 2, 3 rispettivamente) sono passatti attraverso i registri
\Reg{1}, \Reg{2} e \Reg{3}
poco prima della chiamata a \printf call, e gli altri 5 valori sono passati tramite lo stack:

\item chiamata a \printf.

\myindex{Function epilogue}
\item epilogo della funzione:

% TODO to be synced with EN version:
L'istruzione \INS{ADD SP, SP, \#0x14} ripristina il puntatore \ac{SP} al suo valore precedente, pulendo quindi lo stack.
Ovviamente quello che era stato memorizzato nello stack rimarra li', e sara' probabilmente riscritto interamente durante l'e
secuzione delle funzioni seguenti.

\myindex{ARM!\Instructions!LDR}
L'istruzione \INS{LDR PC, [SP+4+var\_4],\#4} carica il valore di \ac{LR} salvato dallo stack nel nel registro \ac{PC}, causando quindi
l'uscita dalla funzione.
Non c'e' il punto esclamativo ---infatti \ac{PC} e' caricato prima dall'indirizzo memorizzato in \ac{SP} ($4+var\_4=4+(-4)=0$, questa istruzione e' quindi analoga a \INS{LDR PC, [SP],\#4}), e successivamente \ac{SP} e' incrementato di 4.
Questo e' detto \emph{post-index}\footnote{Maggiori dettagli: \myref{ARM_postindex_vs_preindex}.}.
Perche' \IDA mostra l'istruzione in quel modo?
Perche' vuole illustrare il layout dello stack ed il fatto che \GTT{var\_4} e' allocata per salvare il valore di \ac{LR} nello stack locale.
Questa istruzione e' piu' o meno simile a \INS{POP PC} in x86\footnote{E' impossibile settare il valore di \GTT{IP/EIP/RIP} usando \POP in x86, ma in ogni caso hai capito l'analogia.}.

\end{itemize}

\myparagraph{\OptimizingKeilVI: \ThumbMode}

\lstinputlisting[style=customasmARM]{patterns/03_printf/ARM/tmp2.asm}

L'output e' quasi identico al precedente esempio. Tuttavia questo e' codice Thumb e i valori sono disposti nello stack in modo differente:
8 per primo, quindi 5, 6, 7, e infine 4.

\myparagraph{\OptimizingXcodeIV: \ARMMode}

\begin{lstlisting}[style=customasmARM]
__text:0000290C             _printf_main2
__text:0000290C
__text:0000290C             var_1C = -0x1C
__text:0000290C             var_C  = -0xC
__text:0000290C
__text:0000290C 80 40 2D E9   STMFD  SP!, {R7,LR}
__text:00002910 0D 70 A0 E1   MOV    R7, SP
__text:00002914 14 D0 4D E2   SUB    SP, SP, #0x14
__text:00002918 70 05 01 E3   MOV    R0, #0x1570
__text:0000291C 07 C0 A0 E3   MOV    R12, #7
__text:00002920 00 00 40 E3   MOVT   R0, #0
__text:00002924 04 20 A0 E3   MOV    R2, #4
__text:00002928 00 00 8F E0   ADD    R0, PC, R0
__text:0000292C 06 30 A0 E3   MOV    R3, #6
__text:00002930 05 10 A0 E3   MOV    R1, #5
__text:00002934 00 20 8D E5   STR    R2, [SP,#0x1C+var_1C]
__text:00002938 0A 10 8D E9   STMFA  SP, {R1,R3,R12}
__text:0000293C 08 90 A0 E3   MOV    R9, #8
__text:00002940 01 10 A0 E3   MOV    R1, #1
__text:00002944 02 20 A0 E3   MOV    R2, #2
__text:00002948 03 30 A0 E3   MOV    R3, #3
__text:0000294C 10 90 8D E5   STR    R9, [SP,#0x1C+var_C]
__text:00002950 A4 05 00 EB   BL     _printf
__text:00002954 07 D0 A0 E1   MOV    SP, R7
__text:00002958 80 80 BD E8   LDMFD  SP!, {R7,PC}
\end{lstlisting}

\myindex{ARM!\Instructions!STMFA}
\myindex{ARM!\Instructions!STMIB}
Quasi lo stesso codice visto prima, ad eccezione dell'istruzione \INS{STMFA} (Store Multiple Full Ascending),
che e' sinonimo di \INS{STMIB} (Store Multiple Increment Before). 
Questa istruzione incrementa il valore nel registro \ac{SP} e solo successivamente scrive il prossimo valore del registro in memoria, invece che operare le due azioni in ordine inverso.

Un'altra cosa che salta all'occhio e' che le istruzioni sono disponste in maniera apparentemente casuale.
Ad esempio, il valore nel registro \Reg{0} e' manipolato in tre posti diversi
agli indirizzi \GTT{0x2918}, \GTT{0x2920} e \GTT{0x2928}, quando invece sarebbe stato possibile farlo in un punto solo.

Ad ogni modo, il compilatore ottimizzante avra' avuto le sue ragioni per ordinare le istruzioni in questa maniera ed ottenere una maggiore efficacia durante l'esecuzione del codice.

Solitamente il processero prova ad eseguire simultaneamentele istruzioni vicine.
Ad esempio, istruzioni come \INS{MOVT R0, \#0} e
\INS{ADD R0, PC, R0} non possono essere eseguite simultaneamente poiche' entrambe modificano il registro \Reg{0}. 
D'altra parte, \INS{MOVT R0, \#0} e \INS{MOV R2, \#4} 
possono invece essere eseguite simultaneamente poiche' l'effetto della loro esecuzione non genera conflitti tra loro.
Presumibilmente, il compilatore prova a generare codice in questo modo (quando possibile).
 
\myparagraph{\OptimizingXcodeIV: \ThumbTwoMode}

\begin{lstlisting}[style=customasmARM]
__text:00002BA0               _printf_main2
__text:00002BA0
__text:00002BA0               var_1C = -0x1C
__text:00002BA0               var_18 = -0x18
__text:00002BA0               var_C  = -0xC
__text:00002BA0
__text:00002BA0 80 B5          PUSH     {R7,LR}
__text:00002BA2 6F 46          MOV      R7, SP
__text:00002BA4 85 B0          SUB      SP, SP, #0x14
__text:00002BA6 41 F2 D8 20    MOVW     R0, #0x12D8
__text:00002BAA 4F F0 07 0C    MOV.W    R12, #7
__text:00002BAE C0 F2 00 00    MOVT.W   R0, #0
__text:00002BB2 04 22          MOVS     R2, #4
__text:00002BB4 78 44          ADD      R0, PC  ; char *
__text:00002BB6 06 23          MOVS     R3, #6
__text:00002BB8 05 21          MOVS     R1, #5
__text:00002BBA 0D F1 04 0E    ADD.W    LR, SP, #0x1C+var_18
__text:00002BBE 00 92          STR      R2, [SP,#0x1C+var_1C]
__text:00002BC0 4F F0 08 09    MOV.W    R9, #8
__text:00002BC4 8E E8 0A 10    STMIA.W  LR, {R1,R3,R12}
__text:00002BC8 01 21          MOVS     R1, #1
__text:00002BCA 02 22          MOVS     R2, #2
__text:00002BCC 03 23          MOVS     R3, #3
__text:00002BCE CD F8 10 90    STR.W    R9, [SP,#0x1C+var_C]
__text:00002BD2 01 F0 0A EA    BLX      _printf
__text:00002BD6 05 B0          ADD      SP, SP, #0x14
__text:00002BD8 80 BD          POP      {R7,PC}
\end{lstlisting}

L'output e' quasi lo stesso dell'esempio precedente, ad eccezione dell'uso di istruzioni Thumb. 
% FIXME: also STMIA is used instead of STMIB,
% which is why it uses LR, which is 4 bytes ahead of SP

\myparagraph{ARM64}

\mysubparagraph{\NonOptimizing GCC (Linaro) 4.9}

% TODO translate to Italian:
\lstinputlisting[caption=\NonOptimizing GCC (Linaro) 4.9,style=customasmARM]{patterns/03_printf/ARM/ARM8_O0_EN.lst}

I primi 8 argomenti sono passati nei registri X- o W-: \ARMPCS.
Un puntatore ad una string richiede un registro a 64-bit, quindi e' passato in \RegX{0}.
Tutti gli altri valori hanno timpo \Tint a 32-bit, quindi sono memorizzati nella parte a 32-bit dei registri (W-).
Il nono argomento (8) e' passato tramite lo stack.
Infatti non e' possibile passare un grande numero di argomenti tramite registri, in quanto il loro numero e' limitato.

\Optimizing GCC (Linaro) 4.9 genera lo stesso codice.
}
\JA{\subsubsection{ARM: 8つの引数}

前のセクションの9つの引数を使って例を再利用してみましょう：\myref{example_printf8_x64}

\lstinputlisting[style=customc]{patterns/03_printf/2.c}

\myparagraph{\OptimizingKeilVI: \ARMMode}

\lstinputlisting[style=customasmARM]{patterns/03_printf/ARM/tmp1.asm}

このコードはいくつかの部分に分けることができます：

\myindex{Function prologue}
\begin{itemize}
\item 関数プロローグ:

\myindex{ARM!\Instructions!STR}
最初の\INS{STR LR, [SP,\#var\_4]!}命令は、このレジスタを \printf 呼び出しに使用する予定であるため、\ac{LR}をスタックに保存します。 
最後の感嘆符は、事前索引を示します。

これは、まず\ac{SP}を4減少させた後、\ac{SP}に格納されたアドレスに\ac{LR}を保存することを意味します。
これはx86のPUSHに似ています。 
もっと読む：\myref{ARM_postindex_vs_preindex}

\myindex{ARM!\Instructions!SUB}
第2の\INS{SUB SP, SP, \#0x14}命令は、スタック上に0x14（20）バイトを割り当てるために\ac{SP}（\gls{stack pointer}）を減少させる。 
実際には、スタックを介して5つの32ビット値を \printf 関数に渡さなければならず、それぞれが4バイトを占めます。これは正確に$5*4=20$です。
他の4つの32ビット値は、レジスタに通される。

\item スタックを介して5,6,7および8を渡す：それらはそれぞれ\Reg{0}, \Reg{1}, \Reg{2} と \Reg{3}レジスタに格納される。\\
次に、\INS{ADD R12, SP, \#0x18+var\_14}命令は、これら4つの変数が格納されるスタックアドレスを\Reg{12}レジスタに書き込みます。 
\myindex{IDA!var\_?}
\emph{var\_14}は、スタックにアクセスするコードを便利に表示するために \IDA によって作成された-0x14に等しいアセンブリマクロです。 
\IDA によって生成される\emph{var\_?}マクロは、スタック内のローカル変数を反映します。

したがって、\GTT{SP+4}は\Reg{12}レジスタに格納されます。\\
\myindex{ARM!\Instructions!STMIA}
次の\INS{STMIA R12, {R0-R3}}命令は、レジスタ\Reg{0}-\Reg{3}の内容を\Reg{12}が指すメモリに書き込みます。 
\INS{STMIA}は、\emph{後に複数のインクリメント}を格納します。 
\emph{Increment After}は、各レジスタ値が書き込まれた後に\Reg{12}が4ずつ増加することを意味する。

\item スタックを介して4を渡す：4が\Reg{0}に格納され、\INS{STR R0, [SP,\#0x18+var\_18]}命令の助けを借りてこの値がスタックに保存されます。
\emph{var\_18}は-0x18なので、オフセットは0になります。したがって、\Reg{0}レジスタ（4）の値は\ac{SP}に書き込まれたアドレスに書き込まれます。

\item レジスタ経由で1,2,3を渡す
最初の3つの数値（a、b、c）（それぞれ1,2,3）の値は、printf（）呼び出しの直前に\Reg{1}、 \Reg{2}、 \Reg{3}レジスタに渡されます。
他の5つの値はスタックを介して渡されます。

\item \printf 呼び出し。

\myindex{Function epilogue}
\item 関数エピローグ：

\INS{ADD SP, SP, \#0x14}命令は\ac{SP}ポインタを元の値に戻して、スタックに格納されているものをすべて取り消します。
もちろん、スタックに格納されているものはそこにとどまりますが、後続の関数の実行中にすべて書き換えられます。

\myindex{ARM!\Instructions!LDR}
\INS{LDR PC, [SP+4+var\_4],\#4}命令は、保存された\ac{LR}値をスタックから\ac{PC}レジスタにロードして、機能を終了させます。
感嘆符はありません。次に、\ac{SP}は（$4+var\_4=4+(-4)=0$に格納されているアドレスから最初にロードされるため、この命令は\INS{LDR PC, [SP],\#4}に似ています）、 \ac{SP}は4だけ増加します。これは\emph{ポストインデックス}\footnote{もっと読む： \myref{ARM_postindex_vs_preindex}}と呼ばれます。
なぜ \IDA はそのような指示を表示するのですか？
なぜなら、スタックレイアウトと、\GTT{var\_4}がローカルスタックの\ac{LR}値を保存するために割り当てられているという事実を説明したいからです。
この命令は、x86の\INS{POP PC}と多少似ています。
\footnote{x86では \POP を使って\GTT{IP/EIP/RIP}の値を設定することは不可能ですが、アナロジーとしてはよいでしょう}

\end{itemize}

\myparagraph{\OptimizingKeilVI: \ThumbMode}

\lstinputlisting[style=customasmARM]{patterns/03_printf/ARM/tmp2.asm}

出力は前の例とほぼ同じです。 ただし、これはThumbコードであり、値はスタックに別々にパックされます。
8が先に進み、次に5,6,7、および4が3番目に進みます。

\myparagraph{\OptimizingXcodeIV: \ARMMode}

\begin{lstlisting}[style=customasmARM]
__text:0000290C             _printf_main2
__text:0000290C
__text:0000290C             var_1C = -0x1C
__text:0000290C             var_C  = -0xC
__text:0000290C
__text:0000290C 80 40 2D E9   STMFD  SP!, {R7,LR}
__text:00002910 0D 70 A0 E1   MOV    R7, SP
__text:00002914 14 D0 4D E2   SUB    SP, SP, #0x14
__text:00002918 70 05 01 E3   MOV    R0, #0x1570
__text:0000291C 07 C0 A0 E3   MOV    R12, #7
__text:00002920 00 00 40 E3   MOVT   R0, #0
__text:00002924 04 20 A0 E3   MOV    R2, #4
__text:00002928 00 00 8F E0   ADD    R0, PC, R0
__text:0000292C 06 30 A0 E3   MOV    R3, #6
__text:00002930 05 10 A0 E3   MOV    R1, #5
__text:00002934 00 20 8D E5   STR    R2, [SP,#0x1C+var_1C]
__text:00002938 0A 10 8D E9   STMFA  SP, {R1,R3,R12}
__text:0000293C 08 90 A0 E3   MOV    R9, #8
__text:00002940 01 10 A0 E3   MOV    R1, #1
__text:00002944 02 20 A0 E3   MOV    R2, #2
__text:00002948 03 30 A0 E3   MOV    R3, #3
__text:0000294C 10 90 8D E5   STR    R9, [SP,#0x1C+var_C]
__text:00002950 A4 05 00 EB   BL     _printf
__text:00002954 07 D0 A0 E1   MOV    SP, R7
__text:00002958 80 80 BD E8   LDMFD  SP!, {R7,PC}
\end{lstlisting}

\myindex{ARM!\Instructions!STMFA}
\myindex{ARM!\Instructions!STMIB}
\INS{STMFA}（Store Multiple Increment Before）命令の同義語である\INS{STMIB}（Store Multiple Full Ascending）命令を除き、既に見てきたものとほぼ同じです。
この命令は、\ac{SP}レジスタの値を増加させ、逆の順序でこれら2つの動作を実行するのではなく、次のレジスタ値をメモリに書き込むだけです。

目を引くもう一つのことは、命令が一見無作為に配置されていることです。
例えば、R0レジスタの値は、アドレス\GTT{0x2918}, \GTT{0x2920} and \GTT{0x2928}の3箇所で操作できます。

しかしながら、最適化コンパイラは、実行中により高い効率を達成するために、命令をどのように順序付けするかに関する独自の理由を有することができる。

通常、プロセッサは、並んで配置された命令を同時に実行しようと試みます。
たとえば、\INS{MOVT R0, \#0}、\INS{ADD R0, PC, R0}などの命令は、
両方とも\Reg{0}レジスタを変更するため、同時に実行することはできません。
一方、\INS{MOVT R0, \#0}、\INS{MOV R2, \#4}命令は、
実行の影響が互いに矛盾しないため、同時に実行することができます。
おそらく、コンパイラはそのような方法でコードを生成しようと試みます（どこでも可能です）。

\myparagraph{\OptimizingXcodeIV: \ThumbTwoMode}

\begin{lstlisting}[style=customasmARM]
__text:00002BA0               _printf_main2
__text:00002BA0
__text:00002BA0               var_1C = -0x1C
__text:00002BA0               var_18 = -0x18
__text:00002BA0               var_C  = -0xC
__text:00002BA0
__text:00002BA0 80 B5          PUSH     {R7,LR}
__text:00002BA2 6F 46          MOV      R7, SP
__text:00002BA4 85 B0          SUB      SP, SP, #0x14
__text:00002BA6 41 F2 D8 20    MOVW     R0, #0x12D8
__text:00002BAA 4F F0 07 0C    MOV.W    R12, #7
__text:00002BAE C0 F2 00 00    MOVT.W   R0, #0
__text:00002BB2 04 22          MOVS     R2, #4
__text:00002BB4 78 44          ADD      R0, PC  ; char *
__text:00002BB6 06 23          MOVS     R3, #6
__text:00002BB8 05 21          MOVS     R1, #5
__text:00002BBA 0D F1 04 0E    ADD.W    LR, SP, #0x1C+var_18
__text:00002BBE 00 92          STR      R2, [SP,#0x1C+var_1C]
__text:00002BC0 4F F0 08 09    MOV.W    R9, #8
__text:00002BC4 8E E8 0A 10    STMIA.W  LR, {R1,R3,R12}
__text:00002BC8 01 21          MOVS     R1, #1
__text:00002BCA 02 22          MOVS     R2, #2
__text:00002BCC 03 23          MOVS     R3, #3
__text:00002BCE CD F8 10 90    STR.W    R9, [SP,#0x1C+var_C]
__text:00002BD2 01 F0 0A EA    BLX      _printf
__text:00002BD6 05 B0          ADD      SP, SP, #0x14
__text:00002BD8 80 BD          POP      {R7,PC}
\end{lstlisting}

Thumb命令が代わりに使用される点を除いて、出力は前の例とほぼ同じです。
% FIXME: also STMIA is used instead of STMIB,
% which is why it uses LR, which is 4 bytes ahead of SP

\myparagraph{ARM64}

\mysubparagraph{\NonOptimizing GCC (Linaro) 4.9}

\lstinputlisting[caption=\NonOptimizing GCC (Linaro) 4.9,style=customasmARM]{patterns/03_printf/ARM/ARM8_O0_JA.lst}

最初の8つの引数は、XレジスタまたはWレジスタに渡されます。 \ARMPCS
文字列ポインタは64ビットのレジスタを必要とするため、\RegX{0}で渡されます。 
それ以外の値はすべてint型32ビット型なので、レジスタ（W-）の32ビット部分に格納されます。 
第9引数（8）はスタックを介して渡されます。 
実際には、レジスタの数が限られているため、多数の引数をレジスタに渡すことはできません。

\Optimizing GCC (Linaro) 4.9 は同じコードを生成します。
}

\EN{\subsection{MIPS}

\subsubsection{3 arguments}

\myparagraph{\Optimizing GCC 4.4.5}

The main difference with the \q{\HelloWorldSectionName} example is that in this case \printf is called
instead of \puts and 3 more arguments are passed through the registers \$5\dots \$7 (or \$A0\dots \$A2).
That is why these registers are prefixed with A-, which implies they are used for function arguments passing.

\lstinputlisting[caption=\Optimizing GCC 4.4.5 (\assemblyOutput),style=customasmMIPS]{patterns/03_printf/MIPS/printf3.O3_EN.s}

\lstinputlisting[caption=\Optimizing GCC 4.4.5 (IDA),style=customasmMIPS]{patterns/03_printf/MIPS/printf3.O3.IDA_EN.lst}

\IDA has coalesced pair of \INS{LUI} and \INS{ADDIU} instructions into one \INS{LA} pseudo instruction.
That's why there are no instruction at address 0x1C: because \INS{LA} \emph{occupies} 8 bytes.

\myparagraph{\NonOptimizing GCC 4.4.5}

\NonOptimizing GCC is more verbose:

\lstinputlisting[caption=\NonOptimizing GCC 4.4.5 (\assemblyOutput),style=customasmMIPS]{patterns/03_printf/MIPS/printf3.O0_EN.s}

\lstinputlisting[caption=\NonOptimizing GCC 4.4.5 (IDA),style=customasmMIPS]{patterns/03_printf/MIPS/printf3.O0.IDA_EN.lst}

\subsubsection{8 arguments}

Let's use again the example with 9 arguments from the previous section: \myref{example_printf8_x64}.

\lstinputlisting[style=customc]{patterns/03_printf/2.c}

\myparagraph{\Optimizing GCC 4.4.5}

Only the first 4 arguments are passed in the \$A0 \dots \$A3 registers, the rest are passed via the stack.
\myindex{MIPS!O32}

This is the O32 calling convention (which is the most common one in the MIPS world).
Other calling conventions (like N32) may use the registers for different purposes.

\myindex{MIPS!\Instructions!SW}

\INS{SW} abbreviates \q{Store Word} (from register to memory).
MIPS lacks instructions for storing a value into memory, so an instruction pair has to be used instead (\INS{LI}/\INS{SW}).

\lstinputlisting[caption=\Optimizing GCC 4.4.5 (\assemblyOutput),style=customasmMIPS]{patterns/03_printf/MIPS/printf8.O3_EN.s}

\lstinputlisting[caption=\Optimizing GCC 4.4.5 (IDA),style=customasmMIPS]{patterns/03_printf/MIPS/printf8.O3.IDA_EN.lst}

\myparagraph{\NonOptimizing GCC 4.4.5}

\NonOptimizing GCC is more verbose:

\lstinputlisting[caption=\NonOptimizing GCC 4.4.5 (\assemblyOutput),style=customasmMIPS]{patterns/03_printf/MIPS/printf8.O0_EN.s}

\lstinputlisting[caption=\NonOptimizing GCC 4.4.5 (IDA),style=customasmMIPS]{patterns/03_printf/MIPS/printf8.O0.IDA_EN.lst}

}
\RU{\subsection{MIPS}

\subsubsection{3 аргумента}

\myparagraph{\Optimizing GCC 4.4.5}

Главное отличие от примера \q{\HelloWorldSectionName} в том, что здесь на самом деле
вызывается \printf вместо \puts и ещё три аргумента передаются в регистрах  \$5\dots \$7 (или \$A0\dots \$A2).
Вот почему эти регистры имеют префикс A-. Это значит, что они используются для передачи аргументов.

\lstinputlisting[caption=\Optimizing GCC 4.4.5 (\assemblyOutput),style=customasmMIPS]{patterns/03_printf/MIPS/printf3.O3_RU.s}

\lstinputlisting[caption=\Optimizing GCC 4.4.5 (IDA),style=customasmMIPS]{patterns/03_printf/MIPS/printf3.O3.IDA_RU.lst}

\IDA объединила пару инструкций \INS{LUI} и \INS{ADDIU} в одну псевдоинструкцию \INS{LA}.
Вот почему здесь нет инструкции по адресу 0x1C: потому что \INS{LA} \emph{занимает} 8 байт.

\myparagraph{\NonOptimizing GCC 4.4.5}

\NonOptimizing GCC более многословен:

\lstinputlisting[caption=\NonOptimizing GCC 4.4.5 (\assemblyOutput),style=customasmMIPS]{patterns/03_printf/MIPS/printf3.O0_RU.s}

\lstinputlisting[caption=\NonOptimizing GCC 4.4.5 (IDA),style=customasmMIPS]{patterns/03_printf/MIPS/printf3.O0.IDA_RU.lst}

\subsubsection{8 аргументов}

Снова воспользуемся примером с 9-ю аргументами из предыдущей секции: \myref{example_printf8_x64}.

\lstinputlisting[style=customc]{patterns/03_printf/2.c}

\myparagraph{\Optimizing GCC 4.4.5}

Только 4 первых аргумента передаются в регистрах \$A0 \dots \$A3, так что остальные передаются через стек.

\myindex{MIPS!O32}
Это соглашение о вызовах O32 (самое популярное в мире MIPS).
Другие соглашения о вызовах (например N32) могут наделять регистры другими функциями.

\myindex{MIPS!\Instructions!SW}
\INS{SW} означает \q{Store Word} (записать слово из регистра в память).
В MIPS нет инструкции для записи значения в память, так что для этого используется пара инструкций (\INS{LI}/\INS{SW}).

\lstinputlisting[caption=\Optimizing GCC 4.4.5 (\assemblyOutput),style=customasmMIPS]{patterns/03_printf/MIPS/printf8.O3_RU.s}

\lstinputlisting[caption=\Optimizing GCC 4.4.5 (IDA),style=customasmMIPS]{patterns/03_printf/MIPS/printf8.O3.IDA_RU.lst}

\myparagraph{\NonOptimizing GCC 4.4.5}

\NonOptimizing GCC более многословен:

\lstinputlisting[caption=\NonOptimizing GCC 4.4.5 (\assemblyOutput),style=customasmMIPS]{patterns/03_printf/MIPS/printf8.O0_RU.s}

\lstinputlisting[caption=\NonOptimizing GCC 4.4.5 (IDA),style=customasmMIPS]{patterns/03_printf/MIPS/printf8.O0.IDA_RU.lst}

}
\IT{\subsection{MIPS}

\subsubsection{3 argomenti}

\myparagraph{\Optimizing GCC 4.4.5}

La differenza principale con l'esempio \q{\HelloWorldSectionName} e' che in questo caso \printf e' chiamata 
al posto di \puts, e 3 argomenti aggiuntivi sono passati attraverso i registri \$5\dots \$7 (o \$A0\dots \$A2).
Questo e' il motivo per cui questi registri hanno il prefisso A-, che implica il loro uso per il passaggio di argomenti di funzioni.

% TODO translate to Italian:
\lstinputlisting[caption=\Optimizing GCC 4.4.5 (\assemblyOutput),style=customasmMIPS]{patterns/03_printf/MIPS/printf3.O3_EN.s}

% TODO translate to Italian:
\lstinputlisting[caption=\Optimizing GCC 4.4.5 (IDA),style=customasmMIPS]{patterns/03_printf/MIPS/printf3.O3.IDA_EN.lst}

\IDA ha fuso le coppie di istruzioni \INS{LUI} e \INS{ADDIU} in una unica pseudoistruzione \INS{LA}.
Questo e' il motivo per cui non c'e' nessuna istruzione all'indirizzo 0x1C: perche' \INS{LA} \emph{occupa} 8 byte.%

\myparagraph{\NonOptimizing GCC 4.4.5}

\NonOptimizing GCC e' piu' verboso:

% TODO translate to Italian:
\lstinputlisting[caption=\NonOptimizing GCC 4.4.5 (\assemblyOutput),style=customasmMIPS]{patterns/03_printf/MIPS/printf3.O0_EN.s}

% TODO translate to Italian:
\lstinputlisting[caption=\NonOptimizing GCC 4.4.5 (IDA),style=customasmMIPS]{patterns/03_printf/MIPS/printf3.O0.IDA_EN.lst}

\subsubsection{8 argomenti}

Usiamo nuovamente l'esempio con 9 argomenti dalla sezione prcedente: \myref{example_printf8_x64}.

\lstinputlisting[style=customc]{patterns/03_printf/2.c}

\myparagraph{\Optimizing GCC 4.4.5}

Solo i primi 4 argomenti sono passati nei registri \$A0 \dots \$A3, gli altri sono passati tramite lo stack.
\myindex{MIPS!O32}

Questa e' la calling convention O32 (che e' la piu' comune nel mondo MIPS).
Altre calling conventions (come N32) possono usare i registri per scopi diversi.

\myindex{MIPS!\Instructions!SW}

\INS{SW} e' l'abbreviazione di \q{Store Word} (da un registro alla memoria).
MIPS manca di istruzioni per memorizzare un valore in memoria, e' quindi necessario usare una coppia di istruzioni (LI/SW).

% TODO translate to Italian:
\lstinputlisting[caption=\Optimizing GCC 4.4.5 (\assemblyOutput),style=customasmMIPS]{patterns/03_printf/MIPS/printf8.O3_EN.s}

% TODO translate to Italian:
\lstinputlisting[caption=\Optimizing GCC 4.4.5 (IDA),style=customasmMIPS]{patterns/03_printf/MIPS/printf8.O3.IDA_EN.lst}

\myparagraph{\NonOptimizing GCC 4.4.5}

\NonOptimizing GCC e' piu' verboso:

% TODO translate to Italian:
\lstinputlisting[caption=\NonOptimizing GCC 4.4.5 (\assemblyOutput),style=customasmMIPS]{patterns/03_printf/MIPS/printf8.O0_EN.s}

% TODO translate to Italian:
\lstinputlisting[caption=\NonOptimizing GCC 4.4.5 (IDA),style=customasmMIPS]{patterns/03_printf/MIPS/printf8.O0.IDA_EN.lst}
}
\FR{\subsection{MIPS}

\subsubsection{3 arguments}

\myparagraph{GCC 4.4.5 \Optimizing}

La différence principale avec l'exemple \q{\HelloWorldSectionName} est que dans ce cas, \printf est appelée
à la place de \puts et 3 arguments de plus sont passés à travers les registres \$5\dots \$7 (ou \$A0\dots \$A2).
C'est pourquoi ces registres sont préfixés avec A-, ceci sous-entend qu'ils
sont utilisés pour le passage des arguments aux fonctions.

\lstinputlisting[caption=GCC 4.4.5 \Optimizing (\assemblyOutput),style=customasmMIPS]{patterns/03_printf/MIPS/printf3.O3_FR.s}

\lstinputlisting[caption=GCC 4.4.5 \Optimizing (IDA),style=customasmMIPS]{patterns/03_printf/MIPS/printf3.O3.IDA_FR.lst}

\IDA a agrégé la paire d'instructions \INS{LUI} et \INS{ADDIU} en une pseudo instruction \INS{LA}.
C'est pourquoi il n'y a pas d'instruction à l'adresse 0x1C: car \INS{LA} \emph{occupe} 8 octets.

\myparagraph{GCC 4.4.5 \NonOptimizing}

GCC \NonOptimizing est plus verbeux:

\lstinputlisting[caption=GCC 4.4.5 \NonOptimizing (\assemblyOutput),style=customasmMIPS]{patterns/03_printf/MIPS/printf3.O0_FR.s}

\lstinputlisting[caption=GCC 4.4.5 \NonOptimizing (IDA),style=customasmMIPS]{patterns/03_printf/MIPS/printf3.O0.IDA_FR.lst}

\subsubsection{8 arguments}

Utilisons encore l'exemple de la section précédente avec 9 arguments: \myref{example_printf8_x64}.

\lstinputlisting[style=customc]{patterns/03_printf/2.c}

\myparagraph{GCC 4.4.5 \Optimizing}

Seul les 4 premiers arguments sont passés dans les registres \$A0 \dots \$A3,
les autres sont passés par la pile.
\myindex{MIPS!O32}

C'est la convention d'appel O32 (qui est la plus commune dans le monde MIPS).
D'autres conventions d'appel (comme N32) peuvent utiliser les registres à d'autres fins.

\myindex{MIPS!\Instructions!SW}

\INS{SW} est l'abbréviation de \q{Store Word} (depuis un registre vers la mémoire).
En MIPS, il manque une instructions pour stocker une valeur dans la mémoire, donc
une paire d'instruction doit être utilisée à la place (\INS{LI}/\INS{SW}).

\lstinputlisting[caption=GCC 4.4.5 \Optimizing (\assemblyOutput),style=customasmMIPS]{patterns/03_printf/MIPS/printf8.O3_FR.s}

\lstinputlisting[caption=GCC 4.4.5 \Optimizing (IDA),style=customasmMIPS]{patterns/03_printf/MIPS/printf8.O3.IDA_FR.lst}

\myparagraph{GCC 4.4.5 \NonOptimizing}

GCC \NonOptimizing est plus verbeux:

\lstinputlisting[caption=\NonOptimizing GCC 4.4.5 (\assemblyOutput),style=customasmMIPS]{patterns/03_printf/MIPS/printf8.O0_FR.s}

\lstinputlisting[caption=\NonOptimizing GCC 4.4.5 (IDA),style=customasmMIPS]{patterns/03_printf/MIPS/printf8.O0.IDA_FR.lst}

}
\JA{\subsection{MIPS}

\subsubsection{3 arguments}

\myparagraph{\Optimizing GCC 4.4.5}

\q{\HelloWorldSectionName} の例との主な違いは、 \puts の代わりに \printf が呼び出され、
さらに3つの引数がレジスタ \$5\dots \$7 (または \$A0\dots \$A2 )に渡されるという点です。 
そのため、これらのレジスタの前にA-が付いています。これは、関数引数の受け渡しに使用されることを意味します。

\lstinputlisting[caption=\Optimizing GCC 4.4.5 (\assemblyOutput),style=customasmMIPS]{patterns/03_printf/MIPS/printf3.O3_JA.s}

\lstinputlisting[caption=\Optimizing GCC 4.4.5 (IDA),style=customasmMIPS]{patterns/03_printf/MIPS/printf3.O3.IDA_JA.lst}

\IDA は、\INS{LUI}および\INS{ADDIU}命令のペアを1つの\INS{LA}疑似命令に統合しました。 
だから、アドレス0x1Cに命令がないのは、\INS{LA}が8バイトを\emph{占有}しているからです。

\myparagraph{\NonOptimizing GCC 4.4.5}

\NonOptimizing GCC はもっと冗長です。

\lstinputlisting[caption=\NonOptimizing GCC 4.4.5 (\assemblyOutput),style=customasmMIPS]{patterns/03_printf/MIPS/printf3.O0_JA.s}

\lstinputlisting[caption=\NonOptimizing GCC 4.4.5 (IDA),style=customasmMIPS]{patterns/03_printf/MIPS/printf3.O0.IDA_JA.lst}

\subsubsection{8つの引数}

前のセクションの9つの引数を使用して、例を再度使用してみましょう：\myref{example_printf8_x64}

\lstinputlisting[style=customc]{patterns/03_printf/2.c}

\myparagraph{\Optimizing GCC 4.4.5}

最初の4つの引数だけが \$A0 \dots \$A3 レジスタに渡され、残りはスタックを介して渡されます。
\myindex{MIPS!O32}

これはO32呼び出し規約（MIPS世界で最も一般的なもの）です。 
他の呼び出し規則（N32のような）は、異なる目的のためにレジスタを使用するかもしれません。

\myindex{MIPS!\Instructions!SW}

\INS{SW}は\q{Store Word}（レジスタからメモリへ）の略語です。 
MIPSには値をメモリに格納する命令がないため、代わりに命令ペア（\INS{LI}/\INS{SW}）を使用する必要があります。

\lstinputlisting[caption=\Optimizing GCC 4.4.5 (\assemblyOutput),style=customasmMIPS]{patterns/03_printf/MIPS/printf8.O3_JA.s}

\lstinputlisting[caption=\Optimizing GCC 4.4.5 (IDA),style=customasmMIPS]{patterns/03_printf/MIPS/printf8.O3.IDA_JA.lst}

\myparagraph{\NonOptimizing GCC 4.4.5}

\NonOptimizing GCC はもっと冗長です。

\lstinputlisting[caption=\NonOptimizing GCC 4.4.5 (\assemblyOutput),style=customasmMIPS]{patterns/03_printf/MIPS/printf8.O0_JA.s}

\lstinputlisting[caption=\NonOptimizing GCC 4.4.5 (IDA),style=customasmMIPS]{patterns/03_printf/MIPS/printf8.O0.IDA_JA.lst}

}


\subsection{\Conclusion{}}

Here is a rough skeleton of the function call:

\begin{lstlisting}[caption=x86,style=customasmx86]
...
PUSH 3rd argument
PUSH 2nd argument
PUSH 1st argument
CALL function
; modify stack pointer (if needed)
\end{lstlisting}

\begin{lstlisting}[caption=x64 (MSVC),style=customasmx86]
MOV RCX, 1st argument
MOV RDX, 2nd argument
MOV R8, 3rd argument
MOV R9, 4th argument
...
PUSH 5th, 6th argument, etc. (if needed)
CALL function
; modify stack pointer (if needed)
\end{lstlisting}

\begin{lstlisting}[caption=x64 (GCC),style=customasmx86]
MOV RDI, 1st argument
MOV RSI, 2nd argument
MOV RDX, 3rd argument
MOV RCX, 4th argument
MOV R8, 5th argument
MOV R9, 6th argument
...
PUSH 7th, 8th argument, etc. (if needed)
CALL function
; modify stack pointer (if needed)
\end{lstlisting}

\begin{lstlisting}[caption=ARM,style=customasmARM]
MOV R0, 1st argument
MOV R1, 2nd argument
MOV R2, 3rd argument
MOV R3, 4th argument
; pass 5th, 6th argument, etc., in stack (if needed)
BL function
; modify stack pointer (if needed)
\end{lstlisting}

\begin{lstlisting}[caption=ARM64,style=customasmARM]
MOV X0, 1st argument
MOV X1, 2nd argument
MOV X2, 3rd argument
MOV X3, 4th argument
MOV X4, 5th argument
MOV X5, 6th argument
MOV X6, 7th argument
MOV X7, 8th argument
; pass 9th, 10th argument, etc., in stack (if needed)
BL function
; modify stack pointer (if needed)
\end{lstlisting}

\myindex{MIPS!O32}
\begin{lstlisting}[caption=MIPS (O32 calling convention),style=customasmMIPS]
LI $4, 1st argument ; AKA \$A0
LI $5, 2nd argument ; AKA \$A1
LI $6, 3rd argument ; AKA \$A2
LI $7, 4th argument ; AKA \$A3
; pass 5th, 6th argument, etc., in stack (if needed)
LW temp_reg, address of function
JALR temp_reg
\end{lstlisting}

\subsection{By the way}

\myindex{fastcall}
By the way, this difference between the arguments passing in x86, x64, 
fastcall, ARM and MIPS is a good illustration of the fact that the CPU is oblivious to how the arguments are passed to functions. 
It is also possible to create a hypothetical compiler able to pass arguments 
via a special structure without using stack at all.

\myindex{MIPS!O32}
MIPS \$A0 \dots \$A3 registers are labeled this way only for convenience (that is in the O32 calling convention).
Programmers may use any other register (well, maybe except \$ZERO) 
to pass data or use any other calling convention.

The \ac{CPU} is not aware of calling conventions whatsoever.

We may also recall how new coming assembly language programmers passing arguments into
other functions: usually via registers, without any explicit order, or even via global variables.
Of course, it works fine.

